\documentclass[11pt]{ctexart}
\usepackage[letterpaper,top=0.95in,bottom=0.78in,left=1.0in,right=1.0in]{geometry}
\usepackage{fontspec}
\setmainfont{Times New Roman}
\setCJKmainfont[ItalicFont=Songti SC]{Songti SC}
\setsansfont{Helvetica}
\setmonofont{Menlo}
\usepackage{microtype}
\usepackage{amsmath}
\usepackage{amssymb}
\usepackage{booktabs}
\usepackage{array}
\usepackage{longtable}
\usepackage{enumitem}
\usepackage{etoolbox}
\usepackage{cite}
\usepackage{xurl}
\usepackage[hidelinks]{hyperref}
\Urlmuskip=0mu plus 1mu
\setlength{\parindent}{2em}
\setlength{\parskip}{0pt}
\setlength{\emergencystretch}{2em}
\linespread{1.08}
\setlist[itemize]{leftmargin=2.3em,itemsep=0.15em,topsep=0.35em,parsep=0pt}
\setlist[enumerate]{leftmargin=2.5em,itemsep=0.15em,topsep=0.35em,parsep=0pt}
\setcounter{secnumdepth}{3}
\ctexset{
  section={format=\Large\bfseries,name={},aftername=\quad,beforeskip=2.1ex plus .4ex minus .2ex,afterskip=1.25ex plus .25ex minus .15ex},
  subsection={format=\large\bfseries,name={},aftername=\quad,beforeskip=1.7ex plus .3ex minus .15ex,afterskip=.8ex plus .2ex minus .1ex},
  subsubsection={format=\normalsize\bfseries,name={},aftername=\quad,beforeskip=1.25ex plus .25ex minus .1ex,afterskip=.55ex plus .15ex minus .1ex}
}
\renewcommand{\tablename}{表}
\renewcommand{\figurename}{图}
\renewcommand{\refname}{参考文献}
\renewcommand{\arraystretch}{0.82}
\setlength{\tabcolsep}{3.5pt}
\setlength{\LTleft}{0pt}
\setlength{\LTright}{0pt}
\setlength{\LTpre}{0.5em}
\setlength{\LTpost}{0.5em}
\pagestyle{plain}
\begin{document}
\begin{center}
  \vspace*{0.18in}
  {\fontsize{22}{30}\selectfont\bfseries Arm 与 RISC-V 机密计算安全机制调研与总结\par}
\end{center}
\vspace{0.28in}

\section{引言}

\subsection{研究背景：从执行域隔离到平台级机密计算}

机密计算的研究对象已经不再只是“在处理器上划出一个安全世界”。早期 TrustZone 或 enclave 设计主要关注受保护代码如何避免被普通操作系统直接读取：可信应用运行在分离的执行域中，普通世界即便被完全攻陷，也无法直接读取安全世界的内存。这一模型在移动设备与嵌入式场景中被验证是有效的，但它默认了一个相对简单的信任结构——设备归单一所有者所有，安全服务由厂商预置，边界两侧的角色长期不变。随着虚拟化、异构加速和可组合互连进入部署场景，这一信任结构被打破：受保护工作负载需要在不信任的 host、hypervisor、设备固件、DMA 引擎和网络端点之间维持可识别、可转换且可证明的所有权。云租户不信任云厂商的特权软件，机密 AI 训练不信任宿主机对 GPU 队列的配置，CXL 内存池和 RDMA 远端分页则让数据在物理上离开本机内存控制器。Arm CCA/RME/RMM 与 RISC-V CoVE/AP-TEE 等架构因此把问题推进到平台层面：谁能够创建执行域，谁可以转换内存状态，谁被允许把设备接入该执行域，以及远程验证方如何获知这些状态。\cite{li2022cca,sahita2023cove}

这条边界本质上是一条跨层状态链。处理器执行态、物理页属性、二阶段地址转换、设备访问权限、链路密钥与 attestation 报告若使用不同的对象标识或生命周期规则，就可能在看似“受保护”的系统中留下重新映射、陈旧授权、异常注入或错误共享的空间。TrustZone 漏洞谱系的研究已经表明，即便硬件隔离原语本身成立，监控器、可信应用、驱动与 API 组合而成的系统仍然会产生成体系的真实漏洞；ret2ns 等工作进一步说明，连“从安全态返回非安全态”这样一个看似平凡的切换路径，都可能因寄存器清理不完整而击穿整个隔离假设。因此，本报告不将某项硬件特性孤立地描述为安全结论，而是把它放入创建、运行、共享、迁移、回收和验证的完整过程之中，追问其机制是否真正闭合了相应的攻击路径。

\subsection{调研问题：保护边界为何难以闭合}

本调研围绕三个彼此依赖的问题展开。第一，CPU 执行域、特权监控器和内存所有权如何共同建立最小可信基：Arm 侧以 Realm、Granule Protection Table 和 RMM 回答，RISC-V 侧以 TVM、MTT 和 TSM 回答，二者都试图在保留 host 资源管理能力的同时剥夺其对受保护状态的读写能力。第二，当数据经由 IOMMU、PCIe/CXL、RDMA 或加速器队列离开 CPU 后，设备身份、DMA 映射、中断和链路保护如何延续同一边界：SPDM、TDISP、IDE、CoVE-IO 与 sIOPMP 分别从协议、链路和硬件检查三个方向补这一段链。第三，证明、控制流完整性、内存安全和微架构防护等纵深机制，如何约束边界内部的错误或被滥用的合法权限：attestation 让远端验证方获得判断依据，内存加密与完整性树限制物理层观察与重放，特权架构与 I/O 保护则限定固件和设备的合法行为范围。它们分别对应 Arm/RISC-V 机密计算机制、I/O 与设备数据路径防御，以及 ISA 和硬件级纵深防御三条主线。

这些问题不能只靠术语对照回答。架构规范能够精确定义状态和权限，却通常不提供真实工作负载的代价；原始论文能够说明原型实现和实验条件，却未必覆盖所有部署假设；厂商材料和标准草案则能够反映工程进展，但其语义、版本或可访问性可能仍受限制。例如 RISC-V AP-TEE 的 ABI 尚未 ratified，Arm CCA 的生产级芯片在若干论文写作时还不可得，部分厂商规范只能核验公开页面而无法核验全文细节。报告因此把“论文实际证明了什么”“哪些推论依赖额外假设”作为每一篇的核心阅读线索，并把未说明或无法核验的部分显式保留，而不是用推测补齐。

\subsection{选文范围与文章结构}

本文以 \texttt{out/survey.pdf} 为文献池，但只展开配套 PPT 在 15 个研究方向中实际选出的 primary 论文，不将所有参考文献机械逐条改写。PPT 给出 45 个选入槽位，去重后为 41 个独立来源；同一来源若服务于多个方向，正文只保留一次完整论述，并在基本信息中说明其出现的研究方向。这样的范围既与汇报材料保持一致，也避免将背景性引用误写成已完成的论文评审。

后文按七个类别组织：Arm CCA、RME、RMM 与 Realm 机制；RISC-V enclave、CoVE 与 AP-TEE 机制；TEE 设计空间与执行环境谱系；启动、远程证明与生命周期安全；机密 I/O、网络与数据路径防御；加速器、SmartNIC、DPU 与设备 TEE；ISA 与硬件设计纵深防御。每个类别内部先说明问题边界与技术脉络，再按论文依次展开。每篇均遵循“基本信息—研究背景与动机—关键挑战与核心思想—系统架构—核心方法设计—安全性与威胁边界—实验环境与证据—实验结果与证据强度—综合评价”的叙述顺序：先交代为什么需要该工作，随后解释机制如何运作，再审视实验、规范或工程证据能支持到何种程度，最后讨论其相对于已有工作的贡献、限制和可继续验证的方向。每个类别末尾新增两节综合性讨论：“开发商安全实践建议”把该类论文与标准的结论转化为芯片厂商、固件与系统软件开发商、设备制造商可执行的安全工程要求；“未来研究方向”则把行文中暴露的证据缺口和机制空白整理为具体、可验证的研究问题。这样安排旨在让论文之间形成可比较的论证链，并让调研结论能够同时服务学术研究选型与产品开发决策，而不是由元数据、要点和讲解稿拼接而成的材料汇编。\cite{li2024sokteechoices,boubakri2025riscvtee,arm_cca_spec,riscv_privileged}

\section{证据边界与阅读方法}

正文优先采用仓库中已核验的本地 PDF、PPT 选题记录和 README 评审记录。为避免把不同成熟度的材料混为一谈，本文把官方规范、架构手册、RFC 和标准记为 E0；同行评审原始论文记为 E1；Survey/SoK 记为 E2；公开 draft 或未 ratified 规范记为 E3；厂商或行业材料记为 E4；仅有元数据、访问受限或缺少 PDF 的来源记为 E5。E0/E1 支撑具体机制语义和论文报告的实验，E2 用于术语和分类；E3 至 E5 只用于其可核验的标准化或工程状态，不能据此扩展出未经证实的安全和性能结论。

这一分级直接约束每篇论文的写法。对 E1 论文，正文展开其威胁模型、机制设计、实现细节与实验数字，并区分“摘要声称”与“评估章节实际测量”两个层次；对 E0/E3 规范，正文只使用本地可读或官方页面可核验的语义，明确标注哪些接口细节未经本地验证；对 E2 综述，正文只用其分类框架和对比维度，不把综述转述的结论当作原始实验结果；对 E4/E5 材料，正文收窄到其标准化或产品状态。实验数字一律标注其测量平台（模拟器、FVP、开发板或生产芯片）与适用边界，避免把原型开销外推为产品开销。

因此，正文中的“论文未说明”“证据不足”并非篇幅占位，而是对可得证据范围的明确标记。所有条目均保留 BibTeX key 和正文引用，便于回溯 \path{survey/reference.bib} 与 \path{reference/} 证据库；对没有本地 PDF 的条目，则仅依据已核验的 PPT 事实说明其研究位置，并把机制细节和复现实验收窄到可支持的范围。
\section{Arm CCA、RME、RMM 与 Realm 机制}
Arm 平台的隔离机制正在从传统 TrustZone 的安全/非安全二分，转向面向虚拟机和云端多方参与者的 Realm 执行域。这一转变的实质，是把“谁拥有这块物理内存、谁有权改变它的归属”从 hypervisor 的软件约定，提升为由 Granule Protection Table（GPT）、Granule Protection Check（GPC）与 Realm Management Monitor（RMM）共同维护的架构级状态机：host 仍然调度 CPU、分配内存、管理设备，但它对每一个 4KB granule 的访问都必须经过硬件路径上的 GPC 判定。本节将先说明 RME、RMM 与物理内存所有权为何成为新的控制核心，再逐篇考察细粒度隔离（Shelter、RContainer、NanoZone）、运行时管理（RMM 规范）和 I/O 与攻击分析（ACAI、Devlore、CAGE）如何限定 host、monitor 与受保护工作负载之间的责任边界。\par
这一组工作共同提醒我们：处理器提供新的执行态并不等于隔离自然成立。Li 2022 的设计与形式化验证证明的是 granule 状态转换与访问控制不变量在架构层面闭合，但页面转换、异常路径、调试接口与软件运行时一旦脱离同一所有权语义，就会在委托（delegate）、回收（undelegate）、跨核 GPT 缓存、设备 DMA 与中断注入等环节留下缝隙。后续论文正是沿着这些缝隙逐一展开：Shelter 与 RContainer 追问隔离粒度能否细到用户态应用与容器，NanoZone 追问进程内部的数据域能否低开销切换，ACAI、Devlore 与 CAGE 追问数据离开 CPU 之后设备侧、中断侧与任务元数据侧的边界如何延续。因此，论文中的机制、验证和负面结果都需要放在“创建—委托—运行—共享—回收—验证”的完整状态迁移过程中理解，而不是孤立地读作某个单点防御。\par
为便于后文对照，表 1 先给出本章九篇材料的谱系定位；表中数字均来自各论文本地 PDF 或 README 核验记录，证据等级沿用本报告第 2 章的 E0--E5 划分。\par
\begin{table}[htbp]
\centering\small
\caption{本章材料谱系总览（数字以后文各节标注的出处为准）}
\begin{tabular}{llll}
\toprule
材料 & 角色/粒度 & 平台与证据 & 关键数字 \\
\midrule
Li 2022\cite{li2022cca} & 架构奠基（Realm/GPT/RMM） & Fast Model + 形式化验证，E1 & 43 层抽象；约 200 行可信 Coq \\
CCA spec\cite{arm_cca_spec} & 官方术语与边界（DEN0125） & 官方页面核验、无本地 PDF，E0 & 不提供性能数字 \\
RMM spec\cite{arm_rmm_spec} & 生命周期状态机（DEN0137） & 官方页面核验、无本地 PDF，E0 & 不提供性能数字 \\
Shelter\cite{zhang2023shelter} & 用户态 SApp（multi-GPT） & FVP + 硬件 SoC，E1 & 真实负载 $<$15\% \\
RContainer\cite{zhou2025rcontainer} & 容器（mini-OS + Shim-GPT） & FVP + ARMv8 SoC，E1 & 应用 4\%--7\%；kill 约 8.5\% \\
NanoZone\cite{liu2025nanozone} & 进程内数据域（POE/PIE/PAS） & 模拟器 + 开发板，E3 preprint & 切换 20 cycles；均值 4.87\% \\
ACAI\cite{acai2023} & 加速器数据路径（device-GPC） & 模拟器 + Armv8 板，E3 preprint & GPU 43.5\%；FPGA 12.1\% \\
Devlore\cite{bertschi2026devlore} & 中断控制面（delegate-but-check） & FVP + Rock5b，E3 preprint & 压力 $\le$1\%；glmark2 0.06\% \\
CAGE\cite{wang2026cage} & 加速器工作流（shadow task） & FVP + Juno R2 + VCU118，E1 & GPU 0.58\%--5.31\%；FPGA 9.61\%--16.30\% \\
\bottomrule
\end{tabular}
\end{table}
从表中可以读出一条贯穿本章的线索：隔离粒度越细、越贴近运行态高频路径，机制设计就越依赖“让强边界被调用得足够少”——RContainer 用 mixed-pagetable 避免页表重建，NanoZone 用域分配把 96.72\% 的切换压进 POE 层，Devlore 把检查集中在配置与记录路径，CAGE 用 GPT 同步优化消除 84.63\%--96.55\% 的同步开销。这一共性比任何单篇的绝对数字更能指导工程实践。\par

\subsection{论文 1：Design and Verification of the Arm Confidential Compute Architecture}
\subsubsection{论文基本信息}
《Design and Verification of the Arm Confidential Compute Architecture》由 Xupeng Li、Xuheng Li、Christoffer Dall、Ronghui Gu、Jason Nieh、Yousuf Sait、Gareth Stockwell 合作完成（Columbia University 与 Arm Ltd），发表于 OSDI 2022。本文将其置于“Arm CCA / RME / RMM 基础架构”的研究脉络中考察，并把它作为该方向的主线原始论文（E1）使用。它所回答的核心问题可概括为：如何把“host 不可信”的云威胁模型落到 Arm 平台的具体机制上——Realm 执行域、GPT granule 所有权与 RMM 软件 TCB，并且这一设计的关键安全性质可以被形式化证明而非仅靠实现约定。当前证据状态为：本地 PDF（22 页）已保存并核验；下文仅在该范围内讨论其机制与结论。\cite{li2022cca}\par
\subsubsection{研究背景与动机}
要理解该方案，首先需要看到它面对的系统矛盾：在传统虚拟化中，hypervisor 管理页表、设备、调度与 VM 生命周期，拥有读取和修改任何租户 VM 内存与寄存器状态的全部能力；而云租户希望敏感 VM 对云管理员和 host 软件完全不透明。论文第 1 节明确指出，hypervisor 与 OS 的巨大代码库必然包含可被利用的漏洞，因此“信任系统软件保护应用数据”这一假设本身已经成为风险源。\par
具体而言，既有路径各有不可接受之处。TrustZone 的双世界划分不是为大规模多租户 CVM 管理设计的：secure world 面向少量可信应用与 trusted OS，无法自然扩展到云中成百上千个机密 VM 的创建、调度与回收。嵌套页表（NPT）可以隔离 VM 之间的物理内存，但 NPT 完全由不可信 hypervisor 控制，对 hypervisor 自身毫无防御。若为每个物理帧单独引入所有权跟踪结构，又会显著影响 TLB 设计与性能，或需要引入复杂的 CISC 式指令与微码。CCA 的动机正是在保留云平台可管理性（host 继续动态分配内存、调度 CPU）的同时，把内存内容与执行状态从 host 的可达范围中剥离。\par
\subsubsection{关键挑战与核心思想}
在上述背景下，作者提出的关键判断是：CCA 不是让 hypervisor 消失，而是让 hypervisor 只能管理“不可读”的 Realm 资源；并且由于硬件只提供最小机制、安全保证依赖固件，固件的正确性必须被形式化验证。\par
具体而言，可以把论文命名的挑战归纳为三条。第一，向后兼容挑战：Arm 架构数十年的设计假设是“更高特权级拥有更大访问权”，而 Realm 要求即使更高特权级的 host 也无法读取 Realm 状态；逐条重定义所有指令在 Realm 上下文中的语义不可行，因此 CCA 选择引入与特权级正交的新物理地址空间（PAS）。第二，内存所有权挑战：ownership 必须按 granule 动态转换且由硬件强制，而不能依赖 hypervisor 维护的数据结构。第三，可验证性挑战：RMM 采用 hand-over-hand locking、动态分配的共享多级页表、运行在 Arm 弱内存模型上、C 与汇编互相调用、且资源分配完全由不可信软件控制——这四个性质叠加使此前的验证方法均不适用，论文第 1 节将其列为验证 RMM 的四个具体困难。\par
\subsubsection{系统架构}
从系统结构看，CCA 在既有 Non-Secure（NS）与 Secure 两个 world 之外增加 Realm world 与 Root world：Realm world 承载机密 VM，Root world 在最高特权级 EL3 运行 EL3 Monitor（EL3M），负责在三个 world 之间做 CPU 上下文切换并管理 GPT（论文 Figure 1）。\par
具体而言，每个 world 拥有自己的 PAS，物理内存按 4KB granule 归属某一个 PAS，且不存在 NS 与 Realm 之间的静态分区——granule 可以在运行中从 NS PAS 委托（delegate）到 Realm PAS。硬件对每次内存访问依据 GPT 执行 PAS 检查，访问策略由论文 Table 1 给出：NS world 只能访问 NS granule；Realm 与 Secure world 可以访问各自及 NS 的 granule，但互相不可访问；Root world 可以访问全部 PAS。软件上，RMM 运行在 Realm world 的 EL2，向 NS world 的 hypervisor 暴露 Realm Management Interface（RMI，论文 Table 2），向 Realm 内部暴露 Realm Service Interface（RSI）；任何 world 的软件都可以通过 SMC 请求 EL3M 改变 granule 的 PAS 归属，但只有 EL3M 能实际更新 GPT。\par
\subsubsection{核心方法设计}
论文的关键不是单一组件，而是“最小硬件机制 + 可验证固件 + 严格接口分工”组合成的一条受控执行路径。\par
\noindent\textbf{（1）Realm 与 Security State。}\par
这一机制的直接作用是：把 tenant VM 从 Normal world hypervisor 中隔离出来，同时保持指令语义与特权级结构不变。\par
具体实现上，Realm world 与 NS world 完全兼容，现有软件栈可以原样运行在 Realm world 内，但 NS 软件无法访问 Realm 的 CPU 状态与内存。Realm 与 Secure world 互相不信任，因此引入第四个、更高特权的 Root world 管理 world 间切换。论文明确将机密性定义为“Realm 对私有数据的任何修改不能被其他 Realm 或不可信系统软件观察”，完整性定义为“Realm 不会观察到非自身造成的私有数据变化”，并明确不保证可用性——host 拒绝服务是威胁模型允许的结果。\par
\noindent\textbf{（2）Granule Protection Table 与 GPC。}\par
这一机制的直接作用是：把物理内存的 ownership 从软件约定变成硬件强制的状态机。\par
具体实现上，GPT 记录每个 granule 的 PAS 归属，GPC 在每次访问时对照 Table 1 的访问控制策略做最终判定；即使 host 在自己的页表中建立了映射，也无法绕过 GPC。论文特别指出，CCA 硬件要求所有 DMA 访问同样经过 GPT 检查，从而在架构层面把基于 DMA 重映射的攻击纳入防御范围。只有 EL3M 可以改变 granule 的 PAS；其他 world 的 delegate/undelegate 请求都经由 SMC 进入 Root world 审查。\par
\noindent\textbf{（3）RMM / RMI / RSI 分工。}\par
这一机制的直接作用是：让 host 保留资源管理能力，但所有影响 Realm 安全的操作都必须经过一个小而可审计的状态机。\par
具体实现上，RMI 面向 host/hypervisor，RSI 面向 Realm payload。论文 Table 2 完整列出 RMI 命令集，可按对象分为六组：
\begin{itemize}
\item 版本协商：Version Query，查询 RMI ABI 版本；
\item granule 委托：Granule.Delegate（NS 转为 Delegated）与 Granule.Undelegate（转回 NS）；
\item Realm 生命周期：Realm.Create（创建 Realm Descriptor，RD）、Realm.Destroy、Realm.Activate（Realm 从 New 转为 Active）；
\item 执行上下文：REC.Create/Destroy/Run，REC 对应 VCPU，REC.Run 即进入 VCPU 运行；
\item 数据 granule：Data.CreateUnknown（内容未知）、Data.Create（从 NS 拷贝内容）、Data.Destroy（清零后转回 Delegated）；
\item Realm Translation Table：RTT.Create/Destroy/MapProtected/UnmapProtected/MapUnprotected/UnmapUnprotected/ReadEntry，管理 Realm 的 stage 2 页表。
\end{itemize}
每条 RMI 命令都是一次 SMC：hypervisor 调用后先 trap 到 EL3M，EL3M 把执行切换到 Realm world 的 RMM 处理，完成后 RMM 再经 SMC 返回。RMM 只做安全仲裁，虚拟化功能（资源分配、调度、设备模拟）仍交还给既有 hypervisor，因此 RMM 可以比裸机 hypervisor 小几个数量级——这是其可验证性的工程前提。\par
\noindent\textbf{（4）委托语义与 GST/GPT 的双表分工。}\par
这一机制的直接作用是：把 granule 生命周期拆成“硬件强制”与“软件簿记”两层，各自保持简单。\par
具体实现上，委托（delegation）是 CCA 的核心概念：hypervisor 把内存 delegate 到 Realm world、undelegate 回 NS world；RMM 自己不维护内存池，Realm 用的所有内存都必须先由 hypervisor 委托；已委托但未被 RMM 使用的 granule 必须只含零，降低复用时的意外信息流。RMM 维护 Granule Status Table（GST）跟踪每个 granule 的委托状态与当前用途——例如 delegate 请求到达时，RMM 先查 GST 确认该 granule 未被委托过，再经 SMC 请 EL3M 检查其当前处于 NS PAS 并更新 GPT，最后更新 GST；重复委托或回收使用中的 granule 一律返回错误码。关键设计是：GST 不被硬件检查，只是软件簿记；GPT 保持简单、只含硬件检查所需信息。两表分工使硬件路径与软件状态机的复杂度各自可控。回收路径同样严格：RMM 在 undelegate 前把 granule 清零，Realm 此后访问该 IPA 会触发 stage 2 abort，且未经 Realm 许可，hypervisor 不能把 granule 映射回 PAR 内曾经被支撑的 IPA。\par
\noindent\textbf{（5）REC、RTT、PAR 与初始状态测量。}\par
这一机制的直接作用是：定义 Realm 运行期的地址空间语义与 verifier 可验证的初始状态。\par
具体实现上，REC 对应普通 VM 的 VCPU，RTT 对应嵌套页表——RTT 使用与 NS stage 2 相同的格式与拓扑，但额外提供一个位，允许 Realm 在 RMM 控制下访问 NS granule（用于与 hypervisor 的虚拟 I/O）。每个 Realm 在其 IPA 空间内拥有一个 Protected Address Range（PAR）：RMM 保证 PAR 内只能映射 Realm PAS granule，PAR 之外 hypervisor 可以自由映射 NS granule 或模拟访问——这使 Realm 内 OS 能可靠区分“私有内存”与“可与不可信方共享的缓冲区”（如 DMA buffer）。测量方面，Realm 创建期间 hypervisor 把 granule 指定到特定 IPA 并从 NS granule 拷入数据，IPA 与数据被密码学哈希进 Realm 的 attestation token；Realm 激活后测量固定，之后只能向未使用的 IPA 添加未知内容的内存。工程细节上，物理连续的委托内存可以按大于 4KB 的块映射以优化 TLB；RMM 为支持大 Realm 采用细粒度锁——每个 granule 一把锁、RTT 每级一把锁并以 hand-over-hand locking 支持多 CPU 并发页表操作（Figure 2）。此外论文记录了明确的防御性编程措施：RMM 与 EL3M 不像 Linux 那样映射全部物理内存，RMM 静态映射仅含自身代码与元数据（GST 与锁）并按需映射 Data/元数据 granule，EL3M 只映射自身代码、小栈与 GPT；SMC 参数在 EL3M 中永远按值解释、绝不作为指针——即便未来固件版本引入未验证的 bug，这些措施也使 ROP/JOP 类攻击难以得手。\par
\noindent\textbf{（6）VIA 验证基础设施与理想/现实范式。}\par
这一机制的直接作用是：把“固件可信”从假设变成定理。\par
具体实现上，论文提出 VIA（Verification Infrastructure for Armv9-A），包含四项技术：mover oracle queries，把其他 CPU 的交错操作封装为可按 mover 类型重排的查询，使 hand-over-hand locking 与动态共享多级页表首次可被模块化验证；把并发代码分解为 data-race-free 与非 DRF 部分，对后者定义 permutation conditions，使顺序一致性模型上的证明可以迁移到弱内存硬件；register accounting，利用 Arm 过程调用标准跟踪 C 与汇编相互调用时的寄存器信息流；ideal/real paradigm，定义一个带 declassification 数据通道的理想化安全机器，证明实现经 43 层抽象精化到规范、规范模拟理想机器。整个证明只需信任约 200 行 Coq 规范，是首个针对机密计算架构安全保证的完整证明。\par
\noindent\textbf{（7）原型实现与 KVM 集成。}\par
这一机制的直接作用是：证明 CCA 设计可以落到真实软件栈上，而不是停留在架构图纸。\par
具体实现上，作者修改 Linux KVM hypervisor 使其在 CCA 上运行并管理 Realm VM：host 侧 KVM 继续负责资源分配与调度，涉及 Realm 安全的操作转为 RMI 请求。Realm 内部的 guest OS 无需理解 CCA——Realm world 与 NS world 的兼容性保证了未经修改的软件栈可以直接迁入。RMM 只做安全仲裁这一设计使其代码量可以比裸机 hypervisor 小几个数量级，这既是性能考虑（减少 world 切换路径上的软件），更是可验证性的前提。\par
\subsubsection{安全性与威胁边界}
论文的威胁模型（第 2 节）假设攻击者无物理访问能力，目标是破坏 VM 数据的机密性与完整性。范围内外可明确列举为：
\begin{itemize}
\item 范围内：控制 hypervisor 或其他软件读取/修改 VM 私有内存与寄存器状态；控制 DMA 设备发起访问；内存重映射与别名攻击；来自其他被攻陷 VM 的机密性/完整性攻击；
\item 范围外：host 发起的可用性攻击（拒绝调度、回收内存属合法行为）；侧信道与已知软件错误注入（需配合架构级缓解单独使用）；冷启动、在线探测、重放等 DRAM 攻击（需要额外硬件）。
\end{itemize}
因此，文中机制只在“granule 所有权、状态转换与受控接口同时成立”时构成完整隔离论证；侧信道与物理攻击需要正交机制补齐。\par
值得单独展开的是架构兼容性与安全的张力。论文第 3 节举例指出，Arm 架构中的调试寄存器本来就是为让 hypervisor 窥视 VM 状态而设计的，这与 Realm 的机密性直接冲突；CCA 的处理方式不是逐条重定义指令，而是在 Realm 上下文中约束这类特性的语义。这个例子说明：Realm 不是在旧架构上“加”出来的世界，而是对“更高特权级必然拥有更大访问权”这一数十年假设的首次系统性反转，其边界条件需要逐特性审查。\par
\subsubsection{实验环境与证据}
论文的证据由两部分构成。其一是功能与初步性能证据：CCA 硬件当时不可用，作者在支持 CCA 的 Arm Fast Model 上演示了修改后的 Linux KVM 管理 Realm VM 并运行多种 VM 工作负载，并把固件移植到当时的 Arm 硬件上获得初步数据，结论是 KVM on CCA 相对 vanilla KVM 在真实应用负载上只有 modest overhead。其二是形式化证据：RMM（含 C 与汇编）实现经 43 层抽象证明精化到规范，规范再被证明模拟理想化安全机器，可信基础约 200 行 Coq。本地可核验材料为 22 页 PDF、Figure 1/2、Table 1（访问控制策略）、Table 2（RMI）及验证章节。\par
\subsubsection{实验结果与证据强度}
对于性能或效果，论文能够支持的结论是：架构定义与设计级验证是强证据，全栈性能是弱证据。“modest overhead”来自早期固件移植的初步测量，不能写成“CCA overhead 已被本文全面证明”；I/O、device assignment、interrupt 与后续 RMM 规范细节均不在本文覆盖范围内，具体 container、accelerator、interrupt 与 memory-protection 开销应分别引用 Shelter、RContainer、NanoZone、ACAI、Devlore、CAGE 等后续论文。\par
\subsubsection{失败模式与边界分析}
论文自身与后续工作共同标定了三条边界。第一，验证覆盖边界：形式化证明的对象是原型固件的实现与规范，证明的终结点是“实现精化到规范、规范模拟理想机器”；它不排除规范之外的行为（如未建模的硬件 errata），也不为任何其他 RMM 实现提供保证。第二，威胁模型边界：可用性被显式放弃——host 可以合法地拒绝调度或回收内存，Realm 必须容忍；这意味着依赖 CCA 的可用性敏感业务需要额外的冗余部署，而非架构内解决。第三，证据外推边界：论文的性能数据来自 CCA 硬件可用之前的早期移植，作者在文中明确这只是初步数据；本章后续六篇系统论文同样受困于生产芯片缺失，这说明整个方向在 2022--2026 年间的性能证据都带有“非生产平台”标签，引用时必须连同平台一起引用。\par
\subsubsection{综合评价}
综合来看，Li 2022 是 CCA 方向的奠基材料：它把 Realm、GPT、RMM、attestation 放进统一设计，并首次给出机密计算架构安全保证的完整机器证明，证明了“granule 状态机 + 小 TCB 固件”这一路线在设计上可以闭合。其局限同样明确：不覆盖后续规范细节、设备路径与商用部署；验证对象是固件实现而非未来所有实现，形式化证据不能外推为“任何 RMM 实现都无漏洞”。从部署角度看，它是 Arm server confidential computing 的基础，但真实落地依赖 RMM 实现、firmware、OS 与 I/O 生态的共同成熟。\par

\subsection{论文 2：Arm Confidential Compute Architecture Specification}
\subsubsection{论文基本信息}
《Arm Confidential Compute Architecture Specification》（DEN0125）由 Arm Limited 发布，官方文档页面记录于 2025 年。本文将其置于“Arm CCA / RME / RMM 基础架构”的研究脉络中考察，并把它作为 spec 类（E0）材料使用。它所回答的核心问题可概括为：把论文中的设计语言固定成官方架构语言，使 Realm、RME、GPT、Root/RMM 等术语具有规范语义。当前证据状态为：官方来源页面已核验，但本地 PDF 不可得——README 记录 2026-05-12 复核时 \texttt{den0125/latest} 与 \texttt{den0125/0400} 均返回 access-denied，未找到稳定公开 PDF 端点；下文仅在公开概念范围内讨论，凡涉及规范内部细节处均显式标注证据边界。\cite{arm_cca_spec}\par
\subsubsection{研究背景与动机}
要理解该材料的角色，首先需要看到它面对的系统矛盾：CCA 的讨论极易混用论文、实现与规范三个层次——论文解释“为什么这样设计”，实现给出“某个原型怎么做的”，而只有规范定义“哪些状态、接口与约束属于架构语义”。如果三者混淆，调研报告就会把某个原型的实现选择误写为架构保证，或把论文推测误写为标准事实。\par
具体而言，Realm、RME、GPT、Root/RMM 等术语在本报告中必须按官方语义使用；而规范 PDF 不在本地这一事实，直接限定了本节的写法：可以引用官方 source page 的存在性与公开架构概念，不能引用未本地读取的具体章节、表格或状态编码。\par
\subsubsection{关键挑战与核心思想}
在上述背景下，该材料承担的核心作用是：spec 是规则书而非实验论文，它告诉我们哪些 CCA 论断可以被官方架构支撑，哪些必须回到实现论文寻找证据。\par
具体而言，规范强于论文之处在于权威性——它把 Realm/REC/granule 的对象定义、访问控制的架构性质与软件/硬件分工固定下来；弱于系统论文之处在于不提供任何 workload、实现开销或部署评估。在本调研中，它承担“术语校准”与“边界校准”两个角色：凡本章其他论文的机制描述与公开架构概念冲突时，以规范语义为准；凡涉及性能、兼容性或部署的结论时，规范不能作为证据来源。\par
\subsubsection{系统架构}
从系统结构看，按公开 CCA 概念可以描述的架构总览与 Li 2022 Figure 1 一致：Realm payload 承载机密工作负载，RMM/Root 管理 Realm 生命周期与 granule 状态，Normal host 继续管理资源但不能读取 Realm 内容，GPT/GPC 在硬件路径上执行最终访问控制，attestation 向远程验证方提供初始状态证据。\par
具体而言，这一总览中每个组件的规范语义都来自官方架构公开概念：Realm 不是 TrustZone secure world 的简单改名，而是与 NS world 指令级兼容、但与特权级正交的独立执行态；REC 是 Realm execution context，承载 VCPU 式执行状态而非普通线程；granule 是 ownership 与状态转换的最小单位，委托语义要求未使用的已委托 granule 只含零。组件间交互的细节（RMI/RSI 的具体命令集、状态编码、版本间差异）在本仓库无本地证据：Li 2022 Table 2 给出的是其原型实现的 RMI 命令集，与 DEN0125 当前版本的接口是否逐条对应，证据不足，本报告不作断言。规范文档编号 DEN0125 本身及其官方页面状态是可核验事实（README 记录），这也是本节唯一可以直接引用的“规范级”信息。\par
\subsubsection{核心方法设计}
作为规范材料，其“方法”体现为它固定下来的三组语义。\par
\noindent\textbf{（1）Realm / REC / Granule 术语固定。}\par
这一机制的直接作用是：统一所有后续工作的对象模型。\par
具体实现上，Realm 是保护域，与 TrustZone secure world 是并列而不同的执行态；REC 是 Realm execution context，承载 vCPU 式的执行状态；granule 是内存 ownership 与状态转换的基本单位。本报告对这三个词的使用均以该规范语义为准。\par
\noindent\textbf{（2）访问控制语义。}\par
这一机制的直接作用是：明确访问控制是架构状态与硬件检查，而非软件约定。\par
具体实现上，granule 的 security state 决定哪个 security state 的执行体可以访问它；host 的管理请求必须经过 RMM/Monitor 的合法状态转换。这解释了为什么本章所有部署型论文（Shelter 的 multi-GPT、RContainer 的 Shim-GPT、NanoZone 的 PAS 层、CAGE 的 accelerator GPT）都围绕 granule、页表与 RMM 接口做设计——它们是在架构状态机上选择不同的策略点。\par
\noindent\textbf{（3）状态机与接口语义的标准化对象。}\par
这一机制的直接作用是：界定规范类材料在调研中的可用范围。\par
具体实现上，按照公开架构概念，CCA 类规范标准化的对象可分为四类：
\begin{itemize}
\item 术语与对象模型：Realm、REC、granule、PAS 等概念的唯一语义；
\item 执行态与安全状态的关系：各 world 对 granule 的访问矩阵；
\item granule 生命周期的合法转换：delegate/undelegate 及其前置条件；
\item 软硬件分工：哪些检查由 GPC 硬件完成、哪些状态由 Monitor/RMM 软件维护。
\end{itemize}
本章其余论文的机制都可以映射到这四类对象上：Shelter 的 multi-GPT 是对访问矩阵的 per-core 实例化，RContainer 的 Shim-GPT 是对 granule 生命周期的容器级复用，ACAI/CAGE 的设备侧检查是对 GPC 适用范围的扩展。但需要强调：以上映射是本报告基于公开概念与本地论文的归纳，规范文本中关于这些扩展是否被明确支持，本地证据不足，不作断言。\par
\noindent\textbf{（4）证据边界管理。}\par
这一机制的直接作用是：防止规范材料被过度引用。\par
具体实现上，本节可以写“官方 CCA source page 已核验（README 记录）”，不能写“本文本地核验了 DEN0125 某条具体表格或某版状态机”；具体接口细节一律回引 Li 2022（本地 PDF）、RMM 规范 README 记录的公开生命周期概念，或 Shelter/RContainer 等本地可读论文。\par
\subsubsection{安全性与威胁边界}
规范本身不定义威胁模型，它定义的是“在架构语义内哪些访问不可能发生”。其边界在于：架构语义不等于实现安全——任何具体 silicon、firmware、RMM 实现的漏洞都不在规范覆盖范围内；也不等于性能可行——规范不承诺任何 granule transition、world switch 或 attestation 的开销量级。这些边界在本章其余论文中分别以验证证据（Li 2022）与原型测量（其余各篇）的形式补充。\par
\subsubsection{实验环境与证据}
这是强规范证据，但不是可渲染的论文证据。可核验材料为：Arm DEN0125 官方文档页面（README 记录已核验）；README 记录的下载状态（PDF 不可得、access-denied、2026-05-12 复核）；公开架构概念。不能支撑的材料为：规范内部的任何具体表格、命令编码、版本差异细节。\par
\subsubsection{实验结果与证据强度}
官方架构规范不提供性能数字。它可以说明哪些机制类别可能影响性能——granule transition、world switch、attestation、RMM 调用——但任何具体 overhead 数字都必须来自实现论文或产品文档。本章的做法正是如此：所有性能结论均标注到具体论文的具体图表，不以规范为性能证据。\par
\subsubsection{综合评价}
综合来看，CCA spec 是本方向的权威边界材料：优势是官方来源、术语与架构边界的最终裁决者；局限是本地 PDF 不可用，不能作为任何细节论断的唯一证据，且按定义不含实验。从部署角度看，它代表 Arm confidential computing 的标准化方向，但标准化文本与真实落地之间隔着 silicon、firmware、RMM 与云软件栈四层实现，每一层都需要独立证据。\par

\subsection{论文 3：Realm Management Monitor Specification}
\subsubsection{论文基本信息}
《Realm Management Monitor Specification》（DEN0137）由 Arm Limited 发布，官方文档页面与公开 CDN 链接记录于 2025 年。本文将其置于“Arm CCA / RME / RMM 基础架构”的研究脉络中考察，并把它作为 spec 类（E0）材料使用。它所回答的核心问题可概括为：把 CCA 中最关键的软件 TCB 讲清楚——谁创建 Realm、谁运行 REC、谁改变 granule state，以及 host 与 Realm 各自能看到什么接口。当前证据状态为：官方文档页面与 DEN0137 CDN URL 已核验，但本地 PDF 不可得（README 记录 2026-05-12 复核时直接 CDN 下载返回 access-denied）；下文仅使用公开的 RMM/RMI/RSI 生命周期概念。\cite{arm_rmm_spec}\par
\subsubsection{研究背景与动机}
要理解该规范的必要性，首先需要看到它面对的系统矛盾：CCA 的硬件隔离必须有一个小而受控的管理者，否则 host 仍可通过生命周期操作（创建、销毁、调度、内存回收）绕过安全目标——硬件只管“当前这个 granule 谁能访问”，不管“把一个 granule 从 NS 变成 Realm 再变回来”这一过程的合法性。\par
具体而言，host 需要请求创建与运行 Realm，但不应直接操控 Realm 私有状态；Realm payload 需要 attestation 等服务，但也不能信任 host 转述的任何信息。RMI/RSI 把这两类请求分开，使 RMM 的 TCB 在接口层面就是可审计的：审计 RMI 即审计 host 能对 Realm 做的全部操作，审计 RSI 即审计 Realm 信任链的全部入口。\par
\subsubsection{关键挑战与核心思想}
在上述背景下，该规范的核心判断是：RMM 是 CCA 安全边界的状态机；安全性来自合法状态转换序列，而不是任何单个函数调用的正确性。\par
具体而言，Realm 创建前后，granule、REC、页表与 measurement 都有严格的前置/后置状态：host 只能通过 RMI 发出请求，RMM 检查前置条件满足后才执行转换；Realm 通过 RSI 获取服务与 evidence，避免把信任传递给 normal world。这一“状态机而非函数集”的视角，是理解后续所有部署论文的钥匙——RContainer 的 con-shim 创建/销毁、CAGE 的 shadow task 同步，本质都是在 RMM 状态机旁路或内部再嵌套一层小型状态机。\par
\subsubsection{系统架构}
从系统结构看，RMM 夹在 untrusted host 与 Realm 之间，向两侧暴露不同接口：host 一侧是 RMI，承担创建、内存委托、REC 运行与销毁；Realm 一侧是 RSI，承担 Realm 可见服务与报告；硬件一侧则由 GPT/GPC 执行最终访问判定。\par
具体而言，这一三面结构决定了 RMM 的失败模式：RMI 侧的前置条件检查若有遗漏，host 可获得非法状态转换；RSI 侧的 evidence 生成若绑定不完整，verifier 会被错误初始状态欺骗；而 GPT/GPC 一侧的正确性则由架构与 EL3M 承担，不在 RMM 自身的可证范围内。由于本地无规范 PDF，以上接口分解依据公开生命周期概念与 Li 2022 Table 2 的本地证据，具体命令集未本地核验。\par
\subsubsection{核心方法设计}
作为规范材料，其“方法”体现为生命周期上的三段公开语义。\par
\noindent\textbf{（1）Realm 创建与销毁。}\par
这一机制的直接作用是：让 host 管理生命周期，但 host 不能越过状态机直接读写 Realm。\par
具体实现上，按公开生命周期概念，一个 Realm 的完整轨迹是：分配并初始化 metadata、委托 granule 到 Realm PAS、测量初始状态形成 measurement、激活 Realm、创建并运行 REC、最终销毁并清理。每一步都有前置条件：例如委托要求 granule 当前处于可转换状态，激活要求测量已完成，销毁要求私有 granule 被清理以避免数据残留。这些步骤是后续 container/accelerator 论文必须兼容的底座：RContainer 的 con-shim termination 之所以需要约 637 微秒清理内存，正是这一“销毁必须 scrub”语义在容器粒度上的投影；CAGE 的加速器环境清理则是同一语义在设备侧的投影。\par
\noindent\textbf{（2）REC 运行与调度。}\par
这一机制的直接作用是：保留云平台调度能力，同时把 REC 状态的安全保存/恢复收进 RMM 边界。\par
具体实现上，host 决定何时运行 Realm，但寄存器状态的进出由 RMM 控制，防止 hypervisor 直接读取机密寄存器与内存。interrupt、device 与 scheduling 方向的论文都在这个接口附近扩展：Devlore 检查的是 REC 运行期间被注入的中断是否合法，CAGE 检查的是 REC 运行期间发起的加速器任务是否合法。\par
\noindent\textbf{（3）Attestation 与 Measurement。}\par
这一机制的直接作用是：让生命周期产生 verifier 可用的初始状态证据。\par
具体实现上，测量绑定 Realm 初始配置、内存内容与平台状态，verifier 依据 evidence 决定是否把 secret 交给 Realm。公开概念与 Li 2022 的本地证据互相印证：Realm 创建期间被赋值的 IPA 与数据经密码学哈希进入 attestation token，激活后测量固定，之后只能添加未知内容——这一“先测量后激活”的次序是 RMM 状态机中最关键的安全性质之一，因为任何激活后的初始状态变更都会使 verifier 的决定失效。这是本章与本报告 attestation 方向的连接点：RMM 生命周期语义决定“证据覆盖到什么状态”，attestation 方向的论文决定“证据如何被验证与传递”。\par
\noindent\textbf{（4）RMM 实现的工程要求。}\par
这一机制的直接作用是：把规范语义转化为对实现者的约束清单。\par
具体实现上，按公开概念与 Li 2022 验证经验可归纳四条要求：
\begin{itemize}
\item TCB 最小化：RMM 只做安全仲裁，虚拟化功能交还 hypervisor，Li 2022 证明这可以使 RMM 比裸机 hypervisor 小几个数量级；
\item 接口审计可行性：RMI/RSI 分离使审计边界清晰——审计 RMI 即审计 host 的全部管理能力，审计 RSI 即审计 Realm 信任链的全部入口；
\item 并发与弱内存正确性：RMM 需要支持多核并发操作（细粒度锁、hand-over-hand locking），其实现正确性不能靠测试覆盖，Li 2022 的 VIA 证明了机器检查的可行性但也展示了其成本；
\item 补丁与调试策略：调试接口不能成为绕过 granule 检查的后门（Li 2022 第 3 节讨论调试寄存器与 Realm 语义的冲突），实现者必须对调试与固件更新路径做显式策略。
\end{itemize}
这些要求超出规范文本本身，是从本章证据链归纳的工程推论，规范是否逐条强制，本地证据不足。\par
\subsubsection{安全性与威胁边界}
RMM 规范的安全性边界与 Li 2022 一致：host 可通过 RMI 拒绝服务（不调度、不委托内存），这属于威胁模型允许范围；host 不能通过任何 RMI 序列获得 Realm 私有状态，这是状态机要维护的不变量。边界在于：规范定义状态机语义，不证明任何具体 RMM 实现满足它——Li 2022 的验证针对的是其原型固件，商用 RMM 实现需要独立的形式化或审计证据；且本地无规范 PDF，具体状态机的完整性与版本演进细节证据不足。\par
\subsubsection{实验环境与证据}
官方 source 已记录，但本地 PDF 不可用。可核验材料为：Arm DEN0137 文档页面与 CDN URL 状态（README 记录）；可支撑的结论是 RMM/RMI/RSI 生命周期的公开概念；不能支撑的是未本地读取的具体表格、命令编码、性能或实现细节。本节所有机制性描述均以公开概念与 Li 2022 的本地证据交叉支撑。\par
\subsubsection{实验结果与证据强度}
RMM 规范不给出 overhead，性能需看 RMM 实现或系统论文。RMI/RSI 调用、granule transition 与 REC exit/entry 都可能影响性能，规范只能说明这些路径存在。具体数据在本章由部署论文给出：RContainer Table VIII 报告了 OS switch 约 492 微秒、con-shim 创建约 530 微秒等 RMM 相邻路径的微基准成本，可作为“RMM 类状态机操作在原型系统中的一阶代价”参考，但不能反推为规范本身的性能属性。\par
\subsubsection{综合评价}
综合来看，RMM spec 是 CCA 可信边界的状态机锚点：优势是清晰分离 host-facing 与 Realm-facing 接口，使 TCB 可审计；局限是本仓库缺本地 PDF，细节论断必须让位于本地可读论文。从部署角度看，RMM 是 Arm CCA 云平台的必需组件，其商业化成败取决于实现质量、TCB 大小控制、补丁与调试策略，以及与 device lifecycle（RME-DA 方向）的集成——这些正是本章后续论文逐一切入的缺口。\par

\subsection{论文 4：SHELTER: Extending Arm CCA with Isolation in User Space}
\subsubsection{论文基本信息}
《SHELTER: Extending Arm CCA with Isolation in User Space》由 Yiming Zhang、Yuxin Hu、Zhenyu Ning、Fengwei Zhang、Xiapu Luo、Haoyang Huang、Shoumeng Yan、Zhengyu He 完成（Southern University of Science and Technology、The Hong Kong Polytechnic University、Hunan University、Ant Group），发表于 USENIX Security 2023（论文页码 6257--6274）。本文将其置于“Arm CCA 细粒度隔离与部署模型”的研究脉络中考察，并把它作为 system 类（E1）材料使用。它所回答的核心问题可概括为：把 Arm CCA 从保护整台 Realm VM 扩展到 Normal world 内的用户态安全应用（SApp），使第三方开发者无需 trusted OS 即可获得硬件级隔离。当前证据状态为：本地 PDF 已保存并核验；下文仅在该范围内讨论其机制与结论。\cite{zhang2023shelter}\par
\subsubsection{研究背景与动机}
要理解该方案，首先需要看到它面对的系统矛盾：CCA 的 Realm 粒度太粗，TrustZone 的 trusted OS 又太重。论文引言给出两个具体观察：其一，TrustZone 对第三方开发者可达性差——TEE 厂商需要严格审核每个 Trusted Application 以防引入漏洞，这与服务的快速上线节奏冲突；其二，商用 TrustZone 系统的攻击面随 TA 与 trusted OS 数量增长而扩大，trusted OS 一旦被攻破，整个 TrustZone 系统随之失守。\par
具体而言，很多设备只需要保护少量用户态安全应用（密钥处理、认证、轻量安全服务），而不是整台 VM。把完整 trusted OS 放进 secure world 会扩大 TCB；把每个安全应用放进独立 Realm VM 又引入 VM 级开销与未成熟的 CCA 软件生态（论文指出当时开发者甚至无法使用 RMM）。Shelter 的动机是：CCA 在 Armv9.2 上提供的 GPC/GPT 原语能否直接保护 Normal world 用户态中的单个应用，从而获得比 TrustZone 更轻、比 Realm VM 更细的部署模型。\par
\subsubsection{关键挑战与核心思想}
在上述背景下，作者提出的关键判断是：只要能让 SApp 的物理页在 GPT 上对 host 不可访问，就能在 Normal world 内构造一个轻量 TEE——Monitor 驻留 Root world，用多份 GPT 为 SApp 与其他代码提供不同的物理访问视图。\par
具体而言，论文明确命名了三个挑战：
\begin{itemize}
\item C1（内存隔离）：CCA 原生 GPT 只把内存划分为 Normal/Secure/Realm/Root 四个 PAS，若 SApp 与 host 同在 Normal world，被攻破的特权软件对 SApp 内存拥有完全访问权——Shelter 的对策是基于“每个 CPU core 可以独立配置 GPT（类似扩展页表）”这一观察，提出 multi-GPT 设计，为每个 SApp 建立 per-core 的地址空间视图；
\item C2（启动延迟）：每次创建 SApp 都要新建 SHELTER GPT 并填入表项，初始化开销大——对策是改进 GPT 管理以加速创建；
\item C3（多核 GPT 缓存一致性）：GPT 信息可被 TLB 缓存并跨核共享，攻击者可借助另一核上缓存的 SHELTER GPT 信息绕过 GPC——对策是专用的多核管理。
\end{itemize}
这三条挑战说明：把架构原语变成可用隔离，难点不在原语本身，而在生命周期、性能与多核一致性三个工程面。\par
\subsubsection{系统架构}
从系统结构看，Shelter 把 Monitor 放在 EL3 的 Root world（论文 Figure 2），利用 RME 使 Root world 对其他 world 天然不可访问；用 multi-GPT 为 host 与 SApp 提供不同的物理访问视图（Figure 3）。\par
具体而言，Host OS 继续管理普通应用与大部分资源，不需要被信任；SApp 在 Normal world 用户态执行，但其内存通过 SHELTER GPT（S.GPT）与 host GPT 的差异受到保护；第三方开发者通过 Shelter SDK 与 Monitor 交互，完成 SApp 的创建、内存申请与退出。Monitor 负责创建、切换、验证与回收 SApp 的隔离状态，是系统中唯一新增的高特权组件。\par
\subsubsection{核心方法设计}
论文的关键不是单一组件，而是把 SApp 抽象、multi-GPT 与系统交互协议组合成一条完整的用户态隔离路径。\par
\noindent\textbf{（1）SApp 抽象与威胁模型。}\par
这一机制的直接作用是：把保护对象从 VM/trusted OS 收窄为单个用户态应用，从而收窄 TCB 与部署成本。\par
具体实现上，SApp 的代码与数据需要防御 host OS、hypervisor 以及其他特权软件（trusted OS、secure/Realm hypervisor）的读取与篡改；DoS 与部分侧信道不作为主要目标。这个抽象适合密钥处理、认证与轻量安全服务，也为“CCA 用户态 SDK”这一产品形态给出了原型。\par
\noindent\textbf{（2）Multi-GPT 与 GPC 强制。}\par
这一机制的直接作用是：同一物理页在 SApp 视图可访问、在 host 视图不可访问，且该判定由 GPC 在硬件路径上最终执行。\par
具体实现上，每个 CPU core 可配置独立的 GPC/GPT base，Shelter 据此在 SApp 进入时切换访问视图：Monitor 检查分配区域不重叠后更新对应 GPT 表项，S.GPT 授权 SApp 访问，host GPT 拒绝 host 访问。由于 GPC 在 MMU 页表转换之外独立检查物理地址，host 即使在自身页表中建立映射也无法绕过。Table 1 给出各 world 的物理访问权限矩阵，Table 2 给出 GPI 编码。\par
\noindent\textbf{（3）Monitor 的职责划分与最小 TCB 设计。}\par
这一机制的直接作用是：让 Monitor 只执行安全策略，把非安全职责全部留给不可信 OS。\par
具体实现上，论文为 Shelter 定了四个设计目标：G1 安全（防御高特权软件）、G2 小 TCB、G3 无硬件修改、G4 性能。围绕 G2，Monitor 只通过有限的 SMC API 支持 SApp 部署与隔离维护；系统调用支持、调度与中断处理全部由 untrusted OS 承担。内存管理采用“转发加校验”模式：Monitor 不自己分配内存，而是把 SApp 的分配请求转发给 Host OS 的 buddy allocator，只记录返回的地址与大小并检查多个 SApp 之间无重叠——朴素方案逐页传递地址会引入多次上下文切换，论文第 4.2 节的 SHELTER Memory Management（含 united page allocation）正是为压缩这条路径。选择把功能放进 EL3 Monitor 而非 RMM 的理由也有明确记录：隔离只依赖 GPT 操作、不需要 RMM 的 stage-2 翻译；只有 EL3 有权改 GPT，RMM 操作还需 SMC 进入 EL3；EL3 固件本就负责上下文切换与 GPT 管理，可复用其库与数据结构。GPT 表项的 GPI 编码（论文 Table 2）采用 4 位取值：0000 禁止一切访问，1000/1001/1010/1011 分别仅允许 Secure/Normal/Root/Realm PAS，1111 全部允许——multi-GPT 隔离正是通过这些编码在不同 GPT 副本中的差异化设置实现的。\par
\noindent\textbf{（4）动态内存与系统交互。}\par
这一机制的直接作用是：让 SApp 仍然能使用 OS 服务，而不把私有状态暴露给 OS。\par
具体实现上，SApp 需要分配/释放内存、处理系统调用、与 normal world 交换非敏感数据；Monitor 验证每次 allocation 结果，防止 host 返回重叠区域或恶意页（Iago 式风险）；性能优化集中在减少 GPT 更新、上下文切换与跨核同步。这一部分解释了为什么 Shelter 的开销主要来自生命周期与交互路径，而非稳定运行态。\par
\noindent\textbf{（5）SApp 创建加速与 GPT 管理。}\par
这一机制的直接作用是：回应挑战 C2，把 SApp 的启动延迟从“每次新建完整 GPT”压到可交互水平。\par
具体实现上，朴素实现需要在每次创建 SApp 时新建一份 SHELTER GPT 并逐项填入表项；Shelter 改进 GPT 管理（论文第 4.3 节），通过预分配与复用 GPT 结构加速创建路径。这一优化针对的是用户态安全应用的真实使用形态——按需启动、短生命周期——若创建开销过大，SApp 抽象再轻也无法替代常驻的 trusted service。\par
\noindent\textbf{（6）多核管理与 GPC 缓存一致性。}\par
这一机制的直接作用是：回应挑战 C3，防止跨核的缓存状态绕过 GPC。\par
具体实现上，GPT 信息在多核间可被 TLB 类结构缓存并共享，攻击者控制的核可能利用另一核缓存的 SHELTER GPT 信息访问 SHELTER 内存；Shelter 采用专用的多核管理（论文第 4.4 节），在 SApp 进入、退出与迁移时同步各核的 GPT 配置与缓存状态。这一机制揭示了一个普遍性教训：CCA 的 per-core GPT 配置能力既是 multi-GPT 隔离的基础，也是跨核一致性攻击面的来源，任何复用该能力的后续设计（包括 RContainer 的 GPT 切换与 NanoZone 的 GPC3 bypass）都必须回答同一个问题。\par
\subsubsection{安全性与威胁边界}
Shelter 的隔离论证建立在三点同时成立上：S.GPT 与 host GPT 的视图差异正确维护、Monitor 对分配结果的验证无遗漏、多核 GPT 缓存被专用管理约束。威胁模型内，攻击者控制 Normal/Secure/Realm world 的任意特权软件（包括 untrusted OS、hypervisor、trusted OS、SPM、RMM），也可部署恶意 SApp——Shelter 以双向隔离保证 SApp 不能访问其他 SApp 或任何软件的私密数据，反之亦然。信任假设被论文显式列出：Monitor 可信（经厂商签名校验并由 secure boot 安全加载）、RME 硬件可信（假定无 bug）；SHELTER 的 I/O 数据与持久存储假定由安全加密通道保护；远程 attestation 假定存在并可嫁接既有方案。威胁模型外包括物理攻击（故障注入、冷启动、总线窃听）、DoS 与侧信道（论文指出既有侧信道与投机执行防御可叠加）。论文第 6.2 节调研了 45 个以控制特权软件为目标的 CVE 报告，论证 Shelter 可防御这些漏洞被利用后对 SApp 的未授权访问——这是“特权软件可被攻破”假设下的实证安全分析，而非形式化证明。\par
\subsubsection{实验环境与证据}
Shelter 的证据包括功能验证与性能评估两部分：两个原型分别实现在支持 RME 的 Arm FVP（功能验证）与硬件 SoC（性能评估）上；证据源为 USENIX Security 2023 本地 PDF。可核验材料包括 Figure 1（CCA 组件）、Figure 2（总体设计）、Figure 3（multi-GPT）、Table 1/2、45 个 CVE 的安全分析与实验章节。\par
\subsubsection{实验结果与证据强度}
论文报告的性能证据可分为四层，每层数字的基线与含义不同，必须分开引用。\par
第一层是应用负载开销（Figure 8，相对 Linux 基线）：纯用户态计算型应用开销极低——OTP 与 LeNet 在执行期无系统调用、开销为零，Squeezenet 开销可忽略，AES 仅因分配对象的系统调用带来 5.2\% 开销；大规模应用开销集中在上下文切换密集的路径——Apache 最高 15.0\%，Nginx 11.8\%，Memcached 8.3\%。论文将三者开销较高归因于复杂操作引发的 SHELTER 与 OS 间频繁上下文切换及附加的微架构维护与安全检查。摘要中的“小于 15\%”即来自这组真实工作负载。\par
第二层是对 Normal world 系统的影响（Figure 9，UnixBench 并发运行 SApp）：SApp 以 1000ms 间隔运行时平均开销 0.37\%、单项最高 1.4\%；100ms 间隔平均 1.87\%、最高 3.29\%；最密集的 10ms 间隔平均 7.68\%、最高 16.9\%——说明 Shelter 的系统级代价与 SApp 生命周期频率强相关。\par
第三层是内存利用率（Table 8）：united page allocation 使小对象密集场景受益最大——Numeration 基准（20 万个 16B--128B 小对象）内存利用率改善 98.5\%，Apache 17.6\%、Nginx 68.6\%、Memcached 34.8\%。\par
第四层是与虚拟化路线的对比（Figure 10，作者在 Armv8-A Juno R2 上用 GPT-analogue 方法实现 CCA 性能原型）：在 Apache、Memcached、Nginx 上，Realm VM 平均开销 32.0\%、Vanilla VM 29.8\%，而 Shelter 平均 11.7\%；论文进一步给出切换成本分解——Vanilla VM 与 Linux KVM 间切换 0.41 微秒，而 Realm VM 与 CCA KVM 间切换达 1.886 微秒（VM 到 RMM 0.813、RMM 到 EL3M 0.872、EL3M 到 CCA KVM 0.201 微秒），解释了 VM 式路线的固有成本。与 Sanctuary 的对比（Figure 7）中，Shelter 在创建与销毁两个生命周期操作上均更快。边界在于：CCA 性能原型与 Shelter 原型均不含 attestation 与硬件加密（当时 FVP 尚无 Memory Protection Engine），对比数字应视为同平台上的相对证据而非绝对预测。\par
\subsubsection{失败模式与边界分析}
论文与后续工作共同暴露了 Shelter 的三条边界。第一，EL3 代码量：RContainer 论文指出 Shelter 的 security monitor 有约 2K 行代码运行在 EL3——EL3 是 Arm 架构最高特权级，其中任何缺陷对整个平台都是灾难性的，因此“EL3 新增代码最小化”成为后续工作的直接动机。第二，语义鸿沟：RContainer 的对比实验观察到，把安全功能集中在 EL3 需要额外做 EL1 上下文分析以弥合异常级间的语义差，这推高了开销。第三，粒度边界：SApp 不是完整容器，系统调用与 I/O 仍需逐一设计 declassification 边界，进程内部的数据域隔离完全不在其范围内。\par
\subsubsection{综合评价}
综合来看，Shelter 是 CCA 细粒度隔离的起点：优势是不依赖完整 trusted OS、直接复用 GPC/GPT 原语、机制边界清晰、SDK 形态贴近第三方开发者；局限是粒度停在单应用、EL3 代码量与语义鸿沟推高成本、容器与进程内场景不适用。从部署角度看，它适合设备内密钥服务、轻量安全模块与 CCA 用户态 SDK 的第一代形态，其设计经验（multi-GPT、per-core 视图、多核 GPT 管理）被后续工作反复引用与修正。\par

\subsection{论文 5：RContainer: A Secure Container Architecture through Extending ARM CCA Hardware Primitives}
\subsubsection{论文基本信息}
《RContainer: A Secure Container Architecture through Extending ARM CCA Hardware Primitives》由 Qihang Zhou、Wenzhuo Cao、Xiaoqi Jia、Peng Liu、Shengzhi Zhang、Jiayun Chen、Shaowen Xu、Zhenyu Song 完成（中科院信工所、国科大、宾州州立大学、波士顿大学），发表于 NDSS 2025。本文将其置于“Arm CCA 细粒度隔离与部署模型”的研究脉络中考察，并把它作为 system 类（E1）材料使用。它所回答的核心问题可概括为：把 CCA 的隔离能力包装成 container-friendly 的运行时——强隔离、轻量、最小高特权代码三个目标同时成立，而不是让每个容器都变成重 VM。当前证据状态为：本地 PDF（19 页正文加附录）已保存并核验；下文仅在该范围内讨论其机制与结论。\cite{zhou2025rcontainer}\par
\subsubsection{研究背景与动机}
要理解该方案，首先需要看到它面对的系统矛盾：容器的安全性与轻量性长期冲突。容器共享 host 内核，恶意容器可通过内核接口、procfs 或 namespace escape 攻击其他容器与 host；把每个容器放进 VM 或 microVM 获得强隔离，则损失容器的快速启动与高密度部署优势，且 hypervisor 本身就扩大 TCB。TrustZone 路线（TZ-container、TrustShadow）则面临所有容器在 secure world 内权限相同、trusted OS 复杂度持续增长两个问题。\par
具体而言，论文明确提出三个同时成立的目标：强隔离（对不可信 OS 与其他容器双向隔离）、轻量（接近原生容器的运行时与生命周期开销）、最小高特权代码（新增 EL3 代码尽可能少）。论文进一步指出 Shelter 式方案虽大幅降低了相对 hypervisor 的开销，但把约 2K 行代码放进 EL3 仍违背第三个目标——EL3 是平台最高特权级，其中问题的影响是全局性的。CCA 的 RME 原语（GPC/GPT、四个 PAS）提供了新的机会：物理地址空间的权限可以由硬件判定，那么容器隔离的主体逻辑能否下沉到 EL1 执行？\par
\subsubsection{关键挑战与核心思想}
在上述背景下，作者提出的关键判断是：RContainer 不把 OS 整体放进 TCB，而是用一个可信 mini-OS 管理安全关键路径，让 deprivileged OS 继续提供普通服务；两者同处 EL1，靠 mixed-pagetable 与 GPT 视图互相隔离。\par
具体而言，论文命名了两个挑战。C1（容器不适合直接放进 Realm world）：多容器共享一个 Realm OS 内核则容器间隔离不足，每容器一个 Realm VM 则退化为 microVM 方案的性能与 TCB 问题。C2（小 EL3 代码量下如何保证其余 TCB 防篡改）：把 TCB 放进 Secure/Realm world 会增加复杂度并引入大量跨 world 操作；把 TCB 放进 Normal world 则与 host OS 同特权级，需要同级隔离——而教科书上两种同级隔离做法（独立页表 + EL3 挂钩页表操作；共享页表 + 页表页只读 + EL3 代理更新）都会因频繁页表/上下文切换与 TLB 刷新带来显著开销。mixed-pagetable 正是为绕开这两个高开销方案而设计。\par
\subsubsection{系统架构}
从系统结构看，RContainer 的 TCB 由 EL1 的 mini-OS 与 EL3 的原生 secure monitor 构成，容器隔离单元是 con-shim（论文 Fig.1）。\par
具体而言，mini-OS 是一个精简基础 OS，负责容器/con-shim 内存管理与控制流保护；服务供给仍由 deprivileged OS 承担。每个容器在内核层获得一个 con-shim 实例——一个隔离的物理地址空间，由独立的 Shim-GPT 定义其访问权限：在该容器的 Shim-GPT 中，mini-OS、deprivileged OS 与其他容器的内存均被剔除，形成用户态与内核空间的双向隔离（Security Insight 1）。mini-OS 与 deprivileged OS 同级运行但由 mixed-pagetable 隔离（Fig.2）；pagefault 处理区分快慢路径，快路径尽量不进入高权限逻辑（Fig.3）。架构设计同时服务 Security Insight 2（最小化最高特权代码：新增 EL3 代码仅 130 行）与 Security Insight 3（安全特性可扩展：mini-OS 处于 EL1，可以集成 Iago/MUMA 等更多防御而不膨胀 EL3）。底层依赖的架构事实是：CCA 把全部物理地址空间划分为 Normal/Secure/Realm/Root 四个 PAS，Root world 可访问全部 PAS，而 GPT 相关操作（加载 \texttt{GPTBR\_EL3}、更新 GPT）只能由 EL3 执行——RContainer 的全部隔离能力最终都锚定在这 130 行 EL3 代码对 GPT 的独占操作上。\par
\subsubsection{核心方法设计}
论文的关键不是单一组件，而是把 TCB 拆分、mixed-pagetable 与生命周期快路径组合成一条可用的容器隔离路径。\par
\noindent\textbf{（1）Mini-OS 与 Con-Shim。}\par
这一机制的直接作用是：只把容器隔离所需的关键能力放进高权限，压缩高特权 TCB。\par
具体实现上，con-shim 是容器在内核空间的轻量隔离实例，承接容器的内核侧数据面；mini-OS 管理 con-shim 内存、容器与 OS 间的控制流以及 pagefault 快路径。deprivileged OS 继续提供调度、驱动与普通服务，但在当前容器的 Shim-GPT 视图中不可访问 con-shim 私有页。由此容器获得双向隔离：既防不可信 OS 读取容器，也防恶意容器经内核数据结构攻击其他容器（对应论文 Security Insight 1 的用户态加内核空间双向隔离）。\par
\noindent\textbf{（2）Mixed-Pagetable 与 Shim-GPT。}\par
这一机制的直接作用是：页表可以共享，物理访问权限由不同 GPT 视图隔离，从而避免独立页表方案的高频切换开销。\par
具体实现上，mini-OS 与 deprivileged OS 使用同一份 MMU 页表但不同的 GPT：EL3 secure monitor 维护 Priv-GPT（mini-OS 的视图，全权限）与 OS-GPT（deprivileged OS 的视图，仅可访问自身内存与自身页表页）；维护 mini-OS 映射的页表页在 OS-GPT 中被标记为 No-access PAS，使 deprivileged OS 运行时无法触及 mini-OS 的页表。一个关键约束是：除 L0 外，任何页表页只能存放 mini-OS 相关或 deprivileged OS 相关映射之一，不能混放；唯一同时映射两者的 L0 页由 shadow L0 机制单独保护，防止 deprivileged OS 篡改。容器隔离侧，每个 con-shim 在 EL3 拥有独立 Shim-GPT，其 GPI 编码（论文 Table II）为：自身容器与自身 con-shim 为 1001（Normal PAS，可访问），其他容器、其他 con-shim、mini-OS、deprivileged OS 均为 0000（No-access PAS），secure monitor 为 1010（Root PAS，天然不可达）。mini-OS 维护 GPT-routing 表记录容器到 GPT 的映射：容器创建时由 mini-OS 通知 EL3 分配 GPT 并登记索引，容器运行前查到索引并通知 EL3 切换 GPT。deprivileged OS 照常维护页表（兼容性好、无 EL3 代理开销），权限差异完全由 GPT 视图表达——这等价于把 C2 中“同级隔离”的成本从页表层移到了 GPT 层。\par
\noindent\textbf{（3）Pagefault 与生命周期快路径。}\par
这一机制的直接作用是：让创建、销毁、OS 切换与 pagefault 这些高频路径足够轻，系统才可用。\par
具体实现上，Fig.3 区分快/慢 pagefault workflow；Table VIII 给出关键微基准：multi con-shim 环境初始化约 45607 微秒（仅系统运行期执行一次）、deprivileged OS 与 mini-OS 间的 OS switch 约 492 微秒、con-shim 创建约 530 微秒、销毁约 637 微秒（需清理全部 con-shim 相关内存以消除信息残留）、4KB 页的 Set GPI 约 530 微秒。论文还分析了 30 个容器与 Linux 相关 CVE（附录 Table XI），其共同前提是“内存隔离弱 + 权限不受限”，并逐一模拟深度攻击验证均触发 GPT fault。\par
\noindent\textbf{（4）控制流保护与异常拦截器。}\par
这一机制的直接作用是：mini-OS 作为容器与不可信 OS 之间的安全监视器，必须同时防数据面篡改与系统调用返回值欺骗。\par
具体实现上，控制流保护由 exception interposer 实现。系统调用路径：控制流先从容器切到 mini-OS，mini-OS 按请求信息拷贝系统调用参数、清除敏感 CPU 寄存器，再切到 deprivileged OS 处理；处理完成后控制流回到 mini-OS，检查结果并验证通过才返回容器。异常与中断路径更危险——它们可不经系统调用直接触发容器与 deprivileged OS 间的异常级切换，造成控制流劫持：RContainer 强制异常/中断先进入 EL3 monitor 做 GPT 切换（Shim-GPT/OS-GPT 到 Priv-GPT），再交 mini-OS 分发给 deprivileged OS 的对应 handler（伴随 Priv-GPT 到 OS-GPT 的切换）；mini-OS 维护一份受 deprivileged OS 不可修改的异常向量表副本，所有 handler 入口地址硬编码，每次异常处理完成后检查结果再返回容器。shim-style 隔离的理论依据是论文的两个观察：容器相关攻击多发源于控制面（控制流劫持、缺陷代码），但最终落在数据面（全局变量、内核数据结构）；数据面被访问更频繁、需要更强隔离。据此 RContainer 在内核数据面识别三类关键数据——内核边界点（系统调用进出点）、容器私有数据结构（如 \texttt{task\_struct}）、共享全局变量（如 \texttt{nr\_files}）——把它们连同共享内存与私有栈一起装进每个容器的 con-shim，而全部 con-shim 都在 TCB 之外。\par
\noindent\textbf{（5）容器生命周期保护与 Iago 防御。}\par
这一机制的直接作用是：把隔离从“运行时的静态视图”扩展为覆盖启动、任务创建、内存分配与回收的全过程。\par
具体实现上，启动采用 secure boot 加 lazy loading：deprivileged OS 先被 EL3 加载度量并正常启动（此时网络等服务未初始化，远程攻击面未打开），mini-OS 启动后立即通知 EL3 初始化 OS-GPT 与 Priv-GPT、把自身内存区设为 OS-GPT 中的 No-access PAS，此后 deprivileged OS 即失去对 mini-OS 的访问能力。任务创建区分三种情形：新进程拥有独立地址空间时，mini-OS 验证并记录新 task 结构与新页表地址，并指示 secure monitor 将相关区域在 OS-GPT 中设为 No-access；新线程共享地址空间时只需记录 task 结构；execve 替换地址空间时还需擦除旧区域并在 Shim-GPT 中设为 No-access。内存分配引入 reservation 机制：启动期预留若干连续物理区，全部页初始在所有 GPT 中 No-access；分配给某容器时由 EL3 把该页在其 Shim-GPT 中标记为 Normal PAS——快 pagefault 路径（Fig.3a）因此无需离开 mini-OS；内存池耗尽时走慢路径（Fig.3b）向 deprivileged OS 申请新内存区，此时 mini-OS 必须做 Iago 检查：确认该区域未被映射、核对进程映射范围链表中的 VA 合法性、清空全部物理页后才回填页表。共享内存与 COW 也走受控路径：系统调用参数经 con-shim 内专门的共享缓冲区（OS-GPT 中 Normal PAS）中转，mini-OS 检查源页属容器、目的页属 deprivileged OS 后才调整 GPT。\par
\noindent\textbf{（6）文件、网络与 IPC 保护。}\par
这一机制的直接作用是：覆盖容器数据离开内存边界的三条通道。\par
具体实现上，文件系统分两类：容器内部文件依赖应用自身的加解密；需由 deprivileged OS 加载执行的二进制采用 scan-encrypt 预处理——\texttt{rc\_image\_binary\_crypto} 脚本扫描镜像、用公钥加密全部 ELF（保留 ELF 头供加载器解析）、哈希并重打包，配对私钥存于 host 安全存储、仅 mini-OS 与 secure monitor 可达；加载时 mini-OS 校验哈希完整性并解密，动态链接库的 FD 在容器创建时已记录于 con-shim。网络 I/O 不做额外系统级保护，依赖 HTTPS/SSL/TLS 等安全传输协议——论文对此边界是显式的。IPC 分两种：共享内存（shmget/shmat）复用 con-shim 共享内存路径；管道、socket、消息队列则按 FD 建立加密通道，使用容器初始化时的密钥由 mini-OS 在系统调用拦截点加解密——这解释了实验中 IPC 密集负载（Hackbench、lmbench IPC 项）开销偏高的直接原因。\par
\noindent\textbf{（7）系统调用兼容与 TCB 量化。}\par
这一机制的直接作用是：证明“安全容器”不以牺牲 Linux ABI 为代价，并把 TCB 从定性口号变成可审计的数字。\par
具体实现上，论文用 Linux Test Project（LTP v20240129）的 1409 个测试用例在原生容器与 RContainer 上对照执行：970 个通过、280 个跳过、159 个失败——两组结果一致，失败项主要源于内核版本（kernel-firefly-4.4.194）不支持 io\_uring 等新特性，与 RContainer 机制无关，即 RContainer 的系统调用覆盖面与原生容器基本相同。TCB 分解（Table VI）：未加密配置 2.6K SLoC，其中 EL3 的 secure monitor GPT 管理仅 130 行 C；启用镜像与 IPC 加密（AES-ECB，Python 镜像解析/扫描/加密组件加约 2308 行 C）后总计约 4955 SLoC，另有约 200 行补丁。作为对照，论文统计原生 EL3 代码总量约 421621 SLoC——RContainer 新增的 130 行相对其基座可忽略，相对 Shelter 式方案约 2K 行的 EL3 代码则低一个数量级。\par
\subsubsection{安全性与威胁边界}
威胁模型（Table I 矩阵）假设攻击者可控制 Secure/Realm/Normal world 的系统软件，可劫持控制流、重映射内存、控制 DMA 设备并发起 Iago/MUMA 攻击；也假设攻击者可部署恶意容器进行提权与资源滥用。容器主动泄露自身数据、侧信道、密码学攻击、物理攻击与一般 DoS 在范围外（但内存相关 DoS，如抽象资源攻击与内存耗尽，被纳入防御）。论文声明系统启动期良性、容器镜像可加密、既有 attestation 可叠加。隔离论证依赖：mixed-pagetable 完整性、Shim-GPT 视图正确、EL3 对 GPT 操作（\texttt{GPTBR\_EL3} 加载与 GPT 更新）的独占。\par
\subsubsection{失败模式与边界分析}
论文与其实验设计暴露了四条边界。第一，平台偏差：性能主要在 ARMv8 机器上模拟 GPT 操作获得，作者明示与真实 ARMv9-A 机器可能存在偏差，所有百分比数字都应连同该前提引用。第二，高开销路径：文件系统类微基准开销约为原生两倍（系统调用参数与数据需来回拷贝）、IPC 加密模式约 20\%、Hackbench 接近 1 倍——“平均 5\%”不适用于 IPC/IO 密集型负载，部署前需按负载画像评估。第三，扩展性天花板：超过 32 个并发容器时原生与 RContainer 均因开发板内存不足无法运行，多容器扩展性结论仅限于该规模内。第四，威胁模型留白：一般 DoS、侧信道与容器内应用自身漏洞在范围外，生产部署需要叠加独立的侧信道缓解与应用层防护。\par
\subsubsection{实验环境与证据}
RContainer 的两个原型分别实现于支持 ARMv9-A 的 FVP 与 ARMv8 硬件 SoC（RK3399 开发板）；评估覆盖 threat model 矩阵、架构、mixed-pagetable、pagefault workflow、lmbench 微基准、六类真实应用容器、docker 生命周期与多容器 kernbench 扩展性。TCB 量化：未启用加密时 2.6K SLoC（其中 EL3 仅 130 SLoC），启用 IPC/镜像加密（AES-ECB 等）时增加约 2.3K SLoC，总计约 4955 SLoC；作为对照，原生 EL3 代码总量约 421621 SLoC。论文明确声明性能评估的局限：因无法获得 ARMv9-A 真机，性能主要在 ARMv8 机器上模拟 GPT 相关操作，与未来真实 ARMv9-A 机器可能存在偏差。\par
\subsubsection{实验结果与证据强度}
论文报告的关键数字（Fig.4--Fig.8、Table VIII/IX）可归纳如下。lmbench：内存类用例几乎无开销（不拦截页表转换），进程类（fork+exec 等）约 10\%（需验证父子进程 task struct 与地址空间并修改 GPT），文件系统类最高约为原生两倍（deprivileged OS 不能直接访问容器内存，需来回拷贝系统调用参数与数据），IPC 类在加密模式下约 20\%。内核编译（linux-4.19.309 defconfig，Table IX）：RContainer 约 4.5\%、加密模式约 6.6\%，远优于虚拟化方案（原生 3248.38s、RContainer 3393.12s、RContainer+IPC 3463.1s、Container+VM 8060.29s）。\par
真实应用（Fig.5）：Apache 与 Nginx 相对原生 docker 开销 4\%--7\%、平均约 5\%，而虚拟化超过 50\%；Memcached 稳定在约 0.3\%；MySQL 约 0.2\%--0.3\%（IPC 加密模式也仅 0.7\%，虚拟化约 34\%）；Netperf 低于 3\%；Hackbench 最差、开销接近 1 倍——论文归因于频繁线程创建引发的内存属性检查/GPT 修改与 IPC 加解密，但仍比虚拟化降低近 70\%。与 Shelter 的对比（Fig.6，作者在 RK3399 上尽力移植 Shelter）：大规模应用（Apache、Nginx、Memcached、MySQL）上 RContainer 平均开销低约 5.7\%。生命周期（Fig.7，busybox 镜像、每命令 10 轮取均值）：多数命令开销低于 5\%——docker create 3.06\%、start 1.17\%、pause 3.62\%、unpause 0.21\%、rm 2.41\%；docker kill 最高约 8.5\%，因需清理容器内存并修改 shim-GPT 与 OS-GPT 权限。多容器扩展性（Fig.8，kernbench allnoconfig）：开销可忽略，因各容器操作限于自身 con-shim/Shim-GPT；超过 32 个并发容器时原生与 RContainer 均因开发板内存不足无法运行。\par
\subsubsection{与相关工作的定量对比}
\begin{table}[htbp]
\centering\small
\caption{RContainer 关键性能数字汇总（据论文 Fig.5/Fig.7/Table VIII/IX）}
\begin{tabular}{lll}
\toprule
路径 & 指标 & 报告值 \\
\midrule
Apache/Nginx（Fig.5） & 相对原生 docker 开销 & 4\%--7\%，平均约 5\%（虚拟化 $>$50\%） \\
Memcached（Fig.5） & 吞吐开销 & 约 0.3\% \\
MySQL（Fig.5） & 吞吐开销 & 0.2\%--0.3\%（IPC 加密 0.7\%；虚拟化约 34\%） \\
Netperf（Fig.5） & 开销 & $<$3\% \\
Hackbench（Fig.5） & 最差路径 & 接近 1 倍，仍比虚拟化低约 70\% \\
内核编译（Table IX） & 时间开销 & 约 4.5\%（加密 6.6\%） \\
OS switch（Table VIII） & 单次成本 & 约 492$\,\mu$s \\
con-shim 创建/销毁（Table VIII） & 单次成本 & 约 530$\,\mu$s / 637$\,\mu$s \\
docker 生命周期（Fig.7） & 各命令开销 & 多数 $<$5\%；kill 最高约 8.5\% \\
对比 Shelter（Fig.6） & 大应用平均开销差 & RContainer 低约 5.7\% \\
\bottomrule
\end{tabular}
\end{table}
该表同时标定了证据边界：所有数字来自 FVP 加 ARMv8 模拟环境，作者明示与真实 ARMv9-A 机器可能有偏差；Hackbench 与 IPC 加密是高开销路径，不能被 0.3\% 类数字掩盖。\par
\subsubsection{综合评价}
综合来看，RContainer 是 CCA 容器化部署的重要一步。在谱系定位上，它相对 hypervisor 类方案（Inktag、X-Container、BlackBox、gVisor）减小了 TCB 并保留容器 UX；相对 SGX 类方案（SCONE）不需要定制 C 库与 enclave 抽象；相对 TrustZone 类方案（TZ-Container、TrustShadow）解决了 secure world 内权限同质与 trusted OS 膨胀问题；相对 Shelter 则把安全功能主体从 EL3 移到 EL1，论文自称隔离强度相当而 EL3 新增代码降低一个数量级。其优势是容器抽象清晰、mixed-pagetable 真正解开了“同级隔离必然高开销”的死结、评估覆盖微基准到真实应用到生命周期与扩展性；局限是仍需修改 OS/runtime 路径、加密模式与 IPC 密集负载开销显著、设备/I/O 与复杂容器生态（编排、镜像仓库、运维工具链）未覆盖、性能证据依赖模拟环境。从部署角度看，它可服务云端机密容器平台，但依赖 Arm CCA silicon、RMM、container runtime 与运维工具链的共同就绪。\par

\subsection{论文 6：NanoZone: Scalable, Efficient, and Secure Memory Protection for Arm CCA}
\subsubsection{论文基本信息}
《NanoZone: Scalable, Efficient, and Secure Memory Protection for Arm CCA》由 Shiqi Liu、Yongpeng Gao、Mingyang Zhang、Jie Wang 完成（华中科技大学、George Mason University），2025 年 arXiv 预印本（arXiv:2506.07034）。本文将其置于“Arm CCA 细粒度隔离与部署模型”的研究脉络中考察，并按其预印本状态作为 E3 级材料使用——机制与性能结论均需等待同行评审，不能当作标准或商用产品证据。它所回答的核心问题可概括为：把隔离粒度推进到 CCA CVM 内部——同一进程内的不同数据域也需要隔离，且切换必须快到可用于每次请求。当前证据状态为：本地 PDF 已保存并核验；下文仅在该范围内讨论其机制与结论。\cite{liu2025nanozone}\par
\subsubsection{研究背景与动机}
要理解该方案，首先需要看到它面对的系统矛盾：CCA 保护整个 CVM，Shelter/RContainer 把粒度收窄到应用/容器/进程，但进程内部的 bug 仍然会泄露同进程内的秘密——Heartbleed 正是这类“非关键代码路径泄露私钥”的代表。论文给出两个量化观察：新版 Linux 超过 4000 万行代码，特权提升漏洞风险随代码量上升；OpenSSL 超过 90 万行代码，但其中处理敏感数据的代码占比很小，把整个程序当作单一保护单元会让非关键路径的漏洞暴露全部秘密。\par
具体而言，既有机制各有缺口。SGX 生态的 intra-enclave 方案不兼容 CVM 式隔离：Nested Enclave 通过硬件修改把 enclave 分为内外两层，但层间切换约 1.1 微秒（与原生 ecall/ocall 延迟相当），对每次请求一次的键值访问类高频场景不切实际；LightEnclave 用 MPK 把用户态域切换压到约 20 cycles，性能达标，但 MPK 长期是 Intel 平台专属——Armv7 曾有类似的 memory domains 特性，但其 Domain Access Control Register 用户态不可访问，且该特性在 AArch64 已被废弃；NestedSGX 类在 SEV CVM 内模拟 SGX 的方案移植到 CCA 需要对 realm world hypervisor 做非平凡修改。CCA 的 PAS/GPT 边界强但切换重。转机是 Arm 在 Armv8.9-A 引入的 Permission Overlay Extension（POE），使用户态权限切换在 Arm 上成为可能——NanoZone 正是以 POE 为支点，组合 PIE 与 CCA 的 PAS/GPC3，构造进程内多域隔离。论文还给出对照坐标：RContainer 完全不信任原生 OS、引入 5K 行以下的可信 mini-OS 管理机密容器，Shelter 不引入任何可信内核态组件、把粒度细化到进程级，但两者都停留在“单体式隔离模型”，未处理进程内威胁——这正是 NanoZone 的切入点。\par
\subsubsection{关键挑战与核心思想}
在上述背景下，作者提出的关键判断是：高安全边界与低切换开销不能只靠一种机制；NanoZone 用 POE 承接高频用户态切换，用 PIE/PAS 提供强边界兜底，用 root world 监控与 Code-Pointer Integrity 防止切换机制本身被滥用。\par
具体而言，论文命名了三个挑战：
\begin{itemize}
\item 可扩展性挑战：POE 最多支持 16 个域且仅 7 个可通用，而服务器按客户端隔离私数据的域数轻易超过 7——NanoZone 观察到 Arm 的间接权限模型还有 PIE（Permission Indirection Extension），把 4 个未用 PIE 索引映射到已有 POE 索引可将可用 POE 域从 7 扩到 28；再观察到 CCA 的 GPT（按 PAS 划分物理页）与 GPC3（允许每核对选定 PAS 绕过 GPC）与 POE/PIE 同构，于是跨 PAS 复用同一套 PIE/POE 索引，形成三层 zone 结构，支持事实上无限数量的域；
\item 效率挑战：PIE 域间或跨 PAS 的切换需要陷入更高特权级，测量显示其延迟约为轻量 POE 切换的 80 倍——对策是优化域分配，尽量把请求路由到同一 PIE 域与 PAS；
\item 安全挑战（domain-switch abuse）：进程内攻击者可劫持控制流滥用切换指令——对策是以 Armv9.4-A 的 Guarded Control Stack（GCS）硬化返回地址，并扩展出 Pointer-Integrity Memory（PIM）与 Code-Pointer Integrity（CPI）机制。
\end{itemize}
\subsubsection{系统架构}
从系统结构看，NanoZone 在用户态提供 PIM 与权限切换能力，在 root world 放置安全模块监控页表与权限状态（Figure 3 展示 permission indirection 与 overlay 的组合）。\par
具体而言，PIM 区域保存全部函数指针的影子副本，只能由一条专用非特权指令写入（二进制扫描把该指令的其他出现替换为普通 store，保证只有插桩代码能写 PIM）；root world 安全模块监控页表更新（防止 OS 修改 POE 域配置与内存映射），并拦截中断控制流、先路由到 root world 保存/恢复敏感上下文（如 POE 权限寄存器），防止被攻破的 OS 借中断路径破坏域状态。\par
\subsubsection{核心方法设计}
论文的关键不是单一组件，而是“快路径切换 + 强边界兜底 + 反滥用监控”的三层组合。\par
\noindent\textbf{（1）POE 快速域切换。}\par
这一机制的直接作用是：把高频域切换留在用户态，避免每次 trap 与 TLB flush。\par
具体实现上，POE 索引（类似 MPK 的 protection key）配置在 PTE 中，域内执行时通过在用户态直接修改 POE 寄存器临时授予敏感内存访问权，最小化暴露时间；POE 是 per-core 权限寄存器，其他核仍受自身配置约束。论文报告这类用户态切换约 20 cycles 量级（与 LightEnclave 的 MPK 切换同量级），这正是它适合 session key、per-client data 等每次请求一次切换的场景的原因。\par
\noindent\textbf{（2）PIE/PAS 强边界与三层 zone。}\par
这一机制的直接作用是：在 POE 域数受限与强度不足处，用 PIE/PAS 组合把权限绑定到 CCA 更强的地址空间语义。\par
具体实现上，PIE 以内核态基权限加用户态 overlay 权限提供间接层（Figure 3：PTE 的 PIIndex/POIndex 分别索引 PIRE0\_EL1 基权限与 POR\_EL0 overlay 权限）；PAS/GPT 提供 CCA 物理访问控制兜底，GPC3 让每核对选定 PAS 绕过 GPC 以降低检查成本。代价是跨 PIE/PAS 切换比 POE 慢约 80 倍，必须靠域分配策略压低其频率。\par
\noindent\textbf{（3）域分配优化。}\par
这一机制的直接作用是：把切换路径分布压向 POE 层。\par
具体实现上，Figure 2 给出 1000 个 Memcached 请求的实测分布：96.72\% 的切换命中 POE 域内，PIE 域转换占 2.73\%，跨 PAS 边界仅 0.55\%；由此平均域切换延迟被压到特权级切换的 4.87\%。这组数字是理解 NanoZone 性能结论的关键——其低开销不是来自强边界变快，而是来自强边界被调用得足够少。\par
\noindent\textbf{（4）GCS+PIM 的 Code-Pointer Integrity。}\par
这一机制的直接作用是：防止进程内攻击者劫持控制流后滥用域切换指令。\par
具体实现上，GCS 硬化返回地址防御 ROP；LLVM pass 自动插桩，把每个函数指针存入 PIM 并在每次解引用时与影子副本比对；整个 CPI 机制运行在用户态，开销很小（论文第 7.2 节）。root world 模块同时监控页表更新并重定向中断控制流，防御内核态对手破坏 POE 配置或借中断腐蚀权限状态。\par
\noindent\textbf{（5）三个案例研究的设计意图。}\par
这一机制的直接作用是：用三个互补场景检验三层结构的真实收益分布。\par
具体实现上，Nginx 案例隔离 per-session 密钥，Memcached 案例保护内存键值数据，NVM 案例做非易失内存数据隔离；三者分别对应“每请求一次切换”“高频键值访问”“持久化数据保护”三种切换频率与数据形态，恰好覆盖 POE 快路径与 PIE/PAS 强路径的不同配比。论文报告三个案例均未造成显著退化，其中 Nginx 的 22.67\% 总开销主要源于特权级隔离（即 Shelter 式基线本身的成本），域切换与 CPI 的增量仅 4.40\%。\par
\subsubsection{安全性与威胁边界}
NanoZone 同时考虑进程内对手与内核空间对手：前者可劫持控制流滥用切换指令，由 CPI/GCS 与 PIM 写保护防御；后者可篡改 POE 配置、内存映射或借中断破坏权限状态，由 root world 监控防御。论文第 6.1 节系统分析攻击面、第 6.2 节结合 CVE 与特权攻击场景检验，结论是两类对手均可被挫败。边界在于：这是预印本的系统性安全分析而非形式化证明；机制依赖 POE/PIE/GCS/GPC3 等新硬件特性的语义正确性；root world 模块本身成为新的 TCB 组件。\par
\subsubsection{实验环境与证据}
证据包括两个原型：Arm 官方模拟器上的功能版（验证对进程内与内核空间对手的防御）与 Arm 开发板上的性能版（session-key 隔离、内存键值保护、NVM 数据隔离三个案例研究）。代码量：软件栈约 4.5K LoC，内核支持为小于 1K LoC 的轻量驱动，编译期支持约 500 LoC（LLVM pass 加运行时库），Monitor 中的核心隔离模块约 3K LoC（TCB）。评估对象覆盖微基准（第 7.2 节）与 Nginx、Memcached、NVM 保护三个真实应用案例（第 7.3 节）。证据状态为 arXiv 预印本，未经同行评审。\par
\subsubsection{实验结果与证据强度}
论文报告的关键数字：用户态域切换约 20 cycles；PIE/PAS 切换延迟约为 POE 的 80 倍；Memcached 场景 96.72\% 切换命中 POE 域、PIE 2.73\%、PAS 0.55\%；平均域切换延迟仅为特权级切换的 4.87\%。真实应用：Nginx 经历 22.67\% 开销，但其中大部分来自特权级隔离本身——相对无进程内隔离的 Shelter 基线，域切换与 CPI 带来的额外开销仅 4.40\%；摘要另报告整体约 20\% 性能开销并保留 95\% 吞吐。这些数字应按预印本证据谨慎使用：它们来自开发板原型与特定工作负载，且“22.67\%”与“4.40\%”是相对不同基线的两个数，混用会严重高估或低估实际成本。\par
\subsubsection{失败模式与边界分析}
NanoZone 的边界有四点必须随数字一起引用。第一，证据等级：arXiv 预印本，未经同行评审，所有机制与性能结论都应标注该状态。第二，硬件不可得：POE（Armv8.9-A）、PIE、GCS（Armv9.4-A）与 GPC3 在论文写作时均无商用硬件支持，功能版跑在官方模拟器、性能版跑在开发板变体上，真实芯片上的行为可能不同。第三，性能悬崖：4.87\% 的平均切换延迟依赖 96.72\% 的 POE 命中率——一旦工作负载的域分配无法把切换压进 POE 层，单次 PIE/PAS 切换约 80 倍于 POE 的成本会立即使开销失控，域分配策略实质上是该系统最关键的性能组件。第四，TCB 迁移：root world 监控模块（约 3K LoC）成为新的高特权组件，其自身正确性未获形式化证明，与李 2022 已验证的 RMM 原型不同，这里的信任是新增的。\par
\subsubsection{综合评价}
综合来看，NanoZone 是本章粒度最细的一篇，设计上抓住了 intra-process 隔离的核心矛盾——切换频率与边界强度的乘积决定开销，并用三层结构把该乘积拆开优化。其优势是机制组合有启发性、反滥用设计完整、真实应用证据具体；局限是 arXiv 预印本状态、依赖 POE/PIE/GCS 等新硬件特性（当时无商用硬件支持）、root world 监控引入新的高特权组件。从部署角度看，它适合机密服务器内的 session key、租户数据与库级沙箱，落地需要编译器、运行时与 OS 的配套改造，其结论的稳固性有待同行评审与真实硬件复现。\par

\subsection{论文 7：ACAI: Protecting Accelerator Execution with Arm Confidential Computing Architecture}
\subsubsection{论文基本信息}
《ACAI: Protecting Accelerator Execution with Arm Confidential Computing Architecture》由 Supraja Sridhara、Andrin Bertschi、Benedict Schlüter、Mark Kuhne、Fabio Aliberti、Shweta Shinde 完成（ETH Zurich），本地 PDF 为 arXiv:2305.15986v2，标注为 USENIX Security 2024 论文的扩展版；README 证据等级按 preprint 处理。本文将其置于“Arm CCA 的 I/O、DMA、加速器与中断”的研究脉络中考察。它所回答的核心问题可概括为：让加速器成为 CCA Realm 可安全使用的一等资源，而不是把数据复制到普通 world 的 bounce buffer 中再做软件加密。当前证据状态为：本地 PDF 已保存并核验；下文仅在该范围内讨论其机制与结论。\cite{acai2023}\par
\subsubsection{研究背景与动机}
要理解该方案，首先需要看到它面对的系统矛盾：CVM 越来越需要 GPU/FPGA/TPU 加速大工作负载，但 TEE 的信任边界停在 CPU——云用户面临两难：把计算放到 TEE 外的加速器上则暴露数据，留在 CPU 上则承受性能损失。论文进一步指出，即使加速器自带设备级 TEE（如近年的 NVIDIA GPU），仍有两个缺口：设备无法判断它连接的主机侧执行体是不是合法 realm VM；CPU 侧 TEE 必须防御 rogue 设备与被同租 VM 控制的设备。\par
具体而言，论文区分了三种访问模式（Figure 1、Table 1）。Integrated 模式让设备直接访问 realm 内存，性能最好，但突破了 CCA 的默认假设（加速器被视为不可信外设），威胁模型不成立；Encrypted/bounce-buffer 模式先把数据拷贝到普通世界缓冲区再做软件加解密，安全性尚可，但引入两次额外拷贝与软件加解密成本；ACAI 模式的目标是保留硬件加密与直接路径的性能优势，同时补上设备侧检查。云场景的形态也支持这一设计：P3/F1 类实例天然把 VM 与专用 GPU/FPGA 打包供给，加速器以 PCIe 设备接入，与移动/桌面端间歇使用的集成 GPU 不同。\par
\subsubsection{关键挑战与核心思想}
在上述背景下，作者提出的关键判断是：CCA 的 granule protection 不能只停在 CPU，必须系统性地扩展到设备侧事务；ACA 的核心做法是把 CCA 用于 CPU 侧保护的安全不变量逐条映射到设备侧。\par
具体而言，论文逐条给出映射：CPU 侧“执行模式决定可否访问某模式内存”映射为“把设备安全地标记为 realm 模式，复用 world 级隔离”；CPU 侧“可信软件控制 MMU 实现 VM 间隔离”映射为“用设备侧 MMU（SMMU）实现设备访问的 VM 级隔离，且保证能编程 SMMU 的不可信 hypervisor 无法篡改该隔离”；CCA 的总线级内存加密映射为复用 PCIe 与设备侧 TEE 的原生硬件加密保护设备链路；VM 身份与加密密钥绑定进 attestation report 映射为把设备 attestation 纳入 VM 创建流程，将设备身份绑定到 realm VM。这组不变量使 VM 与设备执行的机密性、完整性可以被系统性推理，而不是逐点打补丁。\par
\subsubsection{系统架构}
从系统结构看，ACAI 修改 CCA 软件栈的三个组件——可信固件 TF-A、可信 hypervisor RMM、不可信 hypervisor（打补丁的 KVM）——并适配 guest/host Linux 与既有设备驱动，使机密应用像在 legacy 环境一样访问加速器（Figure 2 展示 CCA GPC 接口与 RMM/hypervisor/monitor 关系，Figure 3 展示设备访问 realm 内存的攻击与防护）。\par
具体而言，原型基于 Arm 公开模拟器（支持 CCA 的 FVP），展示 GPU 与 FPGA 两类加速器；兼容性被显式列为设计目标——对可信 RMM、monitor、guest 内核与不可信 hypervisor 的改动都要最小，因为把按层级信任模型写的软件栈改造成可信执行模型本身就容易引入错误。\par
\subsubsection{核心方法设计}
论文的关键不是单一组件，而是把访问模式分析、设备侧检查与 attestation 绑定组合成一条端到端的加速器数据路径。\par
\noindent\textbf{（1）访问模式对比与设计空间界定。}\par
这一机制的直接作用是：先把 integrated、encrypted、ACAI 三条路径的安全/性能权衡讲清楚，说明 bounce buffer 为什么慢、integrated 为什么不安全。\par
具体实现上，论文 Figure 1 与 Table 1 把三种模式按拷贝次数与加密方式逐项对照。Integrated 模式让 PCIe 加速器直接访问 realm 内存，没有额外拷贝，但 CCA 默认把外接加速器视为不可信外设——允许其直接访问将绕过 granule protection checks 的系统性强制，威胁模型不成立。Encrypted 模式先把数据从 realm 拷贝到普通世界 bounce buffer、软件加密后再送设备，返回路径对称，共两次额外拷贝加软件加解密，开销随数据量线性增长。ACAI 模式用 PCIe IDE 类硬件加密与设备侧检查取代软件路径，目标是让“受保护的直接访问”成为默认路径而非例外。这一对比同时解释了为什么 ACAI 把兼容性列为一等目标：bounce buffer 方案之所以流行，恰恰因为它对驱动与应用完全透明，任何替代方案若要求重写软件栈就难以落地。\par
\noindent\textbf{（2）Device-Side GPC 与安全不变量。}\par
这一机制的直接作用是：设备访问 realm 内存前必须像 CPU 一样通过 granule protection 检查。\par
具体实现上，设备被标记为 realm 模式后，其 DMA 与 MMIO 访问受 GPT 约束；SMMU 实现设备侧的 VM 级隔离，且 SMMU 配置路径本身被防护，使可编程 SMMU 的 hypervisor 无法借此篡改隔离；设备不能越权访问其他 realm 或 host 内存。Figure 3 给出设备侧攻击面与对应防护的完整枚举。\par
\noindent\textbf{（3）Attestation、密钥绑定与兼容性。}\par
这一机制的直接作用是：让 verifier 不仅相信 VM 初始状态，也相信设备侧配置与密钥绑定正确。\par
具体实现上，设备 attestation 被纳入 VM 创建流程，设备身份与 device-side 加密密钥绑定到 realm VM 的报告链；链路数据由 PCIe 原生硬件加密保护，避免软件加解密与拷贝；既有加速器驱动与应用尽量不改。原型在 GPU 与 FPGA 两类设备上演示了该流程。\par
\noindent\textbf{（4）兼容性工程与软件栈改造。}\par
这一机制的直接作用是：让安全改造不以重写加速器生态为代价。\par
具体实现上，ACAI 修改的三个可信/半可信组件（TF-A、RMM、patched KVM）与两个内核（realm guest Linux 与 host Linux）都遵循最小改动原则；既有设备驱动与用户应用保持在 legacy 形态下可用，机密应用像在非 TEE 环境一样发起加速器请求。论文把兼容性显式列为与设计安全并列的目标，理由是：把按层级信任模型编写的软件栈改造为可信执行模型，本身就是错误高发区，改动面越大，新引入的漏洞越可能抵消机制收益。\par
\subsubsection{安全性与威胁边界}
ACAI 的威胁模型覆盖：不可信 hypervisor 操纵 I/O 子系统硬件、rogue 设备与被同租 VM 控制的设备发起的越权访问、物理对手窃听/篡改 PCIe 通信。防御由三层叠加而成：world 级设备标记、SMMU 设备侧隔离（配置路径受保护）、总线级硬件加密。边界在于：设备本身与设备侧 TEE 被信任；CCA 生产 CPU 当时不可用，所有性能估计来自模拟器与 Armv8 开发板移植，属原型证据；设备固件可信度、TDISP/SPDM 式标准化设备 attestation 不在其范围。\par
\subsubsection{实验环境与证据}
实验环境为 CCA 模拟器原型加真实加速器工作负载：基于 Arm 公开模拟器、TF-A、RMM 与打补丁的 KVM 修改实现，展示 GPU 与 FPGA 两类加速器；性能数字来自模拟器加 Armv8 开发板移植（论文明示因 CCA CPU 不可用）。本地可核验材料为 20 页 PDF、Figure 1（访问模式）、Table 1（拷贝/加密对比）、Figure 2（GPC 接口）、Figure 3（攻击与防护）与评估章节。\par
\subsubsection{实验结果与证据强度}
论文报告：相对 CPU 与加速器均无 TEE 保护的执行，ACAI 在 GPU 工作负载上平均引入 43.5\% 开销、FPGA 上平均 12.1\%；对系统其余部分（normal world 执行）的影响为 GPU 3.8\%、FPGA 1.9\%，其中 hotspot 最高 11.8\%。评估的方法学本身值得展开：FVP 是指令精确而非周期精确的模拟器，论文以指令数为代理指标，并在 Cortex-A53（1.2GHz）开发板上测量关键操作做投影，作者明确声明这些数字不反映真实 CCA CPU。\par
成本分解（论文第 10 节）揭示了开销来源：生命周期上，系统启动时把 SMMU 数据结构委托到 root world 增加 11.42\%（一次性），realm 创建增加 5.34\%（额外的 RSI/SMC 调用），删除几乎无开销（删除 SMMU stream table 项仅 69 条指令）；运行态上，GPU 与 FPGA 负载的上下文切换分别增加 103.8\% 与 62.2\%（内存分配触发的 RSI 与 timer interrupt 是主因，大内存分配在碎片化时最坏可达每页一次 RSI），RMI 调用增加 7.7\%；微基准上，kernel 到 hypervisor 切换 0.89 微秒、hypervisor 到 monitor 0.77 微秒，经 monitor 的 SMC 加单页写为 2.33 微秒（直接执行 1.75 微秒），MMIO 寄存器写 0.60 微秒（直接 0.13 微秒）；双 GPT 原型中设备 attach 的额外 GPT 操作使 realm 启动指令数增加 99.15\%，但单次 GPT 变更仅 10 条额外指令，GPT Setup 与 Delegate 在开发板上分别增加 20.8\% 与 0.7\% 周期。与 encrypted（bounce buffer）模式的对比（第 10.3 节）是 ACAI 价值的最强证据：encrypted 模式在 GPU 上平均需要 26.8 倍于 ACAI 的指令数，数据传输量最大的 backprop 达 47.8 倍；FPGA 的 SVD 因小数据高频轮询达 3271.7 倍；按开发板实测与 Graviton 数据估算，ACAI 相对 encrypted 模式的加速约 10--19 倍。优化变体 ACAI\_OPT 再带来 GPU 26.05\%、FPGA 3.1\% 的平均加速。\par
这组数字的正确读法是：加速器路径的安全不是免费的，且开销与设备类型、工作负载、数据搬移量强相关——GPU 的 43.5\% 说明早期 device-side 检查与保护路径在高吞吐设备上代价显著；它们来自原型评估，不能外推到所有 PCIe 加速器或未来量产 CCA 平台。\par
\subsubsection{失败模式与边界分析}
ACAI 的边界有四条。第一，原型证据等级：CCA 生产 CPU 不可得，43.5\%/12.1\% 等数字来自模拟器加 Armv8 开发板移植，属可行性证据而非产品预测；论文自身也仅以“demonstrate the feasibility”表述结论。第二，GPU 开销读数：43.5\% 的平均开销说明早期 device-side 检查路径在高吞吐 GPU 上代价显著，它既是 ACAI 的诚实负面结果，也是后续工作（如 CAGE 的 GPT 同步优化）的改进空间——引用该数字时应同时引用其原型前提，避免误读为架构固有的性能上限。第三，信任假设：设备本体与设备侧 TEE 被信任，设备固件的可信度、供应链与 fake device 问题不在防御范围。第四，标准化缺口：设备 attestation 依赖 PCIe 生态的 SPDM/TDISP/IDE 成熟，这些在论文范围外，是商用化必须补齐的层。\par
\subsubsection{综合评价}
综合来看，ACAI 是 CCA 加速器数据路径的奠基论文：优势是问题抓得准——把设备侧访问纳入 CCA 不变量体系，终结 bounce buffer 的默认地位，且把 attestation 绑定与兼容性作为一等设计目标；局限是原型与设备生态仍早期，生产硬件、标准化设备 attestation（TDISP/SPDM）、驱动 TCB 控制均需补齐，GPU 路径 43.5\% 的开销也说明机制本身还有优化空间。从部署角度看，它面向机密 AI/GPU/FPGA 云场景，依赖厂商加速器与 PCIe trusted I/O 生态的成熟；其“设备侧复用 world 级隔离”的思想被后续 CAGE 等工作继承并细化。\par

\subsection{论文 8：Devlore: Device Interrupt Protection for Confidential VMs}
\subsubsection{论文基本信息}
《Devlore: Device Interrupt Protection for Confidential VMs》由 Andrin Bertschi、Supraja Sridhara、Mark Kuhne、Benedict Schlüter、Friederike Groschupp、Clément Thorens、Nicolas Dutly、Srdjan Capkun、Shweta Shinde 完成（ETH Zurich），本地 PDF 为 arXiv:2408.05835v2（2026 年 3 月）；README 证据等级按 preprint 处理。本文将其置于“Arm CCA 的 I/O、DMA、加速器与中断”的研究脉络中考察。它所回答的核心问题可概括为：补上 CCA 设备路径的控制面缺口——恶意 hypervisor 可以伪造或操纵中断，而内存隔离对此完全无效。当前证据状态为：本地 PDF 已保存并核验；下文仅在该范围内讨论其机制与结论。\cite{bertschi2026devlore}\par
\subsubsection{研究背景与动机}
要理解该方案，首先需要看到它面对的系统矛盾：内存隔离与中断隔离是两类不同的问题。设备一旦通过内存设置绑定到某个 CVM，硬件即可在无需软件介入的情况下维持内存隔离；但中断控制器及其状态是 hypervisor 与所有 CVM 共享、复用的同一份资源，设备在运行期持续产生中断，每一次投递都经过 host 控制的簿记。已有工作表明，来自恶意 hypervisor 的中断注入攻击可以破坏 CVM 执行。\par
具体而言，hypervisor 通常管理 GIC 配置、虚拟中断注入与优先级；攻击者可注入伪造的物理/虚拟中断影响 CVM 控制流——例如让 CVM 相信某个设备操作已完成。CVM 需要设备完成通知，但不能信任 host 的中断簿记。论文强调，要强制的是中断隔离：攻击者注入的恶意中断永远不应被投递给 CVM，而合法设备中断必须被正确投递；这不仅涉及中断配置隔离，还涉及运行期投递与确认的隔离。与既有 TrustZone 工作的对比能说明 CCA 场景的新约束：TrustZone 类方案把 GIC 配置标记为 secure world、把中断标记为 secure Group 0 以 trap 到 EL3；但 CCA 的威胁模型不信任 secure world——secure world 软件可以直接编程 GIC 恶意触发中断——因此 Devlore 必须把 GIC 配置标记为 root world，这是它与全部前作在信任锚点上的根本差异。论文同时指出，TDX-Connect 与 SEV-TIO 等 TDISP 兼容实现也用可信固件检查虚拟中断注入，但只允许“全部允许或全部拒绝”设备中断，缺乏 Devlore 这样的细粒度中断隔离。\par
\subsubsection{关键挑战与核心思想}
在上述背景下，作者提出的关键判断是：不必把整个中断管理搬进 TCB；可以让 hypervisor 继续管理中断，但由可信侧记录并检查其行为——delegate-but-check。\par
具体而言，论文的设计探索表明若干“显而易见”的方案并不充分，因为中断隔离需要对整个中断生命周期（注册、物理到虚拟转换、投递、优先级、确认）逐阶段推理，每个阶段都要权衡设计选项。Devlore 的核心思想是最小化改动：把中断管理与投递继续委托给不可信 hypervisor，可信软件只负责记录设备/中断分配关系并检查 hypervisor 是否违规——物理中断与虚拟中断两层都被隔离，且不需要修改任何既有设备驱动。\par
\subsubsection{系统架构}
从系统结构看，Devlore 在 CCA/RMM/GIC 路径上增加中断隔离状态，不改应用与设备驱动（Figure 1 展示 Arm 中断架构与攻击面，Figure 2 展示 Devlore 设计：隔离、物理中断记录、虚拟中断检查，Figure 3 展示 GIC 配置流）。\par
具体而言，Table 1 给出设备生命周期 API，Table 2 给出驱动兼容性结果：对 hypervisor、固件与内核（host 与 guest）的改动合计约 7k 行代码，设备驱动零改动。原型在 Arm FVP（CCA 使能模拟器）上验证功能、正确性与对硬件规范/软件栈的兼容性；因 CCA 生产 CPU 尚不可得，性能估计在 Armv8.2-A 的 RK3588 Rock5b 开发板上回植测量。\par
\subsubsection{核心方法设计}
论文的关键不是单一组件，而是把中断生命周期建模、双层检查与兼容性工程组合成一条完整的控制面防御路径。\par
\noindent\textbf{（1）中断生命周期建模。}\par
这一机制的直接作用是：把“防中断注入”从单点检查扩展为从注册到投递的全过程推理。\par
具体实现上，生命周期被拆为五个阶段：注册（设备与中断号绑定）、物理到虚拟转换（physical IRQ 与 virtual IRQ 的映射被记录）、投递（运行期检查是否为合法设备产生的中断）、优先级、确认。每个阶段的可信侧记录与检查点相互衔接，使攻击者无法在任一阶段注入未被记录的状态。\par
\noindent\textbf{（2）物理与虚拟中断双层检查。}\par
这一机制的直接作用是：攻击既可发生在 physical IRQ 层（设备/路由伪造），也可发生在 virtual IRQ 层（hypervisor 注入），两层都要检查。\par
具体实现上，物理层的关键观察是：集成设备只能拉自己的中断线，无法伪造其他设备的物理中断；hypervisor 伪造物理中断的唯一途径是直接写 GIC。Devlore 因此把 GIC 配置内存标记为 root world——hypervisor 的每次 GIC 访问触发 Granule Protection Fault，trap 处理程序重构出错指令并经 SMC 请求 Monitor 代为执行，Monitor 完成三项检查（地址/长度落在 GIC 配置区、不涉及受保护中断等）后才放行（Figure 3）。物理中断的记录则复用 GIC 既有特性，把所有已注册中断路由到 EL3 Monitor：中断到达后 Monitor 将其写入 realm 侧数据结构（ground truth）、转发给 hypervisor 并恢复原执行上下文（Figure 2.b 第 1--5 步）。虚拟层的检查放在 RMM：CCA 本就不允许 hypervisor 直接编程 CVM 的 vGIC 系统寄存器，hypervisor 只能请求 RMM 注入虚拟中断；RMM 在注入前对照 Monitor 记录的 ground truth，区分“对应真实物理中断的合法虚拟化”与“恶意注入”（第 6--9 步），CVM 进入时 GIC 才触发该中断（第 10 步）。设备注册依赖平台厂商提供的可信映射——（devID, function）二元组在每个平台上唯一，RMM 据此把虚拟中断号绑定到平台特定的物理中断号，且先校验设备确实挂在发起请求的 CVM 上；论文举例 Mali GPU 会产生三个中断（电源管理、页错误、任务完成），说明单设备多功能场景下该映射的必要性。\par
\noindent\textbf{（3）兼容性与最小改动。}\par
这一机制的直接作用是：让防御可以落到真实软件栈上。\par
具体实现上，delegate-but-check 策略使 hypervisor 保留全部管理能力，可信侧只做记录与验证；改动合计约 7k LoC，覆盖四类不同设备的附加场景，应用与设备驱动均无需修改（Table 2）。这与“把中断管理整体移入 TCB”的方案相比，在 TCB 膨胀与生态兼容性之间取了明确的工程平衡点。\par
\noindent\textbf{（4）设计空间探索与被否定的方案。}\par
这一机制的直接作用是：说明 delegate-but-check 不是唯一选项，而是权衡后的结果。\par
具体实现上，论文的设计探索逐一审视了更“直觉”的方案并给出否定理由。物理中断记录的三个候选：直接投递给 CVM 不可行——CVM 运行在虚拟化环境，收到物理中断会被硬件强制退出并 trap 到 EL2 的 RMM；固定一个核专跑 RMM 并配置 GIC 只向该核投递可行但浪费一个核且成为瓶颈；修改 GIC 硬件使其能像对 secure world 一样只向 realm 核投递则需要不可得的硬件与工具，且仍不能解决虚拟中断层问题——Devlore 最终选择复用 GIC 既有的“特定中断仅投递到 EL3”特性，无硬件修改。虚拟中断管理的候选：把全部中断管理移入 RMM 会让 RMM 处理与 hypervisor 调度深度耦合的 timer/IPI 等非注册中断，且 vGIC 系统寄存器数量有限，会造成 RMM 与 hypervisor 的寄存器争用和排队管理复杂度，论文明确判定其“可行但不实用，且破坏 CCA 的 delegate-but-check 策略”——最终选择只让 RMM 做注入前检查。生命周期的五个阶段各有不同的可信边界与成本结构——注册与映射是低频配置路径，可以承担较重验证；投递是高频运行路径，检查必须极轻。Devlore 的最终设计正是按这个成本结构分配检查强度，这也是其低开销的来源。\par
\subsubsection{安全性与威胁边界}
Devlore 的隔离目标是：恶意中断永不投递、合法中断正确投递。威胁模型内是控制 hypervisor 与 host 软件的对手，可操纵 GIC/vGIC 配置与中断注入；威胁模型外包括 DMA 机密性（属 ACAI/CAGE 一类数据面问题）、设备 attestation、TDISP/SPDM、恶意设备固件与侧信道。其安全论证依赖可信侧分配记录的完整性与每个生命周期阶段检查点无遗漏；作为 preprint，尚无形式化证明。\par
\subsubsection{实验环境与证据}
实验环境为双层：CCA-enabled Arm FVP 用于功能、正确性与兼容性验证（四类设备、interrupt stress、glmark2 集成 GPU 案例、驱动兼容性）；RK3588 Rock5b（Armv8.2-A）用于性能估计。本地可核验材料为 21 页 PDF、Figure 1--3、Table 1--3 与评估章节。论文明示 CCA 生产 CPU 不可得，性能为移植估计，属原型证据。\par
\subsubsection{实验结果与证据强度}
论文按五个问题组织评估，数字层次分明。功能层（Q1，FVP）：Table 2 报告 GPU（G76 kbase driver，Midgard test suite）、UART（Amba-pl011，nano/浏览器/IRC 客户端）、PS2 键鼠（Amba-kmi）、LED 与按钮（不产生中断，无需保护）四类设备均无需驱动修改；四项对抗性检查全部通过——良性中断全部正确投递、hypervisor 注入的 1 万个恶意虚拟中断全部被 RMM 拒绝、经 GIC 写伪造物理中断的尝试全部 trap 到 Monitor 并被拦截、SMMU test engine 模拟的恶意设备 MMIO 访问与 hypervisor 的 MMIO 访问均触发 GPT violation。\par
平台成本层（Q2，Rock5b）：整机启动开销 2.4\%（8.9s 增至 9.1s，多 221ms，其中约 80\% 即 177.7ms 来自 SMMU 内存隔离设置，GIC 保护在 BL31 阶段仅贡献 0.05\%/0.1ms）；Linux 内核启动期间共 1816 次 GIC 配置 trap（445 读、1371 写），平均每次 2.5 微秒，相对 4.8 秒的内核启动可忽略；sysbench 非中断负载在 VM 与 CVM 上均无可见影响。\par
中断路径层（Q3/Q4，Rock5b 上的 Mali G610 GPU，驱动栈约 90K LoC）：端到端中断投递加确认时延从基线 669 ticks（27.9 微秒）增至 703 ticks（29.3 微秒），约 5\% 开销；分解显示 Monitor 记录阶段增加 2.2 微秒、RMM 检查阶段增加 5\%（0.4 微秒）、退出确认阶段增加 18\%（1.1 微秒），而 guest 调度阶段反而快 12\%（1.6 微秒）、收尾簿记快 38\%（0.7 微秒）——部分 hypervisor 簿记被前移到 Monitor 阶段。持续压力测试（Q4，GPU Midgard Ktuft 与 UART 内部回环，10k/100k/1M 中断）显示开销稳定在 1\% 以内（标准差低于均值 1.2\%），即单中断成本不随规模累积。\par
真实应用层（Q5）：Rock5b 上 glmark2（2023.01）全部 17 个基准、每个负载产生 2.5k--32k 次 GPU 中断，启用保护后帧率几乎不变，几何平均比值 +0.06\%；FVP 上以未修改的终端浏览器、IRC 客户端、文本编辑器三个 UART 应用做交互式案例。这组数字说明“delegate-but-check”的中断检查可以非常轻，因为检查点主要在配置与记录路径而非每次投递的关键路径上。边界在于：数字只覆盖中断路径，不能外推为完整设备机密 I/O 的开销；性能来自 Armv8 板卡移植估计，非 CCA 生产芯片。\par
\subsubsection{失败模式与边界分析}
Devlore 的边界有四条。第一，覆盖面边界：它只解决中断控制面，DMA 数据面的机密性与完整性完全不在范围——一个同时需要两类防护的系统必须叠加 ACAI/CAGE 类机制，0.06\% 与 1\% 的数字不能相加或直接外推为完整设备机密 I/O 的成本。第二，设备信任假设：恶意设备固件不在威胁模型内，可信侧记录假设设备按分配产生中断；fake device 或固件被篡改的设备需要设备 attestation 层。第三，证据等级：preprint 加 FVP/Armv8 板卡移植估计，无生产 CCA 芯片复现。第四，检查完备性依赖建模：安全性取决于生命周期五阶段检查点无遗漏，该完整性目前是设计论证而非机器证明，新的 GIC 特性或中断虚拟化路径出现时需重新审视。\par
\subsubsection{综合评价}
综合来看，Devlore 把 CCA I/O 讨论从 DMA 数据面扩展到中断控制面，抓住了“共享复用资源无法靠一次性绑定防御”这一本质差异：优势是生命周期建模完整、delegate-but-check 兼顾兼容性与低开销、四类设备与真实 GPU 案例的验证具体；局限是不处理 DMA 机密性、设备 attestation 与恶意设备固件，仍是 preprint 阶段的早期原型。从部署角度看，它适合手机、边缘与云端集成设备场景，落地依赖 GIC/RMM/OS 上游化与硬件实现；其生命周期建模方法可直接迁移到任何“host 代管共享控制面资源”的场景。\par

\subsection{论文 9：Building Confidential Accelerator Computing Environment for Arm CCA}
\subsubsection{论文基本信息}
《Building Confidential Accelerator Computing Environment for Arm CCA》（系统名 CAGE）由 Chenxu Wang、Kun Lu、Fengwei Zhang、Yunjie Deng、Kevin Leach、Jiannong Cao、Zhenyu Ning、Shoumeng Yan、Tao Wei、Zhengyu He 完成（SUSTech、PolyU、Vanderbilt、Hunan University、Ant Group），发表于 IEEE Transactions on Dependable and Secure Computing 23(1)，2026 年；论文注明为 NDSS 2024 会议版的扩展期刊版。本文将其置于“Arm CCA 的 I/O、DMA、加速器与中断”及“加速器机密卸载”的研究脉络中考察，并把它作为同行评审 SOTA（E1）使用。它所回答的核心问题可概括为：把 Realm 的机密边界延伸到加速器工作流——任务元数据、队列状态、MMIO 与 DMA buffer——而不是让不可信 driver 直接处理敏感任务。当前证据状态为：本地 PDF 已保存并核验；下文仅在该范围内讨论其机制与结论。\cite{wang2026cage}\par
\subsubsection{研究背景与动机}
要理解该方案，首先需要看到它面对的系统矛盾：CCA 的 GPC/GPT 能保护内存所有权，但加速器工作流还有大量内存之外的攻击面——代码、任务描述符、metadata、MMIO 寄存器状态、DMA buffer 与完成状态都经过不可信 driver/runtime 之手。论文指出，Arm 官方的 RME Device Assignment（RME-DA）目前仍是高层概念、没有完整硬件实现，CCA 设计把通用加速器视为不可信外设；既有加速器 TEE（HIX 依赖 SGX、StrongBox 信任 secure world 软件）的硬件前提与 CCA 的 realm 式架构不兼容；近期类似 RME-DA 的设计则需要对 hypervisor 做非平凡修改并把重型加速器软件纳入 realm TCB。\par
具体而言，威胁模型假设攻击者控制 host OS、hypervisor、加速器 driver/runtime 与外设软件栈，可以伪造任务元数据与地址；可信的是 CCA 硬件原语、Monitor/RMM 逻辑与真实加速器硬件本体。物理攻击、侧信道、rollback 被标为 out-of-scope 或需正交机制。GPU 与 FPGA 的工作流差异（unified memory 对 dedicated memory、任务队列对 XDMA 描述符）意味着不存在通用的拷贝包装器，保护工作流必须按设备类型重画。在相关工作坐标系中，论文区分了三类前作：非 Arm GPU TEE（HIX、GEVisor 依赖 SGX，Honeycomb 依赖 SEV-SNP，CURE、HETEE 依赖总线控制器）难以迁移到 Arm；Arm 端侧 GPU TEE（StrongBox 等）依赖 TrustZone 与传统安全虚拟化，尚未覆盖 CCA 的新对手模型；直接在 GPU 硬件内建 TEE 的路线（Graviton、NVIDIA H100）则严重影响硬件兼容性。CAGE 的位置是第四条路：复用 CCA 既有原语、不改硬件、把驱动留在 TCB 外。\par
\subsubsection{关键挑战与核心思想}
在上述背景下，作者提出的关键判断是：让不可信软件栈继续做复杂工程（内存分配、任务调度），但在 Monitor 内重建并验证一个可信的 shadow task——不可信栈创建与真实加速器任务对应的 stub，Monitor 验证并同步这些操作到真实任务。\par
具体而言，论文命名了三个挑战。C1：host 侧加速器软件并非为管理 realm 内应用设计，频繁交互必然引入与 hypervisor 通信和 world 切换的开销——对策是 shadow task 机制，且该机制与具体加速器硬件架构解耦。C2：在缺少 RME-DA 支持的情况下，通用 CCA 设计不允许加速器访问 realm——对策是 two-way isolation，复用 GPC 同时保护 CPU 与外设侧，为每个 realm 提供不同的加速器内存视图，并保护加速器执行环境本身。C3：基于 GPC 的保护需要为不同 GPC 创建多份 GPT 并在保护环境时同步，产生额外开销——对策是 GPT 维护优化，重点是同步不可信组件的 GPT 与为 realm 初始化加速器 GPT。\par
\subsubsection{系统架构}
从系统结构看，CAGE 在 Realm 与加速器软件栈之间插入 Monitor 侧验证与 GPT/GPC 保护：安全模块部署在 root world 的 Monitor 中，防御不可信软件（OS、hypervisor、secure world 组件）与外设。\par
具体而言，不可信 driver 发起 GPU/FPGA 任务；Monitor 从 stub 中提取安全关键字段构建 shadow task，验证元数据、地址与状态转换后同步到真实任务；受保护 buffer 经加速器 GPT/GPC 管理，host 无法访问敏感数据；任务完成后 Monitor 恢复并清理环境（MMIO/寄存器状态、专用内存）。GPU 侧处理 unified memory 与 GPU MMIO/寄存器状态，FPGA 侧处理 dedicated memory、XDMA 元数据与 PCIe 工作流。\par
\subsubsection{核心方法设计}
论文的关键不是单一组件，而是 shadow task、双向隔离、GPT 优化与按设备类型适配的组合。\par
\noindent\textbf{（1）Shadow Task。}\par
这一机制的直接作用是：在可信 Monitor 内保存可验证的加速器任务镜像，避免把庞大 GPU/FPGA driver 纳入 TCB。\par
具体实现上，真实任务仍由不可信栈操作（它创建与真实任务对应的 stub）；Monitor 根据 shadow task 检查 metadata 与地址，验证通过后才同步并启动加速器。数字对比说明了这一抽象的价值：CAGE 的 GPU 安全模块约 1301 行代码、FPGA 扩展额外约 140 行，而作为对照的 Midgard driver 约 3 万行——把 driver 留在 TCB 外是 TCB 控制的决定性因素。\par
\noindent\textbf{（2）GPC/GPT 双向 buffer 保护。}\par
这一机制的直接作用是：用 CCA 既有 GPC 同时约束 CPU 与外设对敏感 buffer 的访问。\par
具体实现上，敏感 buffer 被切换到 realm/加速器可用状态；Monitor 为对应 GPC 创建并同步多份 GPT；错误地址或恶意 metadata 的后果被设计为仅导致 DoS，而不泄露或篡改秘密——这是“fail to denial, never to disclosure”的明确工程选择。\par
\noindent\textbf{（3）GPU 与 FPGA 工作流的分型适配。}\par
这一机制的直接作用是：承认两类设备的保护对象不同，保护模块也随之不同。\par
具体实现上，GPU（Arm Mali，unified memory）侧重任务/队列/buffer 元数据与 MMIO 寄存器状态检查；FPGA（Xilinx VCU118，dedicated memory）侧重 bitstream/任务元数据与 XDMA 类路径及专用内存清理。论文明确指出，扩展到 DPU/SmartNIC 时同样需要重画工作流特定的状态——shadow task 的内容由设备工作流决定，不存在一次写就的通用版本。\par
\noindent\textbf{（4）GPT 维护优化。}\par
这一机制的直接作用是：把 C3 中多 GPT 同步的成本压下去。\par
具体实现上，优化集中在不可信组件 GPT 的同步与 realm 加速器 GPT 的初始化；论文报告优化缓解了 84.63\%--96.55\% 的 GPT 同步开销。这一数字解释了 CAGE 最终开销为何能落在 moderate 区间——未优化的 naive 多 GPT 方案在该路径上代价要高出一个量级。\par
\noindent\textbf{（5）安全分析的攻击清单与防御映射。}\par
这一机制的直接作用是：把“安全”从口号落成可逐条核查的攻击—防御表。\par
具体实现上，论文安全分析覆盖的攻击包括：未授权内存访问（由 GPC/GPT 拦截）、恶意 metadata 与恶意任务提交（由 shadow task 验证拦截）、GPU 页表操纵（由 Monitor 对页表状态的检查拦截）、fake accelerator（设计目标为仅导致 DoS）、SMMU/GPC 绕过尝试（由 Monitor 对 MMIO/寄存器状态的检查与环境清理拦截）、FPGA dedicated-memory 泄漏（由任务结束后的专用内存清理拦截）。每条攻击都映射到具体防御组件，使读者可以审计“哪个组件失效会打开哪条攻击路径”——这种映射关系正是后续做 Monitor 形式化验证时的规格雏形。\par
\subsubsection{安全性与威胁边界}
CAGE 的安全分析覆盖：未授权内存访问、恶意 metadata、GPU 页表操纵、恶意任务提交、fake accelerator、SMMU/GPC 绕过与 FPGA dedicated-memory 泄漏；防御来自 Monitor 对 metadata/代码签名的检查、GPC 对 CPU/GPU/外设访问的限制、MMIO/寄存器状态检查与环境清理。强假设是加速器硬件可信、物理攻击排除、侧信道与 rollback 需正交方案；安全分析系统性较强，但没有对完整状态机的形式化证明，Monitor 中加入的加速器特定逻辑也增加了 TCB。\par
\subsubsection{实验环境与证据}
实现基于 Arm 官方双平台：软件模拟 CCA 特性的 FVP（Base RevC，Linux 5.3.0，TF-A 2.8 Monitor 修改）与 Juno R2 开发板（Mali-T624 GPU 加 PCIe 连接的 Xilinx Virtex Ultrascale+ VCU118 FPGA）。TCB 增量为 GPU 1301 LoC 加 FPGA 140 LoC。GPU 评估使用六个 Rodinia benchmark 与神经网络推理模型；FPGA 评估使用 FFT、矩阵乘等六个 C benchmark 与 Xilinx 工作流。实验覆盖功能验证、安全分析、GPT 优化消融与两类加速器。局限：硬件环境仍非最终 CCA 生产平台，GPU 类型偏 Arm 端侧 unified-memory GPU，不能直接代表数据中心离散 GPU/DPU。\par
\subsubsection{实验结果与证据强度}
论文报告：GPU Rodinia benchmark 开销 0.58\%--5.31\%；FPGA benchmark 开销 9.61\%--16.30\%；GPT 维护优化缓解了 84.63\%--96.55\% 的 GPT 同步开销。对比 ACAI 报告的 GPU 平均 43.5\%，CAGE 在 GPU 路径上的开销低了约一个数量级——这一差距主要来自两方面：CAGE 面向 unified-memory 端侧 GPU 且做了 GPT 同步优化，ACAI 面向 PCIe 外接加速器且包含设备侧检查与链路保护的完整路径；两者数字不可直接互相归一，但共同说明加速器机密计算的开销高度依赖设备形态与数据搬移路径。FPGA 侧 9.61\%--16.30\% 的开销则提醒：dedicated-memory 设备的元数据验证与清理路径代价显著高于 unified-memory GPU。\par
\subsubsection{失败模式与边界分析}
CAGE 的边界有五条。第一，设备身份：fake accelerator 与真实 device identity 未被完整解决，论文显式将 SPDM/TDISP/IDE 列为商业化必须补齐的层；当前设计的底线是“伪造设备仅导致 DoS”。第二，硬件代表性：GPU 评估基于 Mali-T624 这类端侧 unified-memory GPU，不能直接代表数据中心离散 GPU/DPU；FPGA 为 PCIe 连接的 VCU118，两类结论向其他设备形态迁移时需重新评估。第三，Monitor TCB 增长：加速器特定逻辑进入 Monitor 增加了高特权代码，且按设备类型增长（GPU 1301 LoC、FPGA 140 LoC），论文未给出形式化最小化方案。第四，多租户与侧信道：加速器时分复用、跨租户状态残留与调度侧信道均未覆盖。第五，无完整状态机证明：安全分析系统但非形式化，shadow task 状态机的 TOCTTOU 风险需后续机器验证。\par
\subsubsection{设计留白与对架构的反馈}
CAGE 的讨论章节罕见地系统交代了设计留白，这些留白本身就是研究方向。第一，代码机密性被显式放弃：CAGE 保护敏感数据但不保护代码，理由有三——许多加速器应用使用通用代码二进制（如通用深度学习模型），敏感的是输入参数；明文代码有助于加速器软件完成功能且不泄露敏感数据；代码机密性会引入额外加解密开销。第二，SMMU 与 GPC 的结构性缝隙：传统机制用 SMMU 的 StreamID 区分各外设的访问视图，但最新 GPC 只为 SMMU 提供单个 MMIO 配置区域，所有连接的外设共享同一 GPC 视图；CAGE 建议 Arm 为未来设备引入 StreamID 级 SMMU GPC 或允许每个外设拥有独立 Root 控制页。第三，FPGA 特有的任务管理攻击：对手可提供错误的 channel-buffer 关系（把数据喂进错误的 H2C/C2H 通道），防御要求用户在数据 buffer 描述中提供详细通道绑定并由 Monitor 验证完整性。第四，物理与正交攻击的嫁接方案：JTAG/I2C 等物理接口攻击可借助检测方案，PCIe/AXI/NVLink 链路与内存攻击可借助 PCIe IDE 与 CCA Memory Protection Engine，GPU 内存 rollback 可用单调计数器。第五，与 RME-DA 的关系：论文指出 RME-DA 是 Arm 对 TDISP 的实现，需要修改 SMMU 硬件；ACAI 需要修改 hypervisor 中 SMMU 相关代码；两者都要求把重型加速器软件放进 realm 以增加 TCB，而 CAGE 的 shadow task 把这部分留在 TCB 外——这是三者的根本设计分歧。\par
\subsubsection{综合评价}
综合来看，CAGE 是当前 Arm CCA 加速器工作流方向最完整的同行评审 SOTA：优势是机制贴合 CCA 原语、shadow task 抽象清晰、GPU/FPGA 双设备评估、TCB 增量极小、GPT 优化有量化消融；局限是 device identity 与 fake accelerator 的完整防御需要 SPDM/TDISP 层、生产 CCA 硬件缺失、多租户加速器调度与侧信道未覆盖、Monitor 内加速器特定逻辑有待形式化与最小化。从部署角度看，其 shadow task 方法可迁移到 DPU/SmartNIC 的队列描述符与报文处理任务，但必须补齐设备身份层；论文与 ACAI、Devlore 共同勾勒出 CCA 设备子线的三块拼图——数据路径（ACAI）、控制路径（Devlore）、任务工作流（CAGE）。\par

\subsection{开发商安全实践建议}
本章九篇材料共同指向一个工程事实：Arm CCA 的安全不是单个硬件特性，而是 granule 所有权状态机、RMM 软件 TCB、部署粒度选择与设备路径防护的组合。以下建议面向芯片厂商、固件与系统软件开发商、设备制造商（含网络设备厂商），每条均回扣本章具体论文的结论。\par
\noindent\textbf{实践一：在设备平台上按 CCA/RMM 的对象模型规划信任边界，而不是复用 TrustZone 时代的二分模型。}Li 2022 证明 Realm world、Root world 与 GPT granule 所有权可以把 host 降级为“不可读的资源管理者”\cite{li2022cca}，CCA 与 RMM 规范则把 Realm/REC/granule 的术语与生命周期语义固定为官方架构语言\cite{arm_cca_spec,arm_rmm_spec}。设备厂商在规划新产品线时，应以 RMM 状态机（创建—委托—运行—回收—验证）为信任边界的设计起点：凡需要 host 参与的管理操作都应对应明确的 RMI 前置条件，凡需要租户信任的服务都应经 RSI 而非 host 转述。对暂无 CCA 芯片的产品，应先用该对象模型审查现有 TrustZone 方案中 trusted OS 的 TCB 暴露面。\par
\noindent\textbf{实践二：把 granule 生命周期审计作为固件验收的必选项。}Li 2022 的验证经验表明，安全性来自合法状态转换序列而非单个接口正确\cite{li2022cca}；RContainer 的 30 个 CVE 深度攻击模拟显示，漏洞利用的共同前提是“内存隔离弱加权限不受限”，而 GPT fault 能在攻击路径上统一截停\cite{zhou2025rcontainer}。固件开发商应对 delegate/undelegate、GPT 更新与销毁 scrub 路径建立审计清单：每次 granule 状态转换可追溯到合法的 RMI/Monitor 调用链；销毁路径必须包含私有内存清理（RContainer 的 con-shim termination 约 637 微秒即此成本的原型参照）；GPT 操作保持 EL3 独占。条件允许时，应对 RMM 类组件做模块化形式化验证，Li 2022 的 43 层抽象与约 200 行可信 Coq 规范证明了这在工程上可行。\par
\noindent\textbf{实践三：按负载粒度选择隔离方案，避免“一个 Realm VM 包打天下”。}Shelter 证明用户态 SApp 可用 multi-GPT 以低于 15\% 的开销获得硬件隔离\cite{zhang2023shelter}；RContainer 证明容器场景下 mini-OS 加 Shim-GPT 可在 Apache/Nginx 上压到平均约 5\%、MySQL/Memcached 低于 0.5\%，且 EL3 新增代码比 EL3 集中式方案低一个数量级\cite{zhou2025rcontainer}；NanoZone 则表明进程内数据域隔离需要 POE/PIE 级别的快路径，特权级切换的频率必须压到个位数百分比才可用\cite{liu2025nanozone}。系统软件开发商应为不同粒度提供分层 SDK：单应用敏感逻辑用 SApp 式用户态隔离，多租户容器用容器级隔离运行时，进程内秘密（session key、租户私数据）预留 POE/PIE 式域切换接口；同时把 EL3 新增代码量作为硬指标纳入设计评审。\par
\noindent\textbf{实践四：加速器与设备路径必须数据面、控制面、工作流三面同时设防。}ACAI 报告 device-side 保护在 GPU 上平均 43.5\%、FPGA 上 12.1\% 的原型开销\cite{acai2023}；Devlore 报告中断检查可在 glmark2 真实负载下压到 0.06\%\cite{bertschi2026devlore}；CAGE 报告 GPU 0.58\%--5.31\%、FPGA 9.61\%--16.30\%，且 GPT 同步优化可消除 84.63\%--96.55\% 的同步开销\cite{wang2026cage}。三组数字说明：设备路径的开销高度依赖设计点，早做架构决策的收益巨大。设备与平台厂商应：(1) 在设备接入流程中规划 device-side GPC/SMMU 隔离与链路加密（PCIe IDE 类硬件加密优先于软件 bounce buffer）；(2) 对中断采用 delegate-but-check 式设计，保留 hypervisor 管理能力但由可信侧记录分配并检查投递；(3) 对加速器任务采用 shadow task 式元数据验证，把 driver 留在 TCB 外；(4) 预研 GPT 同步优化，避免 naive 多 GPT 方案在生产负载上失控。\par
\noindent\textbf{实践五：把设备身份与 attestation 绑定作为商用化的前置条件，而非后置补丁。}RMM 生命周期产生的 measurement 是 verifier 决策的唯一依据\cite{arm_rmm_spec}；ACAI 与 CAGE 都指出设备身份、fake accelerator 与 SMMU 路径需要 SPDM/TDISP/IDE 类机制补齐\cite{acai2023,wang2026cage}。平台厂商应在产品路线图中把设备 attestation 纳入 Realm 创建流程：设备身份与 device-side 密钥必须在 attestation 报告链中可见；引入加速器、DPU 或 SmartNIC 前，先确认其 TDISP/SPDM 支持状态；对尚无标准支持的设备类型，以 CAGE 的“fake accelerator 仅导致 DoS 而非泄露”为底线设计要求。\par

\subsection{未来研究方向}
本章证据链在架构层（Li 2022 的设计与验证）与原型层（六篇系统论文）之间仍然断裂：所有性能数字来自 FVP、模拟器或 Armv8 开发板回植，所有规范细节依赖无本地 PDF 的官方页面。以下五个方向按“问题—为什么未解决—技术路线—需要的验证实验”展开，每个都对应本章的具体证据缺口。\par
\noindent\textbf{（1）量产 CCA 硬件上的全栈性能复现与方法学。}问题：本章全部性能结论——Shelter 的 $<$15\%\cite{zhang2023shelter}、RContainer 的 Fig.5/Fig.7 数字\cite{zhou2025rcontainer}、NanoZone 的 4.87\%/22.67\%\cite{liu2025nanozone}、ACAI 的 43.5\%/12.1\%\cite{acai2023}、Devlore 的 0.06\%\cite{bertschi2026devlore}、CAGE 的 0.58\%--5.31\%/9.61\%--16.30\%\cite{wang2026cage}——都来自非生产平台，RContainer 作者已明示模拟环境与真实 ARMv9-A 机器可能存在偏差。为什么未解决：CCA 量产芯片长期不可得，各论文只能各自搭建不可比的模拟环境。技术路线：在首批 CCA silicon 上建立统一基准套件，覆盖 granule 转换、RMI/RSI 调用、world switch、GPT 同步、设备路径五类微基准与本章六类应用负载。需要的实验：同一负载在 FVP 与真实芯片上的双测量，量化模拟偏差的方向与幅度，形成可引用的校准系数。\par
\noindent\textbf{（2）RMM 及 Monitor 扩展组件的形式化验证复用。}问题：Li 2022 验证了原型固件，但 Shelter 的 EL3 Monitor、RContainer 的 mini-OS、NanoZone 的 root-world 模块、CAGE 的 shadow task 逻辑都在向 TCB 添加未验证代码\cite{li2022cca,zhang2023shelter,zhou2025rcontainer,liu2025nanozone,wang2026cage}。为什么未解决：VIA 的验证面向特定代码基，迁移成本高；新增组件含弱内存并发与 C/汇编混合，恰是验证最难的部分。技术路线：把 VIA 的 mover oracle、permutation conditions 与 register accounting 工具化，以 RMM 状态机为共享规范层，让各扩展组件以插件形式证明“不破坏 granule 不变量”。需要的实验：对至少一个扩展组件（如 shadow task 状态机）完成无 TOCTTOU 的机器检查证明，并测量验证成本相对 Li 2022 的变化。\par
\noindent\textbf{（3）设备身份、attestation 与 CCA 生命周期的端到端绑定。}问题：ACAI、Devlore、CAGE 都假设设备本体可信或把 SPDM/TDISP 留作未来工作\cite{acai2023,bertschi2026devlore,wang2026cage}，而 RMM 规范的公开概念中设备生命周期（RME-DA）仍是高层概念\cite{arm_rmm_spec,arm_cca_spec}。为什么未解决：设备 attestation 横跨 PCIe 标准、设备固件与 CCA 软件栈三个生态，单一论文无法闭环。技术路线：把设备身份凭证纳入 Realm 创建 measurement 链，设计设备加入/移除时 granule 与中断分配状态的联动状态机。需要的实验：fake accelerator、设备热替换、重放旧凭证三类攻击下的端到端测试，验证 verifier 能否在 secret 释放前发现设备层异常。\par
\noindent\textbf{（4）多租户加速器共享与调度的机密性。}问题：本章加速器工作均按“设备专属单 Realm”建模，CAGE 明确指出多租户 accelerator scheduling 未覆盖\cite{wang2026cage}，ACAI 也按云实例打包形态假设设备专用\cite{acai2023}；而真实 GPU/DPU 云需要时分复用。为什么未解决：复用引入跨租户状态残留、调度侧信道与 shadow task 生命周期交叉三个新问题。技术路线：扩展 shadow task 为 per-tenant 状态机，定义加速器上下文切换时的 GPT/MMIO 清理不变量，对调度元数据做 declassification 审计。需要的实验：双租户交替负载下的残留状态扫描、调度时序侧信道测量与性能开销分解。\par
\noindent\textbf{（5）规范细节的可公开核验与版本追踪。}问题：本调研中 CCA 与 RMM 规范的本地 PDF 均不可得，所有规范级论断只能停在公开概念层\cite{arm_cca_spec,arm_rmm_spec}；DEN0125/DEN0137 的版本演进（如 RME-DA 从概念到实现的落地）缺乏可引用的稳定文本。为什么未解决：官方文档的分发方式限制了自动获取与存档。技术路线：建立规范快照的可核验存档流程（下载状态、哈希、版本号记录），并把每条架构语义论断绑定到具体版本。需要的实验：跨版本的概念差异对照表，检验本章各论文的设计假设在最新规范版本下是否仍然成立——例如 Shelter 的 per-core multi-GPT 与 RContainer 的 GPT-routing 所依赖的 GPT 配置语义是否被规范明确支持\cite{zhang2023shelter,zhou2025rcontainer}。\par
\section{RISC-V enclave、CoVE 与 AP-TEE 机制}
RISC-V 的开放性使 enclave 原型、监控器实现和 ISA 扩展能够并行演进，但也意味着保护边界不能仅靠某个实现约定维持：同一时期内，社区既在产出 Keystone、Penglai、SPEAR-V 这样各自修改硬件或监控器的研究原型，也在推进 IOMMU、AIA 这样已批准（ratified）的基础规范，以及 CoVE、AP-TEE、CoVE-IO 这样仍处于草案状态的机密虚拟机标准。本节以早期 enclave 系统为起点，沿"基础 I/O 规范—enclave 系统—机密虚拟机架构—机密 I/O 草案"的顺序展开，重点辨析每条路线把哪些责任交给硬件、可信监控器、宿主软件与设备。\par
阅读这些工作时，需要区分三类证据强度截然不同的材料：系统论文（Keystone、Penglai、SPEAR-V）给出原型与实测数字；已批准规范（IOMMU、AIA）给出稳定语义但无实验；架构提案与草案（CoVE、AP-TEE v0.7、CoVE-IO v0.3.0）给出生命周期与 ABI 设计，但其接口和状态机仍可能变化。只有把每篇材料的责任划分映射到启动、运行、共享和销毁的全过程，并把每个数字放回其原始评估环境，才能判断开放 ISA 的可裁剪性是否真正转化为可审计的安全性。\par
\subsection{论文 10：A Survey of RISC-V Secure Enclaves and Trusted Execution Environments}
\subsubsection{论文基本信息}
《A Survey of RISC-V Secure Enclaves and Trusted Execution Environments》由 Marouene Boubakri and Belhassen Zouari 完成，发表于 2025 年的 Electronics 14(21):4171，本地 PDF 与 MDPI HTML 均已核验。本文将其置于"RISC-V CoVE / AP-TEE 机密虚拟机、硬件辅助可信执行环境分类"的研究脉络中考察，并把它作为综述类辅助材料使用。它所回答的核心问题是：当 RISC-V TEE 的论文、开源项目与规范在同一时期快速涌现时，如何为它们建立一张不混淆抽象层级的谱系地图。需要在一开始就明确：该综述的作用是画地图，而不是替代原始系统论文；其证据等级为同行评议综述，机制与性能论断必须回引原始论文或官方规范。\cite{boubakri2025riscvtee}\par
\subsubsection{研究背景与动机}
综述面对的系统矛盾是"同名不同物"：同样叫 TEE，Keystone、Penglai、SPEAR-V 保护的是进程或地址空间级 enclave，CoVE 与 AP-TEE 保护的是整台虚拟机，而 IOMMU、AIA 这类规范根本不提供 TEE 本体，只提供被 TEE 复用的 I/O 机制。论文第 1 节指出，RISC-V 的开放、可扩展与模块化特性推动了机密计算基础研究，但已有方案在隔离强度、可扩展性、硬件改动量、性能开销、编程模型与真实验证程度上取舍各异，社区缺少一个整合视图。\par
具体而言，综述在第 2.3 节给出四个纳入标准——架构意义、实现保真度、设计目标多样性与社区/工业影响——覆盖 2016 至 2025 年的工作；Figure 1 把 Sanctum、MI6、Keystone、CURE、SPEAR-V、CoVE、WorldGuard、AnyTEE、TIMBER-V、Elasticlave、VirTEE、AP-TEE、DORAMI、Penglai 放在同一时间轴上。这一时间轴本身就是动机论证：2016 年前后的工作围绕 PMP 与监控器，2021 年前后转向可扩展内存保护，2023 年之后重心明显移向机密虚拟机与标准化。\par
\subsubsection{关键挑战与核心思想}
在上述背景下，综述提出的关键判断是：RISC-V TEE 之间的差异不是细节差异，而是四个设计轴上的系统性分歧——信任放置（trust placement）、内存原语（memory primitive）、证明（attestation）与 I/O 支持。\par
具体而言，信任可以放在 M-mode 监控器（Keystone、Penglai）、HS/S-mode 的 TSM（CoVE/AP-TEE）、硬件 tag（SPEAR-V、TIMBER-V）或 SoC 可信根上；内存原语可以是 PMP/ePMP、Guarded Page Table、Mountable Merkle Tree、Memory Tracking Table 或 DRAM tag store。综述的价值在于把这些轴合并成 Table 1（p.19）的比较矩阵，使"哪个系统解决了哪个问题"变成可逐项核对的事实，而不是叙述印象。\par
\subsubsection{系统架构}
综述不给出单一系统图，而是给出一张从 ISA 原语到 TEE 系统的分层视图（第 2.2 节）。\par
具体而言，底层是特权级架构（U/S/M 与 hypervisor 扩展）、PMP/ePMP、MMU、IOMMU 与 AIA；中层是安全监控器、TSM、enclave runtime 与内存元数据（GPT/MMT/MTT/tag store）；上层是 enclave 应用、serverless 函数、TVM 工作负载与远端验证者。这一分层视图的实际用途是定位证据：当一篇系统论文声称"隔离"时，应追问隔离发生在哪一层、由哪个组件执行。\par
\subsubsection{核心方法设计}
综述的方法不是提出新组件，而是把文献覆盖、分类矩阵与开放挑战组合成一条可复用的分析路径。\par
\noindent\textbf{（1）原语栈梳理。}\par
这一机制的直接作用是把 RISC-V 安全原语当作系统设计的底座来比较。具体实现上，论文第 2.2 节逐一说明特权级、PMP/ePMP、虚拟内存、hypervisor 扩展、密码指令与内存标记各自决定什么：特权级决定谁能管理隔离，PMP/MMU/MTT/tag store 决定内存边界，IOMMU/AIA/CoVE-IO 决定 I/O 方向——而 I/O 方向在多数被综述系统中仍不成熟。\par
\noindent\textbf{（2）谱系与矩阵比较。}\par
这一机制的直接作用是把代表系统放在同一条演进线与同一张表上。具体实现上，Table 1（p.19）按隔离架构、enclave 类型、特权级、内存隔离、secure I/O、cache 侧信道防护、硬件改动、实现验证、SDK、TCB 规模、密码支持与合规性横向比较；Sanctum/Keystone 落在开放 enclave 分支，Penglai/CURE/SPEAR-V 落在不同内存原语分支，CoVE/AP-TEE 落在机密虚拟机分支。需要强调，Table 1 把不同威胁模型、不同验证平台、不同成熟度的系统放进同一矩阵，读者不应把表格共现误解为同等级安全比较。\par
\noindent\textbf{（3）开放挑战归纳。}\par
这一机制的直接作用是支撑"为什么还要继续研究"的问题意识。具体实现上，第 4 节（p.20--29）指出标准化、硬件采用、形式化验证、侧信道边界与 I/O/设备信任均未完全解决；Table 1 的 secure I/O 列显示多数系统只有部分支持或不支持；第 4.16 节进一步把 SDK、工具链、标准符合性、ASIC 量产验证与生态整合列为采用条件。换言之，综述自身就声明：单项机制成立不等于平台可商用。\par
\subsubsection{安全性与威胁边界}
综述不提供新的安全证明，也不复现攻击实验；其安全分析全部来自对被综述系统威胁模型与机制覆盖的比较。可以支撑的结论是：RISC-V TEE 确实分裂为 PMP/监控器轻量 enclave、硬件修改型 enclave、可扩展内存保护与机密虚拟机规范四条路线，且多数系统对 secure I/O、cache 侧信道与合规性的支持不完整（Table 1 与第 4 节）。不能支撑的结论是：不能仅凭本综述证明任何单个系统的机制安全性，也不能把威胁模型、攻击者能力与物理/侧信道假设各不相同的系统合并到同一安全等级上。\par
\subsubsection{实验环境与证据}
综述无新系统实验，证据来自文献覆盖与分类一致性。从可核验的材料看，本地 PDF 与 MDPI HTML 已核验；能够直接支撑的结论是谱系、分类与挑战归纳，具体载体为 Figure 1 时间线、Table 1 比较矩阵与第 4 节讨论。这些材料不足以支撑任何单一系统的真实性能数字、安全证明或 ABI 细节。\par
本节判断以 Boubakri 综述 PDF/HTML 及本地 README 核验块为依据；超出这些材料覆盖范围的实现细节、性能结论或攻击面，本文不作延伸推断。\par
\subsubsection{实验结果与证据强度}
综述能够为本文提供的"结果"是结构性的，而非数值性的：它证明 RISC-V TEE 设计空间可以沿四个设计轴被系统化，并证明 secure I/O 与量产验证是共同弱项。\par
具体而言，Table 1 引用了若干系统的 TCB 规模（如 Keystone 约 8.4 k LoC、Sanctum 约 5 k LoC、Penglai 约 10.2 k LoC 等），但综述 README 明确标注这些是二手汇总，口径与原始论文不完全一致（例如 Keystone 原文按 SM 增量约 1.6 KLoC、含 runtime 最大约 15 KLoC 的口径统计），引用时必须回到原始论文核验。本节后续条目引用性能数字时一律回到 Keystone、Penglai、SPEAR-V 原文，正是遵循这一证据纪律。\par
\subsubsection{综合评价}
综合来看，Boubakri 综述是合格的谱系辅助材料：其优势在于新（覆盖至 2025 年，含 AP-TEE 草案）、维度全（Table 1 的比较轴足够细）、问题意识清楚（第 4 节把标准化与量产验证单独成章）。其局限在于综述粒度天然有限，且 AP-TEE、CoVE-IO、IOMMU、IOPMP 等规范状态在其截稿后仍可能变化，比较矩阵需要持续维护。从使用角度看，它适合作为路线图与差距分析工具；任何具体产品决策仍须回到标准原文与系统实测。\par
\subsection{论文 11：RISC-V IOMMU Architecture Specification}
\subsubsection{论文基本信息}
《RISC-V IOMMU Architecture Specification》由 RISC-V IOMMU Task Group 完成，本地保存的是 RISC-V Ratified Specifications Library v1.0.1 / 20260222 发布版，基础架构与列出的标准扩展均已批准（ratified），PDF 约 109 页。本文将其置于"RISC-V CoVE-IO / TEE-I/O、RISC-V 基础安全机制"的研究脉络中考察，并把它作为规范类材料使用。它所回答的核心问题是：设备 DMA 不走 CPU 页表，那么机密内存的 CPU 侧保护如何延伸到设备侧。规范给出的答案是 device context、process context、两级地址转换、MSI 翻译与 fault/queue 接口——它是 I/O 隔离的必要底座，但不是可信 I/O 的充分条件。\cite{riscv_iommu_2023}\par
\subsubsection{研究背景与动机}
规范面对的系统矛盾是：CPU 侧的 enclave 或 TVM 隔离做得再完整，设备仍可通过 DMA 直接读写任意物理内存，绕过全部页表与 PMP 检查。\par
具体而言，虚拟化场景需要把设备安全直通给某个 VM，要求设备只能访问该 guest 的地址空间；机密虚拟机场景进一步要求 IOMMU 的配置权本身受到约束——host 作为不可信的资源管理者可以发起配置，但 device context 与 TVM 内存所有权的绑定必须由 TSM 级别的可信组件裁决。规范在动机层面只解决前一个问题，后一个问题正是 CoVE-IO 存在的理由；IOMMU 解决地址，不解决设备是否可信。\par
\subsubsection{关键挑战与核心思想}
规范的关键判断是：IOMMU 是 trusted I/O 的地址层，信任层必须在其上另建。\par
具体而言，它管地址转换与权限检查，不自动证明设备身份；它管 DMA，不自动保护 PCIe 链路的机密性与新鲜性；它提供 fault 报告，但不定义 fault 信息对 TVM 的泄露边界。因此在本节的叙事中，IOMMU 的定位是 CoVE-IO 的底层积木之一，与 sIOPMP/IOPMP 的物理访问控制互补：IOPMP 做物理区间放行，IOMMU 做地址翻译与设备上下文管理。\par
\subsubsection{系统架构}
规范的组织方式是：每个设备请求先经 device directory table（DDT）定位 device context，context 指向翻译根（可选 process directory table 与第一/第二级页表），翻译结果经权限检查后形成主机物理地址；运行期管理由 command queue、fault queue 与 page-request queue 支撑，软件通过 MMIO 寄存器与 IOMMU 交互。\par
具体而言，规范 Figure 1 展示非虚拟化场景下用 IOMMU 隔离设备，Figure 2 展示把设备直接直通给 VM 的配置；Table 1 给出术语表。核心对象还包括 IOHGATP（第二级翻译根）、MSI page table（把设备 MSI 写入重映射到指定 guest interrupt file）以及 ATS/PRI（地址翻译服务与页请求接口）处理逻辑。\par
\subsubsection{核心方法设计}
规范的关键不是单一组件，而是把上下文查找、两级翻译、MSI 重映射与失效/故障语义组合成一条受控的 DMA 路径。\par
\noindent\textbf{（1）Device Context 定位。}\par
这一机制的直接作用是：每个 DMA 请求先按 requester/source ID 定位 device context，决定走哪个翻译根。具体实现上，DDT 是隔离的入口——context 配错会让设备进入错误的地址空间，因此 TEE 场景下 context 的创建、更新与销毁必须纳入可信流程，而不能完全交给不可信 host；规范本身定义接口语义，配置权的可信裁决属于 AP-TEE/CoVE-IO 层。\par
\noindent\textbf{（2）两级翻译与设备直通。}\par
这一机制的直接作用是：IOMMU 支持第一级（guest 页表）与第二级（host 页表）翻译的组合，使设备直通给 VM 成为可能。具体实现上，被直通设备的 DMA 必须被限制到已分配的内存；非法访问产生 fault 而不是静默越界。机密虚拟机场景还要求两点规范未覆盖的增强：fault 报告不能向 host 泄露过多 TVM 状态，解除分配后必须清理陈旧映射，防止残留翻译被复用。\par
\noindent\textbf{（3）MSI 翻译。}\par
这一机制的直接作用是：IOMMU 通过 MSI page table 把设备发出的 message-signaled interrupt 写入重映射到正确的 interrupt file，与 AIA 的 IMSIC 模型衔接。具体实现上，这使 host 可以在不知道 guest 中断布局的情况下管理设备，同时把"设备能向哪个 guest 发中断"变成可配置、可检查的策略对象；机密场景下该表的配置同样属于可信边界。\par
\noindent\textbf{（4）失效、故障与队列语义。}\par
这一机制的直接作用是：IOMMU 不是一张静态表，翻译缓存（IOTLB）需要正确的 invalidation，fault queue 让软件观察违规，command queue 承载上下文与缓存管理命令。具体实现上，invalidation 顺序错误会产生陈旧翻译窗口；fault queue 既是可观测性来源，也构成 DoS 与侧信号面——高频注入 fault 可干扰宿主导调。这些运行时语义是未来 CoVE-IO 实现必须显式定义策略的地方。\par
\subsubsection{安全性与威胁边界}
规范的安全边界是设备地址转换与权限检查。其强假设包括：IOMMU 硬件正确实现、host/TSM 配置正确、device ID 不被伪造、缓存失效顺序正确。它不能独立解决恶意设备、PCIe 链路中间人、TDISP 接口状态或内存加密问题；配置错误、陈旧翻译、ATS/PRI 缓存与 confused deputy 都是规范自身无法排除的风险。仅证明 CPU 执行域不可直接读取，并不能自动说明 DMA 路径同样安全——IOMMU 正是把这一缺口显性化的规范。\par
\subsubsection{实验环境与证据}
规范无实验。从可核验的材料看，本地 PDF 约 109 页，Figure 1/2 与 Table 1 能够直接支撑 DMA 隔离、两级翻译与设备直通的术语和语义；这些材料不足以支撑设备信任、SPDM/TDISP 集成或任何性能开销结论。\par
本节判断以 IOMMU 规范 Figure 1/2 与 Table 1 及 README 核验块为依据；超出这些材料覆盖范围的实现细节、性能结论或攻击面，本文不作延伸推断。\par
\subsubsection{实验结果与证据强度}
规范不提供 benchmark。性能相关的设计点可以识别——IOTLB 命中率、页表遍历深度、invalidation 频率、fault 处理路径都会影响实际开销——但具体数字必须来自实现评测，规范文本不支持任何开销结论。\par
本节判断以 IOMMU 规范证据边界为依据；超出这些材料覆盖范围的实现细节、性能结论或攻击面，本文不作延伸推断。\par
\subsubsection{综合评价}
综合来看，IOMMU 是 RISC-V 可信 I/O 方向不可缺的底座，其优势在于以批准状态提供了稳定的 DMA 翻译与设备直通语义，且通过 MSI page table 与 AIA 形成了中断路径的接口闭环。其局限在于不证明设备可信、不处理链路加密、不定义配置的可信裁决。从部署角度看，虚拟化、机密虚拟机与 DPU/加速器直通的任何 RISC-V 平台都必须以它为起点，并与 IOPMP、AIA、CoVE-IO、SPDM/TDISP/IDE 组合使用。\par
\subsection{论文 12：RISC-V Advanced Interrupt Architecture Specification}
\subsubsection{论文基本信息}
《RISC-V Advanced Interrupt Architecture Specification》由 RISC-V International 发布（编辑 John Hauser），v1.0 于 2023 年 6 月批准，本地保存的批准库版本含 20250312 澄清修订，PDF 约 90 页。本文将其置于"RISC-V 基础安全机制、RISC-V CoVE-IO / TEE-I/O"的研究脉络中考察，并把它作为规范类材料使用。它所回答的核心问题是：当平台走向多核、多 guest 与设备直通时，传统 wired IRQ 模型如何被可扩展、可虚拟化的中断架构取代。规范给出的答案是 APLIC、IMSIC、MSI 与 guest interrupt file 的组合。\cite{riscv_aia_2023}\par
\subsubsection{研究背景与动机}
规范面对的系统矛盾有两层。第一层是可扩展性：传统 PLIC 式有线中断在多核、多虚拟机场景下路由复杂、状态集中，难以支撑大规模虚拟化。第二层是可信边界：中断本身就是攻击面——设备和 host 可以通过伪造、错投或滥用中断影响 TVM/CVM 的执行节奏与状态。\par
具体而言，MSI 让设备通过一次内存写触发中断，天然适合与地址翻译基础设施结合；但"设备能向谁发中断"必须由策略约束。机密计算场景进一步要求把中断所有权与 TVM/设备生命周期绑定：设备直通给一个 TVM 期间，其中断投递路径不应被 host 任意重定向。规范解决第一层问题并给出第二层的机制零件，可信策略则由上层定义。\par
\subsubsection{关键挑战与核心思想}
规范的关键判断是：I/O 安全不只看 DMA，还要看中断是否可伪造、错投或被 host 滥用；而机制与策略必须分层——AIA 定义机制，CoVE 的 COVI 接口与 CoVE-IO 这类上层定义可信管理。\par
具体而言，APLIC 与 IMSIC 让中断虚拟化有了清晰的结构，但规范文本本身不判断哪个设备可信、不定义 interrupt file 的所有权转移协议，也不承诺 confidential interrupt isolation。把 spec mechanism 与 trusted policy 分开，是阅读本规范时最重要的纪律。\par
\subsubsection{系统架构}
规范的组织方式是：APLIC 处理有线中断源，按 domain 路由到 machine/supervisor/guest 上下文；IMSIC 接收 MSI，把每个 hart（及每个 guest）的中断状态组织为独立的 memory-mapped interrupt file；hart 通过新增 CSR（Smaia/Ssaia 扩展）取中断。规范同时为 IOMMU 定义了向 VM 投递 MSI 的支持路径。\par
具体而言，Figure 1 展示传统有线中断投递，Figure 2 展示 MSI 与 IMSIC 模型，Table 1 给出规模限制参数。虚拟化支持使 guest 中断投递不再依赖 host 软件注入的单一通路，这是它在机密计算叙事中的价值起点。\par
\subsubsection{核心方法设计}
规范的关键是把中断源、投递目标与虚拟化状态拆成可独立管理的对象。\par
\noindent\textbf{（1）APLIC 与有线中断。}\par
这一机制的直接作用是：APLIC 负责传统外部中断源的域划分与路由。具体实现上，中断源被路由到 machine、supervisor 或 guest 上下文，支持优先级与 pending/enabled 状态管理；TEE 场景下，APLIC 的 domain 配置若由不可信管理者控制，就可能被错配，因此配置权归属需要上层策略裁决。\par
\noindent\textbf{（2）IMSIC 与 MSI 投递。}\par
这一机制的直接作用是：IMSIC 把 MSI 转化为 per-hart、per-guest 的 interrupt file，天然适配虚拟化与设备直通。具体实现上，设备通过 message write 触发中断，写操作本身可经 IOMMU 的 MSI page table 重映射；每个 hart/guest 拥有独立 interrupt file，使中断状态的隔离粒度与 CPU/VM 隔离粒度对齐。CoVE-IO 需要在这一模型上绑定可信设备分配，规范提供了挂接点但未提供绑定协议。\par
\noindent\textbf{（3）虚拟中断边界。}\par
这一机制的直接作用是：AIA 提供 guest external interrupt 的路由状态与虚拟化机制，但可信边界要靠 TSM/CoVE-IO 层检查。具体实现上，host 管理便利与 TVM 安全存在结构性冲突：host 需要能为调度和错误处理介入中断路径，TVM 需要保证关键中断不被伪造或截流。规范的风险面包括 IMSIC ownership 错配、MSI 重映射错误、host 注入与 rebind 竞态——这些都是 README 核验块明确列出的局限，也是后续研究入口。\par
\subsubsection{安全性与威胁边界}
AIA 是功能/架构规范，其安全性不来自文本本身，而来自上层如何分配和保护 interrupt file、MSI 页与注入路径。对 CoVE/CoVE-IO 而言，trusted MSI 需要 AP-TEE 的 COVI 与 CoVE-IO 额外绑定 TVM、设备身份与 IMSIC ownership；错误的 rebind、陈旧配置或 host 注入都可能破坏机密 I/O 语义。仅证明 CPU 执行域与 DMA 路径受控，并不能自动说明中断路径同样安全。\par
\subsubsection{实验环境与证据}
规范无实验。从可核验的材料看，本地 PDF 约 90 页，Figure 1（传统投递）、Figure 2（MSI/IMSIC）与 Table 1（规模限制）能够直接支撑中断原语术语；同时需要看到其适用边界：不证明任何 TEE 的中断隔离。\par
本节判断以 AIA Figure 1/2 与 Table 1 及 README 核验块为依据；超出这些材料覆盖范围的实现细节、性能结论或攻击面，本文不作延伸推断。\par
\subsubsection{实验结果与证据强度}
规范不提供中断延迟或投递开销的 benchmark。IMSIC/APLIC 的结构会影响中断可扩展性，但具体延迟与开销必须引用实现或系统论文；本材料只说明性能相关路径。\par
本节判断以 AIA 规范证据边界为依据；超出这些材料覆盖范围的实现细节、性能结论或攻击面，本文不作延伸推断。\par
\subsubsection{综合评价}
综合来看，AIA 是 RISC-V I/O 控制面的关键规范：其优势在于以批准状态提供了 MSI 与虚拟中断的可扩展模型，并与 IOMMU 的 MSI 翻译形成接口闭环；其局限在于不定义可信设备生命周期，中断路径的安全语义完全依赖上层。从部署角度看，多核、虚拟化与设备直通的 RISC-V 平台都必备 AIA；机密计算平台还必须在其上叠加 COVI 与 CoVE-IO 的可信管理。\par
\subsection{论文 13：Keystone: An Open Framework for Architecting Trusted Execution Environments}
\subsubsection{论文基本信息}
《Keystone: An Open Framework for Architecting Trusted Execution Environments》由 Dayeol Lee, David Kohlbrenner, Shweta Shinde, Krste Asanovic, and Dawn Song（UC Berkeley）完成，发表于 2020 年的 EuroSys '20（article 38，16 页），本地保存作者站点 PDF 并已核验。本文将其置于"RISC-V TEE 谱系：Sanctum / Keystone / Penglai / SPEAR-V"的研究脉络中考察，并把它作为系统类基础文献使用。它所回答的核心问题是：商用 TEE 把设计权衡固化在平台里，研究者与设备厂商能否拥有一个可按威胁模型自由裁剪 TCB 与功能的开放框架。Keystone 的答案是"安全监控器只做不可替代的事，其余交给 enclave 内 runtime"的分层框架。\cite{lee2020keystone}\par
\subsubsection{研究背景与动机}
论文面对的系统矛盾是"一个 TEE 不可能同时适合所有场景"：有的应用要最小 TCB，有的要完整 syscall 兼容，有的要 secure I/O，有的要 cache 分区或内存加密；而 SGX、TrustZone、SEV 各自把一组固定权衡锁死在硬件与固件中——论文引言以 SGXv1 为例：静态 enclave 尺寸、无 secure I/O、无 syscall 支持。使用者一旦选定平台，就被锁定在该平台的设计限制内，与自身应用需求无关。\par
具体而言，RISC-V 提供了开放的硬件接口（PMP、特权级、可修改的开源核），Keystone 把这种开放性转化为 TEE 框架：监控器、runtime、内存模型与安全功能都成为可替换件。这一动机直接解释了为什么 Keystone 不追求"最强"的单一设计点，而追求"可组合"的设计空间。\par
\subsubsection{关键挑战与核心思想}
在上述背景下，论文的关键判断是：让 Security Monitor（SM）只做不可替代的隔离与生命周期管理，把可变策略全部下放给 enclave 内的 runtime。\par
具体而言，SM 运行在 M-mode，负责 PMP 配置、enclave 度量、进入/退出与销毁清理；runtime（Eyrie 或 seL4 移植）运行在 enclave 内部，负责页表、syscall 代理、SDK 与策略模块。这一拆分的收益是双向的：固定 TCB 被压缩到 SM 级别（约 1.6 KLoC 增量），而不同 enclave 可以在安全与性能之间做不同选择——但代价同样清晰，runtime 进入 enclave 的 TCB，功能越多 TCB 越大。\par
\subsubsection{系统架构}
论文的基本组织方式是把 untrusted host、SM、runtime、eapp 与远端验证者串成一个完整生命周期（Figure 1/2）。\par
具体而言，host OS 负责资源分配与启动请求，但不能读写受 PMP 保护的 enclave 内存；SM 在创建时检查并度量 host 准备的内存、配置 PMP，运行期处理 enclave 的陷入，销毁时清理内存并释放 PMP entry；远端验证者依据平台、SM 与 enclave 的度量链决定是否交付秘密。Eyrie runtime 让未修改的普通应用获得 libc、syscall 代理与可选模块；论文还演示了 seL4 作为 runtime 的移植，证明 runtime 层的可替换性。\par
\subsubsection{核心方法设计}
论文的关键不是单一组件，而是把 PMP 隔离、runtime 抽象、端到端证明与可选功能模块组合成一个可实例化的框架。\par
\noindent\textbf{（1）SM 与 PMP 生命周期。}\par
这一机制的直接作用是：Keystone 的隔离基础是 SM 动态配置 PMP，把 enclave 内存从 host OS 视野中切走。具体实现上，创建时 SM 检查并度量 host 准备的内存区域；运行期 PMP 允许 enclave 访问自己的 region、拒绝 host 与其他 enclave；销毁时 SM 清理内存并回收 PMP entry（Figure 3，Table 1 的 create/run/resume/destroy SBI）。PMP entry 数量有限，是这一路线最硬的规模约束。\par
\noindent\textbf{（2）Runtime 抽象。}\par
这一机制的直接作用是：把"要不要页表、syscall、SDK、cache 策略"变成 enclave 级选择。具体实现上，Eyrie 提供原生 runtime（基础 1.8 KLoC，全部模块最多 3.6 KLoC，见 Table 3）；syscall 代理与 edge call 处理 enclave 与不可信 host 的交互；论文同时给出 seL4 runtime 移植，展示同一 SM 上可运行微内核级 runtime。\par
\noindent\textbf{（3）端到端证明。}\par
这一机制的直接作用是：Keystone 不只隔离内存，还把平台、SM 与 enclave 的度量链交给远端验证者。具体实现上，平台可信根提供启动度量，SM 哈希与 enclave 哈希进入证明报告，验证者依据已知 SM/enclave 度量决定是否释放秘密（Figure 2 与第 4.5 节）。\par
\noindent\textbf{（4）可选功能模块与权衡。}\par
这一机制的直接作用是：把 cache 分区、片上内存、self-paging 与内存加密做成可展示、可取舍的模块。具体实现上，Figure 5 比较不同内存模型；模块增强威胁模型，同时增加实现复杂度与开销。这为后续 Penglai（专注可扩展内存保护）与 SPEAR-V（专注统一 tag 原语）的专门化设计提供了基线。\par
\subsubsection{安全性与威胁边界}
威胁模型以不可信 OS 为核心，依赖 RISC-V PMP 与 SM 的正确性。论文展示了内存加密、cache 侧信道防御等可选扩展，但明确不声称解决所有微架构侧信道；安全强度取决于具体实例化——框架本身只保证"隔离与度量路径可信"，不替使用者选择威胁模型。因此把 Keystone 的实验结论外推到"RISC-V TEE 已解决侧信道"是不成立的。\par
\subsubsection{实验环境与证据}
实验环境是真实 RISC-V 板卡与开源核模拟的组合。硬件为 HiFive Freedom Unleashed（FU540）以及 Rocket/BOOM 等开源核设置；软件为 buildroot Linux 4.15、SM、Eyrie runtime、Linux 驱动与 seL4 移植；评估覆盖 CoreMark、Beebs、RV8、IOZone、Torch/FANN 应用、enclave 生命周期操作与 TCB 代码量统计（第 7.1 节，Table 2 硬件平台，Table 3 TCB 分解）。\par
本节判断以 Keystone 第 7 节、Figure 6、Table 2/3 及本地 PDF 核验数字为依据；超出这些材料覆盖范围的实现细节、性能结论或攻击面，本文不作延伸推断。\par
\subsubsection{实验结果与证据强度}
论文摘要与第 7 节给出的可核验数字（均经本地 PDF 复核）如下：CoreMark、Beebs、RV8 开销小于 1\%；IOZone 文件 I/O 开销可达约 40\%；Torch on Eyrie 最大 7.35\%（Densenet）；FANN with seL4 runtime 约 0.36\%。TCB 方面，Keystone 为 SM 增加的代码约 1.6 KLoC；M-mode 组件合计 10.7 KLoC（密码库 4 KLoC、既有 trap/boot/工具代码 4.7 KLoC、SM 基线 1.6 KLoC）；Eyrie RT 基础 1.8 KLoC、含全部模块最大 3.6 KLoC；最大 eapp 配置总 TCB 约 15 KLoC。\par
\begin{table}[htbp]
\centering
\caption{Keystone 可核验评估数字（论文第 7 节与 Table 3）}
\begin{tabular}{lll}
\toprule
维度 & 数字 & 来源 \\
\midrule
SM 增量代码 & 约 1.6 KLoC & 论文第 2 页/Table 3 \\
M-mode 组件合计 & 10.7 KLoC & 第 7 节 TCB 统计 \\
Eyrie RT（基础/最大） & 1.8 / 3.6 KLoC & Table 3 \\
最大 eapp 总 TCB & 约 15 KLoC & 第 7 节 \\
CoreMark/Beebs/RV8 开销 & $<$1\% & 摘要/第 7 节 \\
IOZone 开销 & 可达约 40\% & 摘要/Figure 7 \\
Torch on Eyrie（Densenet） & 7.35\% & 第 7 节 \\
FANN with seL4 & 0.36\% & 第 7 节 \\
\bottomrule
\end{tabular}
\end{table}\par
这些数字的正确读法不是"Keystone 很快"，而是"开销高度依赖路径"：纯 CPU 负载几乎免费，一旦落到经由不可信 host 的文件 I/O，开销上升一个数量级以上。Figure 6 对 enclave 生命周期的拆分进一步说明 create/destroy 成本受内存清理与证明影响。这正是"框架让用户显式选择功能与开销"论断的实验注脚。\par
\subsubsection{失败模式与边界分析}
论文与 README 核验块共同标定的边界有三条。其一，PMP entry 数量有限，enclave 数量与内存区域数量受硬件资源硬约束——这直接催生了 Penglai 的工作。其二，I/O 路径经不可信 host 代理，Iago 式返回值攻击与 40\% 级开销都没有被框架默认解决。其三，侧信道与大规模云部署不是默认能力，需要可选模块或后续工作。换言之，Keystone 的实验证据支撑"框架可行、开销可选"，不支撑"开箱即用的云级 TEE"。\par
\subsubsection{综合评价}
综合来看，Keystone 是 RISC-V TEE 谱系的基础 SOTA：其优势在于框架结构清楚、SM/RT 分工明确、全部组件开源可改，为新硬件原语研究提供了可复用起点。其局限在于 PMP 数量与物理 region 管理的规模瓶颈明显，secure I/O、侧信道与大规模 enclave 均非默认解决。从部署角度看，它适合教学、原型与定制设备；量产需要硬件供应商、工具链、认证与长期维护配套。它也是本节后文所有系统论文的对照基准：Penglai 解决它的规模问题，SPEAR-V 替换它的内存原语，CoVE 则把保护对象从 enclave 提升为整台 VM。\par
\subsection{论文 14：Scalable Memory Protection in the Penglai Enclave}
\subsubsection{论文基本信息}
《Scalable Memory Protection in the Penglai Enclave》由 Erhu Feng et al.（上海交通大学 / 上海人工智能实验室）完成，发表于 2021 年的第 15 届 USENIX OSDI，本地 PDF 已核验。本文将其置于"RISC-V TEE 谱系：Sanctum / Keystone / Penglai / SPEAR-V"的研究脉络中考察，并把它作为系统类同行评议 SOTA 使用。它所回答的核心问题是：PMP"少量连续 region"的模型无法支撑云原生 enclave 时，能否用软硬件协同把安全内存做成动态、细粒度、大规模的资源。Penglai 的答案是两个新硬件原语（GPT 与 MMT）加一套 fork 式快速初始化机制。\cite{feng2021penglai}\par
\subsubsection{研究背景与动机}
论文面对的系统矛盾是：serverless 与微服务把计算单元变成大量短生命周期实例，而既有 enclave 的内存保护机制是静态、受限、初始化昂贵的。论文第 1--2 节明确指出三类痛点：SGX 式 EPC 有固定保护内存容量与高昂的初始化成本；PMP 式 enclave 受寄存器数量与连续内存管理限制；而目标场景要求数千个并发 enclave、大容量安全内存与秒级以下的函数级创建销毁。\par
具体而言，论文把目标量化为三个可检验指标：支持 1,000 量级并发 enclave、安全内存扩展至 512GB（含加密与完整性保护）、消除安全内存的高成本初始化。这三个指标贯穿全文评估，是判断其贡献是否成立的标尺。\par
\subsubsection{关键挑战与核心思想}
在上述背景下，论文的关键判断是：不要在每次内存访问时做昂贵的所有权检查，而应把检查提前到页表映射与安全内存元数据两个结构性位置。\par
具体而言，Guarded Page Table（GPT）保护 host 页表本身，使不可信 OS 无法把 secure page 映射给自己；Mountable Merkle Tree（MMT）把完整性验证组织成可挂载子树，使验证路径长度与安全内存总量解耦；shadow enclave 把代码/数据的度量预处理为可复用模板，把 serverless 场景的重复初始化变成 fork 式创建。三者分别对应容量、规模与启动延迟三个瓶颈。\par
\subsubsection{系统架构}
论文的基本组织方式是：监控器/驱动/SDK 加硬件 GPT/MMT 扩展，组成一个可扩展的 enclave 系统（Figure 1）。\par
具体而言，host 应用经驱动与库发起 enclave 操作；secure monitor 管理 enclave 生命周期与硬件配置（monitor 约 6,399 LoC，论文第 3 页）；host 仍负责普通资源管理。Penglai 把 HPT Area（host 页表区）、enclave 页表与 MMT meta-zone 全部纳入可信管理；硬件上修改了 Rocket Core、内存控制器与 SoC 相关路径，并在 FPGA、QEMU 与 Gem5 上实现。\par
\subsubsection{核心方法设计}
论文的关键是把页表所有权检查、可扩展完整性、快速初始化与 enclave 间服务化组合成一条完整路径。\par
\noindent\textbf{（1）Guarded Page Table。}\par
这一机制的直接作用是：GPT 用硬件与 HPT Area 约束 host 页表，让页所有权成为硬件可检查的属性（Figure 2/3，第 4.1 节）。具体实现上，host 页表必须位于受保护的 HPT Area；MMU 扩展在页表遍历路径上检查 PTE 是否非法指向 secure page；由此避免每次内存访问都查询大型元数据结构，把所有权检查的成本集中到有硬件加速的翻译路径上。\par
\noindent\textbf{（2）Mountable Merkle Tree。}\par
这一机制的直接作用是：MMT 让安全内存完整性保护随容量扩展，而不是让树节点吞噬 cache 与带宽（第 4.2 节）。具体实现上，完整性树被组织为可挂载/卸载的 SubTree，热点安全内存的验证路径被缩短；树根与关键元数据留在 SoC/监控器控制范围内；论文的设计目标是支撑最高 512GB 安全内存的加密加完整性保护。\par
\noindent\textbf{（3）Shadow Enclave 快速启动。}\par
这一机制的直接作用是：把 serverless 场景的重复初始化变成 fork 式创建（第 4.3 节，Figure 11）。具体实现上，shadow enclave 不直接运行，只保存可复用的代码/数据模板；度量可提前计算并密封，fork 时避免重复高成本的度量与初始化；运行期再按需动态添加内存。论文摘要报告该机制把安全内存初始化延迟降低三个数量级，Figure 11 的评估在 16MB enclave 配置下展示这一结论。\par
\noindent\textbf{（4）Server Enclave 与 IPC。}\par
这一机制的直接作用是：把 enclave 组织成服务链，缓解所有功能都信任 host OS 的问题（第 5.2 节）。具体实现上，server enclave 可承载文件系统等服务，降低 Iago 攻击风险；relay page 支持 host-enclave 与 enclave-enclave 通信。这使 Penglai 的抽象更接近微服务组合，而非单个安全进程。\par
\subsubsection{安全性与威胁边界}
威胁模型包括不可信 OS 与对 enclave 内存的访问/篡改；安全性依赖监控器、GPT/MMT、内存加密与完整性树。侧信道不是主要解决目标，README 核验块明确指出证据不足以覆盖 cache 与 controlled-channel 攻击。同时应注意：GPT/MMT 是研究原型的硬件扩展，与后续 CoVE 的 MTT 不是同一标准化路径，不能把 Penglai 的实验结论直接当作标准化机密内存的证据。\par
\subsubsection{实验环境与证据}
实验环境覆盖硬件扩展原型、Linux 修改与 serverless/微服务负载。实现上修改了 RISC-V 核/MMU/SoC 路径并增加 GPT/MMT 引擎，配套监控器、驱动、SDK 与库；负载包括 CPU 密集型 benchmark（RV8、CoreMark）、内存密集型 Redis、MapReduce 与 serverless 应用；平台为 FPGA、QEMU 与 Gem5。评估证明的是可扩展 enclave 原型的可行性，不证明与 AP-TEE/CoVE 标准路径的兼容性。\par
本节判断以 Penglai 摘要、实现与评估章节、Figure 1/11 及本地 PDF 复核数字为依据；超出这些材料覆盖范围的实现细节、性能结论或攻击面，本文不作延伸推断。\par
\subsubsection{实验结果与证据强度}
论文摘要与评估章节给出的可核验数字（均经本地 PDF 复核）：支持 1,000 量级 enclave 实例并发运行；安全内存可扩展至 512GB 并同时提供加密与完整性保护；GPT 对内存密集型负载（Redis）开销约 5\%，对 CPU 密集型负载（RV8、CoreMark）可忽略；安全内存初始化延迟降低三个数量级；真实应用（MapReduce）经 shadow fork 获得约 3.6 倍加速。\par
这组数字的证据强度在于它们全部来自摘要与评估章节的一手实验；其边界在于平台是研究原型（修改核与内存控制器），数字不代表任何量产 SoC 上的预期表现，也不代表标准化路径下的成本。\par
\subsubsection{与相关工作的定量对比}
相对 Keystone 基线，Penglai 的改进是规模维度的：PMP 路线受 entry 数量约束，Penglai 用页表级所有权把并发 enclave 推到 1,000 量级；固定 EPC 式容量被推到 512GB；创建路径从"每次完整度量初始化"变成 fork 模板复用。相对 SPEAR-V，两者都在改内存原语，但方向不同：Penglai 解决容量与初始化，SPEAR-V 解决页表攻击面与嵌套隔离。这三篇放在一起，恰好构成"同一 enclave 抽象、三种内存原语取舍"的对照组。\par
\subsubsection{综合评价}
综合来看，Penglai 是 RISC-V enclave 规模化的强 SOTA：其优势在于直接命中 PMP region 数量、安全内存容量与启动延迟三个真实瓶颈，且全部结论有原型实验支撑。其局限在于硬件改动较大（核、MMU、内存控制器），商业化依赖硬件扩展被供应商采纳；与后续 CoVE TVM 标准不在同一抽象层。从部署角度看，serverless enclave、边缘云与云原生微隔离是其目标场景；主要风险在 ISA/SoC 标准化与生态支持。\par
\subsection{论文 15：SPEAR-V: Secure and Practical Enclave Architecture for RISC-V}
\subsubsection{论文基本信息}
《SPEAR-V: Secure and Practical Enclave Architecture for RISC-V》由 David Schrammel et al.（Graz University of Technology）完成，发表于 2023 年的 ACM ASIA CCS，本地 PDF 已核验。本文将其置于"RISC-V TEE 谱系：Sanctum / Keystone / Penglai / SPEAR-V"的研究脉络中考察，并把它作为系统类同行评议 SOTA 使用。它所回答的核心问题是：能否用一个统一的硬件原语同时解决页表攻击面、双向隔离、共享内存与任意嵌套，而不堆叠多种机制。SPEAR-V 的答案是 DRAM tag store 加不可变页表。\cite{schrammel2023spearv}\par
\subsubsection{研究背景与动机}
论文面对的系统矛盾是：页表自身会成为攻击面。如果页表由不可信 OS 管理，enclave 需要额外防御 remap、unmap 与 fault pattern 观察（controlled-channel 攻击）；如果把页表复制或移入安全区，性能与软件复杂度又会上升。论文摘要同时指出，既有方案或有性能问题与 controlled-channel 风险，或只支持嵌入式受限场景，或对软件施加不现实约束（如固定数量的物理内存区间）。\par
具体而言，现代云应用需要的是灵活内存边界、共享内存与嵌套能力——例如 enclave 内再隔离一个不可信插件或 JIT 组件。SPEAR-V 选择给页与页表加 tag，让硬件翻译路径直接检查合法性，而不是在软件层补救页表的不可信性。\par
\subsubsection{关键挑战与核心思想}
在上述背景下，论文的关键判断是：把页表标记为不可变（immutable）之后，OS 可以继续管理普通页表，但不能偷偷修改 enclave 映射——页表安全由此从"软件审计"变成"硬件状态"。\par
具体而言，tag store 为 DRAM 中每个 4KiB 页记录 owner/enclave ID、页类型、immutable 位、页级别等字段；PTW 与 TLB 在地址翻译时携带 tag 做权限检查。嵌套 enclave 与双向 sandbox 都来自同一条语义——"这页属于谁、谁能访问这页"——因此共享内存与任意嵌套不需要额外机制，这是"单一硬件原语"论断的核心。\par
\subsubsection{系统架构}
论文的基本组织方式是：tag store 位于 DRAM 中专区，SM 在 M-mode 管理 tag，PTW/TLB 在硬件翻译路径上检查 tag（Figure 1）。\par
具体而言，Figure 3 展示 tag 字段如何进入 TLB entry；tag store 自身由 PMP 保护为仅 M-mode 可访问（不变式 ℐ1：U/S-mode 不可读写 tag store，论文还禁读以防未来侧信道）；SM 提供 E\_CREATE、E\_ADD/E\_REMOVE、E\_ENTER/E\_EXIT、SHARE/REVOKE 等 API。enclave ID 空间支持超过 1600 万个 enclave，嵌套层级不受结构限制。\par
\subsubsection{核心方法设计}
论文的关键是把页级 tag、不可变页表、嵌套语义与 SM API 收敛到一个硬件原语上。\par
\noindent\textbf{（1）页粒度 Tag Store。}\par
这一机制的直接作用是：把"这页属于谁"放进硬件可查的元数据（Figure 1，第 5.1--5.3 节）。具体实现上，每个 4KiB 页对应 tag store 中一条记录，含 owner/ID 与页类型；tag store 自身经 PMP 与 SM 自保护；非 taggable 内存可跳过 tag 读取，使普通程序不为此付出代价——这是 unprotected 应用零开销设计的结构基础。\par
\noindent\textbf{（2）不可变页表。}\par
这一机制的直接作用是：immutable 位是对抗页表攻击的关键点（Figure 2，第 5.4 节）。具体实现上，SM 在向 enclave 加入页时验证映射唯一性，并把相关页表标记为 immutable；此后对 immutable 页表的写访问只能经 SM 进行，违规写会陷入 SM；硬件页表遍历器同样尊重 immutable 属性。由此 OS 无法通过修改 PTE 制造 alias 或 remap，也无法通过 unmap 制造 controlled-channel 故障序列。\par
\noindent\textbf{（3）嵌套与双向沙箱。}\par
这一机制的直接作用是：SPEAR-V 不只保护 enclave 免受 host 攻击，也保护 host 免受 enclave 攻击，并允许 enclave 内再隔离子区域。具体实现上，嵌套 enclave 页与父 enclave 页走同一 tag/owner 规则，父 enclave 不能直接读子 enclave 私有页；共享内存通过 SHARE/REVOKE 管理的共享密钥显式建立。这一双向语义对插件、库与 JIT 等不可信组件隔离场景直接可用。\par
\noindent\textbf{（4）SM API 与生命周期。}\par
这一机制的直接作用是：把 tag 写入与页所有权变更收敛到可信路径（第 5.5 节）。具体实现上，E\_CREATE/E\_DESTROY 管理 enclave 生命周期，E\_ADD/E\_REMOVE 管理 enclave 页所有权，E\_ENTER/E\_EXIT 管理执行切换，E\_SHARE/E\_REVOKE 管理共享内存；所有 tag 变更都经 SM 验证，不存在绕过路径。\par
\subsubsection{安全性与威胁边界}
威胁模型覆盖不可信 host/OS，并以缓解 controlled-channel 攻击为明确目标。完整安全性仍依赖硬件 tag 检查的正确实现、SM/runtime 与 tag 管理策略；物理攻击与全部微架构侧信道不在完全解决范围内。README 核验块特别提示：若 tag 元数据不能贯穿 DMA/I/O/fabric 路径，整体机密边界仍不完整——tag 保护的是 CPU 翻译路径，不自动约束设备。\par
\subsubsection{实验环境与证据}
实验环境是研究原型与模拟/硬件评估，基于 CVA6 CPU（32KB cache，TLB 覆盖约 64KiB 内存），目标是验证 tag 扩展的开销与面积/性能权衡。评估区分 unprotected 与 protected 应用，并探索 0/32/64/128-bit tag 等设计点：0-bit 时无法运行 enclave 但也无额外开销，评估主要使用 64 与 128-bit 配置；负载包括 LMbench（lat\_mem\_rd 测延迟、bw\_mem 测带宽）与 Embench（含嵌套 enclave 场景，重复 50 次取 95 百分位）。证据边界：证明原型可行，不代表供应商量产采用。\par
本节判断以 SPEAR-V 摘要、第 8 节评估、Figure 3/6/7、Table 1 及本地 PDF 复核数字为依据；超出这些材料覆盖范围的实现细节、性能结论或攻击面，本文不作延伸推断。\par
\subsubsection{实验结果与证据强度}
论文可核验数字（均经本地 PDF 复核）：摘要报告 unprotected 应用零开销、protected 应用平均约 1\% 开销；第 8 节细分显示，Embench 上仅内存标记造成的开销为 1.1\%，嵌套 enclave 场景总开销 1.2\%（主要开销来自 timer interrupt 时 SM 保存与切换状态）；LMbench lat\_mem\_rd 在 256KiB 以上工作集的平均延迟增加约 1.5\%（128-bit tag）；bw\_mem（4MiB 工作集）带宽影响总体小于 1\%，其中写带宽仅下降 0.09\%、读带宽最高下降 0.95\%。\par
机制层面，对 tagged 内存，TLB 在每个翻译级别为 64-bit tag 增加一次额外内存访问、为 128-bit tag 增加两次——这解释了为什么 128-bit 配置开销略高（PTW 内存接口每周期只能读一个字）。这组数字的证据强度在于低开销结论有逐 benchmark 细分支撑；边界在于全部数字来自研究原型与模拟，嵌套/共享语义虽有实现但跨真实云负载的长期验证不足。\par
\subsubsection{综合评价}
综合来看，SPEAR-V 是 RISC-V enclave 原语方向的强 SOTA：其优势在于把页表安全、双向隔离、共享内存与任意嵌套装进单一硬件语义，且 protected 路径平均约 1\% 的开销有细分实验支撑。其局限在于需要非标准硬件扩展与 DRAM tag store，生态采纳不确定；侧信道、I/O 与标准化仍需其他机制补全。从部署角度看，它适合嵌入式、边缘与插件级隔离场景；主要风险在 ISA 扩展接受度、验证成本与工具链支持。\par
\subsection{论文 16：CoVE: Towards Confidential Computing on RISC-V Platforms}
\subsubsection{论文基本信息}
《CoVE: Towards Confidential Computing on RISC-V Platforms》由 Ravi Sahita et al.（Rivos）完成，2023 年公开预印本（arXiv:2304.06167），本地 PDF 已核验。本文将其置于"RISC-V CoVE / AP-TEE 机密虚拟机"的研究脉络中考察，并把它作为架构提案类基础文献使用。它所回答的核心问题是：RISC-V 如何从进程级 enclave 上升到与 Arm CCA、Intel TDX、AMD SEV-SNP 同层的机密虚拟机。CoVE 的答案是 TEE VM（TVM）+ TEE Security Manager（TSM）+ Memory Tracking Table（MTT）+ 三组 ABI（COVH/COVG/COVI）的参考架构。\cite{sahita2023cove}\par
\subsubsection{研究背景与动机}
论文面对的系统矛盾是：多租户平台中 firmware、host OS、VMM 与运营者全部进入租户 TCB，而租户真正需要的只是继续使用完整 guest OS 的同时，把 host 降级为不可信的资源管理者。\par
具体而言，传统 VM 完全信任 hypervisor；机密 VM 要保留 hypervisor 的调度与资源管理能力，却剥夺其读写私有内存的权限。RISC-V H 扩展提供虚拟化底座，但缺少 confidential 模式、内存归属跟踪与证明机制；Keystone/Penglai 系 enclave 保护的是进程，不是整台 VM。CoVE 在概念层与 CCA/TDX/SNP 对齐，填补的正是 RISC-V 在这一层的空白。\par
\subsubsection{关键挑战与核心思想}
在上述背景下，论文的关键判断是：CoVE 不是把 Keystone 式 enclave 放大，而是用 TSM 与 MTT 重新定义 VM 的内存所有权。\par
具体而言，TVM 包含 guest firmware、guest OS 与应用；MTT 逐页跟踪物理内存属于普通世界、TSM 还是某个 TVM，host 只有在 TSM 允许时才能 donate、reclaim 或 share；三组 ABI 按角色切分接口——COVH 面向 host 的 TVM 生命周期管理，COVG 面向 guest 的证明与共享请求，COVI 面向安全中断设施。host 继续拥有管理便利，TSM 独占安全裁决，这是整套设计的骨架。\par
\subsubsection{系统架构}
论文的基本组织方式是：TSM 作为 TEE 与 non-TEE 世界之间的 TCB 中介，host 继续调度资源（Figure 1 参考架构）。\par
具体而言，SoC 可信根度量 TSM-driver 与 TSM，构成证明链底座；TSM-driver 负责机器相关的 bootstrap 与上下文工作，TSM 负责 TVM 内存与状态不变量的强制；host/VMM 经 COVH 创建 TVM、添加页、调度 vhart；TVM 经 COVG 请求证明、内存共享与半虚拟化 I/O。Figure 2 展示 MTT 强制路径，Figure 3 给出 SoC 视图，Table 1 给出特权级安排。\par
\subsubsection{核心方法设计}
论文的关键是把监控器拆分、内存所有权表、TVM 生命周期与证明/设备边界组织成一个可对齐产业标准的参考架构。\par
\noindent\textbf{（1）TSM 与 TSM-driver 拆分。}\par
这一机制的直接作用是：把安全强制集中在 TSM，把机器相关的 bootstrap 与上下文工作放到 TSM-driver（Figure 1，第 3.1 节）。具体实现上，TSM-driver 为 TSM 建立空间与时间隔离；TSM 聚焦 TVM 内存/状态不变量。拆分的收益是允许实现差异，代价是证明必须覆盖的 TCB 随之扩大——验证者要同时信任 TSM-driver 与 TSM 的度量。\par
\noindent\textbf{（2）Memory Tracking Table。}\par
这一机制的直接作用是：MTT 是 CoVE 内存隔离的核心，逐页记录物理页归属（Figure 2，第 4 节）。具体实现上，host 经 TSM 向 TVM donate 内存；MTT 强制阻止 host 访问 confidential 页；共享/非机密页必须显式标记，服务半虚拟化 I/O；reclaim 路径要求清理后再复用。这与 Arm CCA 的 GPT 在概念上同层，但元数据组织与执行者不同。\par
\noindent\textbf{（3）TVM 生命周期原语。}\par
这一机制的直接作用是：把 VM 创建拆成 allocation、measurement、vhart 创建、finalization 与 execution 等可验证阶段。具体实现上，初始被度量页进入 TVM measurement；demand-zero 机密页可由 VMM 按需加入；执行、fault 与 exit 仍由 host 调度，但机密寄存器状态由 TSM 保存与恢复。生命周期被显式分阶段，是为了让"host 在哪一步还能影响什么"成为可推理的问题。\par
\noindent\textbf{（4）证明与设备边界。}\par
这一机制的直接作用是：明确证明与 I/O 是 CoVE 能否安全商用的关键边界（第 6.4 节与证明章节）。具体实现上，DICE/SoC 可信根建立 TSM/TVM 的证据链；设备挂接需要 SPDM/TDISP/IDE/IOMMU 等额外机制；论文只提出边界要求，不等于完整 TEE-I/O 标准——这正是 CoVE-IO 草案后续要补的内容。\par
\subsubsection{安全性与威胁边界}
攻击者模型是不信任 host/运营者与部分平台软件；目标是保护 TVM 内存与寄存器状态。论文主要是设计讨论，不提供完整形式化证明；侧信道、物理攻击与设备攻击需要额外机制。因此 CoVE 的安全论证足以支撑"架构方向成立"，不足以支撑"任何 CoVE 实现安全"——实现层的 TLB shootdown、页转换竞态与中断绑定都需要独立验证。\par
\subsubsection{实验环境与证据}
CoVE 是架构提案，无完整系统实验。从可核验的材料看，本地 PDF 含参考架构（Figure 1）、MTT 强制（Figure 2）、SoC 视图（Figure 3）与特权级表（Table 1）；能够直接支撑 TVM/TSM/TSM-driver/MTT/COVH/COVG/COVI 的概念与分工。这些材料不足以支撑 AP-TEE 最终 ABI、生产实现开销或设备 TEE-I/O 的完整安全。\par
本节判断以 CoVE Figure 1--3、Table 1 及 README 核验块为依据；超出这些材料覆盖范围的实现细节、性能结论或攻击面，本文不作延伸推断。\par
\subsubsection{实验结果与证据强度}
论文没有给出系统开销数字。可以识别的是未来开销的来源——MTT 查询、TSM 进入/退出、度量与内存转换都可能产生成本——但任何具体数字都必须等待 AP-TEE 实现或系统论文。与 CCA/TDX/SNP 的比较是设计对照，不是基准测试。\par
本节判断以 CoVE 的材料范围与"无评估 benchmark"这一事实为依据；超出这些材料覆盖范围的实现细节、性能结论或攻击面，本文不作延伸推断。\par
\subsubsection{综合评价}
综合来看，CoVE 是 RISC-V 机密虚拟机的入口材料：其优势在于把 TVM 模型、TSM/MTT 分工与三组 ABI 写成与 CCA/TDX/SNP 同层对齐的完整概念框架，直接催生了 AP-TEE 标准化工作。其局限在于提案性质强、无实现无评估；设备 I/O、debug、PMU 与中断边界都留给后续规范。从部署角度看，它对开放机密 VM 生态至关重要；主要风险在标准收敛速度、firmware 质量与验证者生态。\par
\subsection{论文 17：RISC-V Application-Processor Trusted Execution Environment Specification}
\subsubsection{论文基本信息}
《RISC-V Application-Processor Trusted Execution Environment Specification》（即 CoVE 规范）由 RISC-V AP-TEE Task Group 完成，本地保存 v0.7 / RC2（ARC 评审候选）版本，PDF 约 80 页以上。必须显式标注：该材料为草案、尚未批准（Draft / not ratified: AP-TEE v0.7 / RC2），后续版本可能修改状态机、接口名称或 ABI。本文将其置于"RISC-V CoVE / AP-TEE 机密虚拟机"的研究脉络中考察，并把它作为本方向的主 SOTA 规范材料使用。它所回答的核心问题是：CoVE 的概念如何变成 host、guest、firmware 与验证者可以互操作的同一套状态机与 ABI。\cite{riscv_ap_tee_2024}\par
\subsubsection{研究背景与动机}
规范面对的系统矛盾是：机密虚拟机只有概念不够——没有标准 ABI，OS/VMM/firmware 之间无法互操作；没有明确的内存生命周期语义，donate/reclaim/share 容易产生数据残留或所有权混乱；没有证明令牌格式，依赖方无法判断 TVM 与 TSM 的真实状态。\par
具体而言，草案第 4 节定义威胁模型与安全需求，第 5 节定义部署模型与隔离，第 6 节定义 TVM 证明与度量，第 7 节定义 TVM 生命周期，第 8 节给出 CoVE SBI 扩展提案。从第 4 节到第 8 节的结构本身就说明动机：先把"要防什么"写清，再把"谁在哪一步能做什么"写成状态机，最后落到可调用的接口。\par
\subsubsection{关键挑战与核心思想}
在上述背景下，草案的关键判断是：把 TVM 当作有状态对象，所有页、vCPU、中断与令牌都围绕生命周期变化。\par
具体而言，TVM build/initialization 决定初始度量；TVM execution 决定 host 与 TSM 如何处理 exit/fault/interrupt；TVM memory management 分别规定 measured、demand-zero、shared、private 页的安全要求。规范引入 Supervisor Domain 抽象组织 TSM 管理的隔离域，并用 THCS/VHCS 等接口组织 host 与 guest 侧调用——概念沿自 CoVE 的 COVH/COVG/COVI 三分法，落到 SBI 扩展上成为可实现 ABI。\par
\subsubsection{系统架构}
规范的主体不是结构图，而是规范化的状态机与 ABI 契约。TSM/TSM-driver 与 host/guest 的边界由 SBI 扩展切开；第 6 节定义 TSM token 与 TVM token 的内容与关系；第 7/8 节把创建、内存管理、执行与关闭对应到运行期接口。\par
具体而言，部署依赖 RISC-V H 扩展、AIA（含 IMSIC 绑定）、IOMMU/IOPMP 与 SoC 级内存隔离；规范不规定单一实现，但规定了任何实现都必须满足的安全需求与接口语义。这意味着"符合 AP-TEE"是一个可逐条核对的属性，而不是品牌声明。\par
\subsubsection{核心方法设计}
草案的关键是把构建度量、内存语义、令牌模型与 SBI 接口写成可实现的合同。\par
\noindent\textbf{（1）TVM 构建与度量。}\par
这一机制的直接作用是：规定 TVM 初始内容如何进入度量，避免 host 在创建后偷换载荷（第 6.1.2、7.1、8.2 节）。具体实现上，多步与单步创建都必须形成可验证初始状态；measured 页的指派影响 TVM token；finalization 之后，初始被度量载荷不应再被 host 任意修改。构建路径的设计目标是让验证者拿到的度量值唯一对应一份初始镜像。\par
\noindent\textbf{（2）内存管理语义。}\par
这一机制的直接作用是：把 private/shared/zero/reclaimed 页的转换写清楚（第 5.2、7.3 节）。具体实现上，TVM 私有页对 host 不可见；共享页服务半虚拟 I/O 但必须显式转换；reclaim/reuse 要求清理，避免秘密残留；页转换伴随 TLB 失效与所有权记录更新。内存生命周期是草案中安全需求密度最高的部分，因为 donate/reclaim/share 的竞态与残留是机密 VM 实现最常见的缺陷来源。\par
\noindent\textbf{（3）证明令牌模型。}\par
这一机制的直接作用是：把 TSM token 与 TVM token 分开，让验证者同时看到平台与工作负载（第 6 节）。具体实现上，TSM token 证明 TSM/TCB 状态，TVM token 证明 TVM 的度量、配置与安全姿态；debug、PMU 等可观测能力被要求进入证明姿态——即可观测性不是隐含默认，而是验证者可裁决的属性。\par
\noindent\textbf{（4）CoVE SBI 扩展。}\par
这一机制的直接作用是：SBI 扩展是 host/guest/TSM 互操作的落地点（第 8 节）。具体实现上，host 运行期接口覆盖能力探测、TVM 创建、内存管理与运行；guest 接口覆盖证明请求与共享内存操作；中断路径经 COVI 语义与 AIA IMSIC 绑定衔接。草案 ABI 可能变化，引用任何接口名时都必须标注 not ratified 状态。\par
\subsubsection{安全性与威胁边界}
威胁模型把 host/hypervisor 视为不可信管理者，目标是保护 TVM 内存、寄存器状态与证明。强假设是 TSM 与底层硬件正确、平台证明密钥与度量正确、内存生命周期无别名与竞态。规范是契约而非证明：实现仍需独立验证 TLB shootdown、页转换、中断绑定、MMIO/共享内存边角案例；host 控制的资源管理带来的竞态、侧信道、DoS 与实现缺陷仍是草案明确承认的风险面。\par
\subsubsection{实验环境与证据}
规范草案无实验。从可核验的材料看，本地 PDF 约 80 页以上，章节覆盖威胁模型（第 4 节）、部署与隔离（第 5 节）、证明（第 6 节）、生命周期（第 7 节）与 SBI（第 8 节）；能够直接支撑 TVM 生命周期、内存映射安全需求、证明令牌与运行期接口的语义。适用边界必须反复强调：v0.7 草案、未批准，不能引用为最终标准或生产 ABI，也不能作为性能证据。\par
本节判断以 AP-TEE 草案 PDF 第 4--8 节及 README 核验块为依据；超出这些材料覆盖范围的实现细节、性能结论或攻击面，本文不作延伸推断。\par
\subsubsection{实验结果与证据强度}
草案不提供 benchmark，但揭示了未来开销来源：TSM 调用、页转换、度量、TLB 失效与中断处理。规范的价值在互操作，不在证明开销；性能结论必须等待后续实现或系统论文。\par
本节判断以 AP-TEE 草案证据边界为依据；超出这些材料覆盖范围的实现细节、性能结论或攻击面，本文不作延伸推断。\par
\subsubsection{综合评价}
综合来看，AP-TEE 是 RISC-V 机密 VM 方向最关键的当前材料：其优势在于把 CoVE 的生命周期、ABI 与证明写成可实现接口，与 Arm CCA/RME/RMM 处于同一抽象层，使跨平台对照成为可能。其局限在于未批准状态、缺生产实现与性能数据、设备 I/O 闭环依赖独立的 CoVE-IO 草案。从部署角度看，一旦稳定，它将成为开放机密 VM 软件栈的对齐点；在此之前，任何基于它的设计都应按"草案依赖"管理版本风险。\par
\subsection{论文 18：RISC-V CoVE-IO Specification}
\subsubsection{论文基本信息}
《RISC-V CoVE-IO Specification》（Confidential VM Extension I/O）由 RISC-V AP-TEE-IO Task Group（Samuel Ortiz、Jiewen Yao 等）完成，本地保存 v0.3.0 版本。必须显式标注：该材料为草案、尚未批准（Draft / not ratified: CoVE-IO v0.3.0）。本文将其置于"机密 I/O 协议、可信设备接口与网络端点、RISC-V CoVE-IO / TEE-I/O"的研究脉络中考察，并把它作为规范类材料使用。它所回答的核心问题是：AP-TEE 保护了 TVM 的内存与执行之后，真实设备经 MMIO、DMA、中断与 PCIe/CXL.io 链路接触 TVM 数据的路径如何纳入同一信任框架。\cite{riscv_cove_io_2026}\par
\subsubsection{研究背景与动机}
草案面对的系统矛盾是：机密 VM 的 I/O 默认会退回共享内存加 bounce buffer——安全边界清楚，但 host 可见中转页，TVM 必须复制、加密、验证，性能与语义都不理想。\par
具体而言，要让 TVM 直接使用设备，需要同时证明四件事：设备身份可信（固件可度量）、设备接口确实被锁定分配给本 TVM、链路上的数据不被窃听篡改、DMA/MMIO/中断被约束在授权范围。任何一个环节缺失，设备都会成为 host 观察或篡改 TVM 的通道；而设备重置、热插拔与错误恢复又随时可能破坏已建立的信任状态。草案把这条链写成生命周期与协议要求，而不是把"设备自证身份"误认为终点。\par
\subsubsection{关键挑战与核心思想}
在上述背景下，草案的关键判断是：可信设备分配是证据、地址策略、链路保护与清理四个机制的交集，CoVE-IO 不是单一模块，而是把多个外部标准与 RISC-V TSM 生命周期绑定起来。\par
具体而言，SPDM 解决设备身份与度量，TDISP 解决设备接口的分配状态，PCIe IDE 解决链路机密性与完整性，IOMMU/IOPMP 限制 DMA，AIA/IMSIC 约束中断投递，TSM 负责把这些证据与 TVM 的内存、中断、DMA 策略绑定。草案引入 TDI（TDISP Device Interface）、TDM（TDISP Device Module）、DSM/RDSM（设备安全管理器及其可信根）等角色，并定义 CoVE-IO manifest 与 COVH/COVG/COVT 函数承载设备侧 ABI。\par
\subsubsection{系统架构}
草案把 TVM、host VMM、TSM、IOMMU、root port 与设备组织成一个可信 I/O 编排：Figure 1 给出 bounce buffer 基线，Figure 2 解释 PCIe 拓扑，Figure 3 给出高层架构，Tables 1--3 列出需求。\par
具体而言，host 仍管理资源，但可信分配状态不能只由 host 决定；设备接口从 unassigned 到 assigned to TVM 的转换必须是可证明的；TSM 记录 TVM 与设备的绑定关系，并参与连接 transcript、IDE 密钥与接口绑定的管理。架构的实质是把"设备可用"与"设备可信"拆成两个独立状态，后者需要证据驱动。\par
\subsubsection{核心方法设计}
草案的关键是把证据绑定、基线/直通模式与 DMA/中断/MMIO 策略写成完整生命周期。\par
\noindent\textbf{（1）Bounce Buffer 基线。}\par
这一机制的直接作用是：先定义安全基线——TVM 私有内存不直接暴露给设备 DMA，I/O 数据经共享/非机密缓冲中转（Figure 1）。具体实现上，这降低了对设备信任的要求，但把复制与加解密负担放在数据路径上；草案明确把 bounce buffer 作为默认安全选项，把直接分配作为需要完整证据链的增强选项，避免把直通误写成默认可用。\par
\noindent\textbf{（2）可信设备分配。}\par
这一机制的直接作用是：直接分配要把设备接口状态绑定到具体 TVM。具体实现上，设备必须经 SPDM 证明身份与度量状态，经 TDISP 证明接口处于安全锁定状态；TSM/IOMMU/IOPMP 配置为只允许访问 TVM 授权内存；unassign、reset 与 cleanup 同样是安全生命周期的一部分——分配是可逆的，逆向路径的清理与正向路径的建立同等重要。\par
\noindent\textbf{（3）DMA/中断/MMIO 策略。}\par
这一机制的直接作用是：可信 I/O 不只看 DMA，还包括 MMIO 寄存器与中断投递。具体实现上，DMA 翻译需 IOMMU/IOPMP 配合；MMIO 配置空间不能让 host 注入危险状态；中断/MSI 路由必须防伪造与错投，与 AIA 的 IMSIC 模型及 IOMMU 的 MSI 翻译衔接。草案列出的威胁包括 trusted MMIO 恶意访问/重映射/重定向、PCIe 链路中间人、PCIe ID 伪造与 confused-deputy 式 DMA 重映射——这些条目为后续实现的安全测试提供了直接清单。\par
\noindent\textbf{（4）证据绑定。}\par
这一机制的直接作用是：让最终验证者能知道 TVM 使用的设备也处于可信状态。具体实现上，平台证明只证明 CPU/TSM 不够；设备证据（SPDM 证书与度量、TDISP 状态）需要与 TVM token 或策略关联，形成"平台令牌—TSM 策略—设备证据—TVM-设备绑定—验证者裁决"的链条。这是 CoVE-IO 与 SPDM/TDISP 原始规范的连接点，也是它相对纯架构提案的增量。\par
\subsubsection{安全性与威胁边界}
草案明确列出 trusted MMIO 恶意访问/重映射/重定向/预配置、PCIe 链路中间人、PCIe ID 伪造与 confused-deputy DMA 重映射等威胁。强假设是设备支持 PCIe/CXL.io、TDISP、DOE 与 IDE，平台 TSM/RDSM 正确，IOMMU/IOPMP/AIA 组合无配置错误。草案不证明具体实现安全，也不解决所有 DoS 与侧信道；multi-function 设备、ATS/PRI 缓存、共享加速器与 CXL 内存设备都是 README 核验块列出的复杂风险点。\par
\subsubsection{实验环境与证据}
规范草案无实验。从可核验的材料看，本地 PDF 的 Figure 1--3 与 Tables 1--3 能够直接支撑可信 I/O 需求分类、生命周期与架构；这些材料不足以支撑性能开销、最终批准状态或任何具体设备的实际安全性。\par
本节判断以 CoVE-IO 草案 PDF 及 README 核验块为依据；超出这些材料覆盖范围的实现细节、性能结论或攻击面，本文不作延伸推断。\par
\subsubsection{实验结果与证据强度}
草案只解释模式，不测模式。可以识别的性能风险来自 bounce buffer 复制与加解密、SPDM/TDISP 建立开销、IOMMU 翻译、IDE 加密与中断虚拟化；直接分配可减少复制，但显著增加可信生命周期复杂度。具体数字必须来自未来实现或设备/厂商报告。\par
本节判断以 CoVE-IO 草案证据边界为依据；超出这些材料覆盖范围的实现细节、性能结论或攻击面，本文不作延伸推断。\par
\subsubsection{综合评价}
综合来看，CoVE-IO 是 RISC-V 机密 I/O 的关键蓝图：其优势在于把 SPDM、TDISP、IDE、IOMMU、AIA 与 TSM 生命周期放进同一框架，使"TVM 直通设备"第一次有了端到端的可核对需求清单；它是 RISC-V 追平 CCA/TDISP/IDE 生态的关键接口。其局限在于草案状态、无实现无性能数据，且对 PCIe/设备生态的依赖很强。从部署角度看，云机密 VM 的设备直通、加速器、存储与 SmartNIC 场景都需要它；在批准与设备生态成熟之前，bounce buffer 基线仍是唯一无附加假设的选项。\par
\subsection{开发商安全实践建议}
本章九篇材料覆盖了从已批准规范到研究原型的完整证据谱系，对芯片厂商、固件与系统软件开发商、设备制造商而言，可操作的结论如下。\par
\noindent\textbf{实践一：按"批准规范优先"原则落地 I/O 隔离底座。}\par
SoC 厂商应优先实现已批准的 IOMMU（v1.0.1）与 AIA（v1.0）语义——device context、两级翻译、MSI page table、IMSIC/APLIC——而不是自研私有等价物 \cite{riscv_iommu_2023,riscv_aia_2023}。理由很直接：CoVE 与 AP-TEE 的全部 I/O 设计都以这两份规范的术语为挂接点 \cite{sahita2023cove,riscv_ap_tee_2024}，自研偏离会把后续 TSM 固件与上游软件栈的适配成本转移到自己头上。同时应以 IOPMP/sIOPMP 补充物理区间访问控制，覆盖 IOMMU 不约束的路径。\par
\noindent\textbf{实践二：把 IOMMU/AIA 配置权纳入可信流程，而非交给 host 默认路径。}\par
IOMMU 规范自身声明配置错误、陈旧翻译与 confused deputy 是固有风险 \cite{riscv_iommu_2023}；AIA 的 IMSIC ownership 错配、rebind 竞态与 host 注入同样是已识别风险 \cite{riscv_aia_2023}。固件开发商应把 device context 更新、MSI 重映射与 interrupt file 绑定/解绑设计为 TSM 裁决的操作，并为 invalidation 顺序、fault queue 处理与 unassign 清理写出可测试的不变量——这些正是 AP-TEE 内存生命周期语义在设备侧的对应物 \cite{riscv_ap_tee_2024}。\par
\noindent\textbf{实践三：用小 TCB、可度量的监控器工程换取可审计性。}\par
Keystone 的 SM 增量约 1.6 KLoC、最大 eapp 总 TCB 约 15 KLoC \cite{lee2020keystone}，Penglai 的 monitor 约 6,399 LoC \cite{feng2021penglai}——两个系统都证明安全关键监控器可以控制在万行量级。固件团队应把 TSM/监控器作为独立工程管理：最小化代码量、公开度量值、纳入安全审计与模糊测试，并把 TSM-driver 与 TSM 的拆分边界写进威胁模型文档 \cite{sahita2023cove}。监控器一旦膨胀到数十万行，开放 ISA 带来的可审计性优势就被固件质量抵消。\par
\noindent\textbf{实践四：按工作负载实测选择隔离原语，把 I/O 路径当作开销热点管理。}\par
三篇系统论文的数字给出了一致的结构性结论：纯 CPU 负载开销可忽略（Keystone CoreMark/Beebs/RV8 $<$1\% \cite{lee2020keystone}；Penglai RV8/CoreMark 可忽略 \cite{feng2021penglai}；SPEAR-V protected 平均约 1\% \cite{schrammel2023spearv}），而 I/O 路径开销上升一个数量级（Keystone IOZone 约 40\% \cite{lee2020keystone}）。设备厂商选型时不应比较抽象的"TEE 开销"，而应在自己的真实负载（尤其是存储与网络路径）上实测，并把初始化延迟纳入 serverless 场景的 SLO（Penglai 的 shadow fork 证明初始化可优化三个数量级 \cite{feng2021penglai}）。\par
\noindent\textbf{实践五：设备直通默认 bounce buffer，直接分配走完整证据链。}\par
CoVE-IO 把 bounce buffer 定义为默认安全基线，把直接分配定义为需要 SPDM 身份、TDISP 接口锁定、IDE 链路保护与 IOMMU/IOPMP/AIA 路径约束同时成立的增强选项 \cite{riscv_cove_io_2026,riscv_iommu_2023,riscv_aia_2023}。设备与平台厂商应按此分层交付：先保证 bounce 路径正确，再把 direct assignment 作为显式 opt-in 发布，并把 unassign/reset/cleanup 的逆向生命周期纳入出厂测试。鉴于 CoVE-IO 仍是 v0.3.0 草案，实现应按草案依赖管理版本，接口封装留出演进空间。\par
\noindent\textbf{实践六：把开放 ISA 的可审计性转化为产品安全资产。}\par
Boubakri 综述把 SDK、工具链、标准符合性、ASIC 量产验证与生态整合列为采用条件 \cite{boubakri2025riscvtee}；Keystone 证明监控器、runtime 与评估可以全部开源复现 \cite{lee2020keystone}。厂商应主动开源安全监控器与度量流程、发布可复现的度量值、引入第三方审计，并持续跟踪 AP-TEE 与 CoVE-IO 的批准状态 \cite{riscv_ap_tee_2024,riscv_cove_io_2026}。开放性本身不产生安全，只有把开放性变成"任何人都能核对度量与代码"的工程流程，它才转化为采购方可验证的安全声明。\par
\subsection{未来研究方向}
\noindent\textbf{（1）AP-TEE/CoVE 参考实现与性能基准。}\par
问题是：CoVE 与 AP-TEE 均无任何系统开销数字，MTT 查询、TSM 进入/退出、页转换与度量的成本完全未知 \cite{sahita2023cove,riscv_ap_tee_2024}。现有工作未解决的原因是规范先于实现。可行路线是在 QEMU/FPGA 上实现 TSM stub 与 COVH 创建/运行/内存转换路径，配套 host KVM/QEMU 补丁；需要的实验包括 TVM 创建/运行/销毁延迟、页转换与 TLB 失效开销、与 Arm CCA 实现的同负载对照，以及对 hostile host 的模糊测试。\par
\noindent\textbf{（2）CoVE-IO 端到端组合验证。}\par
问题是：SPDM/TDISP/IDE/IOMMU/AIA/TSM 的组合语义只在草案层面存在，没有任何材料证明组合后无竞态、无遗漏 \cite{riscv_cove_io_2026}。未解决的原因是涉及五个独立演进的标准与设备生态。可行路线是构建模拟 PCIe 设备加 TSM/RDSM 测试框架，建立连接 transcript、TDISP 绑定、MMIO 映射与 DMA 策略的形式化模型；需要的实验覆盖设备连接/断开、SPDM 密钥更新、IDE 刷新、恶意重映射与中断注入，验收标准是 host 无法重映射 trusted MMIO/DMA、TVM 只接受已证明接口。\par
\noindent\textbf{（3）内存所有权原语的统一与标准化。}\par
问题是：PMP、GPT、MMT、MTT 与 tag store 解决的是同一问题的不同切面——Keystone 受 PMP entry 数量约束 \cite{lee2020keystone}，Penglai 用 GPT/MMT 解决容量与初始化 \cite{feng2021penglai}，SPEAR-V 用 tag 解决页表攻击面与嵌套 \cite{schrammel2023spearv}，AP-TEE 用 MTT 语义规范 VM 级所有权 \cite{riscv_ap_tee_2024}——但四者之间没有统一的抽象与兼容性分析。未解决的原因是研究原型与标准化沿不同路径演进。可行路线是以 AP-TEE 内存生命周期为参考模型，逐项映射 GPT/MMT/tag 的语义差异；需要的实验包括同一负载下不同原语的开销与表达能力对照，以及原语组合（如 tag store 加 MTT）是否引入语义冲突。\par
\noindent\textbf{（4）机密 VM 的侧信道边界。}\par
问题是：Boubakri 综述的 Table 1 显示多数 RISC-V TEE 对 cache 侧信道只有部分或不支持 \cite{boubakri2025riscvtee}，而 CoVE/AP-TEE 的威胁模型把侧信道显式留给额外机制 \cite{sahita2023cove,riscv_ap_tee_2024}；SPEAR-V 虽以 controlled-channel 缓解为目标，但未覆盖全部微架构侧信道 \cite{schrammel2023spearv}。未解决的原因是侧信道防御与 TVM 抽象层的结合点尚不清楚。可行路线是把 cache 分区、页表故障屏蔽与中断时序防护设计为 TSM 可配置策略，并进入证明姿态；需要的实验包括在 TVM 模型下复现已知 controlled-channel 与 cache 攻击，评估防御开销与证明表达。\par
\noindent\textbf{（5）跨平台证明令牌的统一验证。}\par
问题是：AP-TEE 定义了 TSM token 与 TVM token 的分层模型 \cite{riscv_ap_tee_2024}，CoVE-IO 又要求设备证据与 TVM 令牌关联 \cite{riscv_cove_io_2026}，但 RISC-V 令牌与 Arm CCA token、EAT 等格式之间没有统一验证者，依赖方无法写一份策略同时裁决多平台证据。未解决的原因是标准化各自为政且设备证据刚进入草案。可行路线是定义平台无关的令牌转换与策略层，把 debug/PMU 等可观测姿态统一编码；需要的实验是用同一验证者裁决 RISC-V 与 Arm 的真实证据链，并测试设备证据更新（固件升级、密钥轮换）下的策略连续性。\par
\noindent\textbf{（6）可信中断路径的安全属性。}\par
问题是：AIA 提供 IMSIC/MSI 机制但不定义可信策略 \cite{riscv_aia_2023}，AP-TEE 的 COVI 语义与 IMSIC 绑定/解绑的竞态、陈旧配置与 host 注入风险只在 README 级被识别 \cite{riscv_ap_tee_2024}，没有任何材料给出 trusted interrupt 的安全不变量。未解决的原因是中断路径横跨 AIA、IOMMU MSI 翻译与 TSM 三份材料。可行路线是为 interrupt file 所有权转移定义状态机与不变量，形式化 rebind/unbind 的竞态自由性；需要的实验包括 host 注入测试、rebind 压力测试与 IOMMU MSI 表和 AIA 状态的一致性检查，验收标准是 host 无法向未授权 TVM 注入 trusted MSI、rebind 后旧路径立即失效。\par
\section{TEE 设计空间与执行环境谱系}
TEE 的实现形态虽然多样，但 CPU、内存、I/O、证明和 TCB 划分始终构成同一设计空间。本节先用一篇 server-side TEE 设计综述和 Arm 官方 TrustZone 白皮书建立比较坐标：前者把``谁管理运行时资源''抽象为四种可比较的管理模式，后者给出双世界模型与 SoC 级安全状态传播的原始机制定义；再把 Cortex-M 快速状态切换漏洞（ret2ns）与 TrustZone 商用 TEE 的漏洞谱系放回这些坐标中，分析设计选择如何影响可验证性和可攻击面。四篇材料恰好构成``框架—机制—攻击—谱系''的闭环：框架告诉我们在哪些维度上比较，机制告诉我们硬件实际提供了什么，攻击与漏洞谱系告诉我们软件层在哪里失守。\par
因此，后续各论文的结论不被孤立解读为某一特性``安全''或``不安全''，而是追问其假设是否覆盖特权软件、外设、异常与生命周期管理。例如，TrustZone 的 NS bit 只能保证硬件事务级隔离，Cerdeira 等人的漏洞统计则说明绝大多数失守发生在 secure world 软件自身的输入校验与接口设计上；反过来，ret2ns 又提醒我们：即使硬件机制本身经过精心设计，为性能引入的``快速路径''也可能改变权限传播语义。这样的对照也为后续 CCA、CoVE 和设备机密计算的跨平台比较提供统一语言。\par

\subsection{论文 19：SoK: Understanding Design Choices and Pitfalls of Trusted Execution Environments}
\subsubsection{论文基本信息}
《SoK: Understanding Design Choices and Pitfalls of Trusted Execution Environments》由 Mengyuan Li、Yuheng Yang、Guoxing Chen、Mengjia Yan、Yinqian Zhang（MIT、上海交通大学、南方科技大学）完成，发表于 2024 年的 ACM Asia Conference on Computer and Communications Security (ASIA CCS 2024)。本文将其置于``硬件辅助可信执行环境分类''的研究脉络中考察，并把它作为本节的主 taxonomy 锚点使用。当前证据状态为：本地 PDF 已保存并核验（ASIA CCS 2024 正式版本）；下文仅在该范围内讨论其机制与结论。\cite{li2024sokteechoices}\par
\subsubsection{研究背景与动机}
论文要解决的问题是 confidential computing 的一个结构性矛盾：cloud TEE 一方面要保护云租户的 data-in-use，另一方面又必须让不可信的 host OS/hypervisor 继续管理 CPU、memory、I/O 等云资源。论文第 1 页摘要和第 2 页 Introduction 明确把问题写成：host OS 原本完全控制资源，TEE 设计必须在 host OS 与 TCB 之间重新划分每一项资源管理职责；否则非本领域读者只能看到 SGX、SEV、TDX、CCA、PEF、Keystone、Sanctum、CURE、Penglai 等产品名列表，难以理解设计选择背后的共性。\par
论文把背景知识压缩为四个概念（第 2--3 页 Section 2）。其一，Secure Remote Execution (SRE)：远程用户希望在不可信云平台上执行敏感负载，并获得 secure measurement、confidentiality、integrity 三类保证。其二，威胁模型：攻击者控制整套软件栈，可执行特权指令、launch/pause/destroy TEE 实例、管理非 TEE 侧的内存映射、I/O 设备与 CPU 调度；商用 TEE 通常把 DoS 与微架构侧信道排除在外。其三，TCB 二分：Manufacturer TCB 包含硬件、微码、固件、security monitor、quoting/provisioning enclave 等平台方组件，ISV TCB 是 enclave/CVM 内的应用或 OS。其四，生命周期抽象：Figure 2 把 TEE 生命周期分为 remote attestation（证明初始状态可信）与 runtime management（决定运行时资源由谁管理）两段。这个划分是本篇 SoK 的解释力来源：安全问题不再按``哪个平台''组织，而按``哪个运行时事件由谁管理''组织。\par
\subsubsection{关键挑战与核心思想}
在上述背景下，作者提出并解决了三个层面的挑战。\par
\noindent\textbf{C1：产品异构导致的不可比较性。}不同 TEE 的文档、论文和规范使用各自术语描述本质上相同的资源管理问题。作者的应对是先把运行时事件归一化：Figure 3（第 5 页）把 CPU、memory、I/O 三组常见事件列成统一清单，使任何平台都能被映射到同一坐标系。\par
\noindent\textbf{C2：资源管理委托的安全—性能—TCB 三角权衡。}把管理权交给 host 性能最好但攻击面最大，完全收进 TCB 边界最清晰但复杂度最高。作者的核心思想是提出 TEE Runtime Architectural Framework (TRAF)，用四种 management mode 刻画``host 管、RTPM（Runtime Protection Module）管、host 管+RTPM 校验、实例自管''这四种基本委托方式（Figure 4，第 5 页），把看似各异的设计选择还原为同一权衡的不同落点。\par
\noindent\textbf{C3：漏洞与设计缺陷之间缺少系统映射。}已知攻击散落在各平台的安全公告与论文中，难以回答``哪类设计选择更容易出错''。作者用 Table 1（第 9 页）和 Appendix Table 2 把已知攻击回扣到具体资源管理模式的弱保护或校验遗漏上，使 pitfall 分析具备可复用的结构。\par
\subsubsection{系统架构}
TRAF 的系统视图包含四类角色：SoC/硬件信任根、RTPM（即运行时保护模块，在不同平台上对应 security monitor、TSM、RMM 或 TDX module 等）、不可信的 host OS/hypervisor，以及 TEE 实例及其内部 runtime。生命周期被抽象为两段：remote attestation 决定``启动时加载了什么''，runtime management 决定``运行中每个资源事件由谁仲裁''。Figure 3 把运行时事件分为三组：CPU 管理（CPU scheduling、context switch、interrupt/instruction emulation）、内存管理（virtual address space、memory allocation、page fault、page table、protected/shared memory）、I/O 管理（I/O data transmission、I/O operations、shared memory、I/O devices）。Figure 5（第 7 页）进一步把多个代表性 TEE 在这八个维度上的模式选择填入统一矩阵，证明该坐标系确实能承载跨平台比较，而不是只按产品名分段描述。\par
\subsubsection{核心方法设计}
论文的关键贡献是四种管理模式及其在逐事件上的落点分析。\par
\noindent\textbf{（1）四种 runtime management mode（Figure 4）。}表~\ref{tab:traf-modes} 按论文 Figure 4 与 Section 4 的论述整理四种模式的管理主体、收益与风险。\par
\begin{table}[htbp]
\centering
\caption{TRAF 四种运行时管理模式（依据论文 Figure 4 与 Section 4 整理）}
\label{tab:traf-modes}
\small
\begin{tabular}{p{0.14\linewidth}p{0.22\linewidth}p{0.28\linewidth}p{0.28\linewidth}}
\toprule
模式 & 管理主体 & 典型收益 & 典型风险 \\
\midrule
Unprotected & 不可信特权软件 & 性能好、云平台可调度 & host 可观察/操控资源，攻击面最大 \\
RTPM-only & 运行时保护模块 & 安全边界最清晰 & TCB 变大，复杂任务代价高 \\
RTPM-guarded & host 配置、RTPM 校验 & 平衡 TCB 与可管理性 & corner case 校验遗漏形成漏洞 \\
Instance-assisted & TEE 实例自身参与 & 灵活，减少 host 观察面 & 恶意实例或复杂 runtime 扩大 TCB/DoS 风险 \\
\bottomrule
\end{tabular}
\end{table}
\noindent\textbf{（2）CPU 管理事件的落点。}CPU scheduling 通常是 unprotected mode，因为 CSP 必须保留调度权，放进 RTPM 会显著增加 TCB 与性能开销（第 6 页 Section 4.3.1）。Context switch 通常采用 RTPM-only mode 以保护寄存器与微架构状态；AMD SEV 早期是著名例外，其未充分保护的寄存器状态后来成为攻击入口，SEV-ES/SEV-SNP 才转向更强保护（第 6 页 Section 4.3.2）。\par
\noindent\textbf{（3）内存管理事件的落点。}Page fault 与 memory allocation 是模式分歧最大的灰区（第 7--8 页 Section 4.4）：RTPM-guarded 可降低 TCB，但 page fault 可被 host 用作 controlled-channel 观察 enclave 的访问模式；RTPM-only 更强但实现复杂；instance-assisted 允许 Keystone、CURE 等框架的 enclave runtime 参与页表管理。\par
\noindent\textbf{（4）I/O 事件的落点与 pitfall 映射。}I/O 在多数平台仍是 unprotected mode，靠应用层加密保护数据；trusted I/O 需要设备侧可信逻辑、PCIe IDE、SR-IOV 等额外机制（第 8 页 Section 4.5）。Pitfall 分析用 Table 1（第 9 页）把 vulnerable design flaws 归因到具体模式：未保护寄存器、TLB 误用、未认证加密、未保护 NPT 等问题大多来自 unprotected 模式或 RTPM-guarded 的校验遗漏；Appendix Table 2 给出已知攻击目录作为对照。\par
\subsubsection{安全性与威胁边界}
论文的威胁模型清晰（第 3 页 Section 2.2）：攻击者控制特权软件、可管理内存/I/O/调度、可启动/暂停/销毁 TEE 实例；物理攻击是否纳入取决于具体 TEE；DoS 与微架构侧信道被多数商用 TEE 明确排除。需要强调的是，论文不声称 TRAF 本身防御任何攻击；它是解释设计选择与漏洞来源的分析框架。Table 1 说明许多 SEV/SEV-ES/SNP 相关攻击可归因于 context switch、instruction emulation、memory allocation 等资源管理模式的弱保护，但论文也说明侧信道、硬件缺陷、推测执行类攻击不属于 runtime design flaw 的主线讨论。范围边界同样明确（第 2--3 页）：本篇 SoK 聚焦通用处理器上的 hardware-based server-side TEE，明确排除没有 SRE 语义的 TrustZone、纯软件 TEE、TEE 的软件扩展与嵌入式 TEE。因此，如果只引用该 SoK 就断言某平台``安全''或``性能好''，证据不足；它只能支撑 taxonomy、design-choice framing 与 pitfall explanation，不能替代 Arm CCA、SEV-SNP、TDX 等原始论文或规范的一手机制证据。\par
\subsubsection{实验环境与证据}
这是一篇 SoK，无新系统实现，也无 benchmark。其证据对象是系统化整理：Figure 1（传统云与 TEE 云的信任关系对比）、Figure 2（generalized TEE lifecycle）、Figure 3（CPU/memory/I/O 事件分类）、Figure 4（四种管理模式）、Figure 5（代表性 TEE 的模式选择矩阵）、Table 1（vulnerable design flaws）与 Appendix Table 2（已知攻击目录）。论文没有代码 artifact；商业平台的 TCB 规模、固件行为与 corner-case 实现受闭源限制，论文未说明也无法独立验证。\par
\subsubsection{实验结果与证据强度}
对``TRAF 能否解释不同 TEE 的设计选择''这一问题，Figure 5 的矩阵提供了较强支撑：SGX、SEV、TDX、CCA、Keystone 等确实可以在同一组维度上被区分。对``TRAF 能否解释已知设计缺陷''，Table 1 与 Section 5.2/5.3 的 case study 支撑较强，但只覆盖论文选定的案例。性能维度上论文没有新测量，任何``某模式开销为 X\%''的说法都不在该文证据范围内，必须回引各平台原始论文。\par
\subsubsection{综合评价}
综合来看，本篇 SoK 的价值在于把 cloud TEE 的核心矛盾抽象为 resource management delegation，而非又一份按平台罗列的产品年表。这个抽象能解释为什么 CPU scheduling 放给 host、context switch 必须强保护、page fault 处于灰区、I/O 仍大量依赖 unprotected path，因此比 schneider2022soktee 式的广义 taxonomy 更适合作为本报告的坐标系。其局限有三：一是范围限于 server-side 通用 CPU TEE，不覆盖 TrustZone 移动/嵌入式主线与 accelerator/device TEE；二是时间截面截至 2024 年，之后的 AP-TEE、CoVE-IO、TDISP/IDE、Arm CCA device assignment 需另行补齐；三是无新实验与形式化证明，模式级别的真实性能开销仍然未知。在本节中，它的角色是为后续三篇 TrustZone 材料提供``设计选择—缺陷''的统一语言。\par

\subsection{论文 20：ARM Security Technology: Building a Secure System Using TrustZone Technology}
\subsubsection{论文基本信息}
《ARM Security Technology: Building a Secure System Using TrustZone Technology》由 ARM Limited 于 2009 年发布（文档编号 PRD29-GENC-009492C），是 TrustZone 架构的官方白皮书。本文将其置于``Arm TrustZone TEE 与漏洞谱系''的研究脉络中考察，并把它作为 vendor 机制定义类材料使用。它所回答的核心问题是：如何把安全边界从 CPU 指令级扩展到整个 SoC——让每个总线事务都携带安全状态。当前证据状态为：本地 PDF 已保存并核验（109 页全文）；下文仅在该范围内讨论其机制与结论。\cite{arm_trustzone_whitepaper}\par
\subsubsection{研究背景与动机}
要理解该方案，首先需要看到它面对的系统矛盾：TrustZone 的诞生背景不是云多租户，而是单设备内把安全服务与 rich OS 分开。普通 OS 代码量大、复杂、易被攻破，但系统仍需要显示、键盘、存储、基带、安全启动等安全敏感服务。更关键的是，安全边界必须横跨 CPU、cache/TLB、内存控制器、DMA、debug 与 interrupt 路径，而不能只在软件进程层做隔离——任何一个路径上的疏漏都会让进程级隔离形同虚设。这解释了 TrustZone 为什么强调 SoC integration，而不是只定义一组 enclave 指令。\par
\subsubsection{关键挑战与核心思想}
白皮书的核心思想是把系统状态二分：Secure 事务可以访问安全资源，Non-secure 事务被硬件拒绝。围绕这一思想，白皮书处理了三个挑战。\par
\noindent\textbf{C1：处理器级的状态二分与受控切换。}每个物理核被虚拟化为 Secure 与 Non-secure 两个虚拟核，二者以时间片方式共享流水线，通过新增的 monitor mode 完成上下文切换。进入 monitor 的路径被严格收紧：Normal world 只能通过 SMC 指令或 IRQ、FIQ、外部 Data Abort、外部 Prefetch Abort 等异常子集进入；Secure Configuration Register (SCR) 中的 NS bit 只能由 monitor mode 软件直接修改，Normal world 软件无法访问 SCR（白皮书 3.3 节）。\par
\noindent\textbf{C2：安全状态的 SoC 级传播。}仅有 CPU 状态位没有意义，必须让内存与外设也知道一次访问来自哪个 world。白皮书的解法是 AMBA3 AXI 总线为读写通道各增加一个控制信号，即 Non-Secure bit（NS bit）：AWPROT[1] 标记写事务、ARPROT[1] 标记读事务的安全属性，低为 Secure、高为 Non-secure（白皮书 3.2 节）。所有 Non-secure master 的 NS bit 在硬件上被固定拉高，使其在地址译码阶段就无法匹配任何 Secure slave。白皮书还指出，NS bit 附加到总线事务与 cache tag 上，相当于提供了``第 33 根地址线''，形成 Secure 与 Non-secure 两套物理地址空间，因此必须小心处理 memory aliasing 问题。\par
\noindent\textbf{C3：控制面的独立处理。}中断与 debug 是两个最容易被忽视的旁路。白皮书推荐 IRQ 作为 Normal world 中断源、FIQ 作为 Secure world 中断源的模型，并由 monitor 在每次世界切换时改写 SCR 中的 IRQ/FIQ 陷阱位；CP15 中还有只能由 Secure world 访问的配置寄存器，可禁止 Normal world 屏蔽 CPSR 的 F 位与 A 位（白皮书 3.3.3 节）。Debug 侧，CoreSight 组件通过 SPIDEN、SPNIDEN、DBGEN 等认证接口信号独立控制 Secure debug 的使能，而不是依赖普通外设保护机制（白皮书 3.4 节）。\par
\subsubsection{系统架构}
架构总览是``双世界 + monitor + 安全外设''：普通 OS 仍运行绝大多数功能，安全服务在 Secure world 中响应请求。CPU 支持安全状态与 monitor mode；Normal world 通过 SMC 请求安全服务；内存与外设由 TZASC/TZPC/TZMA 等控制器按安全属性隔离；DMA 引擎（如 PrimeCell DMAC）必须支持并发 Secure/Non-secure 通道并遵守同样的安全状态约束。白皮书 Figure 2-1 给出一个蜂窝手机 SoC 的完整示例，Figure 4-1 展示用 TZPC 动态控制外设安全属性的接法。\par
\subsubsection{核心方法设计}
\noindent\textbf{（1）双世界与 monitor mode。}Secure world 运行 trusted OS/firmware/TA，Non-secure world 运行 rich OS 与应用。Monitor mode 软件是两世界之间的守门人：保存当前世界状态、恢复目标世界状态，然后执行 return-from-exception。世界切换不是进程切换，而是更底层的安全状态切换；SMC 还可以在寄存器中携带消息负载，配合 world-shared memory 构建高效的跨世界通信协议；对 VFP/NEON 等协处理器还支持 lazy context switching 以降低切换成本（白皮书第 5 章）。\par
\noindent\textbf{（2）NS bit 与总线/外设检查链。}AXI 主总线携带 NS bit，APB 低速外设总线不携带；因此 AXI-to-APB bridge 承担 APB 外设的安全裁决——桥内为每个外设提供 TZPCDECPROT 输入信号，决定该外设是 Secure 还是 Non-secure，并拒绝安全属性不匹配的访问（每个桥最多管理 16 个 APB 外设）。AHB 总线同样无安全状态，AXI-to-AHB/AHB-to-AXI 桥只能把整个 AHB 域映射为单一安全属性，除非增加自定义逻辑按 master ID 生成 NS bit。\par
\noindent\textbf{（3）内存分区控制器：TZASC、TZMA 与 TZPC。}TrustZone Address Space Controller (TZASC, PL380) 是 AXI 组件，可将其 slave 地址范围划分为多个区域并由 Secure 软件逐区域配置安全属性，主要用途是把单颗片外 DRAM 划分为多个安全域——这比放两颗小容量内存更便宜；TZASC 只能分区 memory-mapped 设备，不能用于 NAND flash 等块设备。TrustZone Memory Adapter (TZMA, BP141) 用于片内 SRAM/ROM，以 4KB 为粒度划出单个 Secure 区域，其 R0SIZE 信号可由 TZPC 的 TZPCR0SIZE 寄存器动态驱动。TrustZone Protection Controller (TZPC, BP147) 是 APB 上的可配置信号控制块，含三组 TZPCDECPROT 寄存器（每组控制 8 路 SoC 信号）；其上电默认状态是所有 DECPROT 位为 0（全部 Secure）、TZPCR0SIZE 为 0x200（TZMA 覆盖的整个 2MB 范围全部 Secure），由启动代码按需放宽——这是一种默认安全的配置哲学（白皮书第 4 章）。\par
\noindent\textbf{（4）中断与 debug 锁定。}在推荐的 IRQ/FIQ 模型下，若中断到来时处理器恰好在对应世界，则无需经过 monitor 本地处理；否则先陷入 monitor 再路由。Debug 路径必须通过 CoreSight 认证接口锁定，TZPC 还可用于动态控制 SPIDEN 等 debug 使能信号，否则攻击者可绕过全部软件隔离直接观察 Secure world。\par
\subsubsection{安全性与威胁边界}
白皮书定义的是机制而非安全证明。硬件能保证的是：NS bit 标记的事务在总线、内存控制器与外设检查点被一致裁决；monitor 是唯一受控的世界切换入口。硬件不能保证的是：Secure world 内部软件（trusted OS、TA、monitor 本身）的正确性——一旦 trusted side 变复杂，TrustZone 的硬件边界无法自动修复软件漏洞，这一点后来被论文 22 的漏洞统计充分证实。此外，白皮书明确要求 SoC 集成者自查每条路径：第 7 章的设计 checklist 逐项追问 APB 外设安全配置、legacy FIQ 路由、Non-secure 软件尝试世界切换（SMC、中断、外部 abort）是否都被正确拒绝，说明厂商早已意识到``机制存在''与``集成正确''之间的鸿沟。\par
\subsubsection{实验环境与证据}
这是 vendor 白皮书，没有实验环境、没有 benchmark。证据源为 Arm 官方 PDF（本地已核验，p.1--p.109），可直接支撑的结论是：TrustZone 基础术语、双世界模型、NS bit 的 SoC 级安全状态传播、TZASC/TZMA/TZPC 的分区机制、monitor/中断/debug 控制面设计。不足以支撑的结论是：CCA/RME/RMM 的隔离语义、任何商业 TEE 实现的安全性、漏洞统计与性能开销。\par
\subsubsection{实验结果与证据强度}
白皮书没有任何性能表格，不能为其补造 latency 或 overhead 数字。可以讨论的只是定性权衡：把检查放进硬件路径（总线信号 + 地址译码）降低了软件逐次检查的复杂度与开销，TZASC 被设计为允许 burst 访问以最小延迟影响通过；但 Secure world TCB 规模与 world switch 的软件成本仍需引用具体 TEE OS/SoC 论文才能量化。\par
\subsubsection{综合评价}
综合来看，TrustZone 是 Arm TEE 的历史根基：双世界模型简单、部署广泛、SoC 集成度高，适合保护设备内密钥与安全服务，并长期支撑了移动与嵌入式生态。其局限同样清晰：双世界粒度粗，Secure world TCB 容易膨胀，TA/driver/API 层的漏洞会直接击穿安全服务（论文 22 的主题）；它也不提供云 confidential VM 所需的多租户隔离与远程证明语义。在本报告中，它的价值主要是作为机制定义的来源，以及与 CCA/RME 形成历史对照的基线。\par

\subsection{论文 21：Return-to-Non-Secure Vulnerabilities on ARM Cortex-M TrustZone}
\subsubsection{论文基本信息}
《Return-to-Non-Secure Vulnerabilities on ARM Cortex-M TrustZone: Attack and Defense》由 Zheyuan Ma、Xi Tan、Lukasz Ziarek、Ning Zhang、Hongxin Hu、Ziming Zhao（University at Buffalo、Washington University in St. Louis）完成，发表于 2023 年第 60 届 ACM/IEEE Design Automation Conference (DAC)。本文将其置于``Arm TrustZone TEE 与漏洞谱系''的研究脉络中考察，并把它作为 system/attack 类材料使用。它所回答的核心问题是：Cortex-M TrustZone 为性能引入的快速状态切换机制是否改变了权限传播语义——答案是肯定的，作者据此提出 ret2ns 攻击，可导致非安全态特权级任意代码执行。当前证据状态为：本地 PDF 已保存并核验；下文仅在该范围内讨论其机制与结论。\cite{ma2023ret2ns}\par
\subsubsection{研究背景与动机}
该工作的动机建立在一个规模事实上：Cortex-M 是最主流的 32 位微控制器架构，仅 2020 年第四季度就出货 44 亿颗基于 Cortex-M 的芯片（论文引 Arm 数据），广泛用于 Fitbit、Oculus VR、车载 ECU、手机蓝牙控制器等产品。与 Cortex-A 不同，Cortex-A 的世界切换必须经过特权 secure monitor mode 这一单一入口（SMC 指令），而 Cortex-M TrustZone 为降低开销允许通过函数调用与返回直接切换状态，入口数量不受限制。同时，Cortex-M 的安全状态划分是基于内存映射的：执行 secure 内存中的代码即处于 secure state；IPSR（中断程序状态寄存器）在两个状态间不分组（not banked），CONTROL 寄存器的 nPRIV 位则分组。论文 Figure 1 明确对比了这两种切换模型。这种``更快但语义不同''的设计此前未被系统研究，构成本文的切入点。\par
\subsubsection{关键挑战与核心思想}
\noindent\textbf{C1：刻画快速状态切换的攻击面。}Cortex-M 的跨状态控制流转移既可能发生在 handler mode 也可能发生在 thread mode，既可能通过间接返回（bxns）也可能通过间接调用（blxns），组合出四种情形。作者按``起始执行模式''和``转移指令''两个维度把攻击系统化，得到 ret2ns 的四个变种：handler-mode-originated 攻击多见于 RTOS，thread-mode-originated 攻击多见于支持特权分离的安全增强裸机系统（如 EPOXY、ACES）。\par
\noindent\textbf{C2：把语义差距转化为权限提升链。}Secure state 程序与 non-secure state 内存之间固有的语义差距导致 confused deputy 问题。此前的 boomerang 攻击（Cortex-A）只能借 secure 程序之手读写 non-secure 内核内存，不能直接获得提权代码执行；ret2ns 更进一步：它是 return-to-user (ret2usr) 类攻击在 TrustZone 上的新形态，把被污染的 secure state 代码指针重定向到 non-secure 用户态代码，而处理器仍以特权级继续执行。论文指出 ret2ns 影响所有带 TrustZone 的 Cortex-M 处理器，且任何``允许 secure state 直接向 non-secure 用户态转移控制流且保持特权级''的 TEE 实现都可能存在同类问题。\par
\noindent\textbf{C3：在没有 PXN 的硬件上防御。}x86 上防 ret2usr 依赖 SMEP，Arm 上对应机制是 Privileged Execute-Never (PXN)；但 PXN 只在规划中的 Cortex-M55/M85（ARMv8.1-M）上才有，市占量大的 M23/M33/M35P 均不支持，且 MPU 区域数只有 8 或 16 个，粒度不足以覆盖复杂 RTOS。防御必须另辟蹊径。\par
\subsubsection{系统架构}
攻击与防御都建立在 Cortex-M 的执行模型上：处理器有 thread 与 handler 两种执行模式，特权级由 IPSR 与 CONTROL 组合决定（IPSR 非零即 handler mode）；从非特权升特权需经 SVC 进入 handler。TrustZone 扩展引入 interstate branch 指令：blxns（间接调用）把 RETPSR 与返回地址压入 secure 栈并把 LR 置为 FNC\_RETURN（0xFEFFFFFF）；bxns（间接跳转）在目标地址非 FNC\_RETURN/EXC\_RETURN 时直接分支到 non-secure 代码。MPU 在两个状态间分组，secure state 可用 tta（test target alternate domain）指令查询 non-secure 地址的 MPU 访问权限——这是论文防御机制的关键指令。\par
\subsubsection{核心方法设计}
\noindent\textbf{（1）威胁模型。}攻击者是 non-secure state 用户态程序，通过正常系统调用路径（svc）与 secure state 程序交互；假设 secure state 固件或库（如 Arm TF-M）存在内存破坏漏洞，攻击者只需污染一个将被 bxns 或 blxns 使用的代码指针——不需要污染其他指针，也不需要（也不追求）在 secure state 执行任意代码；不假设 non-secure 内核本身有漏洞。\par
\noindent\textbf{（2）四变种攻击链与 walking example。}论文用一个 secure display 服务完整演示 handler-mode-originated + bxns 变种（Listing 1、Figure 3）：non-secure 用户程序调用库函数 print\_LCD()，经 svc \#0 进入 handler mode（IPSR=11）；SVC handler 调用 secure state 的 non-secure callable 函数 print\_LCD\_nsc()；由于 IPSR 不分组，secure state 程序仍在 handler mode 以特权执行。该 NSC 函数内 sprintf 向 128 字节栈缓冲区 buf 拼接用户消息，超长输入溢出并覆盖栈上保存的 LR；函数以 bxns 返回时，处理器从被污染的返回地址——攻击者在 non-secure 用户态的 attacker\_controlled()——继续执行，且仍处于 handler mode 特权级。论文强调被利用的内存破坏漏洞不必在 NSC 函数本身，任何能污染其栈帧保存 LR 的漏洞（如 format string）都可用。thread-mode 变种则利用 CONTROL\_NS.nPRIV 被清除后的特权 thread 执行路径，经 blxns 回调或 bxns 返回提权。\par
\noindent\textbf{（3）防御一：MPU-assisted address sanitizer。}作者首先形式化``非法转移''：两类控制流转移被禁止——(i) 从 secure handler mode 到任何 non-secure 用户态地址（无论 non-secure 特权级）；(ii) 从 secure thread mode 到 non-secure 用户态地址且 non-secure 处于特权级。防御在所有 NSC 函数尾声（bxns 之前）与所有 blxns 之前插桩（Listing 2）：依次读取 IPSR 与 CONTROL\_NS.nPRIV 判断转移是否合法；若存疑则用 tta 查询目标地址的 MPU 属性，经 MRVALID 位与区域号进一步读 MPU\_NS.RNR（地址 0xE002ED98）取得 RBAR 中的非特权读权限位；若目标区域允许非特权执行则判定为 ret2ns 并关中断停机。该方法不侵入 non-secure 软件，无需重编译。\par
\noindent\textbf{（4）防御二：address masking。}在 non-secure 用户态与内核态地址范围已知且分离的假设下，插桩简化为：检查 IPSR 与 CONTROL\_NS 后，对 bxns 的返回地址先做 EXC\_RETURN（0xFFFFFF**）特判，再做位掩码强制落入内核区域并写回栈（Listing 3）。EXC\_RETURN 特判是实测 ARMClang 生成代码所必需：non-secure handler 经 bx 调用 NSC 函数时 LR 与返回地址可能恰为 EXC\_RETURN。该方法更轻量但依赖更强的内存布局假设。\par
\subsubsection{安全性与威胁边界}
防御的安全性依赖三个前提：secure firmware 中所有 bxns/blxns 跳转点被完整识别并插桩；non-secure MPU 配置本身可信且正确实现了用户/内核分区（对 sanitizer）；地址布局假设成立（对 masking）。攻击面边界也清晰：ret2ns 不依赖篡改跨状态消息或冒充通信方，因此 REE-independent 调用者鉴别、ShieLD 等跨世界通信保护机制无法防御它；反之，这些机制与 ret2ns 防御正交。论文未声称覆盖编译器自动生成的全部跨世界代码路径或 RTOS 集成场景，这些被留作后续工作。\par
\subsubsection{实验环境与证据}
实验在两块真实/原型硬件上进行：(1) Microchip SAM L11 评估平台（Cortex-M23）；(2) 基于 ARM MPS2+ FPGA 原型板的 AN505 IoT kit 镜像（Cortex-M33）。作者修改 Microchip Studio IDE 与 Keil IDE 自带的示例 TrustZone 工程，在 NSC 函数中注入如 Listing 1 所示的内存破坏漏洞，验证四种攻击变种。防御评估在 Cortex-M33（配置为 20MHz）上进行；由于缺少面向 Cortex-M TrustZone 跨世界性能的现成 benchmark，作者改造 Keil 自带的 Blinky 跨世界工程（3 个 LED + 1 个 UART），用 SysTick 周期参数 T 模拟常规双向通信频率、用每次 NSC 调用前的 nop 数 N 模拟服务请求频率，取四组 (T, N) 组合，主循环执行 10 次、每组重复 5 次取平均；全部实验启用 non-secure MPU。项目已开源，并已向 Arm 负责任披露。\par
\subsubsection{实验结果与证据强度}
四种攻击变种在两个硬件平台上均被实验确认可实现提权任意代码执行；两种防御均能阻断全部四个变种。最坏路径（插桩全部指令执行）下，MPU-assisted sanitizer 对 bxns/blxns 分别增加 32/30 个 CPU 周期，address masking 分别增加 18/12 个周期。Table I 的端到端结果显示：在最密集的跨世界负载（T=$10^5$、N=$10^7$，约每 5ms 一次跨状态调用）下，两种防御的相对开销分别不超过 0.0381\% 与 0.0201\%；在最稀疏负载下低至 0.0004\% 与 0.0002\%。证据强度的边界同样明确：作者承认因 Cortex-M 安全固件生态尚新，未在量产系统中发现真实 ret2ns 漏洞，实验验证的是攻击可行性与防御有效性，不能据此推导所有 Cortex-M 产品均已暴露。\par
\subsubsection{综合评价}
综合来看，这篇工作的价值在于揭示了 Cortex-M TrustZone 为性能引入的直接状态切换并非只是实现细节，而是改变了权限传播语义：IPSR 不分组、CONTROL 分组、入口不受限这三件事叠加，使单一指针污染即可完成从 non-secure 用户态到特权执行的跨越。它把 confused-deputy 风险具体化为可执行的权限提升链，并给出两种可部署、开销可忽略（最坏 32 周期）的插桩式缓解。局限在于防御仍依赖跳转点的完整识别与 MPU/地址布局的可信假设；编译器生成代码、RTOS 集成与真实固件中的覆盖率仍待检验。在本节坐标系中，它是``机制设计缺陷''类的代表：问题不在软件 bug 本身，而在快速路径的权限语义。\par

\subsection{论文 22：SoK: Understanding the Prevailing Security Vulnerabilities in TrustZone-assisted TEE Systems}
\subsubsection{论文基本信息}
《SoK: Understanding the Prevailing Security Vulnerabilities in TrustZone-assisted TEE Systems》由 David Cerdeira、Nuno Santos、Pedro Fonseca、Sandro Pinto（Universidade do Minho、INESC-ID/IST Lisboa、Purdue University）完成，发表于 2020 年的 IEEE Symposium on Security and Privacy (S\&P)。本文将其置于``Arm TrustZone TEE 与漏洞谱系''的研究脉络中考察，并把它作为 sok 类材料使用。它所回答的核心问题是：TrustZone 提供了硬件隔离，为什么商用 TEE 系统仍然漏洞频出——答案是漏洞绝大多数不在 NS bit 失效，而在 secure world 软件把不可信输入当成可信上下文处理。当前证据状态为：本地 PDF 已保存并核验（全文 18 页）；下文仅在该范围内讨论其机制与结论。\cite{cerdeira2020trustzone}\par
\subsubsection{研究背景与动机}
该 SoK 的动机是一个规模与风险的反差：数亿乃至数十亿设备依赖 TrustZone TEE 保护 keystore、DRM、支付等敏感资产——论文引用估计称 Trustonic 的 TEE（Kinibi，前身 Mobicore）运行在约 17 亿台设备（多为三星）上，华为 2018 年售出 2 亿台手机——但针对 Qualcomm QSEE、Trustonic、华为 Trusted Core 的完整攻击链不断公开。作者要回答的问题是：这些漏洞是偶发实现失误，还是源于架构层面的系统性缺陷。\par
\subsubsection{关键挑战与核心思想}
\noindent\textbf{C1：闭源生态下的系统化研究。}多数商用 TEE 不公开源码，二进制难以获得且常有混淆，同一厂商还存在多代遗留版本并存，TrustZone 硬件又高度异构。作者的方法是五源并用（Table II）：CVE 数据库、三星 SVE 数据库、近十年安全会议论文（SP）、网络杂项报告（MR）、以及开源 OP-TEE 的源码与 changelog（SC，含对系统设计者的访谈）；并对三个代表性固件做逆向——三星 Galaxy S7 (Exynos) 固件 G930FXXS1APG3（Trustonic）、Pixel XL 固件 PQ2A.190205.003（Qualcomm）、华为 P8-Lite 系统镜像 ALE-L21C432B603（华为）——以量化各系统 TCB 并确认架构细节（如华为 monitor 基于 Arm Trusted Firmware (ATF)，高通为自研）。\par
\noindent\textbf{C2：把漏洞归一到组件与根因。}不同厂商、TEE OS 与版本的漏洞命名混乱。作者把 2013--2018 年间的报告手工归一化：在 207 份各来源报告中识别出 124 个 TEE 漏洞（Table III），按 CVSS 分为 critical（$\geq$9）、severe、medium、low 四档，再把每个漏洞映射到组件（secure monitor、TA、trusted kernel、boot loader 等）与根因类别，最终得到``架构—实现—硬件''三大类共 23 项编号问题（I01--I23）。\par
\noindent\textbf{C3：从个案到攻击链与设计教训。}SoK 的价值不止于分类，还在于说明小漏洞如何沿 TEE 组件串成高影响攻击链（Table I 的 E1--E6），并据此推导设计教训：缩小 TCB、收窄接口、隔离 TA、限制共享内存、强化输入校验。\par
\subsubsection{系统架构}
论文 Figure 1/2 给出 TrustZone-assisted TEE 的典型软件结构：normal world 应用经特权 daemon 调用 TrustZone driver 发出 SMC；secure world 侧自下而上为 secure boot loader（分 EL3 与 EL1 两段，实现信任启动）、secure monitor（五个被研究系统中四个基于 Arm ATF 参考实现）、trusted OS（内核）与若干 TA；安全外设（指纹、显示、SE 等）经 SW 驱动接入。分析对象覆盖 Qualcomm (QSEE)、Trustonic (Kinibi)、华为 (Trusted Core)、Nvidia 与 Linaro (OP-TEE) 五个现役商用/开源系统；研究原型与未规模部署的产品被排除。\par
\subsubsection{核心方法设计}
\noindent\textbf{（1）漏洞统计与严重度分布（Table III）。}六个系统 2013--2018 年共 124 个 CVE 类 TEE 漏洞：Qualcomm 92、Trustonic 5、华为 3、Nvidia 10、Linaro 3、其他 11；其中 critical 53 个（42\%）、severe 27 个。每个被分析的系统都至少有一个非 low 级漏洞。作者特别提醒：Qualcomm 占 74\% 不能解读为它最不安全，差异可能仅来自报告勤勉度；这些数字是漏洞量的下界。横向对照下，同期代码量大几个数量级的 Linux 有 1647 个 CVE，FreeRTOS/VxWorks 分别 13/10 个，且 TEE 的 critical/severe 占比更高——作者据此认为部分主流 TEE 的开发方法论落后于其他系统软件。\par
\noindent\textbf{（2）架构层问题（I01--I09）。}架构缺陷是最具设计教训价值的部分。I01：SW 驱动运行在 TEE 内核态（S-EL1）——Qualcomm 与 Linaro 为宏内核，Trustonic 与 Nvidia 采用微内核把驱动放到 S-EL0，华为把 TA 生命周期管理委托给特权 TA GlobalTask。I02：子系统间接口过宽——Android 至少四个 daemon 有 TrustZone driver 特权访问权；Trustonic TEE 有 32 个 TA；Widevine TA 实现 70 条命令且大量处理复杂媒体结构；Qualcomm TEE 向 TA 暴露 69 个系统调用且权限粒度粗（TA 可混杂访问全部 syscall）；Trustonic 的指纹驱动 TA 几乎对所有 TA 开放。I03：TCB 过大（Table IV）——Qualcomm TEE 核心二进制 1.61MB（monitor 96.2KB + QSEOS 1.50MB）外加 2.71MB TA（venus 924KB、fingerprint 600KB、Widevine 391KB、keymaster 332KB 等），Trustonic 的 TA 合计达 5.02MB；对照组 seL4 仅 166.5KB/19Kloc。I04：TA 可映射 NW 物理内存——Qualcomm 的 syscall 允许任意 TA 映射包括 REE 内核在内的 NW 物理页，构成 E6 攻击（经被控 TA 直接改写 Linux 内核）。I05：调试通道向 NW 泄密——华为 TEE 允许 TA 把栈轨迹 dump 到 NW 内存，泄露 GlobalTask 物理地址布局；Trustonic 的调试日志同样暴露 TA 内部信息。I06：ASLR 缺失或羸弱——Trustonic 所有 TA 加载于固定地址 0x1000、公共库 mcLib 固定于 0x7D01000；Qualcomm 虽有 TA ASLR 但 TA 全部加载于约 100MB 的连续物理区域，熵约仅 9 bit；五个系统均无 KASLR。I07：栈 cookie、guard page、执行保护不全（Table V）——Trustonic 无栈 cookie 且全局区与栈相邻排布，Qualcomm 有随机化栈 cookie 但无 guard page，华为（基于 Micrium $\mu$/OS）无栈 canary、无数据执行保护、.text 段未写保护。I08：缺少软件无关的 TEE 完整性报告——TrustZone 没有硬件机制向远程方报告软件测量，远程证明只能由 SW 软件实现，信任链含 EL3 全部软件。I09：TA 降级与撤销支持不良——TA 通常从 REE 文件系统加载，旧版已知漏洞的 Widevine TA 在 Qualcomm 与 Trustonic 上均被成功降级加载。\par
\noindent\textbf{（3）实现层缺陷分类（Table VI）。}实现 bug 分三类。Validation bugs 共 121 个，占绝对多数：TA 内 62 个（33.16\%）、trusted kernel 内 52 个（27.81\%）、secure boot loader 5 个、secure monitor 2 个。代表性案例：ATF 用于检测指针溢出的宏 check\_uptr\_overflow 自身存在缺陷（CVE-2017-9607），AArch32 下无法检测和落入 ($2^{32}$, $2^{64}-1$) 区间的溢出，波及多个 monitor 入口；华为 RTOSck 一个无任何输入校验的系统调用允许向 NW/SW 任意地址写（Listing 2）；Qualcomm TEE 内核完全缺少指针校验代码；启动加载器 X.509 证书解析器栈溢出可替换 TEE 镜像（CVE-2017-7932）；TA 校验漏洞可被系统性利用（如 CVE-2016-5349 经 boomerang 攻击直接提权至 Linux 内核）。Functional bugs：内存保护 32 个（如 ATF 翻译表配置错误使 S-EL1 只读区域恒可执行、OP-TEE ARMv7 FIQ 寄存器保存/恢复错误可致 REE 提权进 TEE）、外设配置 8 个（如 Qualcomm SPI 总线非独占访问 CVE-2016-10423 允许一个 TA 读另一个 TA 的 SPI 数据、OP-TEE 伪随机数生成器熵源不足）、安全机制 11 个（如 Amlogic S905 安全启动只验完整性不验签名、OP-TEE LibTomCrypt 的 Bellcore 攻击漏洞 CVE-2017-1000412 危及 RSA 私钥、RPMB 硬编码密钥泄露）。Extrinsic bugs：并发 11 个（如三星 TIMA 驱动两个竞态 SVE-2017-8974/8975、Nvidia DRM TA 的 TOCTOU CVE-2017-6296）、软件侧信道 4 个（如 LibTomCrypt 模幂优化泄露指数的时序信道 CVE-2017-1000413）。\par
\noindent\textbf{（4）硬件层问题（I19--I23）与微架构攻击。}I19：FPGA 可重构逻辑可破坏安全启动、NS bit 向 FPGA 传播的研究暴露六种攻击。I20：CLKSCREW 借恶意 non-secure 内核驱动超频超压诱导计算错误，击穿 TrustZone 边界提取密钥、绕过代码签名。微架构侧（Table VII）：cache 争用支撑 Prime+Probe 类攻击，ARMageddon、TruSpy 均恢复出完整 128 bit AES 密钥；分支预测器 BTB 在 NW/SW 间共享，Prime+Probe 式探测从 Qualcomm 硬件 backed keystore 完整恢复 256 bit 私钥；Rowhammer 在 Trusty TEE 上通过跨安全边界翻转位腐蚀 RSA 私钥并借错误签名推导密钥。\par
\noindent\textbf{（5）攻击链分析（Table I，E1--E6）。}以 QSEE 为例的完整链：E3 利用 Android Mediaserver 漏洞（CVE-2014-7920/7921）获得 TEE 接口访问；E2 利用 TrustZone Linux 驱动 bug（CVE-2014-4322）拿到 Linux root 与 SMC 能力；E1 利用 QSEOS 输入校验弱点实现 SW EL1 任意执行，完全控制 TZ 内核，进而可劫持 TA 提取密钥破解 Android 全盘加密或解锁 bootloader。第二条链只需 TA 接口：E4 经 Widevine TA 提权漏洞（CVE-2015-6639）控制 TA，再经 E5 利用 Qualcomm trusted OS 系统调用缺失校验（CVE-2016-2431）写入内核任意地址劫持 TEE 内核；E6 则借被控 TA 与``TA 可映射 NW 物理内存''的系统调用反向夺取 Linux 内核。这组链说明：normal world 立足点不必直接拿到密钥，它只是进入 TEE API 的起点。\par
\subsubsection{安全性与威胁边界}
论文的攻击者模型限定在 NW：攻击者目标为窃取 TEE 或 REE 秘密、提权至 REE 内核或 TEE；访问 SMC 接口的途径是 N-EL1 代码执行（直接构造 SMC）或 N-EL0 应用经 TA 命令间接触达；所有 NW 组件均不可信。有效性威胁作者写得很坦白：多数漏洞无公开 PoC、CVE 描述信息不足，分类可能不精确；闭源系统信息不全；某一系统报告量远超其他系统时存在过度代表风险；且只分析已披露漏洞，未知漏洞类型可能存在。边界上，本文证据只支撑 TrustZone 漏洞谱系本身，不能据此推导 CCA/RME 一定存在相同漏洞——它只能作为迁移动机与设计教训的来源。\par
\subsubsection{实验环境与证据}
这是同行评议 SoK，``实验''是公开漏洞语料库（207 份报告、124 个 CVE 类漏洞）加三个商用固件的逆向分析，而非新 TEE 实现。证据源为 IEEE S\&P 2020 正式版本 PDF（本地核验，p.1--p.18），核心可引用对象为 Figure 1/2（架构）、Table I（QSEE 攻击链）、Table II/III（语料来源与严重度分布）、Table IV（TCB 规模）、Table V（内存保护矩阵）、Table VI（实现 bug 分类）、Table VII（微架构攻击）及 I01--I23 编号问题。\par
\subsubsection{实验结果与证据强度}
论文不提供任何 latency/throughput 结论，其 claim 强度分布为：漏洞分类与组件绑定强（每个编号问题都有公开报告或逆向证据），攻击链复现描述强（E1--E6 均有原始 exploit 文献），聚合趋势判断中等（作者自行加注报告勤勉度偏差），跨系统安全水平横向比较弱（明确不可做），性能维度无证据。商业风险的结论主要来自 patch lag、闭源 TCB 与厂商私有 TA 生态三者叠加。\par
\subsubsection{综合评价}
综合来看，这篇 SoK 是 TrustZone 方向最重要的``反面教材''：它用 124 个漏洞与 23 项编号问题证明，硬件隔离必须与小微内核化 TCB、窄接口、完整内存保护、可靠撤销机制和严格输入校验配套，否则 secure world 只是另一个会被攻破的内核。其优势在于把真实漏洞绑定回架构组件，``错在哪里''一目了然；其局限在于依赖公开披露，闭源系统与未公开漏洞覆盖不足，横向比较有系统性偏差。从部署角度看，它可直接转化为 TEE 审计 checklist、TA 沙箱化需求与向 CCA/RME 迁移的论据。在本节坐标系中，它与论文 21 形成互补：一个统计``软件层失守''的谱系，一个深挖``机制层快速路径''的缺陷。\par

\subsection{开发商安全实践建议}
本章四篇材料分别提供了框架（\cite{li2024sokteechoices}）、机制定义（\cite{arm_trustzone_whitepaper}）、机制级攻击（\cite{ma2023ret2ns}）与漏洞谱系（\cite{cerdeira2020trustzone}）。面向 TEE 系统开发商、芯片厂商与固件团队，可提炼出以下五条可操作建议。\par
\noindent\textbf{实践一：用管理模式矩阵审计运行时资源管理面。}参照 TRAF 的事件—模式坐标（\cite{li2024sokteechoices}），为自有 TEE 产品建立一张``CPU scheduling、context switch、interrupt/instruction emulation、memory allocation、page fault、I/O''逐事件的责任矩阵，逐项标明由 host、RTPM/monitor 还是实例 runtime 管理。审计重点放在 RTPM-guarded 模式的校验完备性上：该 SoK 的 Table 1 表明未保护寄存器、TLB 误用、未保护 NPT 等缺陷多来自 guarded 模式的 corner case 遗漏。每个 guarded 检查点都应能回答``host 能否通过 present bit、中断注入、I/O emulation 或共享内存影响 secret-dependent 行为''。\par
\noindent\textbf{实践二：以微内核化和接口收缩控制 secure world TCB。}Cerdeira 等人的 TCB 测量（\cite{cerdeira2020trustzone}）显示，采用宏内核、驱动跑在 S-EL1 的系统（Qualcomm 1.61MB 内核 + 2.71MB TA）漏洞暴露面显著大于微内核架构（Trustonic、Nvidia 将驱动置于 S-EL0）。开发商应：(i) 把驱动与非核心服务移出 TEE 内核态；(ii) 压缩单个 TA 的命令面（反面教材：Widevine TA 70 条命令）与 TEE 系统调用面（反面教材：69 个混杂授权的系统调用）；(iii) 对 TA 间服务调用实施细粒度授权并控制白名单规模（反面教材：TIMA 驱动 TA 的 34 项硬编码白名单）。\par
\noindent\textbf{实践三：把内存保护与生命周期管理纳入 TA 强制基线。}依据漏洞谱系（\cite{cerdeira2020trustzone}），TA 输入校验缺陷占全部实现 bug 的最大份额（62/187 个实现类 bug），而 ASLR、栈 cookie、guard page、执行保护在各商用 TEE 上普遍缺失或羸弱。SDL 基线应强制：全量 TA 与内核的高熵 ASLR/KASLR（避免 Qualcomm 约 9 bit 熵的反面案例）、栈 cookie 与 guard page、WXN/UXN/PXN 执行保护、调试通道出厂关闭（避免栈轨迹泄露物理布局）、TA 版本防降级与撤销机制（避免 Widevine 降级重放）。\par
\noindent\textbf{实践四：对 Cortex-M 产品执行跨世界控制流切换专项检查。}凡使用带 TrustZone 的 Cortex-M（M23/M33/M35P 等）的固件团队，应把 ret2ns 检查纳入发布门禁（\cite{ma2023ret2ns}）：枚举固件中全部 bxns/blxns 跳转点（含编译器生成的 NSC 尾声），在跳转前校验 IPSR、CONTROL\_NS.nPRIV 与目标地址的非安全 MPU 属性；无法完整插桩时至少在地址布局已知的路径上启用 address masking；迁移到 M55/M85 后应配置 MPU 的 PXN 属性从硬件层禁止特权级执行用户页。同时注意该检查不假设 non-secure 内核无漏洞，secure 固件中的任何内存破坏漏洞（含 NSC 调用链下游函数）都可能成为入口。\par
\noindent\textbf{实践五：将 NS bit 传播完整性作为 SoC 集成的签收项。}TrustZone 白皮书（\cite{arm_trustzone_whitepaper}）第 7 章的 checklist 本质上是一份集成验收表：每条 master/slave 路径是否携带并检查 NS bit；TZASC 区域划分、TZPC 上电默认（全 Secure）是否被启动代码按最小权限放宽；APB/AHB 桥的安全裁决是否覆盖全部外设；DMA 通道、ACP 一致性端口、CoreSight debug 认证信号（SPIDEN/SPNIDEN/DBGEN）是否锁定；legacy FIQ 路由与 Non-secure 世界切换尝试（SMC、中断、外部 abort）是否都被正确拒绝。芯片厂商应把这份清单形式化为流片前验证项，固件团队应在每次启动配置变更后回归验证。\par

\subsection{未来研究方向}
\noindent\textbf{方向一：面向 trusted I/O 与设备直通扩展 TRAF 坐标。}问题：TRAF 的 I/O 维度在 2024 年截面上仍主要落在 unprotected 模式（\cite{li2024sokteechoices}），而 CoVE-IO、TDISP、PCIe IDE、SPDM 与 Arm CCA device assignment 正在快速改变这一格局。现有工作未解决：新机制的管理主体（设备侧可信逻辑、IDE 密钥、RMM）无法落入原四模式。技术路线：把设备队列、设备内存、固件、驱动 TCB 映射为 host/RTPM/device/instance 四方分工，扩展事件—模式矩阵。验证实验：对至少两个已实现的 trusted I/O 路径（如 CoVE-IO 原型与 CCA RME-DA 文档）做模式标注，并检查矩阵能否预测其已知缺陷。\par
\noindent\textbf{方向二：TrustZone 漏洞的自动化发现与可复现基准。}问题：Cerdeira SoK 依赖公开披露语料，存在系统性偏差，闭源 TCB 的真实漏洞密度未知（\cite{cerdeira2020trustzone}）；ret2ns 作者也未在量产固件中发现真实漏洞样本（\cite{ma2023ret2ns}）。现有工作未解决：缺少面向 TA 接口、共享内存与跨世界控制流的统一 fuzzing/静态分析基准。技术路线：以 I01--I23 编号问题为标签体系构建漏洞基准，结合固件重托管对 SMC/NSC 接口做结构化 fuzzing。验证实验：在 OP-TEE（开源可控）与至少一个商用固件镜像上对比工具发现率，并人工确认误报/漏报。\par
\noindent\textbf{方向三：Cortex-M 跨世界控制流的编译器级与形式化防护。}问题：ret2ns 防御依赖人工或插桩识别全部 bxns/blxns 跳转点，编译器自动生成代码（如 ARMClang 的 EXC\_RETURN 复用）已被证明会引入边界情形（\cite{ma2023ret2ns}）。现有工作未解决：没有编译器 pass 或形式化模型能保证``任何跨状态转移不携带越权特权''这一性质的完备性。技术路线：在编译器 IR 层标注安全状态与特权级，插入经验证的 sanitizer；或对状态切换语义（IPSR 不分组、CONTROL 分组、EXC\_RETURN/FNC\_RETURN 协议）建形式模型并做模型检验。验证实验：对 TF-M 等真实固件全量插桩后测量覆盖率与开销，并对四变种攻击做回归。\par
\noindent\textbf{方向四：从双世界到多隔离域的漏洞谱系迁移评估。}问题：Cerdeira 的 23 项编号问题中，哪些是 TrustZone 双世界模型固有（如过宽接口、TA 可映射 NW 物理内存、缺少硬件证明），哪些会被 CCA/RME 的 granule 委托、attestable lifecycle 与更小管理面消除（\cite{cerdeira2020trustzone,arm_trustzone_whitepaper}），论文未说明。现有工作未解决：缺少把 TrustZone 漏洞类逐一映射到新架构的系统性核查。技术路线：以 I01--I23 为检查表，对照 CCA RMM 规范与参考实现逐项判定``消除/缓解/仍存在/不适用''。验证实验：在 Arm FVP 上对``仍存在''项构造概念验证，区分架构结论与实现结论。\par
\noindent\textbf{方向五：模式感知的远程证明语义。}问题：当前 attestation 只证明初始镜像与配置，不表达运行时资源管理模式与被排除的攻击类（\cite{li2024sokteechoices}）；TrustZone 侧则连软件无关的完整性报告都缺失，只能依赖 SW 软件自证（\cite{cerdeira2020trustzone}）。现有工作未解决：verifier 无法从证据中判断``page fault 由谁仲裁、I/O 是否离开受保护路径''。技术路线：把 TRAF 模式标注编码进证明证据格式，使 verifier 策略可按模式拒绝高风险配置。验证实验：在一个开源 TEE（如 Keystone/OP-TEE）上扩展证据格式与 verifier，展示对两个模式不同但镜像相同的配置的区分能力。\par
\noindent\textbf{方向六：闭源 TEE 的 TCB 度量与漏洞情报持续化。}问题：Cerdeira 的 TCB 测量（Table IV）依赖一次性逆向，无法跟踪版本演进；2020 年之后的漏洞谱系缺少同等严谨度的更新（\cite{cerdeira2020trustzone}）。现有工作未解决：闭源二进制 + 混淆 + 多代并存使自动化 TCB 度量与漏洞归因困难。技术路线：发展面向固件镜像的组件识别与跨版本 diff 技术，建立可持续更新的 TEE 漏洞与 TCB 数据集。验证实验：对同一厂商连续多个固件版本做 TCB 规模与接口面（syscall/TA 命令数）追踪，并与公开 CVE 时间线交叉验证。\par
\section{启动、远程证明与生命周期安全}
证明机制的目标不是简单地输出一份签名，而是让验证方能够把设备身份、受保护代码、运行状态和挑战新鲜性联系起来。本节沿着软件 attestation、启动测量、远程证明和运行时证据的演进展开：先以一篇跨平台综述建立比较维度（论文 23），再以 IETF RATS 架构固定角色与数据流词典（论文 24），随后检查一个把硬件/软件协同证明做到机器可验证的低端设备方案（论文 25），最后落到证据与传输通道绑定这一部署环节（论文 26）。对每篇工作，我们都逐项检查 verifier 实际能据此判断什么、仍然无法判断什么。\par
特别需要避免把“得到证明报告”直接等同于“完整隔离”：报告的覆盖范围取决于测量根、密钥生命周期、证据绑定对象和传输路径。一份签名正确的 quote 既不能证明测量覆盖了全部运行态，也不能证明接收 quote 的通信端点就是产生 quote 的那个 TEE 实例；这两个缺口分别对应本节论文 25 讨论的实现级绕过（DMA、中断、密钥残留）和论文 26 讨论的 relay 攻击。各论文将在这一问题框架下讨论其协议设计、形式化假设与部署限制。\par

\subsection{论文 23：An Exploratory Study of Attestation Mechanisms for Trusted Execution Environments}
\subsubsection{论文基本信息}
《An Exploratory Study of Attestation Mechanisms for Trusted Execution Environments》由 Jämes Ménétrey、Christian Göttel、Marcelo Pasin、Pascal Felber 和 Valerio Schiavoni（University of Neuchâtel）完成，发表于 2022 年的 5th Workshop on System Software for Trusted Execution (SysTEX，与 ASPLOS 同期），正文共 7 页（arXiv:2204.06790v2）。本文将其置于“远程证明、启动与生命周期安全”的研究脉络中考察，并把它作为 survey 类材料使用。它所回答的核心问题可概括为：Intel SGX、AMD SEV、Arm TrustZone 与新兴 RISC-V 方案各自的远程 attestation 机制在哪些维度上存在结构性差异，验证方据此能建立何种信任。当前证据状态为：本地 PDF 已保存并核验；下文仅在该范围内讨论其机制与结论。\cite{menetrey2022attestation}\par
\subsubsection{研究背景与问题定义}
该工作的问题定义并非抽象的安全愿望，而是一个具体的系统矛盾：多平台 TEE 的 attestation 机制各异，缺少统一框架来比较证据语义、freshness 和 verifier policy。论文（第 2 节）首先区分本地证明与远程证明：本地证明发生在同一硬件上的两个软件环境之间，依赖一个硬件绑定的 root of trust 派生签名密钥，其结果可用于在两个环境之间建立新的秘密（如安全通道）；远程证明则跨设备完成，论文明确采用 IETF 术语——relying party 借助 verifier 与 attester 建立信任关系，attester 收集一组 claims（例如应用代码的密码学哈希测量值）并签名形成 evidence，verifier 断言或否认该 evidence 之后，relying party 才决定交互或传输机密数据。\par
论文进一步把已有远程证明实现分为三类（第 2 节）：software-based attestation 不依赖特定硬件、适配低成本设备；hardware-based attestation 依赖防篡改硬件（如 TPM）或 PUF（利用制造差异产生的唯一硬件标记防假冒），也包括把片上熔断密钥（die-fused secret）仅暴露给可信环境的做法；hybrid 方案结合两者优点。此外论文指出，可信应用常需要双向保证——例如从传感 IoT 设备取回机密数据时，设备要认证远程方、远程方也要确认设备未被伪造或篡改——因此 mutual attestation 被单列为一项比较特性。这套“角色—claims—证据类型—双向性”框架解释了为何签名本身不能替代完整的信任决策。\par
\subsubsection{核心方法设计}
这一部分的核心结论是：按 attestation 原则、平台实现、证据格式和验证策略建立 taxonomy，覆盖 SGX、SEV、TrustZone 与 RISC-V。论文第 3.1 节提出一组“基石特性”（cornerstone features）作为比较框架，每项可为缺失、部分满足或完全满足：内存完整性（防止 TEE 实例 DRAM 被篡改的主动机制）、新鲜性（防重放与回滚）、加密、可并发的隔离域数量（部分满足表示域数量有上限）、开源程度、本地证明、远程证明（部分满足表示无内建支持、由文献扩展）、应用可用的证明 API、双向证明、用户态支持、工业成熟度（区分量产工业 TEE 与学术原型），以及隔离和证明的粒度与隔离的系统支撑机制。Table 1 依此对 SGX、TrustZone、SEV 三个版本（vanilla/SEV-ES/SEV-SNP）和 Keystone、Sanctum、TIMBER-V、LIRA-V 打分；该框架的作用不是排出绝对名次，而是把不同 TEE 的安全属性、可部署性和证明能力拆到同一组可观察维度中。\par
论文第 3.2 节给出可信应用的通用部署与证明流程（Figure 1）：应用先在开发方安全场所被编译和测量，随后被传输到不受信任的部署环境并在 TEE 中执行；可信应用与 verifier 建立受信通道时，TEE 环境向其暴露一份 evidence，应用把自己的密钥材料附加进 evidence 以防攻击者窃听通信；verifier 将 evidence 与参考值列表比对，识别可信应用的正品实例。这条流程揭示了实际部署中的关键连接点：代码测量、设备根信任、证据签名、通道密钥以及 verifier 策略必须覆盖同一实例，任一环节未绑定都可能被拼接攻击利用。\par
\noindent\textbf{（1）Intel SGX：enclave、quoting enclave 与 EPID/DCAP。}\par
SGX（2015 年 Skylake 起引入）以 enclave 为隔离粒度：启动时保留的 Enclave Page Cache (EPC) 存放加密的 enclave 代码与数据，OS 和 hypervisor 均不可访问，CPU 与内存之间的流量由 Memory Encryption Engine (MEE) 保密，EPC 同时存放校验码以检测外部软件对内存的修改。本地证明由 \texttt{EREPORT} 指令产生 report，内含身份、属性、TCB 可信状态、目标 enclave 附加信息和 MAC；远程证明则依赖内建的 quoting enclave——它通过 \texttt{EGETKEY} 独占访问设备特定私钥，先与应用 enclave 做本地证明，验证 report 后用设备私钥签名替换 MAC，形成可被远程验证的 quote。quote 可以把应用公钥等附加信息绑定进证据，从而同时证明“正品 Intel SGX 处理器”和“加载进 enclave 的应用测量值”。远程协议基于 SIGMA 并扩展为 Enhanced Privacy ID (EPID)：EPID 不标识唯一实体而标识签名者所属群组，quote 用与处理器固件版本绑定的 EPID 密钥签名，验证方用群公钥和撤销信息或 Intel attestation service 验证；数据中心场景另有支持第三方证明的 DCAP。论文 Figure 3 给出完整流程：verifier 提交带 nonce 的 challenge，应用 enclave 准备 manifest（一组 claims）与用于回传机密信息的公钥，manifest 哈希作为本地证明 report 的附加数据，quoting enclave 验证后签发 quote。论文同时指出 SGX 的局限：quoting enclave 与微码闭源（已有形式化分析工作），驱动与 SDK 开源；MAGE 等工作进一步扩展出无可信第三方的 enclave 群组双向证明。\par
\noindent\textbf{（2）Arm TrustZone：无内建证明，靠安全启动与根密钥补足。}\par
TrustZone 把处理器分为 secure world（TEE）与 normal world，通过 SMC 指令切换，secure world 运行 OP-TEE 等可信 OS 并以隔离进程方式运行 Trusted Application (TA)。尽管商业上广泛部署，TrustZone 缺乏内建 attestation 机制，relying party 无法直接远程验证 TrustZone TEE 的状态；文献中的单向与双向证明协议因此都要求额外条件：secure world 内具备 root of trust 的额外硬件、用于密码操作的安全随机源，以及安全启动机制。设备可用片上熔断密钥作为 root of trust 派生证据签名私钥，安全启动则逐级测量启动阶段、阻止被篡改系统启动。论文以 Shepherd 等人的方案为案例：secure world 内的 Trusted Measurer (TM) 组件由安全启动认证，基于 OS 与 TA 状态生成证据，私钥被 provisioning 后封存（seal）在 TEE 安全存储中、类似 firmware TPM 的使用方式；双向证明分三轮——attester 发起含双方身份与密钥材料的握手，verifier 应答并附上基于 Diffie-Hellman 共享秘密签名的己方 TEE 证据，attester 再回送基于同一共享秘密签名的己方证据，双方确认证据为真后派生共同秘密建立受信通道。\par
\noindent\textbf{（3）AMD SEV：以虚拟机为粒度的启动测量与运行时报告。}\par
SEV 以虚拟机为隔离与证明粒度：嵌入式 AES 引擎配合多密钥加密内存，每 VM 与 hypervisor 被分配带 Address Space Identifier (ASID) 标签的密钥以防跨 TEE 攻击，处理器封装外的代码与数据以 128 位密钥 AES 加密；密码材料由集成的闭源 Arm 安全协处理器生成并保存在 CPU 内。证明的核心是 Chip Endorsement Key (CEK)——AMD 签发、熔断在处理器 die 上的秘密。启动时，AMD 安全处理器把 VM 内容密码学测量为 launch digest（claim），SEV-SNP 还测量内存页元数据，使 digest 同时覆盖初始客座内存的布局。论文 Figure 4 给出六条 launch 命令构成的流程：\texttt{LAUNCH\_START} 建立带客座属主公钥的 guest context，\texttt{LAUNCH\_UPDATE\_DATA} 在 hypervisor 装载 VM 时加密内存并累计测量，\texttt{LAUNCH\_UPDATE\_VMSA}（仅 SEV-ES）初始化寄存器保存区，\texttt{LAUNCH\_MEASURE} 产出带签名的测量值供客座属主核对，核对通过后 \texttt{LAUNCH\_SECRET} 注入镜像解密密钥等敏感数据，最后 \texttt{LAUNCH\_FINISHED} 允许 VM 执行。三个版本的证明能力递进：SEV 与 SEV-ES 只支持客座 OS 启动阶段的远程证明，SEV-SNP 则引导私有通信密钥，允许 guest VM 在任意时刻请求 attestation report 并获取用于数据封存（sealing）的密码材料；SEV-ES 通过加密 VMCB 中的寄存器状态（VMSA）并经 GHCB 共享内存与 hypervisor 通信，修复了 SEV 在中断时经寄存器向 hypervisor 泄漏信息的问题；SEV-SNP 针对的是恶意云厂商物理接触机器、安装恶意固件读取系统的回滚攻击。论文同时指出代价：整 VM 保护免去了 SGX 式应用拆分，但 TCB 更大、攻击面更宽，guest OS 需支持 SEV，不能使用 PCI passthrough，且第一版 SEV 最多支持 16 个虚拟机。\par
\noindent\textbf{（4）RISC-V：PMP/MPU/标签内存上的四种学术取舍。}\par
RISC-V 已批准的 Physical Memory Protection (PMP) 指令提供了与前述技术类似的内存保护能力，论文覆盖四个带远程证明的代表性设计。Keystone 在 M-mode 实现 secure monitor、完全基于 PMP 而不改动硬件，允许自选安全原语（内存加密、动态内存管理、cache 分区）；其安全启动测量 secure monitor 镜像、生成 attestation key 并用硬件可见秘密（root of trust）签名，secure monitor 经 SBI 向 enclave 暴露证据签发接口，证据覆盖 secure monitor、runtime 与应用三层测量，且允许附加任意数据以与 verifier 建立标准安全通道。Sanctum 用两个开源组件——measurement root (mroot) 与 secure monitor——替代 Intel 不透明的微码，设有专用签名 enclave 从 secure monitor 获得派生私钥以签发证据，证据 claims 含请求方 enclave 代码哈希与密钥交换消息信息；后续工作还加入安全启动与基于 PUF 的免 root of trust 身份派生。TIMBER-V 面向小型嵌入式 RISC-V，用硬件辅助内存标签隔离执行，secure supervisor mode 的 TagRoot 管理标签、secure user mode 可承载多个隔离 enclave，配合 MPU 支持任意数量进程；其证据基于 enclave 身份、平台秘密密钥与应用提供的任意标识符，但用 MAC 认证——对称密码在远程证明中不常见，因为 verifier 必须持有同一秘密密钥，作者称可扩展为公钥密码。LIRA-V 面向受限边缘设备的双向证明，不允许在 TEE 内执行任意代码，claims 由 ROM 中的程序对物理内存区域计算，协议与 TrustZone 案例类似、需预置密钥作为 root of trust。论文明确排除了缺乏远程证明机制的 RISC-V TEE（Hex-Five MultiZone、HECTOR-V、Lindemer 等的 PMP 线程隔离方案），以及无 RA 的 IBM PEF 和写作时尚未落地的 Intel TDX 与 Arm CCA Realm。由此可以看出，隔离模型和证明接口并不天然一一对应。\par
\subsubsection{安全性与威胁边界}
从安全边界看，证明机制必须同时回答证据由谁产生、是否新鲜、是否绑定到正确的实例和端点，以及 verifier 据此能做出什么决策。论文的比较框架本身提示了若干边界：测量值本身并不等同于运行期隔离（内存完整性、新鲜性、加密在 Table 1 中是三个独立特性）；TrustZone 类无内建证明的平台，其证据安全性强依赖安全启动与根密钥 provisioning 的正确性；TIMBER-V 式对称 MAC 证据把秘密共享给 verifier，密钥分发性本身构成新的信任假设。论文未声明覆盖物理、侧信道或拒绝服务攻击，也未评估任何平台证据格式对 CCA/CoVE 等后续架构的适用性。\par
\subsubsection{实验设计与证据分析}
从评估设计看，这是一篇七页的综述性工作，不提出新的 TEE 原型或性能基准，因此不存在 benchmark、开发板或模拟器实验。其主要证据是三件可核验材料：Table 1 的跨平台功能矩阵（13 个特性维度、9 个平台/方案）、Figure 1 的通用部署与证明流程图，以及第 3.3–3.6 节对 SGX、TrustZone、SEV 和四种 RISC-V 方案的机制归纳（含 SGX Figure 3、SEV Figure 4 两条具体协议流程）。因而它能够支持的是比较维度和研究谱系层面的结论，而不是某个平台在特定 workload 上的延迟、吞吐或端到端安全结论；凡涉及具体平台性能或部署状态的判断，须回到相应平台的一手规范核对。\par
\subsubsection{综合评价}
综合来看，这是当前跨平台 TEE attestation 覆盖最系统的学术 survey 之一。它的贡献不在于给出新的证明协议，而在于将证明能力从处理器营销特性中拆解出来，说明成熟 TEE 之间仍可能在新鲜性、互证、应用 API、隔离粒度和工业可用性上存在结构性差异：SGX 提供完整内建远程证明协议，SEV 面向虚拟化环境提供启动测量与（SNP 起）运行时报告，TrustZone 与原生 RISC-V 则把证明留给社区用软件方案补足。它适合作为本报告中 Arm、RISC-V 和云端 TEE 证明机制的比较起点；但由于写作时间早于 CCA、CoVE 等后续架构的成熟（论文明确将 Realm、TDX 列为省略对象），具体接口和现行部署状态仍需回到相应规范与原始系统论文核对。本节判断以本地已核验 PDF（arXiv:2204.06790v2，第 1–4 节及 Table 1、Figure 1/3/4）为依据；超出这些材料覆盖范围的实现细节、性能结论或攻击面，本文不作延伸推断。\par

\subsection{论文 24：Remote ATtestation procedureS (RATS) Architecture}
\subsubsection{论文基本信息}
《Remote ATtestation procedureS (RATS) Architecture》由 Henk Birkholz、Dave Thaler、Michael Richardson、Ned Smith 和 Wei Pan 完成，2023 年 1 月作为 IETF RFC 9334 发布。本文将其置于“远程证明、启动与生命周期安全”的研究脉络中考察，并把它作为 spec 类材料使用。它所回答的核心问题可概括为：如何把 TPM、TEE、机密虚拟机与嵌入式设备形态各异、术语不通用的 attestation 统一为一套角色、制品与拓扑语言。当前证据状态为：本地 RFC PDF 已保存并核验；下文仅在该范围内讨论其机制与结论。\cite{rats_rfc}\par
\subsubsection{研究背景与动机}
要理解该架构，首先需要看到它面对的系统矛盾：attestation 经常被误解为“单个签名报告”——拿到一份带签名的证据就当作信任建立完毕。RATS 的基本立场是，Evidence 只说明 attester 声称了什么，并不自动等于可信：Verifier 还必须取得 endorsements（背书）、reference values（参考值）并套用 appraisal policy 才能裁决；Relying Party 拿到 attestation result 后还要按自己的策略决定是否授予访问。RFC 第 2 节给出七个参考用例（网络端点评估、机密机器学习模型保护、机密数据保护、关键基础设施控制、TEE provisioning、硬件 watchdog、FIDO 生物认证），说明这套多方模型必须覆盖从数据中心到受限设备的跨度。\par
\subsubsection{关键挑战与核心思想}
在上述背景下，RATS 提出的关键判断是：信任决策不在 attester，而在 verifier 的 appraisal policy。同一份 Evidence 对不同 relying party 可能得出不同结论，因为裁决是“证据×背书×参考值×策略”的函数，而非证据本身的属性。其中，Endorsement 是 Endorser（通常是制造商）为 attester 各项能力（如 claims 收集与证据签名的完整性）作保的安全声明，用于证明证据签名密钥或组件来源；Reference Values 是用于与 claims 比对的基准值（其他文档中称 known-good values、golden measurements 或 nominal values，RATS 强调比对不限于相等性判断）。这一分工把“设备是否为真”“测量值是否属于可接受基线”“结果是否满足本应用的安全要求”拆成三个可独立演进的问题。\par
\subsubsection{系统架构}
从系统结构看，RATS 把证明过程拆成 evidence generation、conveyance、appraisal 和 result use 四个环节，其角色与制品定义（RFC 第 4 节）带有明确的产生/消费关系。角色方面：Attester 产生 Evidence；Verifier 消费 Evidence、Reference Values、Endorsements 与 Appraisal Policy for Evidence，产出 Attestation Results；Relying Party 消费 Attestation Results 并套用 Appraisal Policy for Attestation Results 做应用层决策；Verifier Owner 与 Relying Party Owner 分别有权配置两类策略；Endorser（典型为制造商）产生 Endorsements；Reference Value Provider 产生 Reference Values。制品方面：Claim 是断言信息（常为名值对），Evidence 是 attester 生成的一组 claims（可含配置数据、测量、遥测或推断），Attestation Result 是 verifier 为结果有效性作保的输出。Attester 可以是 device、TEE、CVM、accelerator 或 service；RFC 第 3.1 节进一步把 Attester 内部拆为至少一个 Attesting Environment 与至少一个 Target Environment（收集证据的可信环境与被执行测量的环境），并讨论分层证明（layered attestation）与复合设备（composite device）两种组合形态。Verifier 可以本地或远端部署；RFC 中立于 CPU 架构、claim 内容与传输协议。\par
\subsubsection{核心方法设计}
RATS 的关键不是单一组件，而是把角色、制品、拓扑与时序组合成一套可落地的架构约束。\par
\noindent\textbf{（1）方法 1：Conceptual Data Flow（Figure 1）。}\par
这一机制的直接作用是：证明不是二方对话，而是多角色数据流。Figure 1 描绘的完整回路是：Endorser 与 Reference Value Provider 分别把 Endorsements 和 Reference Values 注入 Verifier，Verifier Owner 配置 Appraisal Policy for Evidence；Attester 把 Evidence 交给 Verifier 完成“appraisal of Evidence”并产出 Attestation Results；Relying Party 在 Relying Party Owner 配置的 Appraisal Policy for Attestation Results 约束下完成“appraisal of Attestation Results”，据此做出授权等应用决策。两类策略都可经协议下发、静态配置或编程固化。这张图是全文其余章节（拓扑、信任模型、conceptual messages）的公共参照系。\par
\noindent\textbf{（2）方法 2：Passport 与 Background-Check 两种拓扑（第 5 节）。}\par
这一机制的直接作用是：RATS 支持不同部署模型，取决于 verifier 是否在线参与 relying-party 交互。Passport 模型（Figure 5）中，Attester 先把 Evidence 交给 Verifier 换取 Attestation Result，将其视为不透明数据（可缓存）并出示给一个或多个 Relying Party；如同护照——公民是 Attester，签发机构是 Verifier，护照是 Attestation Result，边检柜台是 Relying Party。该模型有三种失败方式：Evidence 未通过 Verifier 策略、Result 未通过 Relying Party 策略、Verifier 不可达；且由于 Attestation Result 随资源访问协议在 Attester 与 Relying Party 之间传递，其序列化格式的互操作与标准化比 Evidence 格式更为关键。Background-Check 模型（Figure 6）中，Attester 把 Evidence 交给 Relying Party，后者不解析、原样转发给可信 Verifier，再按返回的 Attestation Result 决策；如同雇主向前雇主核实候选人履历。此时 Evidence 只需能被封装进资源访问协议，Relying Party 无需为其准备解析器——对受限节点可显著缩小代码体积与攻击面。第 5.3 节还给出组合形态：Relying Party 与 Verifier 同机部署、同一设备对不同 Relying Party 混用两种模型，以及 TEEP 架构采用的变体（Relying Party 把 Attestation Result 回传给 Attester 供其向其他 Relying Party 复用）。两种模型影响隐私、延迟、可缓存性与受限节点适配，是部署期的显式设计选择。\par
\noindent\textbf{（3）方法 3：Freshness 与事件时序（第 10 节、Table 1）。}\par
这一机制的直接作用是：RATS 关注证据随时间的有效性，适合连接 boot、runtime 和 device lifecycle。RFC 第 10 节给出三种确定证据“epoch”的常用途径：其一，可信同步时钟加签名时间戳（可细化到 per-claim），但许多 TEE 的时钟在 TEE 之外、不可信，未必可用；其二，由裁决方发出不可预测 nonce 并要求其被签名进证据，可建立“粗略 epoch”——粒度是整组 claims，且 claims 产生时刻与收集时刻不可区分；其三，由 epoch ID 分发者周期性向收发双方派发 epoch ID（可复用、可多方共用、不必单调递增），需处理纪元切换时的竞态。RFC 明确两点边界：新鲜性裁决永远存在竞态（证据生成后 attester 状态或策略可能立即变化），目标只是把“新近程度”收窄到策略可接受的范围；新鲜性要求的弹性恰是缓存与复用证据/结果的前提，对计算与能耗预算紧张的受限设备尤其重要。Table 1（Relevant Events over Time）把相关事件编号为 NS（nonce 发出）、NR（nonce 被转发）、IR（epoch ID 接收）、EG（Evidence 生成）、ER（Evidence 被转发给 Verifier）、RG（Result 生成）、RR（Result 被转发）、RA（Result 被裁决）、OP（Relying Party 执行所请求操作）、RX（Result 过期），并以带本地时钟下标的时序图讨论各拓扑下的假设解法。Boot measurement 只是初始状态；配置更新、固件升级、device assignment 都可能要求产生新证据。\par
\subsubsection{安全性与威胁边界}
从安全边界看，RATS 把安全考虑集中在第 12 节：Attester 与证明密钥保护（第 12.1 节）要求密钥材料只能被有效建立；密钥 provisioning（第 12.1.2 节，又称 personalization/customization）区分两条路线——off-device 生成时，生成器与注入路径需要机密性保护，且保护密钥本身的密钥又需先行安全注入，形成递归问题，实践中由物理与密码措施组合解决；on-device 生成时秘密部分不出设备，问题减轻，但必须保证公钥保管链（chain of custody）的完整性，否则攻击者可让自己控制的密钥获得背书。任何 conceptual message 的传输方案（第 12.2 节）必须支持端到端完整性保护与防重放，通常还需要端到端加密、拒绝服务防护、认证与审计。需要强调的是，RATS 是架构与术语文档：它规定“谁裁决、依据什么裁决”，不规定某个 TEE 的 claim 内容、签名格式或隔离强度，测量值本身也不等同于运行期隔离；物理、侧信道与拒绝服务攻击不在其证明范围内。\par
\subsubsection{实验环境与证据}
RATS 是标准架构文档，无新实验；它是本方向的语言和模型锚点。从可核验的证据看：本地 RFC 9334 PDF 已验证；核心材料是 Figure 1 的 conceptual data flow、第 4 节的角色/制品定义、第 5 节 Figure 5/6 的两种拓扑、第 10 节的新鲜性三方案与 Table 1 的事件时序。同时需要看到其适用边界：它不规定 Arm CCA、RISC-V CoVE 或 EAT 的具体字段，也不评估任何实现的性能。\par
\subsubsection{实验结果与证据强度}
对于性能或效果，RATS 能够支持的结论是：它不提供 latency/throughput，只提供架构模型。任何 attestation 延迟都来自具体协议、密码运算、网络与 verifier 部署方式；RATS 只解释延迟与可缓存性会受 passport/background-check 模型选择影响（例如 passport 允许 Attester 缓存结果、background-check 要求 verifier 在线）。本节判断以 RFC 9334 的架构范围为限，不作 benchmark 层面的延伸。\par
\subsubsection{综合评价}
综合来看，RATS 是所有后续 TEE 与 device attestation 讨论的共同词典。它的优势在于角色、对象、流程清晰：七类角色、六种制品、两种拓扑与一套时序事件足以把 Arm、RISC-V 与嵌入式设备的证明方案放进同一叙事，论文 23 采用的 attester/verifier/relying party 术语即源于此。它的局限在于抽象层高——遵循 RATS 并不说明某份 evidence 本身安全，appraisal policy 也只有角色归属而无标准格式。从部署角度看，它是供应链审计、confidential VM 准入、device onboarding 与 zero trust 架构的基础语言；引用时只应把它作为术语与模型锚点，平台级安全结论须回到各平台规范与系统论文。\par

\subsection{论文 25：VRASED: A Verified Hardware/Software Co-Design for Remote Attestation}
\subsubsection{论文基本信息}
《VRASED: A Verified Hardware/Software Co-Design for Remote Attestation》由 Ivan De Oliveira Nunes、Norrathep Rattanavipanon（UC Irvine）、Karim Eldefrawy（SRI International）、Michael Steiner 和 Gene Tsudik 完成，发表于 2019 年的 28th USENIX Security Symposium。本文将其置于“远程证明、启动与生命周期安全”的研究脉络中考察，并把它作为 system 类材料使用。它所回答的核心问题可概括为：能否让低端嵌入式设备的远程证明从抽象安全定义到 RTL 实现都可机器验证，且硬件代价足够低。当前证据状态为：本地 PDF 已保存并核验；下文仅在该范围内讨论其机制与结论。\cite{nunes2019vrased}\par
\subsubsection{研究背景与动机}
要理解该方案，首先需要看到它面对的系统矛盾：RA 的安全性要求同时相信硬件规则、软件 MAC 代码和二者组合，而早期方案很少对实现级安全性质做机器可检查的证明。论文第 2 节把已有 RA 分为三类并指出各自死结：software-based（基于计时的）RA 无法保证密钥 K 不被恶意软件读取，依赖通信时延阈值，不适合多跳、抖动网络或有更强同谋设备协助的场景；hardware-based RA（TPM 或 SGX/TrustZone 式 CPU 改动）在物理面积、能耗与成本上对低端设备过于昂贵；hybrid HW/SW 协同（SMART、HYDRA、Trustlite、Tytan）以最小硬件支持逼近硬件方案的安全性，是 IoT 场景最有希望的路线，但包括最早的 SMART 在内都按 ad hoc 方式设计，缺少从定义到实现的形式化保证。只证明协议不够——硬件实现可能允许 DMA 或中断绕过；只证明 HMAC 不够——attestation key 与代码内存也要保护。VRASED 的目标因此是端到端：定义、设计与实现全部可验证，且硬件开销可负担。\par
\subsubsection{关键挑战与核心思想}
在上述背景下，作者提出的关键判断是：可验证 RA 必须把硬件状态机和软件 attestation 代码放进同一个证明边界。论文为此固定了一组显式假设（第 3.1 节）：攻击者可控制 prover 的全部软件状态（读写任何未被 HW-Mod 显式保护的内存、在内存段间转移恶意软件以躲避检测），并完全控制所有 DMA 控制器；验证依赖七条公理——A1（PC 始终指向当前指令）、A2（内存读写时 Daddr 与 Ren/Wen 信号正确）、A3（DMA 访问时 DMAaddr/DMAen 正确）、A4（复位后所有寄存器清零，复位逻辑不可被软件修改）、A5（中断置位 irq），A1–A5 由 OpenMSP430 设计满足；A6（被调函数返回时清理其触碰的寄存器）与 A7（编译器语义保持）则把 msp430-gcc 编译器纳入信任基——论文明确指出当前没有面向 MSP430 这类低端 MCU 的形式化验证编译器。物理硬件攻击明确不在范围内。\par
\subsubsection{系统架构}
从系统结构看，VRASED 由 HW-Mod 与 SW-Att 两部分组成：verifier 发出含挑战 Chal 的请求（Figure 1 的通用 RA 协议：请求—带认证的完整性检查—报告—验证四步），prover 的 HW-Mod 保证 SW-Att 在受保护条件下计算 HMAC。内存被划分为若干区域（Figure 3、Table 1）：CR（code ROM，存放 SW-Att）、KR（key ROM，存放主密钥 K）、XS（exclusive stack，SW-Att 专用的安全 RAM 栈）、MR（MAC RAM，存放 HMAC 结果，同区也用于传入挑战）、AR（attested region，被证明内存，可预定义或由经认证的 verifier 请求指定）。HW-Mod 由若干子模块 FSM 构成，实现为 Mealy 机，输入为 PC、irq、Ren、Wen、Daddr、DMAen、DMAaddr 的子集，唯一输出 reset——任何安全性质被违反即触发 MCU 硬复位。HW-Mod 以极小化为设计原则：受验证 FSM 的状态空间保持最小，以压低硬件成本。\par
\subsubsection{核心方法设计}
论文的关键不是单一组件，而是把访问控制、原子执行与形式化证明组合成一条从定义到实现的受控路径。\par
\noindent\textbf{（1）方法 1：形式化规约——LTL 性质与端到端定义。}\par
这一机制的直接作用是：VRASED 把 RA 安全性质写成 LTL 公式并由 NuSMV 模型检查，再用定理证明器（经 LTL 等价）说明子性质蕴含端到端目标。论文 Definition 1 给出 RA soundness：当 PC 进入 CR 首地址（CRmin）、AR 值为 M、MR 值为 Chal 且在 SW-Att 执行完成（PC 到达 CRmax）前不发生 reset，则最终 MR 必然等于 HMAC(KDF(K, Chal), M)——即被测内存在 HMAC 计算期间不可变化（temporal consistency）。Definition 2 用 RA-game 定义安全性：攻击者可任意修改 AR、可多项式次调用 SW-Att，但若它输出的 (M, σ) 满足 σ = HMAC(KDF(K, Chal), M) 且 M ≠ AR(t)（不等于任一调用时刻的真实内存），则攻击者获胜；论文经规约证明，赢得该游戏意味着打破 HMAC 的存在不可伪造性。支撑这些定义的是七条子性质 P1–P7：密钥保护性质 P1–P3 中，正文明确展开 P1（K 只能被 SW-Att 访问，LTL 规范 2：¬(PC∈CR) ∧ Ren ∧ (Daddr∈KR) → reset）与 P2（密钥机密性——attestation 结束与复位后不得有残留泄漏，寄存器清理由公理 A4/A6 保证，RAM 残留由独占栈或擦除解决）；P4 是 SW-Att 功能正确性（由 HACL* 验证蕴含）；P5 是 SW-Att 不可变性（由 ROM 存储直接获得）；P6 是执行原子性（执行期间的中断、复位或数据篡改必然触发 reset）；P7 是受控调用（进入 SW-Att 的唯一途径是从 CRmin 首指令开始执行）。针对 DMA，论文第 4.7 节追加三条 LTL 规范（Table 2 中编号 8–11 的一部分），因为被攻陷的 DMA 控制器可绕过基于 Daddr 的访问控制直接读 K 或独占栈（破坏 P1/P2），也可把恶意软件移出测量范围（破坏 P6）。\par
\noindent\textbf{（2）方法 2：Verified SW-Att——HACL* HMAC 与编译器信任缺口。}\par
这一机制的直接作用是：软件侧不是随便写一个 HMAC，而是使用经过验证的实现并显式处理编译信任问题。SW-Att（Figure 7 给出全部 C 代码）先用 HKDF 从主密钥 K 与挑战 Chal 派生一次性上下文密钥 key，再用 HACL* 的 HMAC-SHA256 对 AR 计算认证值；密钥派生可扩展为纳入 verifier 指定的测量边界。HACL* 代码在 F* 中被验证为功能正确、内存安全、秘密无关（无依赖秘密的分支，缓解计时侧信道），且全部内存分配在栈上、行为确定——后者使“把执行约束在固定保护区”或“退出前擦除栈”两种设计都成为可能。论文同时坦承两个信任缺口：HACL* 为性能不提供栈擦除，故 P2 不直接从 HACL* 成立，低端裸机 MCU 又没有 OS 提供进程间隔离；A6/A7 把编译器纳入信任基，这是验证链条上未被机器检查的一环。\par
\noindent\textbf{（3）方法 3：HW/SW Composition 与可替换设计。}\par
这一机制的直接作用是：最终安全性来自硬件保护、软件 HMAC 与组合证明三者的合取，而非任何单一组件。HW-Mod 限制密钥/代码访问并保证原子执行，SW-Att 计算挑战绑定的 HMAC，composition proof（附录 A）证明子性质合取足以推出 Definition 1/2 的端到端 soundness 与 security。论文第 5 节还系统比较了保证 P2 的三种可替换设计及其代价：默认设计是 2.3KB 独占安全栈（约占 MSP430 16 位地址空间的 3.5%）；erasure-on-SW-Att 变体取消保留栈，但要求硬件在启动时对整块 RAM 做可验证擦除（否则执行中途复位会让 K 残留在 RAM 中）；compiler-based clean-up 变体借助 zerostack 式编译器插桩（论文指出理想情况下应作为 CompCert 类验证编译器套件的一部分实现并验证，当前无 MSP430 验证编译器，评估数字为手工插桩估计）；double-HMAC 变体利用 HACL* 确定性栈分配的性质，用独立常量做第二次“空”HMAC 覆盖第一次调用的全部秘密变量，代价是运行时间约翻倍，适合内存紧张但无法修改 HACL* 或编译器的场景。\par
\subsubsection{安全性与威胁边界}
从安全边界看，VRASED 的威胁模型覆盖“软件全攻陷 + DMA 全控制”的强攻击者，但有三条显式边界：其一，物理硬件攻击（探针、篡改封装）不在范围内，K 与 SW-Att 的 ROM 保护假设硬件不被物理绕过；其二，A6/A7 把编译器语义保持纳入信任基，验证链在机器码层面存在人为假设；其三，它证明的是“证明行为本身正确”，即报告忠实反映调用时刻的被测内存——它不断言被测软件无漏洞，也不覆盖测量之后运行态的变化（TOCTOU 类问题）。这些边界决定了它适用于设备 RoT 与固件完整性证明，而非云 CVM 的完整证据链。\par
\subsubsection{实验环境与证据}
VRASED 是 peer-reviewed 的系统+验证论文，包含形式化证明与 FPGA 实现，证据链完整可复核：硬件基于开源 OpenMSP430（Verilog HDL，可近周期精确执行任意 MSP430 工具链生成的软件），修改内容包括增加存放 K 与 SW-Att 的 ROM、加入 HW-Mod 并适配内存骨干；用 Xilinx Vivado 将 HW-Mod 的 RTL 综合到 Artix-7 类 FPGA（README 记录为 Basys3 平台）；软件用原生 msp430-gcc 编译、以 linker script 匹配 Figure 3 的内存布局；硬件性质用 NuSMV 在 3.40GHz 桌面机上模型检查，软件 HMAC 复用 HACL* 的 F* 验证结果；实现已公开。\par
\subsubsection{实验结果与证据强度}
实验给出三组可核验数字。第一组是验证复杂度（Table 2）：Key AC 子模块验证 LTL 规范 2、11，耗时 0.02s、内存 7.5MB；Atomicity（规范 3–5、11）0.05s/8.5MB；Exclusive Stack（规范 6、7、11）0.03s/8.1MB；DMA Support（规范 8–11）0.04s/8.2MB；HW-Mod 整体合取（规范 2–11）0.28s/13.6MB——全部性质合计内存 <14MB、时间约 0.3s 即完成，说明“最小状态空间”设计使验证在商品桌面机上可行。第二组是成本与性能（Table 3，以 secure stack 默认设计为例）：相对未修改的 OpenMSP430 核（1842 LUT、684 寄存器、3044 逻辑单元、4034 行 Verilog），VRASED 增加 122 个 LUT（+6.6\%）、37 个寄存器（+5.4\%）、共 193 个逻辑单元（+6.3\%），Verilog 增加约 587 行（约 15\%）；ROM 约 4.5KB，其中 96\% 是 HACL* HMAC-SHA256 代码；在 8MHz 下对 4KB RAM（部分 MSP430 型号的典型内存容量）做一次证明约需 360 万周期、约 0.45s，double-HMAC 变体约 0.90s；作为参照，MSP430 上一次优化的 Curve25519 ECDH 密钥交换约需 900 万周期，即 4KB 证明比一次 ECDH 快约三倍，且证明耗时随被测内存线性增长。第三组是同类对比（Table 4）：VRASED 为 hybrid HW/SW、HMAC-SHA256（256 位摘要）、需 ROM 存放证明代码、支持 DMA 防护、经过形式化验证；SMART 为 hybrid、HMAC-SHA1、不支持 DMA、未经验证；SANCUS 为纯硬件、SPONGENT-128/128/8、无需 ROM、不支持 DMA、未经验证——后两者成本可作参照，但因 ad hoc 设计无法提供 VRASED 的可验证保证。论文自评开销“可负担”（含寄存器、逻辑门与独占内存在内仅 3–6\% 的增幅）。需要强调的是，这些数字是低端 MCU 场景的实现成本，不能当作云端 attestation 的延迟基线。\par
\subsubsection{失败模式与边界分析}
论文显式讨论了若干会让保证失效的路径，构成本节的失败模式清单：若缺少启动时的 RAM 擦除，取消独占栈的变体在执行中途复位会把 K 留在 RAM（第 5.1 节）；若 DMA 控制器被攻陷而没有第 4.7 节的附加规范，攻击者可直读 K、读栈或帮恶意软件移出测量范围（第 4.7 节）；若编译器不满足 A6/A7，F* 层面的证明无法传导到机器码（第 3.1、5.2 节）；HACL* 不提供栈擦除，使 P2 必须靠硬件或额外设计补足（第 4.3 节）。这些条目说明 VRASED 的保证是“假设—性质”条件句，部署时必须逐条核对前提。\par
\subsubsection{综合评价}
综合来看，VRASED 是 RA 方向的强 SOTA，因为它把设计、实现和证明放在一起：proof-driven 的 HW/SW co-design、从 LTL 子性质到端到端 soundness/security 的组合证明、公开可用的 FPGA 实现，三者在同方向工作中罕见地齐全。它的局限同样清晰：MCU 威胁模型不含物理攻击，编译器信任缺口（A6/A7）未被消除，现代 TEE 的生命周期议题（运行时测量、撤销、密钥轮换）不在其范围。从部署角度看，它可直接转化为设备 RoT、IoT RA 与小型安全控制器的设计方法——尤其“最小 FSM + 单 reset 输出 + 违反即硬复位”的硬件模式，以及“验证子性质合取推出端到端定义”的证明组织方式；其 6\% 量级的硬件增幅也说明，对低端设备而言可验证性并非不可负担的奢侈品。\par

\subsection{论文 26：Separate but Together: Integrating Remote Attestation into TLS}
\subsubsection{论文基本信息}
《Separate but Together: Integrating Remote Attestation into TLS》由 Carsten Weinhold、Michael Roitzsch（Barkhausen Institut）、Muhammad Usama Sardar（TU Dresden）、Ionut Mihalcea 与 Yogesh Deshpande（Arm）、Hannes Tschofenig（University of Applied Sciences Bonn-Rhein-Sieg）、Yaron Sheffer（Intuit）和 Thomas Fossati（Linaro）完成，发表于 2025 年的 USENIX Annual Technical Conference (USENIX ATC 2025)。本文将其置于“远程证明的通道绑定与网络端点”的研究脉络中考察，并把它作为 system 类材料使用。它所回答的核心问题可概括为：如何在不改动 TLS 证书体系与 TEE 证明体系各自部署的前提下，证明“这条 TLS 连接”确实终止在“这个 attested TEE/workload”内。当前证据状态为：本地 PDF 已保存并核验；下文仅在该范围内讨论其机制与结论。\cite{weinhold2025tlsra}\par
\subsubsection{研究背景与动机}
要理解该方案，首先需要看到它面对的系统矛盾：机密 workload 经常需要远程注入 secret 或建立业务连接，但两种现有保证各管一段。TLS 1.3 的证书断言的是端点身份与行政控制——CA 在申请者证明控制某 FQDN 后签发证书，把域名与公钥绑定；remote attestation 断言的是 TEE 内部状态——RoT 对含 nonce 与软硬件测量值的 quote 签名。只看 TLS，不知道服务器是否在 TEE 中运行；只看 RA，不知道后续 TLS channel 是否被 relay 到 TEE 之外。论文 Figure 1 给出 relay 攻击：攻击者拿到 TLS 证书私钥后搭建运行任意软件栈的假端点，把客户端的 attestation nonce 转发给真实 TEE 取回“看起来有效”的 quote 再转交客户端——通道表面上同时通过证书校验与证明校验，实际却终止在 TEE 之外。这说明并排（side-by-side）使用 TLS 与 RA 时两者性质并不独立：证明的可信度意外地被挂靠到 TLS 私钥的保护上，若 TLS 私钥由云厂商基础设施代管，则 TEE 的承诺被架空。TLS+RA 要解决的正是二者的密码学绑定。\par
\subsubsection{关键挑战与核心思想}
在上述背景下，作者提出组合协议应满足五条性质（第 2.2 节）：其一，保留安全通道性质（机密性、完整性、防重放）并额外传输 attestation report，且必须保证隧道终止在产生报告的 TEE 内；其二，不给握手增加额外网络往返；其三，通道建立后不增加加密层（沿用普通 TLS 载荷加密即可）；其四，部署独立——证书与 PKI 的签发、过期、撤销、透明度日志实践一概不改，证书不得携带 attestation 信息（否则同一域名由不同机器、不同证明技术服务的负载均衡会被破坏），attestation 部署只涉及 TEE 厂商；其五，失败独立——TLS 私钥或 attestation RoT 任一侧被攻破（如 Heartbleed 式侧信道泄漏私钥、或 CA 被攻破签发假证书），另一侧的断言必须仍然成立。作者指出，既有组合方案没有一方同时满足这五条，尤其是失败独立这一此前未被考虑的性质；TLS 1.3 强制使用临时 Diffie-Hellman（DHE）密钥协商与 extension 机制，是重新审视该问题的两个使能条件。\par
\subsubsection{系统架构}
从系统结构看，client 同时执行 TLS 证书验证与 attestation 验证，两者都对同一次握手成立后才信任通道。Figure 2 给出支持双向证明的握手流程：客户端先向自己的 attestation verifier（本地库或远程服务，与协议设计正交）查询支持的证据格式并取得 nonce，然后在 ClientHello 的 attestation request extension 中携带支持的格式列表与该 nonce；服务器若在 TEE 内运行 TLS+RA 库且支持其中某种格式，就让 RoT 生成 attestation report，并把它放进 TLS certificate 消息的 extension 返回；格式不匹配则握手终止。客户端随后验证证书链与 TEE 证据——只有二者都与握手 transcript 一致，secret 才应被释放。整个设计中，TEE 类型、RoT 与证据格式对协议不透明：TLS+RA 把 report 当作 extension 里的不透明数据搬运。\par
\subsubsection{核心方法设计}
论文的关键不是单一组件，而是把“双链接”密码结构、TLS 1.3 扩展机制与平台插件接口组合成一条兼容现有部署的受控路径。\par
\noindent\textbf{（1）方法 1：Relay 攻击防御与双链接（double linking）。}\par
这一机制的直接作用是：用两个相互独立的对称链接把证据钉在本次握手上。标准 TLS 已有第一重链接：对截至当前握手消息日志计算 transcript hash（涵盖双方 hello、nonce 与 DHE 公钥），用 TLS 证书私钥签名，从而把会话与端点身份绑定。TLS+RA 增加第二重链接：计算 linking hash——覆盖同样的握手消息，但额外纳入 DHE 协商出的共享秘密与 attestation nonce——并把它作为 challenge 交给 attestation RoT，使证据内含该 linking hash。关键设计点在于：若 linking hash 只覆盖公开可见的握手消息，攻击者可以在另一条连接上重放同样的公开信息完成 relay；纳入只有通信双方才知道的 DHE 共享秘密后，转发证据必然导致 linking hash 不匹配。持有窃取的 TLS 私钥的攻击者架设假服务器时，无法自证，只能分别与客户端和真实服务器各建一条通道；由于两条通道的 DHE 秘密不同，客户端算出的 linking hash 与真实服务器证据内的 linking hash 不一致，握手被中止（Figure 3）。两重链接对称（源自同一消息日志）但独立（不同签名密钥、不同部署基础设施），这正对应失败独立性；协议同时继承 TLS 1.3 的前向保密。\par
\noindent\textbf{（2）方法 2：TLS 1.3 集成与相关工作对照。}\par
这一机制的直接作用是：设计目标是兼容 TLS 生态，而不是重写安全通道协议。attestation request 与 report 分别作为 hello 消息与 certificate 消息的 extension 传输，与常规握手消息同帧，不产生额外往返；证书 PKI 与 attestation verifier 保持独立，失败处理得以区分 TLS failure 与 attestation failure。论文 Table 1 把相关工作按五条性质对照：HTTPA 在外层 TLS 内嵌套内层 attestation 通道，满足失败独立但握手与加密双重开销；Goldman 等的 platform certificate 把 TLS 密钥对与 RoT 静态绑定，而证书与连接无关，无法阻止连接级 relay；RA-TLS 用含报告的自签名临时证书，脱离 Web PKI，其 CA 实时签名变体增加延迟且自动化所需的长寿命秘密或域控证明可被复制，仍不满足失败独立；Aziz 等的方案把 RoT 绑定到握手 nonce 但增加消息与一整个往返；RATLS 与 Trusted Channels 用 TLS 私钥链接报告，私钥泄漏即可在 TEE 外重放链接；作者参与的 IETF 草案（draft-fossati-tls-attestation）为与相关标准互操作同样选择绑定 TLS 私钥，安全保证与后两者类似。结论是：只有代价高昂的嵌套或额外往返方案才满足失败独立，TLS+RA 以零额外往返、单层加密取得该性质。\par
\noindent\textbf{（3）方法 3：OpenSSL 实现与平台插件。}\par
这一机制的直接作用是：证明该协议能以接近标准 TLS 工程的方式落地。实现基于 OpenSSL：用 \texttt{SSL\_CTX\_add\_custom\_ext} 注册回调在握手中创建/解析 extension，用 \texttt{SSL\_CTX\_set\_verify} 在收到 certificate 消息 extension 中的 report 时触发验证，用 \texttt{SSL\_CTX\_set\_ex\_data} 在会话上下文中嵌入自定义数据，从而以一组与 TLS 握手并行运行的回调实现 attestation 状态机。唯一需要修改 OpenSSL 代码的地方：linking hash 要在握手进行中、DHE 共享秘密刚算出时计算，而 RFC 5705 的标准 exporter API 只能在握手完成后调用——太晚；作者因此增加了一个非标准 exporter 函数，该改动不涉及线上协议，仅从 TLS 实现内部状态派生 linking hash。平台侧提供插件 API，核心为 \texttt{remote\_attest} 与 \texttt{check\_report} 两个函数；论文实现了 TPM 与 AMD SEV-SNP 两个插件用于定量评估，另有 Arm CCA 插件——因 CCA 硬件尚未公开可用，只在 Arm 模拟器上做功能测试、不参与性能基准。代码已开源。\par
\subsubsection{安全性与威胁边界}
从安全边界看，威胁模型（第 4 节）假设 TLS PKI 与 TEE 厂商的证明基础设施独立运营，并要求 TLS+RA 库、使用它的应用及其安全关键依赖都运行在 TEE 内（使组合握手代码被 attestation report 覆盖，部分 TLS+RA 状态也需纳入测量）。攻击者可攻破 TLS 私钥或 attestation RoT 之一，但不同时攻破两者；RoT 密钥被伪造而 TEE 隔离未受损时，可伪造证据但 TEE 内 TLS 私钥不被假定泄漏。边界同样明确：TLS+RA 不替代底层 TEE 的证明语义（quote 解释仍属 verifier 与策略），不证明应用代码无漏洞；证据格式协商失败即终止握手，可用性依赖双方格式交集；对 DPU/SmartNIC 等设备端点，还需设备身份与平台证明机制另行配合。\par
\subsubsection{实验环境与证据}
论文是 peer-reviewed 的系统实现，评估覆盖三种证明技术并显式对应三类部署形态：AMD SEV-SNP（云服务器 TEE）跑在 Amazon c6a.large 实例（Amazon Linux 2023，内核 6.1）；firmware TPM（终端设备）为 13 代 Intel Core i5-13400 台式机（Ubuntu，内核 6.5）上实现 TPM 2.0 的 fTPM；离散硬件 TPM（IoT 设备）为 Infineon Optiga SLB 9670，接在四核 Cortex-A72 的 Raspberry Pi 4 Model B（内核 5.15）上；Arm CCA 用 Arm 模拟器做功能验证。为测最坏处理开销，客户端与服务器同机不同核、经 loopback 通信并关闭核的 idle sleep 以减少噪声；网络延迟用 Linux netem 模拟单向最高 80ms（代表跨洲距离）；每项测量在两次热身后重复 100 次（吞吐实验为 10 次）取均值与标准差。互操作性方面，四种技术上均成功执行了组合的 TLS+RA 流程，且协议支持服务器向客户端、客户端向服务器与双向三种证明方向。\par
\subsubsection{实验结果与证据强度}
实验支持三层结论。其一，握手开销（Figure 4）：TLS+RA 相对标准 TLS 的增量由证据生成与验证主导；载荷方面，基线 TLS 握手约 1.5KiB，fTPM 插件的 TLS+RA 需传输 6.4KiB——主要来自 quote 与 TPM attestation identity key 的公钥部分（原型插件将其序列化为 JSON），SEV 插件则直接转发约 1KiB 的二进制报告与宿主平台密钥证书；所有报告都是粗粒度的整机软件栈描述。其二，网络距离下的摊薄（Figure 5）：SEV 生成报告仅 6ms，在 80ms 网络延迟下 TLS+RA 相对开销收窄到 5\%；dTPM 生成报告需 210ms，同距离下握手开销为 1.3 倍。与之对照，论文复现了“先 TLS 后证明”的 post-RA 顺序组合（HTTPA 等方案的结构）：对 SEV，网络延迟 10ms 及以上时，顺序执行两次握手带来超过 100\% 的基线开销。其三，数据面（Figure 6）：通道建立后 TLS 与 TLS+RA 吞吐几乎相同，因为二者执行完全相同的密码操作、远程证明不引入握手后开销，吞吐不由 TEE 技术决定；嵌套通道方案（复现 HTTPA 式双层 AES-GCM，发送 1GiB 载荷）因冗余加密整体吞吐更低。论文强调，这些数字支撑的是“握手/控制面开销”结论，不应外推为所有网络 I/O 的性能结论；真实部署还需 verifier 策略、证据格式、证书生命周期与多厂商互操作的进一步评估。\par
\subsubsection{与相关工作的定量对比}
论文在同一 TLS 版本、同一实现、同一硬件上复现了替代组合路线，使对照可定量：post-RA 顺序组合在 SEV、10ms 延迟处越过 +100\% 握手开销，而 TLS+RA 在 80ms 处仅 +5\%；嵌套通道在数据面因双层 AES-GCM 吞吐下降，TLS+RA 与基线 TLS 持平。定性维度（Table 1）上，TLS+RA 是所列方案中唯一同时满足防 relay、零额外往返、零额外加密层、部署独立与失败独立的设计；Walther 等的 RATLS 与 Armknecht 等的 Trusted Channels 同样做到前三项加部署独立，但因用 TLS 私钥链接证据而不满足失败独立。\par
\subsubsection{综合评价}
综合来看，TLS+RA 是 endpoint/channel binding 方向的强 SOTA，特别适合 confidential service 的 secret provisioning 场景：问题切得准——它把“TLS 与 RA 并排使用会被 relay”这一已知但长期未被正确处理的漏洞，落成五条可检验的设计性质，并用双链接一次性满足。它的局限在于仍依赖底层 attestation 的正确性与应用层策略（接受哪些 measurement 由应用 policy 决定，机制本身不判断 TEE 语义），且要求应用与 TLS 库整体进入 TEE 的测量范围。从部署角度看，它直接适用于 API service、KBS（密钥代理服务）、CVM secret injection 与 DPU/NIC 控制通道；与本节其他工作的关系是：RATS 给它提供 verifier 侧的角色语言，论文 23 的比较框架解释其插件为何能跨 SEV-SNP、TPM 与 CCA，VRASED 式验证则提示了把“双链接”本身形式化的下一步。\par

\subsection{开发商安全实践建议}
本节四篇工作从比较框架、架构标准、可验证实现到通道绑定给出了完整证据链，对芯片厂商、固件与系统软件开发商、设备制造商（含网络设备厂商）有以下可操作含义。\par
\noindent\textbf{实践一：把测量根与签名密钥放进不可绕过的硬件边界，并逐条核对假设。}\par
设备厂商应确保 attestation key 只存在于软件（含 DMA）不可直读的存储中，且证明代码的执行不可被中断、篡改或非受控入口调用。VRASED 给出了可直接复用的最小模式：K 存 KR、证明代码存 CR（ROM）、专用栈 XS、违反任何访问规则即硬复位，全部性质可用 NuSMV 在桌面机上秒级验证，硬件增幅仅约 6\%\cite{nunes2019vrased}。对不具备内建证明的 TrustZone 平台，论文 23 的归纳表明必须先落实安全启动、片上熔断根密钥与安全随机源，再由 TM 类组件派生证据密钥\cite{menetrey2022attestation}。采用此类设计时应把公理清单（如 A1–A7）写入安全需求文档逐条核对，特别注意 DMA 路径——VRASED 的分析显示被攻陷的 DMA 可绕过基于 CPU 地址信号的访问控制。\par
\noindent\textbf{实践二：按 RATS 角色拆分证据生产、裁决与消费，建立 endorsement 与 reference value 供应链。}\par
固件与平台开发商不应把“验证 quote 签名”当作证明流程的终点。正确做法是按 RATS 拆分：设备侧只负责产生 Evidence；独立的 Verifier 持有 Appraisal Policy for Evidence，并接入制造商（Endorser）签发的 Endorsements 与 Reference Value Provider 维护的参考值基线；业务侧 Relying Party 消费 Attestation Results 并保留自己的授权策略\cite{rats_rfc}。这意味着产品立项时就要规划三条供应链：密钥/组件背书链、固件版本的 known-good 测量值发布链、策略下发链。对多型号产品线，参考值应覆盖配置、固件升级与 device assignment 等生命周期事件，而非只覆盖出厂启动测量。\par
\noindent\textbf{实践三：凡经网络释放 secret，必须把 evidence 密码学绑定到会话，禁止 TLS 与 RA 并排裸用。}\par
机密注入、远程运维与设备接入场景中，只做“TLS 通道 + 单独校验 quote”会留下 relay 缺口：攻击者可用真实 TEE 的证据包装一条终止在 TEE 外的通道\cite{weinhold2025tlsra}。可操作的做法是：在握手内完成证明（利用 TLS 1.3 extension），挑战包含新鲜 nonce，证据内含覆盖 DHE 共享秘密的 linking hash；同时保持证书 PKI 与 TEE 证明基础设施独立部署，使任一侧私钥泄漏不连坐另一侧。对暂缓改造的系统，至少应盘点现有部署中“先 TLS 后证明”的顺序组合——论文 26 的测量显示该结构在中等网络延迟下握手开销即翻倍，且若链接依赖 TLS 私钥则不具备失败独立性。\par
\noindent\textbf{实践四：把密钥生命周期当作一等工程问题管理。}\par
制造商在 provisioning 阶段就应决定密钥生成路线：off-device 生成需保护生成器与注入路径的机密性，并接受“保护密钥的密钥也要安全注入”这一递归问题的工程解法（物理与密码措施组合）；on-device 生成可让秘密部分不出设备，但必须保证公钥保管链的完整性，防止攻击者让自己控制的密钥获得制造商背书\cite{rats_rfc}。选型时还应注意密钥体系的隐私与撤销语义：SGX 式 EPID 群签名不标识唯一设备、quote 与固件版本绑定，适合需要匿名性的场景，但撤销依赖群公钥与撤销信息的分发\cite{menetrey2022attestation}；对称 MAC 证据（TIMBER-V 式）则要求 verifier 持有同一秘密，密钥分发范围天然受限，跨域验证慎用。密钥轮换、固件升级后的证据失效（RFC 9334 Table 1 的 RX 事件）与缓存策略应一并设计。\par
\noindent\textbf{实践五：MCU 与 IoT 设备优先选择可验证的 hybrid RA，并按内存预算选变体。}\par
对低端设备，纯硬件 RA 面积/能耗/成本过高，纯软件 RA 无法保护密钥\cite{nunes2019vrased}。可操作的选型顺序是：先选 hybrid HW/SW 架构（HMAC 类证明 + 最小硬件访问控制），再按内存约束选 P2 保证方式——RAM 充裕用独占栈（约 2.3KB），RAM 紧张用 double-HMAC（时间约翻倍）或启动擦除变体（要求硬件支持）；证据强度上优先 256 位摘要（VRASED 的 HMAC-SHA256）而非 SHA1 级 MAC。若有形式化验证条件，应把 HW-Mod 的 FSM 设计与 LTL 规约纳入评审材料——VRASED 的 Table 2 证明这在工程上完全可行。\par
\noindent\textbf{实践六：按部署拓扑显式选择 passport 或 background-check 模型，并为受限节点优化。}\par
平台开发商应在架构评审中显式记录证明拓扑：passport 模型允许 Attester 缓存 Attestation Result 并向多个 Relying Party 出示，适合 verifier 可及性不稳定或需要离线准入的场景，但要求 Result 格式标准化；background-check 模型让 Relying Party 原样转发 Evidence，免去其解析器，适合代码体积与攻击面敏感的受限节点，但要求 verifier 在线\cite{rats_rfc}。同一设备可对不同业务混用两种模型（如网络准入用 background-check、数据访问用 passport），并为 Evidence/Result 的新鲜性选择 nonce、时间戳或 epoch ID 方案——注意许多 TEE 的时钟在 TEE 之外不可信，盲目采用时间戳会引入隐性信任假设。\par

\subsection{未来研究方向}
\noindent\textbf{方向一：CCA/CoVE 代际的证据语义与跨平台 verifier。}\par
问题：论文 23 的比较矩阵写作于 Arm CCA Realm 与 Intel TDX 落地之前（论文明确省略二者），RISC-V CoVE/TVM 的证据格式更晚出现；现有跨平台比较维度（Table 1 的 13 个特性）未覆盖 realm 级测量、可扩展 TCB 版本声明等新 claims。为什么未解决：这些架构的 token 格式与验证服务在 2022 年后才陆续可用，学术比较滞后于规范迭代。技术路线：以 RATS 角色语言为底座，为 CCA token、CoVE evidence 与 SGX/SNP quote 建立 claim 级映射表，扩展论文 23 的特性矩阵；论文 26 的 Arm CCA 插件（当前仅经模拟器功能测试）可作为首个对接点。需要的实验：在真实 CCA 硬件可用后复测 TLS+RA 式握手的报告生成/验证开销，并与 SEV-SNP 的 6ms 量级基线对照。\par
\noindent\textbf{方向二：证明-通道绑定的形式化验证。}\par
问题：TLS+RA 的双链接以非形式化论证给出安全性（第 5 节），其安全依赖 linking hash 覆盖 DHE 秘密、nonce 与消息日志的精确组合，以及 OpenSSL 回调状态机与 TLS 握手的并行正确性；目前没有任何一方对该组合给出机器检查的证明。为什么未解决：VRASED 式验证面向单设备 RA 原语\cite{nunes2019vrased}，TLS 1.3 的形式化分析（既有学术工作）不含 attestation 扩展，两个验证社区的对象尚未合并。技术路线：用 Tamarin/ProVerif 对双链接握手建模，把 relay 攻击者形式化为可转发消息但不知 DHE 秘密的 Dolev-Yao 攻击者；在实现侧，借鉴 VRASED 的“子性质合取推出端到端定义”方法，为 TLS+RA 状态机的 linking hash 计算点（DHE 完成与 report 生成之间）建立不变量。需要的实验：对模型做攻击自动搜索（relay、降级、格式混淆），并在原型上变异 linking hash 的覆盖范围验证模型预言的失效模式。\par
\noindent\textbf{方向三：从启动证明走向事件驱动的连续证明。}\par
问题：SEV/SEV-ES 只支持启动时测量，SEV-SNP 的运行时报告仍是按需快照；RATS Table 1 枚举了 EG/ER/RG/RA/OP/RX 等完整事件链，但现有 TEE 实现几乎没有把配置更新、固件升级、device assignment 映射为自动触发的新证据。为什么未解决：运行时测量与性能开销、证据缓存策略、TOCTOU 窗口相互纠缠，VRASED 也只保证“调用时刻”的内存忠实性。技术路线：在 RATS 事件模型上定义触发-证据-缓存契约（哪些事件使缓存证据失效、RX 如何传播到 Relying Party），结合轻量运行时测量（如论文 23 所述 TrustZone 单向/双向协议族）实现增量证明。需要的实验：测量事件触发频率与证明开销的帕累托前沿，以及缓存证据在策略更新后的误接受率。\par
\noindent\textbf{方向四：消除低端设备验证链上的编译器信任缺口。}\par
问题：VRASED 的 P2/P4 在机器码层面依赖 A6（寄存器清理）与 A7（语义保持）两条人为公理，论文明确指出没有面向 MSP430 级 MCU 的验证编译器；这使得“从定义到实现全验证”的承诺在最底层断开。为什么未解决：CompCert 类验证编译器不覆盖低端 MCU 指令集，而 zerostack 式插桩尚未被形式化验证。技术路线：为 MSP430/RISC-V 低端核构建验证的子集编译器，或把栈擦除、寄存器清零下沉为硬件行为（扩展 HW-Mod 的 reset 例程）从而把 A6/A7 转化为可模型检查的硬件性质。需要的实验：对比三种方案（验证编译器、硬件擦除、double-HMAC）在面积、周期数与证明复杂度上的开销，复用 VRASED Table 2/3 的评估方法。\par
\noindent\textbf{方向五：Appraisal policy 的标准化、分发与可验证性。}\par
问题：RATS 把信任决策定位于 appraisal policy，却刻意不规定策略格式——策略由 Verifier Owner/Relying Party Owner 以任意机制配置；论文 26 也把“接受哪些 measurement”留给应用策略。结果是安全关键决策落在最缺乏标准化的环节：策略写错，证据再正确也会误判。为什么未解决：策略天然与应用语义绑定，标准化收益直到多厂商 verifier 生态出现后才显现。技术路线：设计可表达 measurement 白名单、TCB 版本下界、新鲜性阈值与撤销状态的声明式策略语言，给出策略与 attestation result 的绑定签名；进一步可对策略满足性做形式化分析（如“该策略是否拒绝所有已知脆弱固件版本”）。需要的实验：用真实固件发布史（含已知 CVE 版本）回放策略判决，测量误接受/误拒绝率，并评估策略分发本身的新鲜性（复用 RFC 9334 第 10 节的三方案）。\par
\noindent\textbf{方向六：双向证明与 verifier 生态的工业化。}\par
问题：论文 23 的矩阵显示 mutual attestation 在工业 TEE 中多为“部分满足”（依赖文献扩展）；TrustZone 的双向协议与 LIRA-V 停留在学术原型；TLS+RA 虽支持双向，但其评估中的 verifier 仅为本地库。为什么未解决：双向证明要求两端都有 RoT、都有 verifier 可达性，且 verifier 本身的可用性、多厂商证据格式互操作与撤销分发构成运营级门槛。技术路线：以 RATS 拓扑为框架研究 verifier-as-a-service 的信任下放（passport 模型的 Result 缓存如何与撤销协调），并以 TLS+RA 的格式协商 extension 为载体做跨厂商互操作测试。需要的实验：多 TEE（SGX/SNP/CCA/TPM）混合部署下的握手成功率、格式协商失败率与端到端延迟分布，以及 verifier 不可达时（passport 模型第三种失败方式）的业务降级行为。\par
\section{机密 I/\allowbreak{}O、网络与数据路径防御}
机密工作负载的数据一旦离开 CPU 封装，保护边界便延伸到内存控制器、PCIe/CXL 链路、DMA 映射和网络端点。这一延伸改变了威胁模型的结构：在 CPU 侧，机密性由内存加密引擎和 TEE 边界保证；而在链路侧，攻击者可以窃听、篡改、重放 CXL Flit 或 RDMA 报文；在设备侧，攻击者可以是控制 NIC 管理操作系统、MNode 软件栈或存储端点固件的特权对手。因此，单点机制无法给出端到端结论——CXL IDE 保护链路却不回答设备身份，SPDM 证明设备身份却不约束 DMA 窗口，SmartNIC 内部隔离不证明链路可信。本节把 CXL、PCIe、RDMA、IOMMU 与端点协议放进同一数据路径，逐篇考察设备身份、访问授权、链路保护和中断投递是否彼此衔接。\par
逐篇评审将特别关注三点：第一，控制面与数据面是否采用一致的对象标识和生命周期，例如 SPDM 的会话密钥与 TDISP 的接口状态、S-NIC 的 function owner 与 Hazel 的 counter lease 是否绑定到同一工作负载身份；第二，性能证据落在哪一层——模拟器、FPGA 原型、真实 CXL 卡还是规范文本——以及该层证据能否外推；第三，论文对兼容性、撤销和多租户的取舍。这样可以避免把单段链路的加密（如 CXL IDE 的 GCM）或单个 DMA 防护，误读为端到端的机密性保证。\par
\subsection{论文 27：AIORE: Efficient Security Support for CXL Memory through Adaptive Incremental Offloaded (Re-)Encryption}
\subsubsection{论文基本信息}
《AIORE: Efficient Security Support for CXL Memory through Adaptive Incremental Offloaded (Re-)Encryption》由 Chuanhan Li（University of California, Santa Cruz）、Jishen Zhao（University of California, San Diego）和 Yuanchao Xu（University of California, Santa Cruz）完成，发表于 2025 年的 58th IEEE/ACM International Symposium on Microarchitecture (MICRO '25, Seoul)，DOI 为 10.1145/3725843.3756119。本文将其置于“内存加密、完整性与重放保护”的研究脉络中考察，并把它作为 system 类材料使用。它所回答的核心问题可概括为：在 XTS TEE 与 CXL IDE 基线把加解密放进每次 CXL 内存读关键路径的前提下，能否通过按页自适应选择 CTR/XTS 加密、增量重加密与卸载重加密，把安全开销压到可接受水平。当前证据状态为：本地已保存论文 PDF 并已核验（README 评审块尚未建立），下文机制描述与全部数字均以论文原文（第 2--6 节、Table 1--4、Figure 6--20）为准。\cite{li2025aiore}\par
\subsubsection{研究背景与问题定义}
该工作的问题定义并非抽象的安全愿望，而是一个具体的系统矛盾：DRAM 扩展受限于缩放瓶颈，CXL.mem 以 load/store 语义把远端内存挂进主机地址空间，但公有云场景要求这些内存处于 TEE 保护之下，而现有 XTS TEE 加 CXL IDE 的组合把密码计算放在了每次读的关键路径上，内存密集型负载的开销可达 14.9\%（论文第 1 节、第 6.1 节）。\par
具体而言，论文的基线是 Intel TDX/AMD SEV 风格的 counterless XTS TEE 与 CXL 2.0 IDE 的组合（论文第 2.5 节 Figure 5）：主机用 key 1 加密本地内存，CXL 模块用 key 2 加密 CXL 内存（密钥经 TEE-IO 与 TDISP 交换，使 CPU 能直接解密 key 2 密文），链路上的 68 字节 Flit（2 字节协议 ID、64 字节载荷、2 字节 CRC）再用 key 3 以 256 位 AES-GCM 做加密与认证，每最多 128 个 cacheline 传输构成一个 MAC epoch。Skid 模式下单边解密即可使用明文，Containment 模式下须先过 MAC 认证。该结构的问题在于：XTS 的第二次 AES 计算必须等待密文到达，无法隐藏；而改用 CTR 模式虽能在 counter cache 命中时让 OTP 计算与取数并行、把解密延迟压到接近不安全基线，却会引入两类新成本——冷页的计数器污染 counter cache，以及 split counter 的 minor counter 溢出触发整页重加密。论文第 3.2 节给出了定量动机：3-bit split counter 平均“其他”安全开销仅 5.8\%，但溢出处理另加 4.3\%，溢出相关流量达到数据访问量的 45\%；morphable counter 在已无完整性树的现代 TEE 假设下仍有 2.4\% 溢出开销和 20.3\% 溢出流量。更激进的混合方案 Counter Light 按 LLC 写回时的带宽阈值在 cacheline 粒度切换 XTS/CTR，论文用两组实验说明其局限：按带宽切换几乎不改变性能（Figure 9），而按页访问频率静态地把最热页面分配给 CTR 的最优配置可减少 28.2\% 安全开销（Figure 10），且 Counter Light 需要额外 16 字节 ECC 总线，与 CXL 的 68 字节 Flit 不兼容。由此论文把问题收敛为：如何按页访问热度动态选择加密模式，并把模式切换与计数器溢出引起的重加密移出关键路径。\par
\subsubsection{核心方法设计}
这一部分的核心结论是：AIORE 通过页热度跟踪器、增量重加密、卸载重加密三个设计，把“用哪种加密”变成运行时决策，把“换加密的代价”变成可被正常访问平摊或卸载到内存节点的后台工作（论文第 4 节 Figure 11）。\par
论文先把内存布局与威胁模型固定在基线之上（第 4.2.1 节 Figure 12）：CXL 内存页分为 CTR 加密、XTS 加密与 Re-encrypting 三类；CTR 页的计数器存放在 CXL 内存元数据区，并复用 Intel SGX v2 的动态计数器映射——每 enclave 一张 secure page metadata table 记录 PFN 到计数器地址的映射，该表与页表均以 per-enclave 密钥 XTS 加密；CXL 内存 PTE 中两个未用位记录页加密状态，使常规页的 LLC miss 可以跳过元数据查找。一个值得注意的工程细节是：每条 64 字节计数器 cacheline 含两个 major counter、各配 64 个 3-bit minor counter，使共享同一计数器行的两个页可以独立重加密。\par
\noindent\textbf{（1）页热度跟踪器（Page Hotness Tracker）。}\par
该机制的直接作用是：动态判定每个 4KB 页应当使用 CTR 还是 XTS，而不是对所有页采用同一策略。具体实现上（第 4.5 节 Figure 15），主机侧放置一个 64 项、8 路 LRU 的检测缓存，每项为 36 位 PFN 加 5 位访问计数器，只统计 LLC miss 与写回（cache 命中与 silent eviction 不计）。计数器溢出热度阈值且页当前为 XTS 时触发中断转入 CTR；项被淘汰时若计数低于冷度阈值且页非 XTS，则触发转回 XTS。为抑制抖动，热阈值与冷阈值初始分设 16 与 8；系统还周期性监测 counter cache 命中率，以 95\% 命中率为目标微调热阈值（评估中固定使用该目标）。切换代价由下面两个机制吸收。\par
\noindent\textbf{（2）增量重加密（Incremental Re-Encryption）。}\par
该机制的直接作用是：把模式切换或 minor counter 溢出引起的整页重加密，拆成随程序固有读写逐步完成的工作，消除传统方案中“读-解密-重加密-写回 64 条 cacheline 并全程阻塞该页”的停顿。具体实现上（第 4.3 节 Figure 13），主机 Root Complex 内集成 Incremental Re-Encryption Bitmap Cache（IREBC，128 项 LRU，访问延迟 30 cycle），每项含 36 位 PFN、3 位重加密状态与 64 位按 cacheline 的完成位图，替补时写回内存中的 IREBT 表。进入重加密状态的页，读访问先查位图：已完成的 cacheline 用新模式（溢出场景用 $(\text{major}+1)\ll3+\text{minor}$）解密，未完成的沿用旧模式；写回与 silent eviction 一律用新模式加密并置位位图。全部位图置位后更新 PTE 状态并回收计数器项。读路径因此不引入额外内存访问，写路径只多一次位图检查。\par
\noindent\textbf{（3）卸载重加密（Offloaded Re-Encryption）。}\par
该机制的直接作用是：当某页在一段时间内未被完整访问、增量重加密迟迟不能收尾时，把剩余工作交给 CXL 内存模块完成，避免 IREBC 项长期占用与主机关键路径受累。具体实现上（第 4.4 节 Figure 14），CXL 模块 Root Complex 内设置 Incremental Re-Encryption Bitmap Buffer（IREBB，32 项，每项 614 位、含该页 256 位 split counter），按先进先出处理，被逐出的 IREBC 项连同计数器行一并送入；CXL 模块复用其既有 XTS/GCM 内存加密引擎（MEE），在引擎空闲时逐条完成剩余 cacheline 的重加密，全部完成后回送 PFN，由主机更新 PTE 与计数器缓存。论文的实测中卸载重加密只占总加密工作量约 4.1\%；仅当访问恰好落在 IREBB 队首页时才短暂阻塞，而这些页本来就是冷页。三者合起来，AIORE 把“重加密”从一次全局维护操作变成了可被访问行为平摊、可被设备侧空闲算力吸收的后台过程。\par
\subsubsection{安全性与威胁边界}
论文第 2.5 节与第 4.8 节明确其威胁模型完全继承基线：TCB 包含处理器、内存控制器、CXL 控制器、CXL/PCIe 组件、enclave VM 与处理器保留内存；攻击者控制特权软件，可对 CPU--内存总线窃听但不能修改，对 CXL 链路可截获、篡改与重放。AIORE 新增的全部组件（IREBC、IREBB、热度检测缓存）都被声明置于基线 TCB 之内，机密性由 XTS/CTR 保证，链路的完整性、唯一性与新鲜性仍由 CXL IDE 的 GCM 保证；论文自认为不引入新攻击面。边界同样清楚：侧信道、CPU 缺陷与拒绝服务在威胁模型之外；论文第 4.8 节坦承自适应加密带来的可变访问延迟可能构成 covert/side channel，但援引工业 TEE 与 CXL IDE 排除时序信道的惯例处理，且未证明其泄露面不比既有 CTR 方案更大。对本文主题而言，关键定位是：AIORE 是 CXL 内存机密性开销的优化，它把完整性与重放保护整体托付给 CXL IDE，本身不提供内存完整性树或故障恢复机制。\par
\subsubsection{实验设计与证据分析}
从评估设计看，作者用模拟器端到端结果、消融、敏感性与带宽分解四组证据回答其核心问题。评估平台为 Gem5 v24.0.0.1（第 5 节 Table 3）：4 核乱序 x86-64、4.00GHz，三级缓存（8MB LLC、30 cycle），64KB 8 路专用 counter cache，AES 操作建模为 56 cycle、GCM 认证 40 cycle、XOR 1 cycle，CXL 往返延迟 70ns，设备侧为 DDR5 4400MHz。负载为 13 个内存密集型程序，取自 SPEC2017、SPEC2006 与 GraphBIG，从 ROI checkpoint 模拟 50 亿条指令；对比 7 个方案，包括 XTS IDE、CTR IDE、Counter Light（含/不含 ECC 总线）、MorphCTR（含/不含 BRRIP）与 ShieldCXL。\par
端到端结果（第 6.1 节 Figure 16--17）支持其设计动机：单核下 XTS IDE 平均比不安全 CXL 慢 10.5\%（内存密集型负载最高 14.9\%），CTR IDE 慢 8.3\%，MorphCTR 与 BRRIP 变体分别慢 8.1\%、7.9\%，ShieldCXL 因插入 dummy Flit 慢 15.2\%；八核带宽竞争加剧时三个 CTR 系方案退化最严重（最高 28\%），XTS IDE 反而以 9.9\% 占优，说明“CTR 恒优于 XTS”在多核下并不成立，这正是自适应切换的价值所在。AIORE 在 1--8 核下开销为 3.2\%--4.3\%（平均 3.7\%），相对全部基线平均减少 62.8\% 的安全开销（区间 35.3\%--78.3\%）。消融实验（Figure 18）把收益拆源：页热度跟踪与自适应加密相对 CTR IDE 贡献约 27.7\% 的开销削减（省去冷页计数器维护），3-bit split counter 相对 7-bit 贡献 15.6\%，增量加卸载重加密贡献 24.1\%；元数据带宽降至 CTR IDE 的约 20.8\%（Figure 20）。页划分结果（Figure 19）显示策略确实因负载而异：605.mcf\_s、619.lbm\_s、433.milc 过半页被划为 CTR 页，而 imagick\_s、kCore、astar 等低复用负载大多数页保持 XTS。硬件成本（Table 2）为主机侧 67496 字节片上存储（IREBC 1632B、热度缓存 328B、64KB counter cache）加设备侧 IREBB 2456 字节，合计 69952 字节；计数器存储约为数据量的 0.78\%；CACTI 估算新增存储泄漏功耗 60.41mW，约占既有片上缓存功耗的 0.83\%。敏感性（Table 4）表明 IREBC 从 64 增至 256 项仅把开销从 3.5\% 降到 3.1\%，热度检测缓存从 32 增至 128 项把开销从 3.7\% 降到 3.0\%，即默认配置已接近收益平台。证据边界必须同时写明：这是周期近似模拟器而非真实 CXL 控制器与内存节点固件，AES/GCM 延迟与 CXL 往返延迟取自先前工作的建模值，真实部署中的中断代价、MEE 空闲率与固件调度均未在论文中测量。\par
\subsubsection{与相关工作的定量对比}
论文 Table 1 把五条技术路线的性能与安全性质并排比较，是理解其设计取舍的最紧凑证据，本文将其整理为表~\ref{tab:aiore-cmp}。可以看到，CTR TEE 加完整性树的路线虽然 LLC miss 延迟低，却以每次写多次元数据访问、大元数据空间和放弃链路唯一性为代价；AIORE 的目标是在保留 CXL IDE 全部链路安全性质的前提下，把计数器溢出开销压到近零。\par
\begin{table}[htbp]
\centering
\caption{安全 CXL 内存各设计路线的性能与安全性质对比（整理自论文 Table 1）}
\label{tab:aiore-cmp}
\footnotesize
\begin{tabular}{@{}lccccccc@{}}
\toprule
设计 & LLC miss 延迟 & 带宽 & 元数据/次写 & 溢出开销 & 元数据空间 & 重放保护 & 唯一性 \\
\midrule
XTS TEE + CXL IDE（基线） & 高 & 零额外 & 0 & 零 & 零 & 是 & 是 \\
CTR TEE + CXL IDE & 中 & 中 & 1 次 & 低 & 小 & 是 & 是 \\
Split CTR TEE + CXL IDE & 低 & 低--中 & 1 次 & 高 & 小 & 是 & 是 \\
CTR TEE + 完整性树 & 低 & 高 & 多次 & 高 & 大 & 是 & 否 \\
AIORE（该论文） & 低 & 低 & 0 或 1 次 & 近零 & 小 & 是 & 是 \\
\bottomrule
\end{tabular}
\end{table}
在与具体系统的对比上：Counter Light 带 ECC 总线时相对 CTR IDE 减少 35.7\% 安全开销，但受 OTP memorization table 命中率限制且与 CXL Flit 不兼容；不带 ECC 总线时与 CTR IDE 持平（8.3\%），因为少数负载才能达到其带宽切换阈值。ShieldCXL 以 13.5\%--15.2\% 的开销换取 obliviousness，与 AIORE 目标正交。相关工作部分还定位了 Toleo（用可信智能内存替代 Merkle 树扩展 freshness）与 Stark 等对 CXL 内存安全粗粒度现状的批评，说明 CXL TEE 的完整性/语义层问题仍未被任何单篇工作闭环。\par
\subsubsection{综合评价}
综合来看，AIORE 是 CXL 内存加密开销优化的当前 SOTA：把安全开销从 XTS 基线的最高 14.9\%（平均 10.5\%）压到平均 3.7\%，且收益的每一部分都有消融数据支撑。\par
AIORE 的创新性在于没有把内存加密视为固定的单一原语，而是把页面热度、重加密时机和计算位置（主机还是内存节点）共同纳入设计，并复用 CXL 模块已有的加密引擎，避免新总线、新协议。其局限有三：其一，全部结果来自 Gem5 建模，论文未在真实 CXL 硬件上验证；其二，机密性之外的安全性质整体依赖 CXL IDE，内存侧的完整性与 freshness 仍是开放问题（CTR TEE 加完整性树的路线又被论文自己证明放弃唯一性）；其三，可变延迟的侧信道含义仅按惯例排除，未做泄露面分析。因此本文将其定位为 CXL 机密性开销优化的强证据，而不外推为完整的 CXL 内存安全方案。\par
\subsection{论文 28：Direct Access, High-Performance Memory Disaggregation with DirectCXL}
\subsubsection{论文基本信息}
《Direct Access, High-Performance Memory Disaggregation with DirectCXL》由 Donghyun Gouk et al. 完成，发表于 2022 年的 USENIX ATC（README 记录来源为 usenix.org ATC22 页面）。本文将其置于“内存与 I/O 互连：CXL、PCIe IDE、RDMA”的研究脉络中考察，并把它作为 system 类材料使用。它所回答的核心问题可概括为：DirectCXL 的核心贡献不是“更快网络”，而是把远端内存变成 CPU 可直接 load/store 的 CXL.mem 地址空间。当前证据状态为：本地 PDF 已保存并核验；下文仅在该范围内讨论其机制与结论。\cite{gouk2022directcxl}\par
\subsubsection{研究背景与动机}
要理解该方案，首先需要看到它面对的系统矛盾：RDMA memory disaggregation 的成本来自软件路径和数据拷贝，而不只是链路带宽。\par
具体而言，论文把既有路线分成三类并逐一指出其代价（README 证据基线为 Section 1--2 与 Figure 1）。Page-based 方案对应用透明，但关键路径上有 page fault、swap、context switch 和 I/O amplification；object-based 方案性能较好，却要求应用改用它自己的接口，语义绑定在特定对象系统上；RDMA 方案则需要预先注册 memory region、依赖 RNIC，并且远端读写往往要远端 runtime 或 CPU 参与。三条路线的共同点是：远端内存对 CPU 而言不是“地址”，而是“消息”，每次访问都要穿越软件栈和拷贝路径。DirectCXL 的赌注是，CXL 的出现使“把远端内存直接变成地址空间”第一次在硬件上可行。\par
\subsubsection{关键挑战与核心思想}
在上述背景下，作者提出的关键判断是：如果 remote memory 能变成系统地址空间，关键路径就从 message/copy path 变成 address path。\par
具体而言，这一判断分解为三个相互支撑的设计前提（Figure 2--Figure 4）。第一，CXL device 暴露 Host-managed Device Memory（HDM），使远端 DRAM 经 PCIe 枚举后进入 host 物理地址空间；第二，多主机、多内存设备场景需要 CXL switch 建立 virtual hierarchy 与路由，而不是点对点线缆，这把 memory pooling/partitioning 转化为 fabric 管理问题；第三，runtime 负责 namespace 管理，不让远端 CPU 参与每次访问——这是与 RDMA 路线的本质区别。\par
\subsubsection{系统架构}
从系统结构看，论文的基本组织方式是：CXL device 是被动内存模块，host 通过 PCIe/CXL fabric 访问它。\par
具体而言，数据通路为：PCIe enumeration 后 HDM 被映射进 host system memory；应用的 load/store 落到 HDM 地址区间后，root port 把请求转成 CXL flit，经 CXL switch 路由到目标 device；device 上的 CXL controller 直接访问本地 DRAM 并返回数据，全程无远端 CPU、无中间拷贝。runtime 侧（Figure 4），驱动暴露 /dev/directcxl 字符设备，ioctl 创建 cxl-namespace，应用以 mmap 按文件式接口映射远端内存，namespace 由 segment table 管理。这条路径的可核验证据范围是 DirectCXL Section 3.1 与 Figure 2--Figure 5。\par
\subsubsection{核心方法设计}
论文的关键不是单一组件，而是地址映射、交换结构、运行时三层机制组合成一条“无远端 CPU”的访问路径。\par
\noindent\textbf{（1）CXL.mem 地址映射。}\par
这一机制的直接作用是：让远端 DRAM 进入 host physical address space。应用 mmap cxl-namespace 后直接 load/store，HDM base address 将请求导向 CXL root port，全程不需要 host DRAM 中转拷贝。这把 page-based 方案的 fault/swap 路径和 RDMA 方案的 MR 注册/verbs 路径同时移出了关键路径。\par
\noindent\textbf{（2）CXL switch 与 virtual hierarchy。}\par
这一机制的直接作用是：多 host、多 memory device 的池化由 switch 的 routing table 完成，而不是点对点连线。CXL switch 建立 virtual hierarchy，把 host 与 logical device 连接起来，使 memory pooling 与 partitioning 成为 fabric 管理层的问题。从安全视角看，这也意味着 tenancy 边界将来必须在 fabric 层表达，但论文本身未处理这一问题。\par
\noindent\textbf{（3）Linux runtime 与 namespace。}\par
这一机制的直接作用是：把硬件内存资源变成应用可用对象。驱动暴露 /dev/directcxl，ioctl 创建 cxl-namespace，segment table 记录分配，mmap 让应用按熟悉的文件式接口使用远端内存。README 记录的接口路径为 /dev/directcxl → ioctl → /dev/cxl-ns0 → mmap。\par
\noindent\textbf{（4）安全边界（论文未覆盖部分）。}\par
这一机制的直接作用是：标出 DirectCXL 不能替代 confidential I/O 证据的位置。论文不处理 device identity，不处理 CXL IDE/链路机密性，不处理 TVM/Realm 的 ownership lifecycle；README 边界明确它是性能/系统论文，不提供 CXL IDE、SPDM/TDISP 或 confidential I/O 安全证明。\par
\subsubsection{安全性与威胁边界}
从安全边界看，DirectCXL 的论证范围严格限于性能。其数据路径把内存边界从单节点扩展到 fabric，这本身放大了机密计算要面对的问题：HDM 映射谁授权、switch 路由是否可被恶意 fabric 管理器重定向、远端 DRAM 内容是否加密、设备被替换后如何发现——论文均未讨论。因此在本节谱系中，DirectCXL 的价值是给出 CXL.mem 数据路径的性能基线与架构词汇，其安全栏应记为“完全留给上层机制”。\par
\subsubsection{实验环境与证据}
论文没有把这部分只写成结论，而是以如下材料界定其证据范围：实验环境是真实/原型 CXL-enabled cluster，比较 Local、RDMA、Swap/KVS 和 DirectCXL 四条路径。\par
具体而言（Section 4、Figure 5、Table 1），软件栈为 RISC-V Linux 5.13.19、DirectCXL runtime，以及移植的 FastSwap/HERD 基线；硬件为 CXL host/device/switch IP 与自建 CXL device 原型；负载包括 microbenchmark、DLRM、MemDB 与图计算等。需要注意的证据边界是：这是原型 IP 而非商用 CXL 芯片，绝对延迟数字不应直接外推到产品形态，但路径长短的相对结论（地址路径 vs 消息路径）由架构决定，较为稳固。\par
\subsubsection{实验结果与证据强度}
对于性能或效果，论文能够支持的结论是：DirectCXL 显著缩短远端内存访问路径，但仍慢于本地 DRAM。\par
具体而言，三组数字构成本文最常引用的证据。Figure 6：64B 访问 RDMA 需 2705 cycle，DirectCXL 只需 328 cycle，快 8.3 倍。Figure 8（含 switch 层次结构）：RDMA 最好情形仍需 2027 cycle，约为 DirectCXL 的 6.2 倍；同图给出本地 DRAM L2 miss 约 60 cycle，说明 DirectCXL 的远端访问仍比本地慢约 5 倍量级。Figure 10 与结论：真实负载（DLRM、MemDB、图计算）上 DirectCXL 约为 RDMA-based disaggregation 的 3 倍性能。这些数字的支撑范围限于论文 Figure 6、Figure 8/Table 2 与 Figure 10，本文不将其扩展为“任何 CXL 内存都快于 RDMA”的一般结论。\par
\subsubsection{综合评价}
综合来看，DirectCXL 强在证明 CXL.mem data path 的价值，弱在完全不覆盖 confidential trust。\par
具体而言，它的优势在于系统原型完整、数字直观，清楚解释了 CXL 为什么改变 memory boundary：远端内存从“消息对象”变成“地址对象”，swap/fault/拷贝/远端 CPU 全部离开关键路径。它的局限在于无 attestation、无 link encryption、无 device lifecycle、无 multi-tenant ownership。从部署角度看，它直接支撑 CXL memory pooling/expansion 场景；在机密计算语境下，它给出的是必须被 SPDM/TDISP/IDE/TEE policy 叠加保护的那条数据路径本身。\par
\subsection{论文 29：Managing Memory Tiers with CXL in Virtualized Environments}
\subsubsection{论文基本信息}
《Managing Memory Tiers with CXL in Virtualized Environments》由 Yuhong Zhong et al. 完成，发表于 2024 年（README 记录来源为 USENIX OSDI24 页面）。本文将其置于“内存与 I/O 互连：CXL、PCIe IDE、RDMA”的研究脉络中考察，并把它作为 system 类材料使用。它所回答的核心问题可概括为：CXL-Tiers 的关键不是“给 VM 加慢内存”，而是让硬件做 cache-line 粒度的 tiering，让软件只处理多租户 outlier。当前证据状态为：本地 PDF 已保存并核验；下文仅在该范围内讨论其机制与结论。\cite{zhong2024cxltiers}\par
\subsubsection{研究背景与动机}
要理解该方案，首先需要看到它面对的系统矛盾：虚拟化环境下的软件 tiering 很难既快又准。\par
具体而言（README 证据基线为 Section 1--2、Figure 1--2），困难来自三处。其一，host 看不到 guest 的真实访存热度，要么依赖高成本的追踪（PTE 扫描、指令采样），要么接受不准确的估计；其二，2MB/1GB huge page 让冷热 cache line 混在同一页内，页粒度迁移天然错配；其三，多租户共享 local DRAM 时，竞争会制造 outlier——个别 VM 的尾部 slowdown 远超平均值。CXL 内存容量便宜但延迟高，这三个困难决定了一个分工问题：细粒度热度判定应该交给谁。\par
\subsubsection{关键挑战与核心思想}
在上述背景下，作者提出的关键判断是：硬件做细粒度 hotness，软件做多租户公平。\par
具体而言（Figure 3--4、Figure 8--9），Intel Flat Memory Mode 在 memory controller 内以 cache-line 粒度在 local DRAM 与 CXL 之间搬运数据，避开页粒度 hotness 错配；Mixed Mode 在此基础上保留一块 dedicated local DRAM，供 latency-sensitive VM 使用；Memstrata 则不进 guest、不追踪页热度，而是用 per-VM performance events 和 ML 估计器预测哪些 VM 已成为 slowdown 超过 5\% 的 outlier，再由动态页分配器在 VM 间转移 dedicated local pages。\par
\subsubsection{系统架构}
从系统结构看，论文的基本组织方式是：VM 看到一个内存池，硬件在 local/CXL 间搬 cache line，Memstrata 在 VM 间分配 dedicated local pages。\par
具体而言，硬件 tiered NUMA node 聚合 local DRAM 与 CXL memory，对 guest 完全透明；Mixed Mode 额外暴露一个 dedicated local NUMA node；Memstrata 作为 host 侧控制器收集性能事件、运行估计器并触发页迁移，guest 不作任何修改。架构上最值得注意的是信任位置的安排：性能遥测和页迁移策略都在 host 侧，这正是它在机密计算语境下需要被重新审视的地方。\par
\subsubsection{核心方法设计}
论文的关键不是单一组件，而是“硬件 tiering + 软件 outlier 管理”这一分工下的三个机制与一个边界教训。\par
\noindent\textbf{（1）Flat / Mixed Memory Mode。}\par
这一机制的直接作用是：Flat Memory Mode 把 local DRAM 当作 CXL tier 的硬件 cache，以 cache-line 粒度换入换出，避免页级热度错配；Mixed Mode 在此之外留出 dedicated local DRAM，使 latency-sensitive VM 能拿到稳定的快页。两者对 guest 均透明，无需修改。\par
\noindent\textbf{（2）Slowdown Estimator。}\par
这一机制的直接作用是：Memstrata 不追踪每页 hotness，而是直接预测“该 VM 是否已成为 outlier”。输入是 per-VM performance events，目标阈值是 slowdown 超过 5\%；估计器刻意避开高成本的 PTE scanning 与 instruction sampling，使 host 开销可控。这一步把问题从“热度测量”换成“受害识别”，是它能同时做到低成本和多租户公平的关键。\par
\noindent\textbf{（3）Dynamic Page Allocator。}\par
这一机制的直接作用是：把 dedicated local pages 从低风险 VM 转移到被预测的 outlier VM。只有被判定为 outlier 的 VM 获得更高的分配优先级；算法在多种负载组合下压低 worst-case slowdown。其内在风险是预测错误和页迁移自身的开销，论文在评估中分别测量了这两点。\par
\noindent\textbf{（4）对机密计算的边界教训。}\par
这一机制的直接作用是：CXL-Tiers 提示资源管理本身也会成为 side/control boundary。per-VM 性能遥测可能泄露 guest 的负载行为；host 的页迁移策略直接决定 CVM 的尾部延迟；论文不提供任何机密性证明，只提供 tiering 机制证据。换言之，CVM 场景下“让 host 管理内存层级”这一看似纯性能的设计，隐含了把一条影响 tenant 体验的控制通道交给了不可信方。\par
\subsubsection{安全性与威胁边界}
从安全边界看，论文的威胁面是性能隔离而非机密性：它要防的是 noisy neighbor 造成的尾部 slowdown，不是恶意 host。因此对机密计算而言，它的结论应被双向解读：一方面，hardware tiering 对 guest 透明、不依赖 guest 配合，与 CVM 模型兼容；另一方面，估计器依赖的性能事件、host 控制的页迁移与 dedicated local 分配策略，在 CVM 威胁模型下都属于不可信控制面，论文未分析这些机制被滥用时的后果，本文据此将其安全栏记为“机制可用、信任未闭环”。\par
\subsubsection{实验环境与证据}
论文没有把这部分只写成结论，而是以如下材料界定其证据范围：实验覆盖真实 CXL 内存、修改过的 Linux/QEMU/KVM 与大规模 VM 负载组合。\par
具体而言（README 证据基线为 implementation/evaluation sections），系统侧为修改的 Linux kernel 与 QEMU/KVM，硬件为 CXL 卡与 Intel memory controller（Flat/Mixed Memory Mode 为真实硬件特性而非模拟）；负载为 115 个 workload，覆盖 web、data、Spark、ML、GAP、SPEC 等类别的组合；指标包括 slowdown、outlier 比例、CPU overhead 与 memory overhead。使用真实 CXL 硬件和真实内存控制器模式，是其证据强度高于纯模拟研究的地方。\par
\subsubsection{实验结果与证据强度}
对于性能或效果，论文能够支持的结论是：硬件 tiering 对大多数负载可接受，但 outlier 需要软件隔离策略来兜底。\par
具体而言（abstract、Figure 5、Figure 10--12），Mixed Mode 下 82\% 的负载 slowdown 不超过 5\%，95\% 不超过 10\%，但 outlier 可达 34\%——这量化了“为什么光靠硬件不够”。Memstrata 把 realistic multi-VM 场景的 worst-case slowdown 从 35\% 压到 6\% 以下，而自身代价为最高占用单核 4\% 的 CPU 与约 110MB 内存。这组数字的读法是：硬件机制解决均值，软件策略解决尾部，两者的开销都被量化；本文不把“worst-case <6\%”外推为任意负载组合的保证，因为估计器本质上依赖负载组合的统计性质。\par
\subsubsection{失败模式与边界分析}
论文自身暴露的失败模式有三类。第一，outlier 残余：Mixed Mode 下仍有负载 slowdown 达 34\%，Memstrata 介入后 worst-case 仍接近 6\%，说明基于预测的分配无法给出硬保证。第二，预测错误风险：allocator 的效果依赖估计器准确性，论文承认预测错误与迁移开销是其内在代价，但未给出对抗性负载（刻意制造误判）下的表现。第三，也是最贴近本文主题的一点：论文的性能遥测与迁移策略都位于 host，对 CVM 而言这是不可信控制面；遥测是否泄露负载指纹、恶意迁移策略能否制造可观测的尾部延迟信道，论文未说明，证据不足。\par
\subsubsection{综合评价}
综合来看，CXL-Tiers 是 CXL 虚拟化内存层级管理的 SOTA，但不是安全机制。\par
具体而言，它的优势在于分工清楚（硬件做粒度、软件做公平）、数据充分（115 个负载、真实硬件）、并准确刻画了 CXL tiering 在云 VM 中的尾部风险。它的局限在于不处理 confidential memory ownership、device identity 与 link security，且其遥测/控制面位置与 CVM 威胁模型存在结构性张力。从部署角度看，它直接适用于云 CXL capacity tier 管理；机密计算场景若借用其机制，必须先回答 telemetry 与 host policy 的信任问题。\par
\subsection{论文 30：ODRP: on-demand remote paging with programmable RDMA}
\subsubsection{论文基本信息}
《ODRP: on-demand remote paging with programmable RDMA》由 Zixuan Wang et al. 完成，发表于 2025 年（README 记录来源为 USENIX NSDI25 页面）。本文将其置于“内存与 I/O 互连：CXL、PCIe IDE、RDMA”的研究脉络中考察，并把它作为 system 类材料使用。它所回答的核心问题可概括为：ODRP 的贡献是把 remote paging 的动态分配逻辑从 MNode CPU 搬到 programmable RDMA 的 WR chain 上。当前证据状态为：本地 PDF 已保存并核验；下文仅在该范围内讨论其机制与结论。\cite{wang2025odrp}\par
\subsubsection{研究背景与动机}
要理解该方案，首先需要看到它面对的系统矛盾：远端内存解耦存在一个三难（trilemma）——内存利用率、MNode CPU 占用、访问性能三者难以兼得。\par
具体而言（README 证据基线为 Section 1--2、Figure 1），static pre-registration 路线（one-sided）不消耗 MNode CPU，但每个 CNode 预注册的大块远端内存大量闲置，碎片与利用率都差；dynamic/two-sided 路线利用率高，但 4KB 粒度的 paging 需要频繁 allocation/free，MNode CPU 被迫处理每个请求而成为瓶颈。三难的本质是：分配决策的“聪明程度”与它消耗的远端算力成正比，而远端算力恰恰是解耦架构想省下的东西。\par
\subsubsection{关键挑战与核心思想}
在上述背景下，作者提出的关键判断是：RDMA work request 可以组合成小型控制程序，把 remote paging 的控制路径卸载到 RNIC。\par
具体而言，ODRP 利用可编程 RDMA 的 WAIT、ENABLE、CAS、FAA 等原语，把“查转换表、从空闲队列取页、更新映射”这一分配序列编成 WR chain，由 RNIC 在数据路径附近执行；MNode 的内存仍保存数据结构，但 CPU 不处理任何请求；CNode 只维护轻量元数据。这一判断的技术含量在于：WR chain 必须在无远端 CPU 协调的前提下，保证多 CNode 并发分配的一致性。\par
\subsubsection{系统架构}
从系统结构看，论文的基本组织方式是：CNode 发起 fault/swap，RNIC 执行 WR chain，MNode 内存保存 translation table 与 free page queue（Section 4、Figure 6）。\par
具体而言，translation table（TT）记录 virtual remote page 到已分配物理页的映射；free queue 支撑按需分配；CNode 侧处理 page fault 后不再走远端 RPC，而是向 RNIC 提交一条预编的 WR chain；RDMA atomic 操作保证并发更新的一致性。架构的关键性质是控制路径与数据路径都终止于 RNIC/内存，MNode CPU 完全出局。\par
\subsubsection{核心方法设计}
论文的关键不是单一组件，而是 WR chain 程序、按需分配与一致性论证三者的组合。\par
\noindent\textbf{（1）WR chain 作为远端控制路径。}\par
这一机制的直接作用是：用 WR chain 完成查表、分配、更新，而不是唤醒远端 CPU。WAIT/ENABLE 原语组织操作间的依赖关系（如前一个 WR 的结果触发后一个 WR），CAS/FAA 完成并发安全更新；典型序列为 READ TT → 若未映射则对 free queue 做 CAS/FAA → WRITE TT → READ/WRITE 目标页。链式执行把分配逻辑放到 RNIC 数据路径附近，省掉了远端 CPU 的调度与上下文切换。\par
\noindent\textbf{（2）按需 4KB 分配。}\par
这一机制的直接作用是：让远端内存利用率接近实际使用量。static 预留把大量空页锁死在各 CNode 名下；ODRP 不要求任何 CNode 预注册大块远端内存，按需从共享 free queue 取页。代价是细粒度分配引入元数据与原子操作成本，这正是其性能 overhead 的来源，论文用 0.8\%--14.6\% 的区间刻画了这笔账。\par
\noindent\textbf{（3）一致性与并发。}\par
这一机制的直接作用是：保证多个 CNode 并发分配不会把同一页分出去。论文用 RDMA atomic 语义加归纳论证（induction argument）分析正确性：TT 与 free queue 的更新必须满足顺序约束，核心不变式是“一页一主”、TT 更新原子、free queue 一致。论文明确列出了不变式被破坏的后果：页别名（alias）、数据损坏或页泄漏。\par
\noindent\textbf{（4）安全边界（论文未覆盖部分）。}\par
这一机制的直接作用是：标明 ODRP 只解决性能/利用率，不解决 remote memory trust。它不证明 MNode 不会读写数据，不提供链路加密或 attestation；机密内存/存储需要 Hazel、SPDM、TEE-I/O 类机制叠加。README 边界同样注明它是 disaggregated memory performance paper，不提供 confidential link/device/TEE proof。\par
\subsubsection{安全性与威胁边界}
从安全边界看，ODRP 的威胁模型是“性能正确性”而非“恶意对手”：其一致性论证针对并发错误，不针对 Byzantine MNode。对本文主题而言这意味着两点：第一，MNode 及其 RNIC 被默认诚实，远端内存内容的机密性、完整性完全不在保护范围；第二，其 WR chain 一致性论证（原子操作 + 归纳）是正确性证据而非安全证据——恶意 MNode 可以直接读写 TT 与 free queue 制造论文自己列出的 alias/corruption/leak。若要把 ODRP 类机制用于机密场景，分配元数据必须移入受信域或加以完整性保护，论文未讨论此方向。\par
\subsubsection{实验环境与证据}
论文没有把这部分只写成结论，而是以如下材料界定其证据范围：实验比较 one-sided static/dynamic、two-sided、Fastswap 等基线（Section 5、Figure 7--9）。\par
具体而言，评估指标覆盖三难的全部三条边：remote memory utilization、MNode CPU usage、execution time/throughput，外加 swap throughput 与 CNode 可扩展性；负载为 Redis、VoltDB 等典型 memory-disaggregation 场景。证据边界：支撑 remote paging 的性能与利用率结论，不支撑任何 confidential trust 结论。\par
\subsubsection{实验结果与证据强度}
对于性能或效果，论文能够支持的结论是：ODRP 用少量性能开销换取远端内存利用率与 MNode CPU 的大幅改善。\par
具体而言（abstract、Figure 7、Figure 9），三组数字分别对应三难的三条边。利用率边：remote memory utilization 提升 1.72--12 倍。CPU 边：MNode CPU usage 为 0，而 two-sided 方法要消耗 86\%--99\% 的 MNode CPU。性能边：overhead 约 0.8\%--14.6\%，部分 throughput 场景约 14.2\%。这组结果的读法是：ODRP 并没有“解决”三难，而是把三难转换成一个明确的报价——用至多约 15\% 的性能换零远端 CPU 和接近按需的利用率；本文不外推其在更高并发或更严格一致性要求下的表现。\par
\subsubsection{失败模式与边界分析}
论文给出的失败模式分析集中在一致性一侧：若 TT/free-queue 更新顺序被破坏，会出现页别名、数据损坏或页泄漏；其归纳论证依赖 RDMA atomic 的顺序语义，论文未说明在 RNIC 固件缺陷或重排序实现差异下论证是否仍成立，证据不足。性能侧的边界是 overhead 上界 14.6\%（吞吐场景 14.2\%）出现在分配密集负载上，说明 WR chain 并非免费。安全侧的失败模式论文完全未覆盖：恶意 MNode 可在不违反任何一致性不变式的情况下窃听或篡改页内容。\par
\subsubsection{综合评价}
综合来看，ODRP 是 RDMA 解耦内存方向三难取舍最清楚的系统 SOTA，但它把安全问题完整留给上层。\par
具体而言，它的优势在于：利用 programmable RDMA 这一真实可得的硬件能力，把三难的每条边都给出量化数字，且 WR chain 设计不依赖远端 CPU。它的局限在于：不处理 confidentiality/integrity/freshness，不证明 remote endpoint trust，一致性论证也未形式化到模型检查级别。从部署角度看，它适用于数据中心远端内存池与云 paging；机密计算若采用，必须在其上叠加加密、attestation 与策略绑定。\par
\subsection{论文 31：Security Protocol and Data Model (SPDM) Specification}
\subsubsection{论文基本信息}
《Security Protocol and Data Model (SPDM) Specification》由 DMTF SPDM Working Group 制定，本文以 DMTF DSP0274 v1.4.0（2025 年发布，本地 PDF 于 2026-05-12 下载并核验）作为本综述的 SPDM 快照。本文将其置于“机密 I/O 协议、可信设备接口与网络端点”的研究脉络中考察，并把它作为 spec 类材料使用。它所回答的核心问题可概括为：SPDM 让平台能问设备三件事——你是谁（证书链与 challenge）、你的固件/配置处于什么状态（measurement）、我们能否建立安全会话（key exchange）。当前证据状态为：本地规范 PDF 已保存并核验；作为标准规范，它不含有实验，下文只讨论其协议语义与覆盖边界。\cite{dmtf_spdm_2025}\par
\subsubsection{研究背景与动机}
要理解该规范，首先需要看到它面对的系统矛盾：当机密计算把边界扩展到真实设备后，TVM/Realm 不能只相信 PCIe requester ID 或驱动自报的配置。\par
具体而言，加速器、NIC、存储端点在被纳入信任边界前，平台需要三类证据（README 证据基线为 SPDM overview）：证书链证明设备身份或其供应链 root；measurement 证明固件/配置/安全状态；secure session 保护后续控制面消息。三者分别对应“是谁”“是什么状态”“接下来怎么谈”，缺一不可——只有身份没有测量，无法发现降级固件；只有测量没有会话，证据本身会被篡改。而数据面大流量的保护（IDE/TLS/IPsec/设备协议）在 SPDM 范围之外。\par
\subsubsection{关键挑战与核心思想}
在上述背景下，规范的核心定位是：SPDM 是 evidence protocol，不是 complete trusted I/O system。\par
具体而言，它把底层安全能力标准化为可被上层机制复用的控制面：TDISP 用它确认设备接口身份，PCIe IDE 的密钥管理、CXL 设备接入、RISC-V CoVE-IO 的设备证据都可以建立在 SPDM 消息之上；但“measurement 是否满足策略”这一判断被刻意留给 verifier/TSM，SPDM 本身不规定每个设备 measurement 的业务语义。这一克制定位既是它成为共同底座的原因，也是它不能单独构成机密 I/O 的原因。\par
\subsubsection{系统架构}
从协议结构看，SPDM 的基本组织方式是：Requester 与 Responder 通过 request-response 消息序列建立身份、measurement 和 session（README 证据基线为 Section 7--11、Figure 2、Table 3）。\par
具体而言，Requester 可以是 host、TSM 或 root port，Responder 是设备端点；一条连接上可存在多个 secure session；消息格式由 SPDM version 与 request/response code 约束，传输层可绑定 MCTP、PCIe DOE 或 TCP（各有独立绑定规范）。标准机制路径为：先经 GET\_VERSION、CAPABILITIES、ALGORITHMS 协商版本与算法，再由 DIGESTS 与 GET\_CERTIFICATE 取回证书链，CHALLENGE 验证私钥持有，GET\_MEASUREMENTS 取回测量，最后经 KEY\_EXCHANGE 与 FINISH 建立会话并转入 secured messages。\par
\subsubsection{核心方法设计}
规范的关键不是单一消息，而是身份、测量、会话三组机制的衔接，以及它们与上层 TEE-I/O 机制的接口。\par
\noindent\textbf{（1）证书与挑战（身份）。}\par
这一机制的直接作用是：GET\_CERTIFICATE 与 CHALLENGE 证明端点身份和私钥持有（README 证据基线为 Section 10.9--10.10、Figure 1）。DIGESTS/GET\_CERTIFICATE 取回证书摘要与证书链；CHALLENGE\_AUTH 要求 responder 对 nonce 与 transcript 签名，证明其持有对应私钥且消息未被中间人替换。这把“设备身份”从一个可伪造的 requester ID 提升为可验证的密码学声明。\par
\noindent\textbf{（2）测量（状态证据）。}\par
这一机制的直接作用是：GET\_MEASUREMENTS 把设备固件/配置状态交给 requester 判断（Section 10.12）。measurement block 可被签名或绑定到 transcript；哪些 measurement 可接受由策略决定；规范刻意不规定各设备 measurement 的业务语义——这意味着“measurement 返回成功”不等于“设备状态可接受”，解释责任在 verifier。\par
\noindent\textbf{（3）安全会话（控制面保护）。}\par
这一机制的直接作用是：KEY\_EXCHANGE/FINISH 建立 SPDM secured messages（Section 7.4、10.17、11）。会话密钥保护后续 SPDM 载荷的机密性与完整性，transcript hash 防止握手被篡改或降级。边界是：会话保护的是 SPDM 控制面消息本身，数据面大流量仍须依赖 IDE、TLS、IPsec 或设备自有协议。\par
\noindent\textbf{（4）在 TEE-I/O 中的使用位置。}\par
这一机制的直接作用是：SPDM 证据必须被 CoVE-IO/TDISP/TSM 绑定到 device assignment 才有机密计算意义。设备身份本身不等于“已安全分配给 TVM”：TDISP 管理 device interface 状态机，IOMMU/IOPMP 管理 DMA 窗口，IDE 管理链路加密，SPDM 提供的是这些机制消费的身份与测量证据。\par
\subsubsection{安全性与威胁边界}
从安全边界看，SPDM 的安全性质依赖证书链、root of trust、measurement 语义、算法协商、transcript binding、随机数质量与实现正确性。它能证明 responder 的身份与所报告的 measurement，但不能单独保证：设备内部 TDI 隔离正确、DMA 权限受限、IDE 链路已启用、TDISP 状态机处于受信状态、TVM 内存所有权正确。README 第 8 节把这一点概括为：SPDM 只解决 control-plane trust bootstrap；若设备内部隔离、IOMMU/IOPMP、IDE、TDISP 状态机或证书供应链任一环节出错，SPDM 握手的成功不能推出端到端机密 I/O 安全。此外，版本协商降级、算法弱化和 responder reset 等路径被列入“验证应覆盖”的范围，说明规范语义正确不等于实现无漏洞。\par
\subsubsection{实验环境与证据}
规范无实验，本文把这一部分写成协议覆盖与证据边界。\par
具体而言，本地证据源为 DMTF DSP0274 v1.4.0 PDF（README 证据基线为 sections 7--11、Figure 1 证书链模型、Figure 2 协议流、Table 3 通用消息格式与请求/响应表）。这些材料能够直接支撑的结论是：SPDM 定义了消息交换、证书、测量与会话的完整语义，并有 libspdm 等开源实现生态。这些材料不足以支撑的结论是：任何具体设备的 assignment 正确性、任何部署的性能开销，本文不作此类推断。\par
\subsubsection{实验结果与证据强度}
对于性能或效果，规范能够支持的结论是：SPDM 定义协议，不提供 workload benchmark。\par
具体而言，其开销来源可确定为 certificate transfer、signature verification、measurement retrieval 与 key exchange，这些都发生在 setup/control path 而非每个数据包；数据面性能取决于 IDE/TLS/设备协议栈，超出本规范范围。README 给出的最小复现验收标准（transcript 可复验、measurement 与设备状态绑定、算法降级被拒绝）可作为实现一致性测试的参照。\par
\subsubsection{综合评价}
综合来看，SPDM 是 device attestation 的必备底座，但不能单独讲成 confidential I/O。\par
具体而言，它的优势在于标准化、可互操作，把 identity、measurement、session 三件事定义为 PCIe/CXL/DPU/accelerator 生态的共同语言；它的局限在于不管 DMA、不管中断、不管 assignment lifecycle、不管数据面链路。从部署角度看，任何“机密 I/O”产品声明若只提 SPDM，覆盖的仅是证据引导段；完整链路还需要 TDISP 的状态机、IDE 的链路与 IOMMU/IOPMP 的隔离共同闭合。\par
\subsection{论文 32：SmartNIC Security Isolation in the Cloud with S-NIC}
\subsubsection{论文基本信息}
《SmartNIC Security Isolation in the Cloud with S-NIC》由 Yang Zhou、Mark Wilkening、James Mickens 和 Minlan Yu（Harvard University）完成，发表于 2024 年的 EuroSys。本文将其置于“SmartNIC、可信网卡与安全存储数据路径”的研究脉络中考察，并把它作为 system 类材料使用。它所回答的核心问题可概括为：S-NIC 让 SmartNIC 上的 tenant network function 像各自拥有独立的 virtual NIC 一样被硬件级隔离。当前证据状态为：本地 PDF 已保存并核验；下文仅在该范围内讨论其机制与结论。\cite{zhou2024snic}\par
\subsubsection{研究背景与动机}
要理解该方案，首先需要看到它面对的系统矛盾：SmartNIC 自身已变成多租户计算平台，而 NIC OS 不能默认可信。\par
具体而言（README 证据基线为 Section 2--3、Figure 1），packet processing、DPI、NAT、ZIP 压缩、RAID、load balancer 等 network function 可能同时运行在同一块 SmartNIC 上，共享 DRAM、可编程核、加速器与管理 OS。论文实证展示了规则泄漏、packet corruption、bus DoS 等攻击：共享 allocator、加速器或总线都可造成跨 function 的信息泄漏与篡改。传统 NIC 的 SR-IOV 只虚拟化 PCIe 接口，完全不覆盖 SmartNIC 内部的可编程资源。对机密计算而言，这个背景的含义是：把工作负载的 packet processing/offload 放到 NIC 上，并不会因为“在主机之外”而天然安全。\par
\subsubsection{关键挑战与核心思想}
在上述背景下，作者提出的关键判断是：NIC 上也需要类似 enclave/VM 的资源所有权语义，但所有权的对象必须是 NIC-local 资源。\par
具体而言，威胁模型允许攻击者控制其他 network function 甚至 NIC management OS，保护目标是 function state 的机密性与完整性，并消除 NIC 内部共享状态侧信道。核心思想被概括为 single-owner semantics：每块 RAM、每条 TLB 项、每个加速器上下文、每个总线时隙在任意时刻只能属于一个 function 或管理 OS。memory denylist 防止 NIC OS 访问 function 私有内存，locked TLB entries 固定 NF 地址翻译，accelerator virtualization 与 bus arbitration 防止共享硬件越权。\par
\subsubsection{系统架构}
从系统结构看，论文的基本组织方式是：S-NIC 在 SmartNIC SoC 内加入硬件隔离机制与 host 可见的管理 API（Figure 2、Table 1）。\par
具体而言，方法管线为：tenant 上传 network function → NIC OS 请求 nf\_launch → 受信 S-NIC 硬件分配 virtual NIC 资源 → 硬件锁定 RAM/cache/accelerator/DMA/bus 所有权 → function 作为隔离的 virtual NIC 执行 → nf\_destroy 时清零状态。关键模块包括 per-core memory denylist、locked-down TLB、dedicated cache allocation、accelerator TLB banks、DMA isolation、trusted bus arbitration、attestation protocol 与 VXLAN/L2 集成。attestation 让 host/tenant 确认 function 由经认证的 S-NIC 以已知初始状态启动。\par
\subsubsection{核心方法设计}
论文的关键不是单一组件，而是四类资源的所有权机制加一条带证明的生命周期。\par
\noindent\textbf{（1）Locked TLB 与 memory denylist（内存所有权）。}\par
这一机制的直接作用是：建立 NIC-local single-owner memory semantics。NF launch 时锁定地址翻译，此后 NIC OS 不能更新该 NF 的 TLB 项；denylist 直接阻止管理软件访问 function memory。两者分别管住“翻译层”和“权限层”，使恶意的 NIC OS 既不能改映射也不能直接读。\par
\noindent\textbf{（2）加速器虚拟化（加速器所有权）。}\par
这一机制的直接作用是：SmartNIC 上的 DPI/ZIP/RAID 等加速器也必须虚拟化并隔离。每个 virtual accelerator 分配独立状态（DPI 的规则/缓存、ZIP 的压缩缓冲、RAID 的元数据），硬件检查每个请求属于哪个 function，避免一个 NF 读取另一个 NF 的加速器元数据——论文演示的规则泄漏攻击正是针对这一层。\par
\noindent\textbf{（3）Cache 与总线划分（微架构所有权）。}\par
这一机制的直接作用是：即使内存权限正确，共享 cache 与总线仍能造成性能与信息干扰。S-NIC 加入 bus arbiter 与 cache partition，评估报告最坏 IPC 退化约 1.66\%。这是把 NIC 内部侧信道当作系统性问题处理，而不是个例修补。\par
\noindent\textbf{（4）带证明的 function 生命周期。}\par
这一机制的直接作用是：S-NIC 提供 launch/attest/destroy 的完整生命周期，而非静态隔离。nf\_attest 证明 function 由经认证的 S-NIC 启动；nf\_destroy 清零资源状态；生命周期 API 让 host 能管理 NF 但不能越权改私有资源。边界是：这仍不等同于 TVM 的 trusted device assignment——它证明的是“function 在真 S-NIC 上启动”，不证明“设备已安全绑定到某个 CVM”。\par
\subsubsection{安全性与威胁边界}
从安全边界看，S-NIC 防御的是：NIC OS 读取 function RAM、其他 function 访问加速器状态、cache/总线侧信道。论文明确的范围外威胁包括：网络观察者侧信道、物理攻击、host 级侧信道与 DoS。强隔离的代价也被如实记录：资源利用率下降，且默认每个 virtual NIC 只跑单个 function，复杂 function chaining 需要 Rust/编译器强制隔离或扩展机制。它也不是完整的机密计算平台：还需要与 TEE evidence、key release、SPDM/TDISP 或 DPU RoT 组合。\par
\subsubsection{实验环境与证据}
论文没有把这部分只写成结论，而是以如下材料界定其证据范围：SmartNIC 架构模拟/建模 + 六类 network function（README 证据基线为 evaluation section、Tables 2--8、Figure 5--8）。\par
具体而言，评估的 function 包括 DPI、NAT、ZIP、RAID、load balancer、monitor 等六类；指标涵盖芯片面积/功耗估计、throughput/IPC 退化、指令延迟与内存占用。证据边界必须写明：这是 gem5 类模拟与硬件成本估算，不是生产硅片，面积/功耗数字是设计评估而非实测；其强项是明确给出硬件成本与侧信道隔离代价的定量关系。\par
\subsubsection{实验结果与证据强度}
对于性能或效果，论文能够支持的结论是：S-NIC 的隔离成本在其评估中较小，主要硬件成本集中在 TLB/cache/bus 支持。\par
具体而言（abstract、evaluation Tables 2--5、Figure 5--6），三组数字构成核心结果。性能面：强隔离下 function throughput 最坏下降小于 1.7\%，高 co-tenancy 时最坏 IPC 退化约 1.66\%。硬件成本面：4 核配置下芯片面积增加最高 8.89\%，功耗增加最高 11.45\%。管理面：nf\_attest 等管理操作的延迟受 RSA 等密码运算影响，不在 packet fast path 上——这说明其设计正确地把密码成本移出了数据路径。这些数字的支撑范围限于论文评估配置，本文不外推到任意 SmartNIC 产品。\par
\subsubsection{失败模式与边界分析}
论文自述的局限构成本节的失败模式分析。第一，资源利用率：single-owner 语义意味着加速器与 cache 分区在 function 空闲时不能被他人使用，强隔离与资源效率存在结构性矛盾。第二，function chaining 受限：默认模型下一条处理链要拆成多个 virtual NIC 或依赖编译器级隔离，复杂链路的代价未量化。第三，无生产硅片：全部成本数字来自建模。第四，覆盖缺口：不处理 SPDM/TDISP、TVM binding 与链路加密；README 明确其边界为“不等同于完整 VM/Realm trusted I/O”。\par
\subsubsection{综合评价}
综合来看，S-NIC 是 SmartNIC 内部多租户隔离的强基线，但它不是完整的 confidential I/O。\par
具体而言，它的优势在于把 NIC 内部资源对象拆得清楚（内存/TLB/cache/加速器/总线五类所有权），硬件成本被量化，问题设定直接命中 DPU/NIC 多租户的现实痛点。它的局限在于无生产 silicon、不覆盖设备级信任链。从部署角度看，它是 DPU tenant NF isolation 的设计参照；落地风险集中在功能链支持、资源利用率与 vendor attestation 的补齐。\par
\subsection{论文 33：TNIC: A Trusted NIC Architecture}
\subsubsection{论文基本信息}
《TNIC: A Trusted NIC Architecture》由 Dimitra Giantsidi、Julian Pritzi、Felix Gust、Antonios Katsarakis、Atsushi Koshiba 和 Pramod Bhatotia（University of Edinburgh、TUM、Huawei Research）完成，发表于 2025 年的 ASPLOS。本文将其置于“SmartNIC、可信网卡与安全存储数据路径”的研究脉络中考察，并把它作为 system 类材料使用。它所回答的核心问题可概括为：TNIC 把跨主机可信分布式系统所需的 root-of-trust 从 CPU TEE 下沉到 NIC hardware。当前证据状态为：本地 PDF 已保存并核验；下文仅在该范围内讨论其机制与结论。\cite{giantsidi2025tnic}\par
\subsubsection{研究背景与动机}
要理解该方案，首先需要看到它面对的系统矛盾：分布式系统需要可信的消息来源（message provenance），但 CPU enclave 路径会拖慢网络。\par
具体而言，BFT/CFT 协议需要 authentication 与 non-equivocation——节点不能对同一序号给出冲突消息，且收到消息的第三方要能验证来源。若把密码操作与网络栈放进 CPU TEE，异构云环境中要面对编程模型不统一、TCB 庞大、网络 I/O 产生显著延迟尖峰等问题。NIC 本来就在数据路径上，且天然与 host CPU 异构无关，这使它成为 host-agnostic trust substrate 的合理位置。\par
\subsubsection{关键挑战与核心思想}
在上述背景下，作者提出的关键判断是：把最小安全原语固化在 NIC 硬件中，应用通过 kernel-bypass API 获得可信网络操作。\par
具体而言，威胁模型继承 Byzantine fault model：云基础设施、机器与网络可表现为 Byzantine，攻击者可控制 host 软件与网络栈，但不能破坏 TNIC 硬件 root 与密码原语。在此模型下，attestation kernel 位于 TNIC 硬件，request/response handler 处理收发路径，Tamarin 模型验证 bootstrapping、remote attestation、消息发送/接收四部分协议。系统设计的目标性质只有两条：non-equivocation 与 transferable authentication，窄目标使形式化验证可行。\par
\subsubsection{系统架构}
从系统结构看，论文的基本组织方式是：TNIC hardware 介于 host application 与 100Gb MAC/RoCE network stack 之间（Figure 1--4、Table 1）。\par
具体而言，Figure 2 展示 Tx/Rx datapath 与 attestation kernel 的位置；Figure 3 展示 remote attestation 协议；Table 1 给出编程 API。方法管线为：TNIC provisioning/attestation → trusted NIC network stack → kernel-bypass RDMA-like API → 每消息 counter/MAC 原语 → 分布式协议变换。应用侧用类似 RDMA verbs 的 API 直接 post 请求，连接建立时分配内存区域，绕开 OS kernel。\par
\subsubsection{核心方法设计}
论文的关键不是单一组件，而是 NIC 级证明、可转移认证、kernel-bypass 栈与形式化边界四者的组合。\par
\noindent\textbf{（1）NIC 级 attestation。}\par
这一机制的直接作用是：TNIC 先证明 NIC 硬件与控制二进制，再让应用相信其网络原语。Remote attestation 协议用 measurement 与 mutual TLS 建立信任，verifier 检查 Ctrlbin/Ctrlpub 等信息。这与 SPDM 的设备端点证明目标一致，但论文实现的是自有协议而非 SPDM 标准——这意味着它与现有设备生态的互操作性需要另行桥接。\par
\noindent\textbf{（2）Transferable authentication。}\par
这一机制的直接作用是：收到的消息证明可以被第三方验证，而非只在点对点连接内有效。该性质正是 BFT/CFT 系统 non-equivocation 的来源：硬件生成/验证消息认证，应用不必把庞大的 CPU TEE 塞进每个网络跳。README 将其安全性质概括为：节点不能对同一 counter 产生冲突消息，第三方能验证消息来源并转移信任。\par
\noindent\textbf{（3）Kernel-bypass 栈。}\par
这一机制的直接作用是：高性能的关键在于让应用绕过 OS kernel 直接 post 请求给硬件。API 为 RDMA/ibv 风格，memory area 在连接创建时分配，网络栈追求低延迟操作；论文报告小包网络路径延迟在 5--5.5 微秒量级。\par
\noindent\textbf{（4）形式化验证的边界。}\par
这一机制的直接作用是：Tamarin 证明强化协议性质，但不证明硬件实现没有 bug。模型覆盖 bootstrapping、remote attestation、transmission、reception 四部分；实现仍面临 FPGA/HLS/HDL 的工程风险。论文把 proof scope 与 system evaluation 明确分开，本文沿用这一区分：形式化证据作用于协议层，系统证据作用于实现层，两者不能互相替代。\par
\subsubsection{安全性与威胁边界}
从安全边界看，TNIC 的安全目标是窄而可验证的：non-equivocation 与 transferable authentication。它不直接保护 application memory 或 host 侧 TEE 状态，不解决物理攻击与 DoS，也不替代 SPDM/TDISP/IOMMU：它回答的是“这条网络消息是否来自一块真的、未被替换的 TNIC 且发送方无法抵赖”，而不是“这台主机的内存是否机密”。其安全性依赖 TNIC 硬件实现与 provisioning 信任；若用于 CVM 网络路径，还需与 VM attestation、key broker、DMA protection、vNIC lifecycle 组合。\par
\subsubsection{实验环境与证据}
论文没有把这部分只写成结论，而是以如下材料界定其证据范围：FPGA 原型、多组 host 基线、分布式系统案例研究与 Tamarin 证明（evaluation Section 8、Figure 5--13、Table 2--5）。\par
具体而言，平台为 Xilinx U280 FPGA 与 Vitis XRT v2022.2；基线覆盖 RDMA-hw、DRCT-IO、DPDK、SSL-lib 与 SGX/AMD SEV 变体；负载包括 attest/verify microbenchmark、send latency/throughput、A2M、BFT/CFT 系统与 chain replication 等。硬件 TCB 约 2114 行 HLS/HDL，对比 TEE-hosted 代码库超过 200 万行的量级；TNIC 设计占 U280 FPGA 约 16.6\% LUT 与 16.3\% FF。artifact 在 GitHub/Zenodo 公开。\par
\subsubsection{实验结果与证据强度}
对于性能或效果，论文能够支持的结论是：TNIC 明显优于 TEE-based 的网络信任方案，但相对纯 SSL/RDMA 仍有安全成本。\par
具体而言（Figure 8--12、Table 3--4），论文报告相比 AMD SEV/SGX 等 TEE-based 竞争者吞吐提升 3--5 倍，相比 CPU-centric TEE 系统最多 6 倍性能提升。A2M 案例（Table 3）的数字最为直观，本文整理为表~\ref{tab:tnic-a2m}：TNIC 吞吐约为 AMD SEV 方案的 5.3 倍，延迟约为其五分之一。TCB 对比（Table 4）则解释了收益来源：2114 行硬件 TCB 对 200 万行级 TEE 代码库。证据边界：全部数字来自 FPGA 原型与所列基线配置，本文不外推到 ASIC 形态或其他协议负载。\par
\begin{table}[htbp]
\centering
\caption{A2M 系统案例：TNIC 与 AMD SEV 基线对比（整理自论文 Table 3）}
\label{tab:tnic-a2m}
\footnotesize
\begin{tabular}{@{}lcc@{}}
\toprule
平台 & 吞吐（ops/s） & 延迟（$\mu$s） \\
\midrule
TNIC & 158K & 6.34 \\
AMD SEV & 30K & 32.37 \\
\bottomrule
\end{tabular}
\end{table}
\subsubsection{失败模式与边界分析}
论文的证据结构自身标出了三条边界。第一，形式化覆盖不全：Tamarin 模型只覆盖协议四段，RTL 级缺陷、HLS 综合问题与 provisioning 过程均不在证明内。第二，目标窄：TNIC 不提供机密性，消息内容本身仍需上层加密；Byzantine 对手仍可实施 DoS。第三，生态缺口：自有 attestation 协议与 SPDM/TDISP 标准栈之间的映射未建立，README 边界明确其“不替代 SPDM/TDISP/IOMMU/link encryption”。\par
\subsubsection{综合评价}
综合来看，TNIC 是 trusted network endpoint 的强 SOTA，最适合支撑“NIC 也可以是 RoT”这一论点。\par
具体而言，它的优势在于把 root-of-trust 下沉到 NIC 后用极小的硬件 TCB（2114 LoC）换来可验证的两个分布式安全性质，形式化验证与系统评估都达到较高强度。它的局限在于不保护 workload memory、不覆盖设备生命周期标准栈。从部署角度看，它适用于 SmartNIC/DPU attested messaging、BFT/CFT 服务与云网络信任基座；与 S-NIC 的对比很能说明问题——S-NIC 解决“NIC 内部谁碰得到谁”，TNIC 解决“网络对端凭什么相信这块 NIC”，两者是互补层而非替代关系。\par
\subsection{论文 34：Hazel: Secure and Efficient Disaggregated Storage}
\subsubsection{论文基本信息}
《Hazel: Secure and Efficient Disaggregated Storage》由 Marcin Chrapek、Meni Orenbach、Ahmad Atamli、Marcin Copik、Mikhail Khalilov、Fritz Alder 和 Torsten Hoefler（ETH Zurich、NVIDIA、University of Southampton）完成，2026 年以 arXiv 预印本（2510.18756，v2）形式发布。本文将其置于“SmartNIC、可信网卡与安全存储数据路径”的研究脉络中考察，并把它作为 system 类材料使用；其证据等级为 draft/not ratified，引用时必须标注 preprint 状态，不能当作已经同行评审或标准化的 trusted I/O。它所回答的核心问题可概括为：Hazel 把 secure storage data path 的三个目标——confidentiality、integrity、freshness——放进同一个 NVMe-oF + DPU 原型。当前证据状态为：本地 PDF 已保存并核验；下文仅在该范围内讨论其机制与结论。\cite{chrapek2026hazel}\par
\subsubsection{研究背景与动机}
要理解该方案，首先需要看到它面对的系统矛盾：远端存储不只要加密，还要防篡改和旧数据回放，而本地方案直接搬到 NVMe-oF 会在性能与扩展性上同时失效。\par
具体而言（README 证据基线为 Section 2、Figure 2），dm-crypt 只提供 confidentiality；dm-integrity/dm-x 虽补上完整性，却带来 IOPS、latency、throughput 的显著开销；在 NVMe-oF/JBOF 场景中，host CPU 与网络路径都可能成为瓶颈。威胁模型包括特权管理员、其他用户、参与节点上的本地对手与远端对手，以及能观察或修改控制面/数据面请求响应的攻击者；论文假设存在 TEE/KBS、trusted Hazel service 与少量受信非易失内存。这个设定把问题精确限定为：在存储端点与网络都可能不可信的前提下，如何给分解式存储同时提供三种安全性质且不牺牲线速。\par
\subsubsection{关键挑战与核心思想}
在上述背景下，作者提出的关键判断是：让 freshness/control path 慢而少，让 data path 快而 offloaded。\par
具体而言，KBS/TEE 负责 counter leasing 与密钥/控制，DPU 负责线速密码运算与 NVMe-oF 路径，Hazel Merkle Tree（HMT）配合 metadata cache 让读路径尽量不走昂贵的树路径。方法管线为：confidential workload → 本地 Hazel service/KBS 租 counter → NVMe-oF 请求封装（metadata 携带 IV/hash/freshness 信息）→ 远端 Hazel 验证/更新 HMT → SSD 存储数据与元数据。核心模块包括 counter leasing、key derivation、NVMe metadata encapsulation、network freshness header、HMT、metadata cache、eventual consistency 写路径与 BlueField-3 GGA/DOCA offload。\par
\subsubsection{系统架构}
从系统结构看，论文的基本组织方式是：Hazel control path 建立信任/密钥/counter，data path 走 DPU/NVMe-oF 并把安全元数据放进 sector layout。\par
具体而言，Figure 3 是 control path：cluster manager、TEE、KBS、Hazel service；Figure 4 是 data path 与 SSD sector layout；Figure 5 是 HMT 设计。实现基于 SPDK 与 NVIDIA BlueField-3 DPU，使用 DOCA SDK 访问 BlueField-3 的密码 Generic Global Accelerators（GGA）；测试环境包含用户端 BF3 DPU、存储端 BF3 NIC、3.84TB PCIe SSD 与 NVMe-oF 路径。控制面与数据面的这一分离，是它能同时声称安全与性能的结构基础。\par
\subsubsection{核心方法设计}
论文的关键不是单一组件，而是四个机制围绕“元数据随数据走、密码运算下沉 DPU”的组合。\par
\noindent\textbf{（1）Counter leasing。}\par
这一机制的直接作用是：Hazel 不为每次 I/O 都访问 KBS，而是租一段 counter range，把控制路径频率降下来（README 证据基线为 Section 3.2、Figure 3）。KBS 管理密钥与计数器，一段 lease 可覆盖大量写操作；租约管理错误会直接损害 freshness，因此必须由 TEE/KBS 策略管控。这是“控制路径慢而少”的具体化。\par
\noindent\textbf{（2）NVMe metadata encapsulation。}\par
这一机制的直接作用是：利用 NVMe 元数据字段携带安全信息，不重写整个 NVMe-oF 协议。sector layout 中放置 counter/IV、MAC/tag 与安全元数据，DPU 在数据路径上完成加解密与校验。这降低了协议改造成本，也保留了高性能存储路径，是 Hazel 与既有 NVMe-oF 生态兼容的关键。\par
\noindent\textbf{（3）Hazel Merkle Tree（HMT）。}\par
这一机制的直接作用是：用 in-memory IV cache、batch update 与 eventual consistency 降低 freshness 开销（Figure 5--7）。设计的两难是：整棵树放磁盘会造成不确定延迟，全树放内存又浪费 DRAM；HMT 选择缓存/批处理/最终一致的折中，在性能与 freshness 时效之间取舍。这一取舍也直接带来了后文要分析的失败模式。\par
\noindent\textbf{（4）DPU offload 的信任边界。}\par
这一机制的直接作用是：Hazel 借助 BlueField-3 降低 CPU overhead，但 DPU 自身的信任需要额外证明。论文把 TDISP/IDE 作为相关工作/未来信任基础讨论；当前 preprint 没有把 DPU attestation 与存储端点身份完整闭环。商业部署必须补上 SPDM/TDISP/NVMe 端点证据，README 边界亦明确标注此缺口。\par
\subsubsection{安全性与威胁边界}
从安全边界看，Hazel 的目标是防明文泄漏、防篡改、防 stale data replay：confidentiality 由唯一 IV 与 counter leasing 支撑，integrity 由 AEAD/hash 与元数据支撑，freshness 由 HMT 与 freshness cache 支撑。正确性风险集中在 metadata cache 污染、并发写、crash recovery 与 eventual consistency；论文讨论了 crash/coherency 问题，但完整形式化证明不足。信任假设包括 KBS/TEE/Hazel service 的正确配置与远端 Hazel 路径可信；DPU 作为数据路径 TCB 的一部分，其 attestation 与 device lifecycle 未在本文闭环。\par
\subsubsection{实验环境与证据}
论文没有把这部分只写成结论，而是以如下材料界定其证据范围：NVMe-oF + BlueField-3 真实原型，负载覆盖 synthetic、IO500、YCSB 与 ML training（implementation/evaluation Section 4--5、Figure 11--13）。\par
具体而言，硬件为高性能 PCIe SSD 与 NVIDIA BlueField-3 DPU；benchmark 包括 synthetic 读写模式、IO500、YCSB/RocksDB、ResNet50/UNet3D ML pipeline；指标覆盖 throughput、IOPS、latency、CPU usage 与 p99 latency。真实 BF3/SSD/NVMe-oF 原型是其证据强度的主要来源；arXiv preprint 状态是其证据等级的上限。\par
\subsubsection{实验结果与证据强度}
对于性能或效果，论文能够支持的结论是：Hazel 在常见大读写与应用路径上开销低，但小随机 freshness 路径仍是弱点。\par
具体而言（Figure 11--13 与评估文本），Figure 11：IO500 各场景开销大多低于 10\%，平均 6.3\%。Figure 12：YCSB 平均 p99 latency 开销 2.2\%，吞吐开销平均 0.6\%。Figure 13 与文本：ML training 的 freshness 开销约 1\%--2\%；大块/顺序路径总体在 1\%--2\% 量级。相对地，小随机 freshness I/O 的开销可显著变差，随机写与 dirty reads 对 metadata-to-data ratio 敏感。这组数字必须与其 preprint 状态一起引用；本文不外推其在其他 DPU 或协议栈下的表现。\par
\subsubsection{失败模式与边界分析}
Hazel 的失败模式在 README 评审中记录得相当完整，本文按机制归类。第一，小随机 I/O 的 freshness 路径：freshness read path 对 dirty metadata cache 敏感，小随机 I/O 是其明确弱点——这恰是 OLTP 类负载的典型模式。第二，eventual consistency 的时效窗口：HMT 的批处理/最终一致设计在性能与 freshness 时效之间留了窗口，crash recovery 的正确性论文有讨论但无形式化证明。第三，信任闭环缺口：DPU attestation、SPDM/TDISP 端点身份、生产级 KBS 运维均未闭环；README 明确其边界为“preprint；未完整闭环 BlueField DPU attestation、SPDM/TDISP/NVMe-oF endpoint identity”。第四，证据等级：arXiv v2 尚未经同行评审，全部性能数字应视为草案证据。\par
\subsubsection{综合评价}
综合来看，Hazel 是 secure disaggregated storage 方向的强草案 SOTA，但必须标注 preprint 与端点信任缺口。\par
具体而言，它的优势在于把 confidentiality、integrity、freshness 与 DPU offload 集成进一个不修改 NVMe-oF 协议的存储原型，并用真实硬件给出了分场景的开销数字；counter leasing 与 HMT 的设计直接回答了“freshness 为什么贵”。它的局限在于：preprint 状态、小随机 I/O 弱点、crash consistency 未形式化、端点信任未闭环。从部署角度看，它面向云机密存储、NVMe-oF/JBOF 与 DPU offload 场景；落地的主要风险在小随机 I/O 负载与证据闭环的补齐。\par
\subsection{开发商安全实践建议}
面向 NIC/DPU 厂商、交换机与 fabric 厂商、存储设备厂商以及平台固件/系统软件开发商，本节从上述八篇材料中提炼六条可操作建议。每条建议都标明其证据来源与适用边界，避免把单篇论文的结论误读为通用工程准则。\par
\noindent\textbf{实践一：把 SPDM 作为设备身份与固件测量的统一控制面底座，并把“降级拒绝”纳入出厂测试。}\par
NIC、DPU、交换机与存储端点应在出厂时内置证书链与 measurement 能力，使平台能以标准方式完成“你是谁、固件什么状态、能否建会话”三步核验（\cite{dmtf_spdm_2025}）。工程要点有三：其一，measurement block 要覆盖固件版本与安全配置，而非只返回身份；其二，版本协商与算法协商的降级路径必须被测试并默认拒绝，README 给出的验收标准（transcript 可复验、measurement 与设备状态绑定、算法降级被拒绝）可直接转化为一致性测试项；其三，要向上层明确声明 SPDM 只覆盖控制面引导——产品文档若把“支持 SPDM”写成“支持机密 I/O”，就是把证据协议误读为完整防御。\par
\noindent\textbf{实践二：链路加密默认开启，并把密码计算移出读关键路径。}\par
CXL 内存模块与 PCIe 设备厂商应默认启用 IDE 类链路保护，但 AIORE 的测量（\cite{li2025aiore}）表明，把 XTS 解密放在每次内存读的关键路径上会在内存密集型负载产生最高 14.9\% 的开销。可操作的结论是：内存扩展模块应内置可复用的加密引擎（MEE），支持把重加密等批量密码工作卸载到设备侧空闲算力；主机侧设计应为 CTR 类可并行 OTP 的方案预留计数器缓存与元数据通道。对交换机/fabric 厂商，这意味着 IDE 的 MAC epoch 与 Flit 处理必须在 switch 路径上保持线速，否则链路安全会成为 fabric 的吞吐瓶颈，论文未覆盖 switch 侧实现代价，厂商需自行实测。\par
\noindent\textbf{实践三：SmartNIC/DPU 内部资源按 single-owner 语义做租户隔离，并把硬件成本预算写进产品规划。}\par
DPU 厂商在设计多租户卸载（vSwitch、DPI、压缩、存储 offload）时，应把 S-NIC 的五类资源对象——内存 denylist、locked TLB、cache 分区、加速器虚拟化、总线仲裁——作为隔离清单逐项覆盖（\cite{zhou2024snic}）。其评估给出可引用的成本量级：4 核配置下面积约 +8.89\%、功耗约 +11.45\%、最坏吞吐损失小于 1.7\%，这组数字可作为“隔离不是免费的、但可负担”的预算基线；同时应如实告知客户强隔离的资源利用率代价与 function chaining 限制。管理面密码操作（如 nf\_attest 的 RSA）应保持在 packet fast path 之外。\par
\noindent\textbf{实践四：为需要跨主机信任的网卡提供 NIC 级 root-of-trust，并把协议核心纳入形式化验证。}\par
面向 BFT/CFT、机密数据库复制等场景的 NIC/DPU 厂商，可参考 TNIC 的论证（\cite{giantsidi2025tnic}）：把 attestation kernel 与消息认证原语固化在网卡硬件，能把可信网络路径的 TCB 压到约 2000 行 HLS/HDL 量级，相对 TEE 托管方案获得 3--5 倍吞吐。可操作的两点是：优先选择窄而可验证的安全性质（non-equivocation、transferable authentication）而非大而全的“可信网卡”；同时把自有 attestation 协议与 SPDM/TDISP 标准栈的映射作为产品化必修课，因为 TNIC 的自有协议证明不了与既有设备生态的互操作。\par
\noindent\textbf{实践五：存储设备与 DPU 厂商应把 confidentiality/integrity/freshness 做进 NVMe-oF 数据路径，而不是留给主机软件层。}\par
Hazel 的原型（\cite{chrapek2026hazel}）证明了三性合一可以在不改 NVMe-oF 协议的前提下实现：用 NVMe metadata 携带 IV/MAC/freshness 信息，用 DPU 的密码加速器承担线速运算，常见应用路径开销可压到 1\%--2\%（IO500 平均 6.3\%）。可操作的边界同样来自该文：小随机 I/O 的 freshness 路径是已知弱点，产品规格书应给出该模式下的实测开销而非只报平均值；counter lease 与 Merkle 树根的崩溃一致性需要写入固件需求。引用该文时必须注意其 arXiv preprint 状态，数字应视为草案证据并在自有平台上复测。\par
\noindent\textbf{实践六：把遥测与资源调度面当作信任边界来管理，而不是纯性能设施。}\par
CXL 内存池化与分层产品（参考 \cite{gouk2022directcxl,zhong2024cxltiers,wang2025odrp}）的共同教训是：fabric 管理器、host 侧调度器与 MNode 软件在机密场景下都是不可信控制面。可操作建议包括：CXL switch 与 pooling 产品的 namespace/segment table 管理接口应支持设备身份绑定与租户授权（DirectCXL 的机制图给出了挂接点，但论文未实现）；tiering 产品的 per-VM 性能遥测应评估负载指纹泄露风险，页迁移策略对 CVM 尾部延迟的影响应纳入威胁模型（CXL-Tiers 未分析这一点）；远端内存分配元数据（如 ODRP 的 TT 与 free queue）在 MNode 不可信时需要完整性保护，否则恶意 MNode 可在不违反任何一致性不变式的情况下窃听或篡改。\par
\subsection{未来研究方向}
以下六个方向直接对应本章评审中记录的证据缺口，每个方向都给出问题、现有工作未解决的原因、可行技术路线与所需验证实验。\par
\noindent\textbf{（1）CXL 内存的低开销完整性与 freshness。}\par
问题：AIORE 把内存侧完整性整体托付给 CXL IDE，而 CTR TEE 加完整性树的路线被证明放弃链路唯一性（\cite{li2025aiore}）；CXL 内存因此缺少一个既保留 IDE 链路性质、又覆盖内存侧篡改/重放的低开销方案。未解决的原因是：完整性树元数据在 CXL 延迟结构下成本过高，现代 VM 型 TEE 又已放弃树结构。技术路线：按页/按区域自适应地在“IDE 保证”与“树保证”之间切换，复用 AIORE 的热度跟踪与卸载机制摊平树维护。所需实验：Gem5 或 RTL 级原型，报告完整性覆盖下的读关键路径开销、元数据带宽与 counter/树缓存命中率，并与 AIORE 的 3.7\% 基线对比。\par
\noindent\textbf{（2）机密 CXL 内存池的多租户所有权与撤销。}\par
问题：DirectCXL、CXL-Tiers、ODRP 分别解决了池化的性能、分层与分配效率，但三者都不处理 device identity、ownership lifecycle 与撤销（\cite{gouk2022directcxl,zhong2024cxltiers,wang2025odrp}）。未解决的原因是：这些工作威胁模型中 fabric 管理器与 host 可信，而 CVM 模型恰恰不信任它们。技术路线：把 namespace/segment table 的每次变更绑定到 SPDM 证据与 TDISP 接口状态，设计池化内存的租户授权与快速撤销协议。所需实验：端到端 assignment 攻击（恶意 fabric 管理器重定向映射）、撤销时延测量、以及撤销后残余数据的清空验证。\par
\noindent\textbf{（3）可编程 RDMA 控制路径的形式化与机密化。}\par
问题：ODRP 的 WR chain 一致性靠原子操作加归纳论证支撑，未达模型检查级别，且 MNode 被默认诚实（\cite{wang2025odrp}）。未解决的原因是：WR chain 是图灵不完备但状态丰富的硬件程序，现有网络验证工具不覆盖；机密化则要求分配元数据离开不可信内存。技术路线：对 WR chain 的不变式（一页一主、TT 更新原子、free queue 一致）做模型检查或定理证明；再将 TT/free queue 移入 RNIC 受信存储或加 MAC。所需实验：Byzantine MNode 故障注入下的别名/泄漏检测率，以及元数据受信化后的 4KB 分配吞吐与延迟。\par
\noindent\textbf{（4）SPDM measurement 的语义、撤销与 fleet 级验证策略。}\par
问题：SPDM 标准化了测量的传输，却刻意不规定测量的业务语义与证书撤销策略（\cite{dmtf_spdm_2025}），导致“measurement 成功”在不同厂商间不可比较，fleet 级 verifier 策略无从落地。未解决的原因是：语义标准化需要跨厂商的固件清单（reference manifest）生态，超出协议工作组的职权。技术路线：为 NIC/DPU/存储端点定义 measurement schema 与撤销分发机制，对 responder 实现做 fuzzing 与 differential testing。所需实验：跨厂商互操作测试（降级协商、错误路径、responder reset）、撤销传播时延、以及 schema 对固件供应链事件的覆盖率。\par
\noindent\textbf{（5）NIC 信任根与设备生命周期标准的联合证明。}\par
问题：TNIC 的自有 attestation 与 Tamarin 证明覆盖协议四段，但不覆盖 device assignment 生命周期，也未与 SPDM/TDISP 会话绑定（\cite{giantsidi2025tnic}）；S-NIC 的 function 证明同样止步于“在真 S-NIC 上启动”（\cite{zhou2024snic}）。未解决的原因是：学术原型各自为政，标准栈的状态机复杂且仍在演进。技术路线：把 NIC RoT 的 quote 作为 SPDM measurement 的一种，将 function launch/destroy 纳入 TDISP 接口状态机，并把联合模型（RoT + assignment + 会话）纳入 Tamarin 类工具。所需实验：联合协议的安全性质证明、assignment/撤销全流程的延迟分解、以及与 CVM 端 quote 的联合 verifier 策略评估。\par
\noindent\textbf{（6）机密分解存储的小随机 I/O 与崩溃一致性。}\par
问题：Hazel 的小随机 freshness 路径开销显著、eventual consistency 与 HMT 根更新顺序未形式化，且 preprint 证据等级有限（\cite{chrapek2026hazel}）。未解决的原因是：freshness 本质要求全局有序，与小随机 I/O 的并发性冲突；DPU 侧的信任闭环又依赖尚未普及的 TDISP/IDE 生态。技术路线：面向 dirty metadata cache 场景设计 DPU 侧批量 freshness 验证器；用模型检查刻画 HMT 批更新与崩溃恢复的顺序约束；把 DPU attestation 与 Hazel service 证据接入 SPDM/TDISP。所需实验：OLTP 型小随机负载下的 freshness miss 开销分布、注入崩溃后的重放检测率、以及对照 Hazel 报告值（IO500 平均 6.3\%、YCSB p99 平均 2.2\%）的独立复测。\par
\section{加速器、SmartNIC、DPU 与设备 TEE}
GPU、FPGA、NPU、SmartNIC、DPU 等加速器与领域专用设备（DSA）把计算和数据管理带到 CPU 之外，也把 TCB、队列、DMA、I/O 链路和资源复用的边界带进更复杂的设备运行时中。CPU 侧 TEE（无论是 enclave 还是 confidential VM）只能保护 CPU 执行上下文与主机内存的一部分；一旦工作负载被卸载，敏感数据就会经过驱动、命令队列、设备内存、PCIe/CXL 链路和设备固件，这些环节中的任何一个缺少身份、访问控制或加密保护，端到端的机密性承诺就会断裂。本节讨论异构机密卸载的共同难题：谁验证设备身份与固件状态、谁配置设备的访问窗口与 DMA 映射、谁处理并发租户之间的资源复用与清理，以及中断和共享缓冲区如何回到同一条可信链。\par
本节选取三类互补材料：一篇覆盖 51 个设计的 accelerator TEE 系统化分析（SoK）\cite{sok-tee}，一个用机架级安全控制器保护 COTS 加速器池的奠基性系统 HETEE\cite{zhu2020hetee}，以及把控制器思路推到数据中心规模、兼容 legacy 非 TEE 设备的 CloudScale\cite{dhar2024cloudscale}。各论文的价值将据此被拆解为可核验的设计选择，而非泛化为``加速器已经可信''；尤其是端点证明、密钥绑定和设备固件假设，将与主机侧 enclave 或 confidential VM 的假设一起审视。凡是论文未给出数字或机制细节的地方，本节明确标注``论文未说明''或``证据不足''，不做延伸推断。\par

\subsection{论文 35：SoK: Analysis of Accelerator TEE Designs}
\subsubsection{论文基本信息}
《SoK: Analysis of Accelerator TEE Designs》由 Chenxu Wang、Junjie Huang、Yujun Liang、Xuanyao Peng、Yuqun Zhang、Fengwei Zhang、Jiannong Cao、Hang Lu、Rui Hou、Shoumeng Yan、Tao Wei、Zhengyu He 完成（机构包括 SUSTech、Hong Kong Polytechnic University、中国科学院、蚂蚁集团等），发表于 2026 年的 Network and Distributed System Security Symposium (NDSS 2026)。本文将其置于``Arm CCA 的 I/O、DMA、加速器与中断、硬件辅助可信执行环境分类''的研究脉络中考察，并把它作为 SoK 类材料使用。它所回答的核心问题可概括为：当机密工作负载越来越依赖 GPU/NPU/TPU/FPGA/DPU 等加速器时，已有的 accelerator TEE 研究能否被放进一个统一的设计框架中比较，其安全机制、TCB 与部署障碍各自呈现什么规律。当前证据状态为：本地 PDF 已保存并核验（NDSS 2026 论文，reference 目录下同目录 paper.pdf）；下文仅在该范围内讨论其机制与结论。\cite{sok-tee}\par
\subsubsection{研究背景与动机}
要理解该方案，首先需要看到它面对的系统矛盾：accelerator TEE 论文数量快速增长，但绝大多数设计只针对特定 CPU 或特定加速器，缺乏可迁移的统一框架。论文摘要（p.1）指出，已有 TEE 综述只部分总结了加速器计算的威胁与防护，既没有给出``如何构建一个 accelerator TEE''的指导框架，也没有比较不同安全方案的优缺点。\par
具体而言，论文在 p.2--p.3 的 Motivation 部分用一组 corpus 统计量化了这一碎片化的严重程度。作者收集并分析了 51 个 accelerator TEE 研究（Table I），覆盖 GPU、NPU、TPU、FPGA、通用加速器与工业设计。其中 41/51 在 2022 年之后提出，说明该方向正在快速升温；42/51 依赖 CPU 侧 TEE 作为信任基础，说明设备侧安全通常不是独立建立的；32/51 面向 x86 平台、30/51 支持 GPU 计算，反映云端异构计算的主流形态；40/51 只支持特定 CPU 架构、43/51 只支持一种加速器类型，且仅 12/51 释放了源代码，意味着跨平台迁移与独立复现都相当困难。这组数字构成本节的第一个可核验事实：设备侧可信边界问题不是个别方案的疏漏，而是整个领域的结构性缺口——AI 与高性能负载把数据推入 device memory、I/O bus 与命令队列后，若设备缺少身份、attestation 或内存保护，即使主机侧 CVM 完好，机密数据路径仍然不完整。\par
\subsubsection{关键挑战与核心思想}
在上述背景下，作者把分析组织为三个 research questions（论文 p.2），它们同时构成本文读者可以使用的问题清单：\par
\begin{itemize}
\item \textbf{RQ1：}构建 accelerator TEE 的典型框架是什么？论文的回答是把现有设计归为 Host-type、Acc.-type、Mix-type 三类架构（详见``系统架构''）。
\item \textbf{RQ2：}如何在不同 CPU/加速器组合上构建具有强安全保证的 accelerator TEE？论文的回答是把安全机制系统化为 access control、memory encryption、attestation 三条主线，并逐条比较现有方案的覆盖与盲区；核心判断是三者必须同时回答，只解决 DMA 隔离远远不够。
\item \textbf{RQ3：}哪些因素影响 accelerator TEE 在真实平台上的部署？论文的回答是 TCB 规模与兼容性（compatibility），而不是单纯的安全机制设计。
\end{itemize}
这三条主线的选择有明确的证据支撑：在本文仓库的语境中，ACAI 偏向 memory/data path 保护，Devlore 偏向 interrupt/control path 保护，而该 SoK 提醒读者还必须同时考察 device HRoT、endorser、reference values、driver TCB 与多设备兼容性；因此 Arm CCA 的 I/O 讨论必须与 TDISP/SPDM/PCIe IDE 等标准生态关联，而不能停留在单一 GPU/FPGA 方案上。\par
\subsubsection{系统架构}
从系统结构看，论文（Figure 1/p.4，Table II/p.4--p.5）把 accelerator TEE 的保护逻辑落点分为三类：\par
\begin{itemize}
\item \textbf{Host-type}：保护逻辑主要驻留在 CPU 侧 TEE、hypervisor 或 firmware 中，复用已有 CPU TEE 能力。优点是部署门槛低、可通过软件补丁升级；代价是容易把庞大的 accelerator driver/runtime 一并塞进 CVM/enclave，导致 TCB 膨胀。SoK 据此建议：使用商用加速器的云厂商（如阿里云类场景）宜采用 Host-type，在 CVM 内保护负载，并通过修改 TSM/firmware 实现设备环境保护。
\item \textbf{Acc.-type}：保护逻辑迁移到 accelerator controller、设备侧加密引擎与设备侧 attestation 模块中。适合有硬件设计能力的加速器厂商（如 NVIDIA、Xilinx 类场景），但需要芯片/固件配合。SoK 建议加速器制造商走此路线，把控制器、定制加密 IP 核与证明模块集成进设备。
\item \textbf{Mix-type}：CPU 侧与设备侧各自持有 root of trust，通过受保护的控制通道与数据通道协同，最接近 TDISP/IDE 这一产业化方向。论文指出目前 Mix-type 设计很少（4/51），其代表是 Arm RME-DA、AMD SEV-TIO、Intel TDX Connect 与 RISC-V sNPU：设备侧由 TDISP 定义的 Device Security Manager (DSM) 管理可分配给 CVM 的 TEE Device Interface (TDI)，DSM 与 TSM 之间的控制路径由 SPDM 保护，数据路径由 PCIe IDE 保护。
\end{itemize}
\subsubsection{核心方法设计}
作为 SoK，论文的关键贡献不是单一组件，而是一套可复用的分析管线：paper collection 得到 51-study corpus，再按 architecture type、attack vector、security mechanism taxonomy 与 deployment factors 逐层分解。以下按机制逐条展开。\par
\noindent\textbf{（1）攻击向量分类（Figure 2/p.4--p.6）。}\par
这一机制的直接作用是：把``保护加速器''这一模糊目标拆成可枚举的失败点。论文将攻击面按两个阶段组织：task preparation/termination 阶段（任务还在主机侧内存，尚未与加速器交互）与 task computing 阶段（任务已进入设备侧内存并通过 MMIO 与驱动交互）。攻击对象分三类：针对 accelerator task 的 AT 类（task code、输入/输出数据、模型参数、page table 与 metadata，如 buffer pointer），针对 accelerator environment 的 AE 类（加速器软件栈、MMIO 寄存器、硬件设备），以及针对 authenticity 的 AA 类（伪造或替换设备）。全文共枚举 AT1--AT8、AE1--AE6、AA1--AA2，攻击来源覆盖被攻陷的 host OS/hypervisor、恶意加速器软件、恶意并发任务、物理攻击者与 I/O bus 窃听/篡改。值得注意的细节是：论文明确指出统一内存架构（如 Arm GPU 与主机共享内存）比独立显存的加速器（如 NVIDIA dGPU）更容易受到 AT4/AT5 类利用主机权限篡改设备侧任务的攻击。\par
\noindent\textbf{（2）Access Control 方案谱系（Table III/p.6，Table IV/p.7，Table V/p.8）。}\par
这一机制的直接作用是：回答``谁能访问 workload、MMIO、DMA 与设备软件栈''。论文把现有方案归纳为六种 SAC：SAC1 基于 CPU TEE 但不保护加速器软件栈；SAC2 基于 CPU TEE 并把加速器软件栈纳入保护；SAC3 基于 TEE manager（如 RMM）或安全 hypervisor；SAC4 基于 firmware（如 monitor）；SAC5 基于 I/O bus 过滤（CPU bus filter、PCIe filter）；SAC6 基于加速器硬件（如 security controller）与经认证的 kernel。Table III 以矩阵形式标注每种方案对 AT/AE/AA 各攻击是``完全防御''、``部分场景防御''还是``无法防御''。Table IV 进一步给出部署偏好统计：在云端 CPU 与加速器分离的场景（34 个相关研究中），SAC6 出现 26/34、SAC1 出现 14/34、SAC2 出现 12/34、SAC5 出现 6/34、SAC3 出现 4/34、SAC4 仅 1/34；而在 CPU-加速器集成的端侧场景（17 个相关研究中），偏好转为 SAC2 与 SAC4 各 11/17、SAC6 6/17、SAC1 5/17。由于一项研究可同时采用多种方案，这些计数是方案出现频次而非互斥占比。结论是：云端离散加速器更依赖 CPU 侧与 bus/device 硬件检查的组合，端侧集成加速器更常依赖 firmware/monitor 与平台硬件。\par
\noindent\textbf{（3）Memory Encryption 谱系（Figure 3/p.9，Table VI/Table VII/p.9）。}\par
这一机制的直接作用是：设备侧机密性必须同时覆盖 host memory、device memory、I/O link 与 metadata，缺任何一环都留下明文窗口。论文把加密落点分为 SME1（基于 CPU TEE 的加解密）、SME2（基于加速器加密引擎/内核）、SME3（基于 I/O bus，如 TDISP/IDE 链路加密）。Figure 3 给出 memory encryption workflow，Table VI 讨论缺失加密环节的含义，Table VII 比较 metadata（完整性树、freshness 计数器等）的粒度选择。论文特别强调 CPU 与加速器两侧加密粒度不匹配（granularity mismatch）会引入显著开销，这一点在 Figure 4 的量化分析中展开（见``实验结果与证据强度''）。\par
\noindent\textbf{（4）Attestation 谱系（Figure 5/p.10--p.11，Figure 6/p.11，Table VIII/IX/p.10--p.11）。}\par
这一机制的直接作用是：证明远端用户连接的是真实设备、运行的是预期固件与软件栈。论文区分 SAT1（基于带 HRoT 的 CPU TEE）、SAT2（基于软件的证明）、SAT3（基于带 HRoT 的加速器）、SAT4（基于带 HRoT 的外接安全硬件）。Figure 5 给出通用 attestation 流程，Figure 6 与 Table VIII/IX 逐系统标注 accelerator HRoT、endorser 与 reference value 的完备性。论文的关键发现是：许多加速器缺少完整的、由厂商支持的 attestation 链，缺 HRoT/endorser/reference values 中任何一环都会导致仿真设备（emulated device）或错误固件被信任——这正是 AA 类攻击能够得逞的根因。\par
\noindent\textbf{（5）TCB 与兼容性维度（Table X--XIII/p.12--p.14）。}\par
这一机制的直接作用是：把``部署障碍''从定性抱怨变成可比维度。Table X 汇总多个系统的 TCB 增量，Table XII 汇总加速器软件栈规模，Table XIII 比较 multi-type 支持与 plug-and-play 能力。论文指出，许多系统只支持单一加速器或特定 CPU；NVIDIA 内核驱动、ROCm、CUDA/OpenCL 等栈可能是闭源且达百万行量级，把这样的栈放进 TEE 会根本削弱``driver in TCB''的安全性论断。商用落地因此需要减少高权限代码、支持即插即用与标准化的设备生命周期管理。\par
\subsubsection{安全性与威胁边界}
从安全边界看，数据一旦离开 CPU，设备身份、DMA 映射、队列、链路和中断投递都会成为新的受控点。该 SoK 的威胁模型（Figure 2）显式覆盖恶意 host OS/hypervisor、恶意加速器软件、恶意并发任务、物理攻击与 I/O bus 篡改，比只考虑恶意 OS 的模型更完整；但它同时明确了自身边界：SoK 的比较矩阵不等于任何一个具体方案的安全证明。access control 只回答``谁能访问''，不自动提供 memory confidentiality/freshness；memory encryption 不证明设备身份；attestation 不阻止运行期 DMA 越界。三类机制必须与工作负载身份和生命周期绑定一致，才能形成端到端的机密数据路径。此外，论文 p.15 明确承认 TDISP 是重要的产业框架但并非终局答案，仍存在 memory-encryption 开销、CPU-加速器信任链、大 TCB 与兼容性盲点。\par
\subsubsection{实验环境与证据}
论文没有把这部分只写成结论，而是以如下材料界定其证据范围：这是一篇 SoK，不包含新的 end-to-end accelerator TEE 系统实现，其证据由系统化统计、比较表格与一个针对 memory encryption granularity mismatch 的量化示例（Figure 4/p.10，LLM inference 叠加 SME 方案）构成。\par
具体而言，可核验的证据对象为：Table I/p.2 的 51-study corpus；Figure 1 与 Table II 的架构分类；Figure 2 的攻击向量；Table III/IV/V 的 access control 比较；Figure 3 与 Table VI/VII 的 memory encryption 分析；Figure 4 的粒度不匹配开销示例；Figure 5/6 与 Table VIII/IX 的 attestation 完备性；Table X--XIII 的 TCB 与兼容性统计。同时需要看到其适用边界：具体性能、协议状态机与安全证明必须回到原始论文或规范（TDISP、PCIe IDE、SPDM、CoVE-IO、Arm CCA RME-DA 等），不能以 SoK 表格替代。\par
本节判断以 Wang/Huang 2026 的 Table I、Figure 1--6、Table III--XIII（p.2--p.15）为依据。超出这些材料覆盖范围的实现细节、性能结论或攻击面，本文不作延伸推断。\par
\subsubsection{实验结果与证据强度}
对于性能或效果，论文能够支持的结论集中在 Figure 4 的 illustrative analysis：在 LLM inference 场景下叠加不同 memory encryption（SME）方案时，论文报告 secure initialization 开销为 47.88\%--73.45\%，secure communication 开销为 40.36\%--44.94\%。其结论是：memory-encryption 的粒度必须贴合加速器的访问模式，否则初始化与通信两段路径都会产生不可忽略的开销。\par
必须强调证据强度的边界：Figure 4 是针对特定 granularity mismatch 问题的示例性分析/实验展示，不是跨系统的通用 benchmark，不能据此推断``所有 accelerator TEE 都有 40\% 以上开销''。论文本身在 p.15 的边界声明中也持同样立场。\par
本节判断以 Wang/Huang 2026 Figure 4/p.10 及 README 实验分析为依据。超出这些材料覆盖范围的实现细节、性能结论或攻击面，本文不作延伸推断。\par
\subsubsection{综合评价}
综合来看，这篇 SoK 是设备侧 TEE 讨论的 taxonomy 补丁，让 Arm CCA I/O、GPU、DPU/SmartNIC  offload 的论述不局限于单个方案，并提供了一张可直接使用的``设备侧 checklist''：access control（SAC1--SAC6）、memory encryption（SME1--SME3）、attestation（SAT1--SAT4）、TCB 增量与软件栈规模、multi-type 与 plug-and-play 兼容性。\par
具体而言，它的优势在于 RQ 结构清晰、corpus 统计扎实（51 个系统、分平台分年代的可量化事实）、表格维度细，适合作为评估任何新设备 TEE 方案的对照基线。它的局限在于：作为 SoK 不能替代 ACAI、Devlore 或 TDISP/SPDM/IDE 规范本身；51 个研究的成熟度不一（学术论文、工业设计、source-limited 系统并存），不能简单同权比较；2026 年之后的厂商与标准状态可能变化。从部署角度看，它支撑 confidential AI、GPU cloud 与 DPU offload 的产品规划，但落地依赖厂商 HRoT、标准化证明链与 driver TCB 缩减。\par

\subsection{论文 36：Enabling Rack-scale Confidential Computing using Heterogeneous Trusted Execution Environment}
\subsubsection{论文基本信息}
《Enabling Rack-scale Confidential Computing using Heterogeneous Trusted Execution Environment》（HETEE）由 Jianping Zhu、Rui Hou、XiaoFeng Wang、Wenhao Wang、Jiangfeng Cao、Boyan Zhao、Zhongpu Wang、Yuhui Zhang、Jiameng Ying、Lixin Zhang、Dan Meng 完成（中国科学院信息工程研究所/计算技术研究所与 Indiana University），发表于 2020 年的 IEEE Symposium on Security and Privacy (S\&P 2020)。本文将其置于``加速器、DPU 与 SmartNIC 机密卸载''的研究脉络中考察，并把它作为 foundational system 类材料使用。它所回答的核心问题可概括为：如何在不修改任何 CPU 或加速器芯片的前提下，把 confidential computing 的边界从单机 CPU enclave 扩展到机架级（rack-scale）的商用异构加速器池。当前证据状态为：本地 PDF 已保存并核验；下文仅在该范围内讨论其机制与结论。\cite{zhu2020hetee}\par
\subsubsection{研究背景与动机}
要理解该方案，首先需要看到它面对的系统矛盾：深度学习训练/推理等 compute- and data-intensive (CDI) 任务依赖 GPU/TPU/FPGA 的高吞吐，但传统 TEE 只保护 CPU enclave；一旦明文被发给 GPU，driver、PCIe 总线与 device memory 都可能暴露给被攻陷的 host 或云管理员。论文在 Introduction 中给出三个递进的失效点：其一，当时的 Graviton、HIX 等异构 TEE 尝试都要求修改 CPU/GPU 芯片，开发周期长且无法使用现有硬件；其二，SGX 只部署到 Xeon E3 级别的处理器，数据中心主流的 Xeon E5 并不支持，把 CDI 任务连同 AI 框架、GPU runtime 和 driver 全部塞进 CPU TEE 既受资源限制（SGX 受保护内存不足 100MB）又会急剧膨胀 TCB；其三，按核心划分 enclave 的方案需要多个 enclave 紧密协作，扩大了资源共享带来的侧信道面。此外，机架级 pooling 还必须解决多任务的资源动态分配与任务结束后的安全清理（cleanup），这是单机 enclave 方案从未面对的问题。\par
\subsubsection{关键挑战与核心思想}
在上述背景下，作者提出的关键判断是：用一个较小的 Security Controller (SC) 作为机架内唯一的可信协调者，把庞大的 GPU 软件栈（CUDA、TensorFlow、driver）移出 TCB，同时用 PCIe 交换结构提供物理级隔离。论文显式列出两个设计挑战：（a）如何在共享计算资源的同时为 enclave 提供强隔离——答案是利用 PCIe ExpressFabric 动态、物理地把 enclave 与其他 enclave 及不可信 OS 隔开；（b）如何最小化 TCB——答案是两级隔离策略：SC 是唯一可信节点，只运行少量集成安全与管理功能的 firmware；proxy node 运行完整的重型软件栈，但在接触敏感数据之前由 SC 验证和保护，其状态恢复通过 secure reboot 完成，避免引入基于软件的 enclave 控制逻辑。\par
\subsubsection{系统架构}
从系统结构看，HETEE 的运行单元是一个置于防拆机箱内的 HETEE box（Figure 1），包含：Security Controller 节点（CPU+FPGA 实现）、若干 proxy node（基于 Intel Xeon E3 的低成本 microserver）、COTS 加速器资源池（原型为 4 块 NVIDIA TITAN X GPU，亦支持 FPGA 加速器），以及连接它们的 PCIe ExpressFabric 交换结构。HETEE box 通过 PCIe Switch Fabric 与同一机架上的其他主机相连；防拆机箱内置 MCU 与压力、振动、温度传感器，用于入侵检测与响应。\par
具体而言，工作流程为：远端用户通过机架上的某台主机转发请求，先与 SC 完成 remote attestation 并协商共享密钥；用户随后通过三类消息与 enclave 交互——configuration 消息（创建指定类型与数量加速器的 enclave，SC 分配唯一 ID）、code 消息（传输待执行程序，可为 ONNX 格式模型或 CUDA 代码）、data 消息（传输加密后的敏感数据）。对非敏感任务，SC 直接配置交换机把加速器连到普通主机；对敏感任务，SC 把加速器划分给与不可信环境物理隔离的 proxy enclave，解密用户数据后在 enclave 内运行核准代码；任务结束后，所有计算单元被 sanitize 并恢复到可信初始状态，再服务下一个任务。结果经加密路径返回用户，降低网络与 host OS 的观察面。\par
\subsubsection{核心方法设计}
论文的关键不是单一组件，而是将若干检查、状态和接口组合成一条受控的执行路径。\par
\noindent\textbf{（1）Security Controller 与两级隔离。}\par
这一机制的直接作用是：SC 是 HETEE 的可信控制面，负责把用户、任务和设备资源绑定。SC 的固件栈只包含四类功能：secure boot、remote attestation、resource assignment \& isolation、task 加解密与调度；它执行 remote attestation 与 key exchange、配置 PCIe switch 的资源分配，并保证不可信 OS 不直接接触敏感明文。proxy node 运行完整 OS 与 AI runtime，但由 SC 在启动时验证，运行时通过 PCIe fabric 的物理隔离与其他任务隔开。这一划分把 TCB 压缩到 SC 一处，是 HETEE 与``把 driver 塞进 enclave''路线的本质区别。\par
\noindent\textbf{（2）基于 PCIe ExpressFabric 的弹性资源分配与物理隔离。}\par
这一机制的直接作用是：机架级 pooling 而非单设备 enclave，是 HETEE 区别于此前工作的独特点。论文利用了 ExpressFabric 芯片的两个特性：software-defined fabric——交换芯片有专用管理端口，SC（作为管理 CPU）可初始化交换机、配置路由表、处理热插拔事件，使所有主机``只看到 SC 允许它们看到的设备''；以及 flexible topology——允许 mesh 等非层次拓扑。由此 SC 可以按需把多块加速器物理划分给同一 proxy enclave（Figure 2(a)），多 GPU 可用于提升 DNN 任务的并行度；资源回收时必须伴随 secure cleanup，proxy node 通过 secure reboot 恢复到安全状态。\par
\noindent\textbf{（3）Encrypted Data Stream。}\par
这一机制的直接作用是：把明文暴露限制在防拆机箱内部的可信边界内。远端用户将数据加密发送至 HETEE box（密文由不可信主机转发），SC 解密后交给 proxy enclave 调度到 GPU 执行，结果加密返回。由于明文只存在于 SC 与受隔离的 proxy/GPU 路径上，网络、主机 OS 与机架内其他租户都只能观察到密文。\par
\noindent\textbf{（4）面向现代标准的边界定位。}\par
这一机制的直接作用是：提醒读者 HETEE 的时代局限。HETEE 发表于 2020 年，早于 SPDM/TDISP/PCIe IDE/CoVE-IO 生态的成熟讨论：其设备身份建立在``COTS GPU 固件可信''的假设上（论文明确假设 GPU 固件无恶意代码且完整性受保护，并信任厂商固件更新的正确性），assignment 生命周期也没有标准化证据链。因此本文把它作为 foundational accelerator TEE，而非现代 trusted I/O 的完整答案。\par
\subsubsection{安全性与威胁边界}
从安全边界看，HETEE 考虑了一个相当强的对手：控制主机全部软件栈（含 OS）的特权软件攻击者，以及可对服务器内存总线进行嗅探的物理攻击者。其防御成立依赖一串显式假设：攻击者无法物理拆改 HETEE box（防拆自毁机箱 + MCU 传感器假设）；侧信道方面，论文把 covert channel 明确排除在范围外，电磁与功耗分析留给未来工作，对 cold boot 只提供部分防护（非法开箱时间长于内存残影保留时间时攻击才失败）；密文需经不可信主机转发，故 DoS 不在防御范围；密码算法本身与 FPGA 综合工具、mCPU 固件被信任。换言之，HETEE 证明的是``通过外部控制器与物理/总线隔离构造的机架级机密执行路径''，并不证明加速器自身成为完整 TEE；SC 仍是集中式 TCB，其供应链与容错风险由部署方承担。\par
\subsubsection{实验环境与证据}
论文没有把这部分只写成结论，而是以如下材料界定其证据范围：真实硬件原型上的 DNN inference/training 评估与硬件成本分析。\par
具体而言，关于硬件，原型由 PCIe ExpressFabric 背板、CPU+FPGA 的 SC、Intel Xeon E3 microserver（proxy node）与 4 块 NVIDIA TITAN X GPU 构成。关于 workload，评估覆盖真实规模的 152 层 ResNet（6000 万参数、约 200MB 模型）在约 138GiB ImageNet 数据集上的训练与推理，另有 VGG16、GoogLeNet、ResNet50/101 的对比（Figure 7--9）。关于指标，论文测量 throughput overhead、latency overhead、多 GPU 扩展性、带宽/延迟与成本（论文声称机密计算硬件开销低于计算单元成本的 5\%）。\par
本节判断以 HETEE 的 abstract、Figure 5--9 与 evaluation 章节为依据。超出这些材料覆盖范围的实现细节、性能结论或攻击面，本文不作延伸推断。\par
\subsubsection{实验结果与证据强度}
对于性能或效果，论文能够支持的结论是：HETEE 的吞吐开销很低，但延迟开销与 batch、模型与数据传输路径明显相关。\par
具体而言，就 Abstract 而言，ResNet152 上 inference 最大 throughput overhead 为 2.17\%，training 为 0.95\%。就 Introduction 的汇总而言，全部评测平均 inference throughput overhead 1.96\%、latency overhead 34.51\%，training 分别为 0.60\% 与 14.24\%。就 Figure 7 而言，batch=8 时 inference throughput overhead 平均 6.95\%、training 平均 0.91\%，多数训练场景低于 5\%。就 Figure 8 而言，batch=8 时 inference latency overhead 平均 42.96\%、training 18.54\%——延迟开销主要来自加密数据流的传输与 SC 中转，是部署实时推理服务时必须计入的成本。\par
本节判断以 HETEE 的 abstract、Introduction、Figure 7 与 Figure 8 为依据。超出这些材料覆盖范围的实现细节、性能结论或攻击面，本文不作延伸推断。\par
\subsubsection{综合评价}
综合来看，HETEE 是 accelerator confidential offload 的奠基性基线：强在完整的机架级系统视角与真实硬件原型，弱在现代协议闭环。\par
具体而言，它的优势在于不改芯片、直接覆盖 GPU DNN workload，并把资源池化、动态隔离与安全清理放进同一设计。它的局限在于：集中式 SC TCB（单点失效与供应链风险）、对 COTS GPU 固件的信任假设、缺少标准化 device identity/lifecycle 证据；后续 CloudScale 论文还指出其 proxy node OS 在 TCB 内、无法多租户复用同一设备、SC 始终位于数据路径上（CloudScale 报告 HETEE 在 GoogLeNet 上开销超过 2 倍）以及需要修改软件栈做代码/数据交接等部署障碍\cite{dhar2024cloudscale}。从部署角度看，HETEE 直接启发了把 security controller 做成 DPU/SmartNIC 形态的路线，但落地需要标准化 device evidence 与失效隔离机制。\par

\subsection{论文 37：Confidential Computing with Heterogeneous Devices at Cloud-Scale}
\subsubsection{论文基本信息}
《Confidential Computing with Heterogeneous Devices at Cloud-Scale》（CloudScale）由 Aritra Dhar、Supraja Sridhara、Shweta Shinde、Srdjan Capkun、Renzo Andri 完成（Huawei Zurich Research Center 与 ETH Zurich），发表于 2024 年的 Annual Computer Security Applications Conference (ACSAC 2024)。本文将其置于``加速器、DPU 与 SmartNIC 机密卸载''的研究脉络中考察，并把它作为 peer-reviewed SOTA system 类材料使用。它所回答的核心问题可概括为：在数据中心同时存在 TEE 节点与大量 legacy 非 TEE DSA（GPU/NPU/FPGA/SSD）的现实下，如何让用户不必在``性能''与``数据保护''之间二选一，并把 HETEE 式控制器思路扩展到跨机架、多租户的云规模。当前证据状态为：本地 PDF 已保存并核验；下文仅在该范围内讨论其机制与结论。\cite{dhar2024cloudscale}\par
\subsubsection{研究背景与动机}
要理解该方案，首先需要看到它面对的系统矛盾：现代数据中心按类型把节点聚成机架、机架间经高速交换互连（Figure 1 的 disaggregated 拓扑），但只有一小部分节点具备 TEE 能力——主流 CPU 已有 SEV/TDX/CCA 级 VM-based TEE，而除 NVIDIA H100 等少数例子外，大多数 GPU、NPU、SSD、FPGA 都不是 TEE 设备。把全部 DSA 替换为 TEE-capable 硬件成本高、周期长；而既有的 host-centric 方案（用 MMU 保护直连 TEE 主机的 MMIO 设备）不适用于 disaggregated 场景。论文进一步逐条指出 HETEE 在云规模下的不足：proxy node 的 OS 在 TCB 内（secure boot 只保证加载了签名镜像，不保证 OS 无漏洞）；proxy 与 GPU 独占分配给单租户，无法多租户；加速器与 proxy 物理绑定，不能动态改派；SC 始终位于数据路径上（带来如 GoogLeNet 超过 2 倍的开销）；且要求修改软件栈完成代码/数据交接。云规模还要求跨机架、多 SC 协同与多租户资源池，这些是单机架方案没有回答的。\par
\subsubsection{关键挑战与核心思想}
在上述背景下，作者先系统化地探索了五种分布式 TEE 设计（S0--S4，Figure 2）：S0 为集中式 SC 托管全部节点（HETEE 即属此类，单点失效、难扩展、无多租户）；S1 为集中式 SC 直连无 TEE 节点（设备仍须独占）；S2 为半集中式 SC 加部分 TEE 节点（SC 仍在用户到 enclave 的数据路径上）；S3 为全节点 TEE 的完全去中心化设计（安全与扩展性最优，但要求所有设备具备硬件 TEE，对 legacy DSA 不可行）；S4 为混合式——TEE-capable 设备按 S3 去中心化管理，legacy 非 TEE 设备挂在机架级 SC 之后由 SC 代理。论文选择 S4，其核心洞察是：SC 作为低 TCB 的专用硬件 TEE proxy，把非 TEE DSA 纳入可证明的物理/逻辑安全边界，而 TEE 节点的 enclave 不经 SC 直接通信，从而避免 HETEE 的全路径瓶颈。\par
由设计空间分析，论文提炼出五个必须解决的挑战（Table I）：恶意重配置（现代 DSA 固件/分区可被 hypervisor 重配，如 NVIDIA A100/H100 的 MIG）、持久安全状态（启动正确不代表运行期不被重刷）、绕过 SC（异步执行的恶意设备可直连其他租户数据）、物理攻击者（控制全部节点与网络设施、可探针总线）、以及 scale-out 不得削弱安全。对应的缓解为：SC+DSA attestation、SC 独占控制面、SC 拦截全部出入流量、AES-GCM 加密通信与防拆容器，以及 SC 芯片级的扩展性设计。\par
\subsubsection{系统架构}
从系统结构看（Figure 3），每个机架部署一个 SC，它是个裸金属逻辑的专用硬件模块：拥有硬件 root of trust，能对自身 firmware 生成带版本号的签名测量；对外分两侧——加密侧（encrypted side）面向 TEE 节点与其他机架的 SC，明文侧（plaintext side）面向与本 SC 同处防拆容器内的非 TEE 节点；容器配 tamper-detecting MCU（压力/振动传感器）。TEE 节点（如 AMD SEV-SNP confidential VM）由设备自身的 security monitor 提供隔离与证明，不经 SC；非 TEE 节点的一切出入流量必须经 SC 中转。SC 内部维护 Enclave Routing Table (ERT)——记录同一 job 的节点集合与 job 级共享密钥 K 的硬件查找表——以及供 DMA 映射的输入/输出缓冲（各 32KB）。跨机架 job 由主 SC 与其他机架 SC 互相证明后共享同一 K，数据只加解密一次。\par
\subsubsection{核心方法设计}
论文的关键不是单一组件，而是将若干检查、状态和接口组合成一条受控的执行路径。\par
\noindent\textbf{（1）Manifest 驱动的资源分配。}\par
这一机制的直接作用是：CloudScale 用一份 enclave manifest（Figure 4）描述用户任务需要哪些设备和安全属性——每个 CPU/DSA enclave 与非 TEE 节点的数量、类型、核数、内存/容量，以及要在其上运行的代码（示例 manifest 同时申请 4 核 64GB 的 CPU-TEE、8 核 32GB 的 NPU enclave 与一块 2TB 非 TEE SSD）。CSP-MP 按 manifest 分配资源，SC 依据 manifest 检查 ERT 中节点未被占用后参与证明；manifest 错误或被篡改会直接改变安全边界，因此 remote verifier 必须按 manifest 核验组合证明报告。资源分配是静态的（类似 AWS F1 的固定 1 VM+1 FPGA 模式）；动态扩容需要重新 attestation 并更新 ERT。\par
\noindent\textbf{（2）分布式远程证明协议。}\par
这一机制的直接作用是：把分布在 TEE 与 SC -backed 非 TEE 节点上的 enclave 聚合成一次可验证的证明。流程为：SC 先生成自身 firmware 测量与版本号的签名报告；收集 CPU enclave 的原生 attestation report；对每个 DSA enclave 收集含配置信息的设备证明；对非 TEE 节点，SC 通过物理接口将其 power-reset 到出厂配置与固件（销毁前一租户状态），并由 tamper-detecting MCU 提供 secure boot 与持续防拆监控的证明，附入组合报告。主 SC 将聚合报告发给 remote verifier，后者先验证主 SC 真实性，再逐项核验各节点报告与 manifest 的一致性，最后与 SC 建立共享密钥 K（节点间密钥交换采用 SIGMA 协议），SC 将 K 写入 ERT 并经安全通道下发。证明通过后 CSP-MP 即无法再篡改该分布式 enclave 的配置。\par
\noindent\textbf{（3）SC 数据面：ERT 访问控制与 AES-GCM 加密代理。}\par
这一机制的直接作用是：SC 不只是管理面，而是同时承担数据面。TEE 与非 TEE 节点之间的数据经 SC 的 staging buffer 换向：离开 SC 的数据用 ERT 中目标 job 的密钥做 AES-GCM 加密认证，进入 SC 的密文用对应密钥解密后放入明文侧 buffer；同一 SC 后的非 TEE 节点之间不走密码操作，而是由 ERT 做访问控制——仅当事务的源与目的属于同一 job 才放行，否则阻断；跨 SC 的非 TEE 节点通信则先用 K 加密再经机架间链路传输。SC 为每个非 TEE 节点分配不重叠的 buffer 区域，防止相互读取明文。原型中设备驱动只需小幅修改，使 CPU-TEE 发往非 TEE DSA 的 DMA 消息被拦截并复制到 SC 的输入 buffer。论文报告 SC chiplet 数据面 RTL 约 2.5 KLoC SystemVerilog，ERT 为 256 bit 宽、1024 项的硬件查找表，配 64KB（2$\times$32KB）SRAM 与单 AES-GCM 引擎。\par
\noindent\textbf{（4）SSD 侧的 TEE 化改造（对照设计）。}\par
这一机制的直接作用是：展示 SC 路线之外``把设备做成 TEE 节点''的增量改造成本。论文基于 SimpleSSD cycle-accurate 模拟器实现 SSD-enclave：用 namespace 对应的连续 NAND 物理地址范围加隔离 DRAM cache 构造 enclave；固件内实现硬件 LUT（namespace id 到 CPU-enclave id 与 AES-GCM key 的映射）与 access control unit（ACU），所有检查在 I/O 命中 DRAM cache 与页翻译之前完成，因此不改动 SSD 控制器、性能特征不变；用现成安全芯片做 measured boot，PCR 度量值用根密钥签名作为证明报告。\par
\subsubsection{安全性与威胁边界}
从安全边界看，CloudScale 的威胁模型（Section II-C）包含恶意 CSP（控制 OS、hypervisor、CSP-MP 与全部网络设施）、互不信任的共驻租户，以及可插拔恶意节点、探针总线的物理攻击者；Iago 攻击、侧信道、硬件木马与 DoS 明确排除。其安全论证（Section V）逐类对应：CSP-MP 错配资源会被组合证明报告发现；hypervisor 的恶意重配置对 TEE 节点可由 SM 配置证明检测，对非 TEE 节点由 SC 的 factory-reset 兜底；非 TEE 节点的越界 MMIO/DMA 被 ERT 阻断，ERT 与密钥对外不可读；物理攻击被防拆 MCU 容器（检测到入侵即复位 SC 与非 TEE 节点）与全链路 AES-GCM 阻断。边界同样清晰：SC 的正确性与供应链可信是硬假设；SC 被攻破时其管辖的所有 legacy 节点的机密性同时受损（blast radius 集中于 SC，这是集中式 TCB 的固有代价，论文未给出 SC 内部分区或冗余设计，此处论文未说明）；设备内部恶意固件、微架构侧信道与调度侧信道需另行处理。\par
\subsubsection{实验环境与证据}
论文没有把这部分只写成结论，而是以如下材料界定其证据范围：FPGA 控制面 + ASIC 综合数据面的 SC 原型、三类真实云负载的端到端评估与芯片级扩展性估算（Table III--V，Figure 5--7）。\par
具体而言，关于实现，SC 控制面（节点增删、attestation、VM 唤起）在 FPGA 上以 Vitis 密码库实现；数据面综合为 28nm HKMG CMOS ASIC（0.8V，TT corner），单 chiplet 面积 0.46mm$^2$、主频 870MHz、端到端延迟 1.15ns、每周期传输 16 字节，实用吞吐 12.96GB/s——论文特意对比 100GbE 以太网的理论吞吐 12.5GB/s，说明单 chiplet 足以线速处理机架级流量；chiplet 间经光学解复用总线互连（论文引用现有近零延迟 12.5ps 的解复用设计）。CPU TEE 为 256GB DRAM 的 AMD SEV-SNP 主机（Ubuntu 20.04 + KVM）；AI 加速器为 32 核 Huawei Ascend 910A NPU 与 Atlas 200 DK（Ascend 310）；SSD 用 SimpleSSD/DRAMSim3 模拟。初始化时 SC 从 VM 与 DSA 的 SW-TPM 收集约 245KB 的签名 secure boot traces。关于 workload，三类实验分别压测不同设备：FIO（4K 块、6 种读写模式）的文件系统路径、YCSB 驱动的 Redis 读写查询、以及 ImageNet 数据经 TEE-SSD 读入后在非 TEE NPU 上跑 ResNet-34/50（batch 1 与 16）的推理路径。\par
本节判断以 CloudScale 的 Section IV--VI、Table III--V 与 Figure 5--7 为依据。超出这些材料覆盖范围的实现细节、性能结论或攻击面，本文不作延伸推断。\par
\subsubsection{实验结果与证据强度}
对于性能或效果，论文能够支持的结论是：CloudScale 把端到端开销控制在实用范围，代价是 SC 成为关键瓶颈点与 TCB。\par
具体而言，就 Abstract 而言，AI、Redis、file-system 三类负载的平均 overhead 约 1.5\%--5\%。就 Figure 7 的端到端评估而言：FIO 场景 SEV VM 路径与 SC 引入的延迟开销分别为 1.84\%（顺序读写）与 1.94\%（随机读写）；论文进一步解释低开销的原因——CPU 侧 I/O 请求以突发形式下发，间隔分别为 10.237$\mu$s 与 6.607$\mu$s，而单个 4K 块在 SC 处的访问控制加 AES-GCM 只需 3.585$\mu$s，故 8 个 I/O（对应 32KB SC buffer）可流水隐藏全部延迟。Redis 场景读/更新查询延迟开销约 4.37\%。AI 场景（README 核验的 evaluation 数据）ResNet-34 batch 16 的最大开销约 5.3\%。\par
就扩展性而言（Table IV/V，芯片级估算）：以 ImageNet 平均 170KB JPEG 计算，单个 SC chiplet 可支持 RetinaNet-RN50（800 分辨率、30.68ms/图、5.41MB/s 饱和速率）达 2236 个并发 NPU，VGG-16 达 483 个、YOLOv3 达 562 个、UNet 达 171 个、ResNet-34 batch 1 达 90 个；新增 job 方面，最大的模型 VGG-16（380MB）也能以 31.85 jobs/s 的速率部署而无性能退化。SSD 场景下，单 chiplet 可支持 117 个并发 Intel 535 240G（随机读写 102.9MB/s）或 100 个（顺序读写 120.23MB/s）；Samsung Z-SSD、983DCT、850Pro 与 Intel 700 系列的并发数在 30--43 之间。论文同时声明扩展上限受 PCIe lane 数、SC 吞吐与负载延迟约束。\par
本节判断以 CloudScale 的 abstract、Section VI（含 Figure 6--7、Table IV--V）与 README 实验分析为依据。超出这些材料覆盖范围的实现细节、性能结论或攻击面，本文不作延伸推断。\par
\subsubsection{与相关工作的定量对比}
CloudScale 与 HETEE 的关系是``同思路、不同时代与规模''：两者都用外部安全控制器保护非 TEE 加速器，但 HETEE 的 SC 始终位于全部数据路径上、proxy OS 在 TCB 内、设备独占单租户且需修改软件栈；CloudScale 只把 SC 放在 TEE 与非 TEE 节点之间的路径上，不允许任何 CPU（因而任何 OS）作为非 TEE 节点接入 SC，从而把 OS/hypervisor 整体移出 TCB，并以 ERT 支持同 SC 后多设备的单租户共享与跨机架扩展。开销数字也呼应这一差别：HETEE 在 batch=8 时 inference latency overhead 平均 42.96\%（Figure 8），而 CloudScale 三类负载的平均 overhead 为 1.5\%--5\%。但两者共享同一硬假设——集中式 SC 的绝对可信与物理边界的完整。\par
\subsubsection{综合评价}
综合来看，CloudScale 是 cloud-scale accelerator/DPU 机密卸载的强 SOTA：它直面 legacy 异构设备这一真实云约束，给出了从设计空间论证、协议、RTL 到芯片扩展性估算的完整证据链。\par
具体而言，它的优势在于： threat model 显式包含物理攻击者与恶意 CSP，比只考虑恶意 OS/VMM 的工作更强；overhead 数据较完整且有机制级解释（3.585$\mu$s 的 SC 处理延迟可被 I/O 突发流水隐藏）；SC 硬件开销小（0.46mm$^2$）。它的局限在于：SC compromise 的 blast radius 集中、物理防拆容器依赖商用 MCU 假设、vendor attestation 与设备生命周期仍复杂；它不替代 SPDM/TDISP/IOMMU/PCIe IDE 标准，而是提供互补的硬件控制器路线。从部署角度看，SC 天然可以落地为 DPU/SmartNIC 形态的 security gateway，支撑 confidential accelerator pool；商业化风险集中在多租户失效隔离、SC 车队级运维与标准互操作。\par

\subsection{开发商安全实践建议}
本节三篇工作从 taxonomy、机架级原型与云规模系统三个角度，共同指向一组可操作的工程实践。以下建议面向 GPU/FPGA/DPU 厂商、SmartNIC/DPU 固件与系统软件开发商、以及运营机密计算池的云厂商。\par
\noindent\textbf{实践一：按``三机制 checklist''审计设备 TEE 设计，拒绝只解决 DMA 的方案。}SoK 的 SAC1--SAC6/SME1--SME3/SAT1--SAT4 矩阵（Table III）可直接转成设计评审清单：任何加速器/DPU 机密卸载方案必须同时回答 access control（谁可访问 workload、MMIO、DMA 与设备软件栈）、memory encryption（host memory、device memory、I/O link、metadata 四处的机密性/完整性/freshness）与 attestation（HRoT、endorser、reference values 三环齐备）\cite{sok-tee}。评审时对照 AT1--AT8/AE1--AE6/AA1--AA2 攻击枚举逐项确认``完全防御/部分防御/无法防御''，只覆盖 DMA 窗口的方案应在文档中显式标注残余攻击面。\par
\noindent\textbf{实践二：为每台设备建立硬件级身份与固件证明链，而非依赖平台侧间接推断。}SoK 的 Figure 6/Table VIII--IX 显示，缺 HRoT/endorser/reference values 的加速器可被仿真设备或错误固件欺骗\cite{sok-tee}；CloudScale 则用 power-reset 到出厂配置 + tamper-detecting MCU 证明 + 组合 attestation report 的方式，为没有原生证明能力的 legacy 设备补出等效证据\cite{dhar2024cloudscale}。设备厂商应实现 SPDM 级身份与固件测量（新设备直接支持 TDISP/DSM；legacy 设备由机架级控制器代理证明），云厂商的 verifier policy 应把``设备证明缺失''视为硬失败而非降级选项。\par
\noindent\textbf{实践三：把可信控制面收敛到一个小型安全控制器，并把 OS/driver/AI runtime 移出 TCB。}HETEE 与 CloudScale 的共同经验是：可行的部署形态是一个固件功能可枚举的 security controller（secure boot、remote attestation、resource assignment、任务加解密调度四类功能）\cite{zhu2020hetee}，其数据面可以小到约 2.5 KLoC SystemVerilog、0.46mm$^2$\cite{dhar2024cloudscale}。DPU/SmartNIC 厂商应把该控制器做成板载安全岛；云厂商不应把百万行级的闭源 GPU driver/runtime 纳入 enclave TCB——SoK 的 Table X--XIII 已量化这一风险\cite{sok-tee}。\par
\noindent\textbf{实践四：为任务结束设计可验证的清理与状态恢复流程。}HETEE 用 secure reboot 把 proxy node 与加速器恢复到可信初始状态\cite{zhu2020hetee}；CloudScale 用 power-reset 到出厂配置销毁前一租户状态，并在 ERT 中删除映射后才允许再分配\cite{dhar2024cloudscale}。云厂商的多租户调度器应把``清理完成且可证明''作为设备重新上架的前置条件，特别是 GPU 显存、FPGA 配置存储与 SSD namespace 的残留数据。\par
\noindent\textbf{实践五：按数据面实测参数设计可信队列与缓冲，使安全处理可被流水隐藏。}CloudScale 的评估给出可复用的设计点：单 4K 块在 SC 的访问控制 + AES-GCM 处理为 3.585$\mu$s，而 CPU 侧 I/O 突发间隔为 6.6--10.2$\mu$s，32KB 缓冲即可让 8 个 I/O 流水化、把端到端开销压到 2\% 量级\cite{dhar2024cloudscale}。设备厂商设计可信队列时应按``安全处理延迟 $\le$ 突发间隔''配平流水线深度与缓冲容量，并把该参数写入数据手册；反之，若安全路径延迟超过负载的 I/O 间隔（如 HETEE 的小 batch 推理路径，latency overhead 平均 42.96\%\cite{zhu2020hetee}），应在产品文档中明示适用负载类型。\par
\noindent\textbf{实践六：将调度隔离建立在硬件路由表而非软件策略上。}CloudScale 的 ERT（256 bit 宽、1024 项硬件查找表，源/目的同属一个 job 才放行）证明租户间隔离可以在数据面硬件内完成，且不随 hypervisor 被攻陷而失效\cite{dhar2024cloudscale}。DPU 厂商应提供类似的硬件级 job/租户路由原语；云厂商应把多租户调度决策（manifest）交由 attested 控制面核验，避免让不可信 CSP-MP 单方决定设备-租户绑定。\par

\subsection{未来研究方向}
以下研究点按``问题—为何未解决—技术路线—所需实验''组织，其入口均可回溯到本章标注的证据缺口。\par
\noindent\textbf{方向一：CPU-加速器一体化证明与统一 verifier 策略。}问题是：现有 attestation 把 CPU TEE report、设备 HRoT 测量、driver/runtime 度量与任务级证据分散处理，SoK 的 Table VIII/IX 显示多数加速器缺少完整证明链\cite{sok-tee}，CloudScale 的组合报告则依赖自研 SC 而非标准格式\cite{dhar2024cloudscale}。技术路线：以 SPDM/TDISP 证据为底，叠加 driver/runtime 的测量与任务 manifest，定义可组合的 verifier policy 与证据聚合格式。所需实验：对 GPU/NPU/DPU 三类设备构建端到端证明原型，测量证明生成与验证延迟，并用仿真设备、降级固件、重放证据三类攻击检验 policy 完备性。\par
\noindent\textbf{方向二：面向加速器访问模式的安全内存元数据粒度。}问题是：SoK Figure 4 显示 memory encryption 粒度不匹配会带来 47.88\%--73.45\% 的初始化开销与 40.36\%--44.94\% 的通信开销\cite{sok-tee}，但该数字来自 illustrative analysis，缺乏跨设备、跨负载的系统化最优粒度理论。技术路线：按 GPU/NPU 的突发大块访问与 SSD 的 4K 随机访问分别设计可变粒度完整性树与 freshness 计数器，并与 IOMMU/SMMU 的页粒度联动。所需实验：在真实 accelerator TEE 原型（或 cycle-accurate 模拟）上扫描粒度参数空间，报告 LLM 推理、CNN 训练与 FIO 三类负载的开销曲线。\par
\noindent\textbf{方向三：driver/runtime TCB 的结构性缩减。}问题是：SoK 指出加速器软件栈可达百万行且常闭源，把其放进 TEE 根本削弱安全论断\cite{sok-tee}；HETEE/CloudScale 用外部控制器绕开了这一问题，但代价是设备功能受限（静态资源分配、无动态 malloc\cite{dhar2024cloudscale}）。技术路线：把 driver 拆为可信薄层（队列、MMIO 门控、页表切换）与不可信富层（编译、调度优化），用硬件强制薄层中介所有特权操作。所需实验：对开源驱动栈（如 ROCm 或 nouveau）做 TCB 划分案例研究，量化薄层 LoC、接口数与性能损失，并做脆弱性对比分析。\par
\noindent\textbf{方向四：机密感知的加速器多租户调度。}问题是：HETEE 与 CloudScale 都要求非 TEE 设备单租户独占、资源静态分配\cite{zhu2020hetee,dhar2024cloudscale}；云厂商真正需要的是带隔离与计费保证的动态共享，而调度本身会引入侧信道（两篇论文均将调度/应用层侧信道排除在威胁模型外）。技术路线：将调度策略纳入 manifest 并由 SC/DSM 执行，结合 MIG 式硬件分区与 ERT 式路由表，对调度时序做常量化或加噪。所需实验：多租户混合负载下的隔离性测试（跨租户残存数据扫描）、吞吐公平性与调度侧信道测量。\par
\noindent\textbf{方向五：集中式安全控制器的 blast-radius 削减与车队级运维。}问题是：HETEE 与 CloudScale 都把机架级机密性押在单一 SC 上，SC 被攻破则其辖内全部 legacy 节点同时失密；两篇论文均未给出 SC 内部分区、冗余或密钥轮换设计（论文未说明）\cite{zhu2020hetee,dhar2024cloudscale}。技术路线：SC 内部按租户/job 做硬件隔离域，ERT 密钥支持按 job 轮换与撤销；跨 SC 用阈值化或分片信任（如多 SC 联合授权敏感操作）。所需实验：注入 SC 固件漏洞与供应链篡改的故障演练，测量受影响 job 比例、密钥轮换开销与撤销传播延迟。\par
\noindent\textbf{方向六：混合 TEE/non-TEE 负载的端到端基准与标准化证据链。}问题是：本章三篇工作的性能数字来自各自原型与负载（HETEE 的 DNN、CloudScale 的 FIO/Redis/ResNet、SoK 的 LLM 示例），缺乏可横向复现的 confidential offload benchmark；且 HETEE 明确缺少 SPDM/TDISP/IDE 时代之前的标准化 device evidence\cite{zhu2020hetee}。技术路线：定义覆盖 AI 训练/推理、存储、网络的公开 benchmark 套件，固定威胁模型与报告字段（开销分解为初始化/通信/调度三段），并把 TDISP 状态机、IDE 链路加密与 SPDM 身份纳入参考实现。所需实验：在至少两类 CPU TEE（SEV-SNP 与 Arm CCA 或 TDX）与两类设备形态（GPU 与 DPU/SmartNIC）上运行套件，公开原始数据供第三方复现。\par
\section{ISA 与硬件设计纵深防御}
平台级隔离还需要 ISA 与微架构提供纵深防御。前几节讨论的 TEE、证明与设备隔离机制，最终都要落在一组更底层的语义上：谁拥有最高特权级、谁处理 trap、物理内存的哪一段对谁可读可写可执行、设备 DMA 请求在哪条路径上被检查、内存密文的新鲜度由什么结构保证、以及虚拟化切换时 CPU 状态向谁暴露。本节讨论权限、I/O 保护、内存加密完整性和 guest 状态保护等机制，说明它们如何约束边界内代码对合法权限的滥用，并与上层 TEE、attestation 和设备隔离形成互补。\par
这里的比较重点是机制的可组合性与证据等级差异：某一项防护可能收紧控制流或内存访问，却不自动处理 DMA、回滚、侧信道或可信软件缺陷。本节四份材料恰好覆盖四种证据形态——一份 ISA 规范（语义权威但无实验）、一篇有完整 RTL 原型与 FPGA 评估的系统论文、一篇 2007 年模拟器时代的经典机制论文，以及一份只有机制图示、没有 benchmark 的工业白皮书。逐篇分析会明确每种方案实际收紧的攻击面、可核验的性能数字，以及仍需由系统软件和部署策略承担的部分；对规范与白皮书不覆盖的内容，本节明确标注证据边界，不作外推。\par
\subsection{论文 38：The RISC-V Instruction Set Manual: Privileged Architecture}
\subsubsection{论文基本信息}
《The RISC-V Instruction Set Manual: Privileged Architecture》由 RISC-V Foundation 完成，本地核验版本为 20260120 Official Release（Volume II，221 页 A4 PDF），具体发表载体在本地元数据中未注明。本文将其置于"RISC-V 基础安全机制"的研究脉络中考察，并把它作为 spec 类材料使用。它所回答的核心问题可概括为：Privileged spec 是所有 RISC-V TEE 的 CPU-side 词典。当前证据状态为：本地 PDF 已保存并核验；下文仅在该范围内讨论其机制与结论，不把它扩展成任何具体 TEE 实现的安全性或性能证据。\cite{riscv_privileged}\par
\subsubsection{研究背景与动机}
要理解该规范的安全意义，首先需要看到 RISC-V 生态面对的系统矛盾：开放 ISA 允许任何实现者自由组合特权级、扩展和平台模块，这带来灵活性，也意味着没有任何一家厂商替整个生态规定"安全边界应该画在哪"。TEE 设计者必须自己明确每个 privilege level 的责任，并证明所选语义足以支撑其威胁模型。\par
具体而言，规范在 Chapter 1 给出三级特权模式（Table 1：U=0、S=1、Reserved=2、M=3），其中 M-mode 是唯一强制实现的特权级，规范原文明确指出"M-mode can be used to manage secure execution environments"。S/HS/VS 管 OS、hypervisor 与 guest state，PMP/Smepmp 与页表是早期 enclave 与后续 CoVE 的隔离底座。规范 Chapter 1 同时承认一个更深层的事实：整套 privileged 设计"could be replaced with an entirely different privileged-level design without changing the unprivileged ISA"——即安全语义的权威来源是这份规范本身而非 ISA 的不可变部分，这决定了它在证据链中的特殊地位。\par
\subsubsection{关键挑战与核心思想}
在上述背景下，可以把规范要解决的设计问题归纳为三条，它们也是所有 RISC-V TEE 论文的分叉点。\par
\noindent\textbf{挑战一：最高权限放在哪一层。} M-mode 权力最大、通常承载 firmware 与 security monitor，但把全部信任压在一层会扩大 TCB；Keystone/Penglai 把信任放在 M-mode monitor 加 PMP，CoVE/AP-TEE 则把 confidential VM 推到 TVM/TSM 与额外硬件扩展。规范本身不裁决这条分叉，只提供两套语义的公共底座。\par
\noindent\textbf{挑战二：trap 与异常的控制权如何分配。} medeleg/mideleg 决定异常与中断流向 M、S 还是 HS/VS；错误的 delegation 会把本应由 monitor 处理的安全事件交给不可信 OS，直接扩大其控制面。\par
\noindent\textbf{挑战三：物理内存与虚拟内存两层防护如何组合。} PMP 以物理地址为对象、由 M-mode 配置，页表以虚拟地址为对象、由 S/HS 配置；二者检查对象、配置者和失败语义都不同，规范用"effective privilege mode"统一描述（页表 walk 按 S-mode 等效检查），但组合的正确性完全落在软件配置上。\par
\subsubsection{系统架构}
从系统结构看，规范定义的是 execution stack 和可用硬件隔离语义，TEE 在其上分配 TCB。Figure 1 展示三类特权执行栈（单应用 AEE、多程序 OS+SEE、hypervisor+HEE），Table 1 给出三级特权编码。安全相关核心章节包括：Chapter 3 的 Machine-Level ISA（mstatus、menvcfg、mseccfg、medeleg/mideleg、PMA 与 PMP），Chapter 6 的 Smepmp 扩展，Chapter 12 的 Sv32/Sv39/Sv48/Sv57 页式虚存，Chapter 22 的 Hypervisor 扩展与两阶段地址翻译，以及 Chapter 23 的 CFI（Zicfilp landing pad 与 Zicfiss shadow stack）。本地 221 页 PDF 中这些模块全部标注为 Ratified。\par
\subsubsection{核心方法设计}
规范的价值不在单一组件，而在于它把若干检查、状态和接口组合成一条受控的执行路径。以下逐条拆解与安全最相关的机制，条目均可在本地 PDF 对应章节核验。\par
\noindent\textbf{（1）Privilege Levels 与 Trap Delegation（Chapter 1、3.1.8）。}\par
这一机制的直接作用是：TEE 要先决定谁处理 trap、谁拥有最高权限。规范规定 M-mode 是 mandatory highest privilege，越权操作触发异常并 trap 到下层执行环境。medeleg/mideleg 按 bit 把同步异常与中断下放给 S-mode，H 扩展下再由 hedeleg/hideleg 继续下放到 VS-mode。规范不给"正确的 delegation 表"，但错误配置的后果是确定的：被下放的 page fault 类异常会让不可信 OS 看到 enclave 的缺页模式，被错误保留的中断会让 monitor 成为所有 I/O 事件的单点。\par
\noindent\textbf{（2）PMP 与 Smepmp 内存边界（3.7、Chapter 6）。}\par
这一机制的直接作用是：PMP 是 RISC-V enclave 谱系最早的硬件隔离抓手。规范定义的关键语义包括：\par
\begin{itemize}
\item 每个 PMP entry 由一个 8-bit 配置寄存器（pmpcfg，含 R/W/X/A/L 字段）和一个地址寄存器（pmpaddr）组成；最多支持 64 个 entry，实现可选 0、16 或 64 个，且必须从低编号连续实现；PMP CSR 只对 M-mode 可访问。
\item 地址匹配模式 A 字段编码为 OFF/TOR/NA4/NAPOT，标准编码支持最小 4 字节区域；entry 静态优先级，编号最小且匹配全部字节的 entry 决定访问成败，部分匹配即失败。
\item 默认语义不对称：S/U-mode 默认无任何物理内存权限，须由 PMP 授予；M-mode 默认拥有全部权限，除非 entry 的 L 位置位。L 位同时是锁——置位后配置与地址寄存器写被忽略，直到 hart 复位，且 R/W/X 权限对 M-mode 同样生效。
\item 检查时机与虚拟化组合：PMP 与 PMA 并行检查，覆盖 S/U 的取指与数据访问、mstatus.MPRV 置位且 MPP=S/U 时的 M-mode 数据访问，以及页表 walk（等效 S-mode）。修改 PMP 后必须执行 SFENCE.VMA（rs1=x0, rs2=x0）同步地址翻译缓存；PMP 违规总是精确陷入。
\end{itemize}
Smepmp（Chapter 6，含 6.1 Threat model 一节）在此基础上解决"M-mode 自身不可信"的问题：新增 mseccfg 的 MML/MMWP/RLB 三个锁定式字段，给出一张 L 位与 R/W/X 组合的锁定规则表（包括 Locked Execute-only、Locked Shared code/data region 等 16 种编码，锁定规则除非 mseccfg.RLB 置位否则到 PMP 复位前不可移除），并要求 BootROM 在早期启动阶段置位 MMWP/MML 以便 firmware 发现扩展。这把"默认拒绝、固件自锁"从软件约定提升为 ISA 语义。\par
\noindent\textbf{（3）Virtual Memory 与 Hypervisor Extension（Chapter 12、22）。}\par
这一机制的直接作用是：CoVE/AP-TEE 需要从 enclave 走向 TVM，这依赖 H extension 和两阶段翻译。规范定义 HS/VS/VU 模式分工、hgatp/vsatp 两套地址翻译 CSR、以及 guest physical 到 host physical 的第二阶段翻译（22.5 节，含 guest-page fault 语义）。对机密计算而言，页表权限位（R/W/X/U/G/A/D）与第二阶段翻译是所有 CVM 内存隔离的底座，但规范只定义翻译与权限语义，不定义 memory tracking 与 ownership——后者正是 CoVE 需要 MTT 一类额外结构的原因。\par
\noindent\textbf{（4）状态收敛与侧通道控制（Chapter 4、7、23）。}\par
这组机制回答"新扩展状态会不会成为跨安全域的通道"。Smstateen（Chapter 4）的动机章节明确以 AIA 新增的至多十个 supervisor CSR 为例：若 hypervisor 不知道这些寄存器存在就不会做上下文切换，它们就成为 guest OS 之间的 covert channel，state-enable CSR 由此提供按扩展封堵的开关。Smcntrpmf（Chapter 7）指出 cycle/instret 计数器默认跨特权级连续计数会"leak information about privileged software execution to user mode"，引入 MINH/SINH/UINH/VSINH/VUINH 按模式过滤。Chapter 23 的 Zicfilp/Zicfiss 给出 landing pad 与 shadow stack 的 ISA 级 CFI 语义（含 ELP 状态、ssp CSR 访问控制与 shadow stack 内存保护）。这些章节说明规范已把微架构状态泄漏纳入 ISA 治理范围，但只提供开关与语义，不提供泄漏分析本身。\par
\subsubsection{安全性与威胁边界}
从安全边界看，规范提供的是底层权限、内存与控制流约束的语义，而不是完整机密计算平台。它能支撑的判断是：PMP/Smepmp 可限制物理内存访问并把 M-mode 自身纳入约束，delegation 语义决定 trap 控制面，H 扩展支撑虚拟化分工，state-enable 与 CFI 收敛扩展状态与控制流。它明确不覆盖：DMA 与 I/O requester（需 IOPMP/IOMMU，见论文 39）、device identity、内存加密与完整性（见论文 40、41）、attestation、侧信道与推测执行。规范的强假设是实现与 privileged firmware 正确；PMP entry 数量有限、配置错误、M-mode TCB 过大都是 README 与本节共同确认的失效路径。\par
\subsubsection{实验环境与证据}
这是官方规范类材料，无实验。本地证据为 221 页 Official Release PDF（pdfinfo 核验），可直接支撑的结论是 privilege level 编码（Table 1）、PMP/Smepmp 规则（3.7、Chapter 6 及其锁定规则表）、trap delegation CSR（3.1.8）、Sv* 页表权限与 H 扩展两阶段翻译（Chapter 12、22）、CFI 与 state-enable 语义（Chapter 4、23）。这些材料不足以支撑任何具体 TEE（Keystone/Penglai/SPEAR-V/CoVE）的性能、TCB 大小或完整安全证明。\par
\subsubsection{实验结果与证据强度}
规范不提供 overhead 数字。trap 处理、PMP 检查、页表 walk 都可能影响性能，但规范不测也不应被引用为性能证据；性能必须引用各系统论文。本材料的 claim strength 在于：凡与规范语义冲突的 TEE 设计叙述都可以被它证伪，这是它在综述中的独特作用。\par
\subsubsection{综合评价}
综合来看，Privileged spec 是 RISC-V 安全的地基，价值高但抽象层低。它的优势在于官方、全面、且 20260120 版本已把 Smepmp、Smstateen、CFI、Pointer Masking、Hypervisor 全部收编为 Ratified 模块，决定了所有后续设计的书面语言。它的局限在于它不是 TEE 方案：不含 verifier、attestation、I/O lifecycle，也不证明任何配置正确。从部署角度看，所有 RISC-V confidential computing 产品都必须兼容这些语义；对开发者而言，真正的风险不在规范缺失，而在把 64 个 PMP entry、menvcfg/mseccfg 与 delegation 表组合成具体平台时的配置正确性。\par
\subsection{论文 39：sIOPMP: Scalable and Efficient I/O Protection for TEEs}
\subsubsection{论文基本信息}
《sIOPMP: Scalable and Efficient I/O Protection for TEEs》由 Erhu Feng、Dahu Feng、Dong Du、Yubin Xia、Wenbin Zheng、Siqi Zhao、Haibo Chen 完成（上海交通大学 IPADS、阿里巴巴达摩院、清华大学、上海人工智能实验室），发表于 ASPLOS 2024（La Jolla，DOI: 10.1145/3620665.3640378）。本文将其置于"RISC-V CoVE-IO / TEE-I/O"的研究脉络中考察，并把它作为 system 类材料使用。它所回答的核心问题可概括为：sIOPMP 的核心贡献是把设备访问控制做成可扩展硬件检查，而不是靠少量粗粒度 IOPMP entry。当前证据状态为：本地 peer-reviewed PDF 已保存并核验；下文机制与数字均出自该 PDF。\cite{feng2024siopmp}\par
\subsubsection{研究背景与动机}
要理解该方案，首先需要看到它面对的系统矛盾：CPU 侧 TEE 边界会被 I/O 打穿，因为 DMA 不经过 CPU page table。enclave/TVM 私有页即使 CPU 访问受限，也可能被设备 DMA 直接读写；而现有 I/O 隔离方案各自失效于不同点。论文第 1 页明确指出：传统 IOMMU 或软件 I/O 在 I/O 密集负载下会造成至少 20\% 的吞吐下降，其原因被拆解为四层——IOTLB 失效走异步命令队列，已有工作测得 IOTLB flush 可造成 20\%$\sim$30\% 的 I/O 吞吐开销；IOMMU 只支持页级隔离，无法处理网络栈中大量 sub-page 报文，需要额外拷贝；IOVA/页表配置复杂且暴露过共享数据结构（descriptor ring）绕过检查的攻击面；多设备多租户下 IOVA 分配与 IOTLB 成为扩展性瓶颈。区域式方案同样受限：TrustZone 只支持 16 个区域与两种设备角色，IOPMP 基线只支持 64 个 source ID；而现代 DMA 控制器的 scatter-gather 模式可支持 512 乃至 1024 个 scatter buffer，纯内存加密方案则挡不住 DMA replay（无完整性树时）且阻断合法 DMA、引入 bounce buffer 拷贝（论文引用 bifrost 的 23\% 开销为例）。\par
\subsubsection{关键挑战与核心思想}
在上述背景下，作者提出的关键判断是：I/O protection 的关键不是"有没有检查"，而是检查是否能随设备和内存规模扩展。论文把设计目标拆成三个可核验的挑战，并各给一个观察支撑的机制：\par
\noindent\textbf{挑战一：超过 1000 个优先级区域的检查不能压低时钟频率。} 优先级区域检查是组合逻辑，线性检查 1000+ entry 会拉长关键路径。论文的观察是：设备对带宽敏感而对几个 cycle 的延迟不敏感（即使在 CXL 场景），因此 CPU 侧不能用的流水线设计在设备侧可用，由此提出 Multi-stage-Tree-based checker（MT checker）。\par
\noindent\textbf{挑战二：有限 SoC 资源如何支持数量不受限的设备。} 虚拟功能（VF）与热插拔使设备总数无法硬编码。论文的观察是：同时活跃的热设备数量受 CPU 核数与总线容量约束，总是有限的，由此提出 mountable entry——冷设备条目放在受保护内存的扩展表中，只在首次访问时装载。\par
\noindent\textbf{挑战三：设备负载是动态的。} 冷热状态需要随 workload 切换。论文利用 CAM 可按内容单周期查找的特性，设计零成本 remapping 机制在冷热状态间切换设备。\par
\subsubsection{系统架构}
从系统结构看，设备请求先被识别为某个 source（SID/DeviceID），再经 sIOPMP metadata 判断能否访问目标物理地址。威胁模型（论文 3.2 节）规定 TCB 只包含 SoC 内部模块（CPU core、sIOPMP 扩展等）与轻量 firmware（secure monitor）；一切片外设备（GPU、加速器、SmartNIC）都可被攻破，REE 侧一切软件（含驱动）不可信；显式不考虑 DRAM 物理攻击（冷冻内存、侦听、splicing）与 SoC 内侧信道，这些被列为正交问题。上层软件栈建立在 Penglai enclave 的 secure monitor 之上，扩展出 sIOPMP 配置能力，并用 capability-based ownership 接口（Create\_TEE()、Device\_map()）把设备与内存区域的 ownership 在 TEE 间转移。\par
\subsubsection{核心方法设计}
论文的关键不是单一组件，而是四个机制组合成的检查路径。以下拆解均可在论文第 4、5 节核验。\par
\noindent\textbf{（1）Multi-stage-Tree-based IOPMP checker（MT checker）。}\par
这一机制的直接作用是：在不降低时钟频率的前提下把可检查 entry 数推到 1000+。具体实现上，MT checker 由两部分组成：IOPMP pipeline 把大检查器拆成多级流水线小检查器，中间结果存寄存器，最后一级给出 allow/deny；tree-based arbitration 把"从低优先级到高优先级逐个比较"的线性逻辑改为逐对比较后按树形电路归约，把检查延迟降到 log(N) 量级，且在 RTL 层（而非依赖 EDA 后端）实现，可按面积/时序需求选二叉树或 N 叉树。流水线引入了新的状态一致性问题：总线上已阻断的 DMA 事务可能仍在检查器流水线中，论文为此加了一个 monitor 维护总线与检查器之间 block 状态的一致视图。\par
\noindent\textbf{（2）Mountable IOPMP 与扩展表。}\par
这一机制的直接作用是：用有限的硬件 entry 支持数量不受限的设备。具体实现上，扩展 IOPMP 表放在受保护内存中（RISC-V 下可由 PMP 保护），无硬件大小限制；每个 mountable entry 含扩展 SID（eSID）、关联 memory domain 索引与附加 IOPMP entry。当请求设备不在 SRC2MD 表中时触发 SID-missing 中断，secure monitor 从扩展表取出 eSID、domain 与 entry 装载进真实硬件寄存器；冷设备总是与最后一个 memory domain（MD62）配对，切换时需 flush 该 domain 的 entry。装载开销只在冷设备首次 DMA 请求时发生一次，之后该设备临时获得快速路径。\par
\noindent\textbf{（3）零成本冷热 remapping。}\par
这一机制的直接作用是：按动态 workload 在冷热状态间切换设备而不阻塞数据通路。具体实现上，利用 CAM 按内容单周期查找返回地址的特性做 remapping，热设备请求不经过任何额外状态转换。论文第 6 节的评估表明该机制的正确使用是性能前提（见"失败模式与边界分析"）。\par
\noindent\textbf{（4）违规处理与配置一致性。}\par
这一机制的直接作用是：发现越权 DMA 后既不能放行数据，也不能让配置更新打开攻击窗口。具体实现上，sIOPMP 提供两种违规处理：packet masking 用 write strobe 与 read clear 信号屏蔽越权请求的读写数据；bus-error handling 直接终止事务。修改 entry 时存在 entry inconsistency（新旧区域同时可访问的窗口），论文用 SID block bitmap 在修改期间按 SID 阻断对应设备；冷设备切换存在 device inconsistency（切换未完成前可能访问上一设备的 domain），论文在全部修改完成前阻断该冷设备的 DMA。两种阻断都是 per-SID 的，不影响其他设备的 I/O 性能。\par
\subsubsection{安全性与威胁边界}
从安全边界看，sIOPMP 提供的是 DMA/MMIO 路径上的访问控制 enforcement，而不是完整 trusted I/O 栈。它能阻止的具体越权路径是：恶意或被攻破设备发起的任意 DMA 读写（含 REE 软件触发的 DMA），并对 DMA replay 给出区域级防护语义（Table 1 中 read/write/replay 一栏）。它没有覆盖的部分同样明确：不证明设备固件身份（需 SPDM/TDISP），不提供 PCIe 链路加密与新鲜度（需 IDE），不处理中断投递（需 AIA 类机制），不覆盖 DRAM 物理攻击与 SoC 内侧信道（论文 3.2 节自述）。sIOPMP 与 TEE-IO 是正交的：论文指出 SEV-TIO/TDX-TEE-IO 解决"可信设备如何直接访问 TEE 明文"，但仍依赖 RMP/GPC 做 I/O 隔离，动态负载下同样受异步失效拖累。\par
\subsubsection{失败模式与边界分析}
论文自身讨论了两类值得写入综述的负面结果。其一是配置错误导致的性能崩塌：在冷设备切换评估中（Figure 17），若把两个并发设备都错误设为冷设备，当二者 DMA 请求比为 1:10 时，频繁的换入换出会浪费"热"设备 85\% 的 I/O 吞吐；设备状态必须按 workload 正确设置（这正是 remapping 机制存在的理由）。其二是机制边界的显式声明：mountable 机制"ideal if there is only one cold device running simultaneously"，多冷设备并发场景的性能不在论文覆盖内；bus-error 与 packet masking 在错误检测延迟上行为不同（后者需设备自行消化被屏蔽的包），选用哪种违规处理会影响错误路径时延。此外论文 Discussion 节确认当前配置（1000 entry、64 热设备）是按当代数据中心机器（CPU 少于 64 核、scatter buffer 少于 1024）标定的，参数需随机器演进调整。\par
\subsubsection{实验环境与证据}
论文的实现与评估环境完整可核验：基于 chipyard（RISC-V SoC 生成器）实现原型，CPU 为 4 核 Boom（乱序）或 4 核 Rocket（顺序），模拟频率 3.2GHz，L1 I/D 各 32KB，L2 512KB 15 路；设备含 IceNet 100Gb/s NIC、NVDLA 深度学习加速器与 dummy DMA 节点；sIOPMP 可配置为 per-device 或集中式、1/2/3 级流水、32$\sim$1024 entry、两种违规处理。由于 RISC-V SoC 当时没有 IOPMP 实现，作者把 Boom/Rocket 核内 PMP 检查器移植为 IOPMP 基线。微基准用 Verilator，应用基准用 FireSim；IOMMU 对比在 Intel Xeon Gold 5317 真实服务器上进行，SWIO（SEV-SNP 式 bounce buffer）对比环境为 Intel Xeon Gold 5317、AMD EPYC 7T83 与 100Gbps NIC。软件栈为 Linux 5.15 + 扩展后的 Penglai monitor。\par
\subsubsection{实验结果与证据强度}
论文的主要定量结论如下，均可对照原文图表核验：\par
\begin{itemize}
\item 时钟频率（Figure 10）：基线 IOPMP 在 FPGA 平台（上限 60MHz）只能维持到 128 entry，1024 entry 无法通过时序分析；纯流水线 2 级只能维持 256 entry，1024 entry 时跌至 10MHz；流水线加树形仲裁可在 512 entry 维持 60MHz，1024 entry 仅轻微退化；3 级流水加树形仲裁可支持 1024+ entry。
\item 检查延迟（Figure 11，64 个连续 burst 请求、每 burst 8 beat 每 beat 8 字节的最坏情形）：读路径基线 1510 cycle，2 级流水 bus-error 为 1575 cycle，packet masking 为 1634 cycle；写路径基线 1081 cycle，两种违规处理下分别为 1175/1189 cycle。吞吐（Figure 12）：Read-Read 场景基线 5.18 byte/cycle，2 级流水 5.08 byte/cycle，Write-Write 与 Read-Write 场景流水线不引入额外开销。
\item 配置修改成本（Figure 13）：阻断机制固定 35 CPU cycle，每 entry 修改 14 CPU cycle，64 个 entry 总修改成本低于 1000 CPU cycle；对照组 IOMMU 的 IOTLB 异步失效可达毫秒级。
\item 硬件资源（Figure 14 及正文）：512 entry 不用树形仲裁需额外 17.3\% LUT 与 1.8\% FF；树形仲裁仅需额外 1.21\% LUT/FF，把 LUT 成本降低 93\%；支持 1024+ entry 的 sIOPMP 整体只多消耗 1.9\% LUT 与 FF。
\item 网络吞吐（Figure 15，iperf）：sIOPMP 与 sIOPMP-2pipe 相比无防护基线影响小于 3\%；IOMMU-strict 单核引入 25\%$\sim$38\% 开销、多核 20\%$\sim$27\%；IOMMU-deferred 开销较小但留下攻击窗口；SWIO（SEV-SNP bounce buffer）牺牲 23\%$\sim$24\% 网络带宽；sIOPMP+IOMMU 组合性能接近 IOMMU-deferred 且相比纯 IOMMU 有 19\% 提升。摘要口径为：相比 IOMMU-based 或软件 I/O 机制提升 20\%$\sim$38\% 网络吞吐。
\item 应用层（Figure 16）：distributed memcached 在相同 50/99 分位延迟要求下 QPS 不受 sIOPMP 影响；冷设备切换单次为 341 CPU cycle（切换 8 个 entry，Figure 17）。
\end{itemize}
本节判断以论文第 6 节与 Figure 10$\sim$17 为依据。证据强度上，这是 peer-reviewed 系统论文，机制、RTL 原型、FPGA 时序与真实系统网络评估齐全；但数字对应其 chipyard/FireSim/特定服务器配置，外推到其他 SoC 时需谨慎。\par
\subsubsection{综合评价}
综合来看，sIOPMP 是 RISC-V trusted I/O 的关键机制论文，但它只解决 data-path protection 的一段。它的优势在于把 DMA/MMIO access control 作为一等硬件 primitive 处理：区域模型匹配 DMA 连续性本质，MT checker 用流水线加树形仲裁绕开组合逻辑瓶颈，mountable entry 与 remapping 把"无限设备"分解为"有限热设备加受控切换"，且全程给出可核验的时序、面积与吞吐数字。它的局限在于设备身份、链路保护、中断与 lifecycle 需要 SPDM/TDISP/IDE/AIA 与 CoVE-IO 组合；多冷设备并发与配置错误的敏感性如上节所述。从部署角度看，它适合 RISC-V SoC、DPU 与加速器平台；风险在标准采纳节奏以及与 IOMMU/TEE-IO 的接口分工。\par
\subsection{论文 40：Using Address Independent Seed Encryption and Bonsai Merkle Trees to Make Secure Processors OS- and Performance-Friendly}
\subsubsection{论文基本信息}
《Using Address Independent Seed Encryption and Bonsai Merkle Trees to Make Secure Processors OS- and Performance-Friendly》由 Brian Rogers、Siddhartha Chhabra、Milos Prvulovic、Yan Solihin 完成（North Carolina State University 与 Georgia Tech），发表于 2007 年的 MICRO 2007。本文将其置于"内存加密、完整性与重放保护"的研究脉络中考察，并把它作为 system 类材料使用。它所回答的核心问题可概括为：BMT 的贡献是把 replay protection 做得更省 cache——树只保护 counter，data 用 MAC 保护。当前证据状态为：本地 PDF 已保存并核验；下文机制与数字均出自该 PDF。需要提醒：这是 2007 年模拟器时代的评估，其数字用于解释机制权衡，不能与现代 CVM 产品横向比较。\cite{rogers2007bonsai}\par
\subsubsection{研究背景与动机}
要理解该方案，首先需要看到它面对的系统矛盾：secure processor 要同时支持 OS 功能和低开销的 memory integrity，而当时的两条技术路线各占一端失效。加密侧，counter-mode 加密的 seed 构造直接决定 OS 兼容性：全局 counter 方案简单，但 64-bit counter 占每块 12.5\% 存储、cache 与预测效果差（32-bit counter 会每几分钟触发全内存重加密）；以虚拟地址为 seed 则不同进程相同虚址会 pad 重用，加 process ID 后又破坏 shared-memory IPC、mmap、fork 的 copy-on-write 与共享库；以物理地址为 seed 则每次 page swap 都要整页解密重加密，且换出到磁盘的页失去保护。完整性侧，传统 Merkle Tree 能防 replay，但树覆盖全部 data block，论文实测在访存密集应用下 MAC/树节点最多可占 1MB L2 cache 的约 50\%，严重挤压数据。\par
\subsubsection{关键挑战与核心思想}
在上述背景下，作者提出的关键判断由两条安全观察支撑。第一，spoofing 与 splicing 只需每个 data block 配一个 MAC 即可检测，Merkle Tree 真正不可替代的职责是防 replay；第二，counter-mode 加密中每个 block 的 counter 随写递增、天然是版本号。由此得到论文的核心 claim（原文第 5.2 节以显式命题形式给出并附论证）：只要（1）每个 block 有以密钥哈希（如 HMAC-SHA-1）计算的 MAC，（2）MAC 输入包含该 block 的 counter 与地址，（3）全部 counter 的完整性有保证，则 data block 无需再由 Merkle Tree 覆盖——攻击者改密文则 MAC 不匹配，连 MAC 一起回滚则因新 counter 参与计算而失败。这把 freshness 的树从"覆盖全部数据"缩小到"只覆盖不可回退的 counter"，而 counter 与 data 的大小比为 1:64，树因此小一到两个数量级，这就是 Bonsai Merkle Tree（BMT）。加密侧的对应物是 Address Independent Seed Encryption（AISE）：seed 不用地址，而用逻辑标识符。\par
\subsubsection{系统架构}
从系统结构看，CPU 片内有可信 root、AES 引擎、counter cache 与 GPC（Global Page Counter），片外内存中存放 data block、counter block、每块 MAC 与 BMT 节点。读数据时并行处理解密、MAC 验证与 BMT 路径验证；写数据时更新 counter、MAC 与树路径。AISE 的 seed 由 LPID（Logical Page Identifier）、page offset、block counter 与 chunk id 拼接成 128 bit（Figure 3）：LPID 取自片内非易失寄存器中的 64-bit GPC，每页首次分配时取新值并递增，跨重启、休眠与断电不重复，从机制上保证 seed 在全系统生命周期、含 swap 空间内唯一。每个 counter block 与一个 4KB 页绑定，存放该页的 64-bit LPID 加 64 个 7-bit block counter，恰好 64 字节、可被 counter cache 直接索引；平均每个 4KB 页只需 64 字节元数据，即 1.6\% 存储开销。block counter 溢出时只为该页取新 LPID 并重加密该页，不触发全内存重加密。\par
\subsubsection{核心方法设计}
论文的贡献由三个相互配合的机制组成，以下拆解均可在原文第 4、5 节核验。\par
\noindent\textbf{（1）Address Independent Seed Encryption（AISE）。}\par
这一机制的直接作用是：把 ciphertext 与固定虚拟/物理地址解绑，消除 OS 管理障碍。具体实现上，page movement 只需连同 LPID 与 counter block 一起搬运，无需重加密，甚至可由 DMA 完成；page sharing、shared-memory IPC、fork 的 copy-on-write 与共享库因 LPID 按物理页唯一而天然可用；换出到 swap 的页保留 LPID，密文不需特殊处理。原文 Table 1 将 AISE 与 global counter、物理地址、虚拟地址三种 seed 方案在 IPC 支持、延迟隐藏、存储开销与其他系统问题上逐一对比：AISE 在全部维度不劣于任一方案，存储开销 1.6\%，为四者最低。\par
\noindent\textbf{（2）Bonsai Merkle Tree。}\par
这一机制的直接作用是：把 replay 保护的树规模从"全部数据"缩减到"counter 集合"。具体实现上，data block 由各自 MAC 覆盖（MAC 输入含密文、counter 与地址，$M=H_K(C, ctr)$），BMT 只覆盖 counter block；片上可信 root 保护整棵树。原文进一步做了两个 cache 策略决定：不缓存 data block MAC（每个 MAC 只覆盖 4 个 block，时间局部性差），只缓存 BMT 节点。论文同时给出把 Merkle 保护扩展到 swap 的方法（5.1 节）：为磁盘上的页维护 Page Root Directory，页根目录本身再被树上的 Merkle 保护，换入换出时按页根校验，避免依赖可信 OS 或每进程一棵树。\par
\noindent\textbf{（3）Cache pollution 治理与存储权衡。}\par
这一机制的直接作用是：把完整性元数据对 cache 与内存的占用量化并压到接近无防护方案。论文 Figure 9 实测：标准 MT 下 data 平均只占 L2 的 68\%（其余 32\% 被树节点占据，art、swim 等极端负载接近 50\%），因为上层树节点时间局部性高于数据；BMT 下 data 占 98\%，BMT 节点只占 2\%。代价是主存存储开销（Table 2）：128-bit MAC 时 AISE+BMT 总开销 21.55\%（其中 BMT 节点 20.02\%、page root 目录 0.31\%、counter 1.23\%），仍显著低于 global64+MT 的 33.51\%。\par
\subsubsection{安全性与威胁边界}
从安全边界看，该方案在其威胁模型内提供 memory spoofing、splicing 与 replay 三类攻击的检测：spoofing/splicing 由 MAC 检出，replay 由 BMT root 检出，论证链条依赖 counter 完整性与密钥哈希不可伪造两个前提。它没有覆盖的部分包括：片内侧信道与推测执行（2007 年威胁模型之外）、故障注入与 Rowhammer 类物理攻击、多租户 CVM 的设备与 attestation lifecycle，以及 crash consistency——断电时 counter/MAC/树路径的原子持久化不在论文范围内。GPC 存于片内非易失寄存器以维持跨重启唯一性，但论文未讨论该寄存器本身被物理篡改的情形。\par
\subsubsection{实验环境与证据}
论文实验（第 6 节 Experimental Setup）使用 SESC 开源 execution-driven 模拟器，建模 2GHz、3-issue 乱序处理器：L1 I/D 各 32KB、2 路、2 cycle 命中；L2 统一 1MB、8 路、10 cycle；L2 级另设 32KB、16 路 counter cache；全部 cache 64B 块、LRU；主存 1GB、访问延迟 200 cycle。加密引擎为 128-bit AES、16 级流水、总延迟 80 cycle；MAC 建模为基于 SHA-1 的 HMAC、80 cycle；默认 MAC 128 bit。工作负载为 21 个 C/C++ SPEC2000 benchmark（Fortran 90 因基础设施不支持而省略），reference 输入，fast-forward 50 亿条指令后模拟 10 亿条；图中只单列 L2 miss 率超过 20\% 的 benchmark，平均值按全部 21 个计算。完整性验证采用 timely but non-precise 策略（块上片即验，但不阻塞触发指令退休），论文声明两种验证策略均兼容。\par
\subsubsection{实验结果与证据强度}
论文的主要定量结论如下，均可对照原文图表核验：\par
\begin{itemize}
\item 总体开销（Figure 6）：global64+MT 平均执行时间开销 25.9\%、最大 151\%；AISE+BMT 平均 1.8\%、最大 13\%。
\item 分解（Figure 7、8）：仅加密的 AISE 平均 1.6\%，优于 32/64-bit 全局 counter 的约 4\%/6\%；固定 AISE 后，标准 MT 把完整性开销推到 12.1\%，BMT 压回 1.8\%。访存密集的 art、mcf、swim 上 BMT 开销低于 15\%，而标准 MT 可超 60\%。
\item 机理证据（Figure 9、10）：标准 MT 使 L2 miss 率从 37.8\% 升至 47.5\%、总线利用率从 14\% 升至 24\%；BMT 仅升至 38.5\% 与 16\%。
\item MAC 位宽敏感性（Figure 11）：MAC 从 32 增至 256 bit 时，MT 平均开销从 3.9\% 近指数升至 53.2\%，BMT 仅从 1.4\% 升至 2.4\%；L2 中 data 占比 MT 从 89.4\% 跌到 36.3\%，BMT 仅从 99.5\% 跌到 94.9\%。
\item 存储开销（Table 2）：128-bit MAC 时 AISE+BMT 总开销 21.55\%，global64+MT 为 33.51\%；256-bit MAC 时分别为 35.03\% 与 55.71\%。AISE+BMT 在 256-bit MAC 下少用 1.6 倍内存，32-bit 下优势扩到 2.3 倍。结论节口径为：执行开销从 12.1\% 降到 1.8\%，存储开销从 33.5\% 降到 21.5\%。
\end{itemize}
证据强度上，消融（加密与完整性分离）、机理（cache 污染、miss 率、总线）与敏感性（MAC 位宽）三层证据互相咬合，机制结论可信；但全部数字来自 2007 年的 SESC 模拟器与 SPEC2000，平台假设（200 cycle 内存延迟、80 cycle AES/HMAC、单核）与现代多核 CVM 差距大，跨时代比较只能引用其趋势而非绝对值。\par
\subsubsection{与相关工作的定量对比}
在论文自身的谱系里，全球 counter 加标准 MT 是唯一与 AISE+BMT 提供同等系统级支持（shared memory、虚存、swap）的对照组，Figure 6 的 25.9\% 对 1.8\% 即二者差距；与同期工作相比，论文指出已有低开销加密完整性方案在访存密集负载上开销仍接近 20\% 且只假设 64-bit MAC，而 BMT 在保持标准 MT 安全强度的前提下把这一区间压到 15\% 以内。Table 2 的横向对比还显示 BMT 的存储优势随 MAC 变长而收窄（32-bit 时 2.3 倍、256-bit 时 1.6 倍），说明其剩余成本主要来自每块 MAC 而非树本身——论文明确表示与"一个 MAC 覆盖多个 block"等降存储技术兼容。\par
\subsubsection{综合评价}
综合来看，BMT 是 memory integrity/freshness 的经典 SOTA 基线，至今仍适合解释为什么 metadata design 决定性能。它的优势在于机制清晰（tree-over-counter 一句话可概括）、安全论证显式、实验量化充分且直接命中标准 MT 的 cache pollution 病灶。它的局限在于年代与平台：模拟评估基于 SPEC2000 与单核乱序假设，不覆盖现代多租户 CVM、CXL 内存池、持久内存 crash consistency 与侧信道；1.6\%/21.55\% 这类数字应读作"机制相对优劣的证据"，而非现代部署的成本预期。从部署角度看，tree-over-counter 思想可用于 CPU/TEE/存储的新鲜度设计；落地必须结合现代 cache 层级、内存控制器与固件持久化策略重新评估。\par
\subsection{论文 41：Secure Encrypted Virtualization Secure Nested Paging (SEV-SNP)}
\subsubsection{论文基本信息}
《Secure Encrypted Virtualization Secure Nested Paging (SEV-SNP)》由 AMD 完成，元数据记录为 2020 年 AMD 白皮书（《Strengthening VM Isolation with Integrity Protection and More》）。本文将其置于"内存加密、完整性与重放保护"的研究脉络中考察，并把它作为 vendor 类材料使用。它所回答的核心问题可概括为：AMD SEV-SNP 方向的贡献是把 VM 机密性从 memory encryption 推向 guest state、ownership 和 hypervisor interaction hardening。当前证据状态需要特别说明：本地保存并核验的 PDF 实为 AMD 2017 年 2 月白皮书《Protecting VM Register State with SEV-ES》（David Kaplan），其图示与机制叙述覆盖 SEV-ES/GHCB/NAE，而 SEV-SNP 的 RMP、VMPL 等具体机制不在本地 PDF 覆盖范围内；下文凡超出本地材料的内容均标注证据不足，细节须以 AMD APM 与 SNP 规范另证。\cite{amd_sev_snp}\par
\subsubsection{研究背景与动机}
要理解该方向，首先需要看到它面对的系统矛盾：hypervisor 即使看不到加密内存，仍能观察或操纵 VM exit 时的 CPU 状态。本地白皮书 Figure 1 给出具体攻击场景：guest 使用 x86 AES 指令时 AES key 常驻 XMM 寄存器，VM 因中断停止运行时寄存器内容被保存到 hypervisor 可读的内存，恶意或被攻破的 hypervisor 读取 XMM 即可窃取密钥、静默外泄数据。写方向同样危险：hypervisor 改写 RIP 等控制寄存器可让 guest 跳过关键安全检查；replay 旧状态则破坏执行新鲜性。简言之，SEV 的内存加密把攻击面从 DRAM 压到了"VM 停止运行的那一瞬间"，SEV-ES/SNP 这条产品线就是围绕这一刻的状态保护逐步收敛的。\par
\subsubsection{关键挑战与核心思想}
在上述背景下，白皮书提出的关键判断是：VM isolation 要保护 data、CPU state 和（在 SNP 方向上）memory ownership，而不只是 DRAM ciphertext。其设计难点在于 hypervisor 并非完全不需要 guest 状态——device emulation、MSR 处理等服务要求受控访问，且这种访问难以审计。SEV-ES 的回答是把"共享什么"的决定权从 hypervisor 移到 guest：VMEXIT 时全部寄存器状态由硬件原子保存并加密，guest 需要 hypervisor 服务时通过协议化 buffer 显式共享最小状态集合，这与 SEV 中 guest 自主决定共享哪些内存页的原则一致。SNP/RMP 方向进一步把 host 对 guest page ownership 与完整性的攻击纳入防护，但本地白皮书只支撑这一方向性叙述，RMP 的表结构、页验证流程与 VMPL 权限分级等细节本地材料未覆盖，证据不足。\par
\subsubsection{系统架构}
从系统结构看，白皮书描述的运行路径为：guest 执行中触发需要 hypervisor emulation 的事件（Non-Automatic Exit，NAE）时，不再发生 world switch，而是向 guest 投递一个新的 \#VC（VMM Communication Exception）异常；guest 的 \#VC handler 决定把哪些寄存器/信息写入位于共享内存页的 GHCB（Guest Hypervisor Communication Block），然后用 VMGEXIT 指令触发一次 AE 式退出把控制交给 hypervisor；hypervisor 读取 GHCB 完成 emulation、把结果写回 GHCB，再以 VMRUN 恢复 guest；guest 仍在 \#VC handler 内，检查 hypervisor 写回的状态是否可接受后才拷贝到真实寄存器并 IRET 返回。Figure 2 给出 VMCB 的两分结构：control area（不加密，hypervisor 管理，含 intercept bits、interrupt delivery、ASID、paging 信息）与 save area（加密且完整性保护，含 segment、control、system 寄存器、GPR 与浮点寄存器）。Figure 3 给出 NAE 示例流（CPUID 场景共享 RAX，OUTW 场景共享 AX 与 DX，返回 RAX/RBX/RCX/RDX）。\par
\subsubsection{核心方法设计}
白皮书的关键不是单一组件，而是三层机制的组合。以下拆解均可在本地 PDF 核验，第四层（SNP）单独标注证据边界。\par
\noindent\textbf{（1）Encrypted State / SEV-ES。}\par
这一机制的直接作用是：把 VMEXIT 时的 guest register state 从 hypervisor 视野里拿掉。具体实现上，SEV-ES 把 legacy AMD-V 中 VMRUN/VMLOAD/XRSTOR 等多步状态切换合并为单条原子 VMRUN：进入 guest 时硬件一次性装载全部寄存器（系统寄存器、GPR、浮点），VMEXIT 时硬件自动把全部状态保存回内存并装载 hypervisor 状态，world switch 不可被中断。保存的寄存器状态用该 VM 的内存加密密钥加密（任何 CPU 软件包括 hypervisor 都拿不到该密钥），且 CPU 计算一个 integrity-check value 存于任何 CPU 软件都不可访问的受保护 DRAM；VMRUN 恢复时校验该值，不匹配则拒绝恢复 guest——白皮书原文的结论表述为：guest 只能被恢复到"它上次退出时完全相同的状态"，由此同时阻断读取、篡改与 replay 三类寄存器攻击。\par
\noindent\textbf{（2）GHCB / \#VC 协议。}\par
这一机制的直接作用是：让 guest 对 hypervisor 的信息共享显式化、最小化、可审计。具体实现上，VMEXIT 事件被分为 Automatic Exit（AE：异步中断、shutdown、某些 page fault，不需要 emulation，走完整 world switch）与 Non-Automatic Exit（NAE：MMIO、MSR 访问等需要 hypervisor emulation 的 guest 行为，不走 world switch 而走 \#VC）。\#VC handler 在运行时按事件类型决定暴露哪些状态：CPUID 共享 RAX，OUTW 共享 AX 与 DX；GHCB 结构本身未被架构显式定义，白皮书建议镜像 VMCB save area 以便双方解析；VMGEXIT 是 guest 主动交权的新指令，返回后 handler 必须校验 hypervisor 写回的值再决定是否接受。白皮书同时指出该设计的软件收益：instruction cracking 等 emulation 工作移入 guest 侧 handler，hypervisor 的 emulation 支持面反而缩小，且 guest 内驱动无需修改，所有 hypervisor 通信收敛到单一 handler；性能侧 handler 可缓存 CPUID 结果、把多个请求合并为一次 VMGEXIT。\par
\noindent\textbf{（3）启动度量与 attestation 接口。}\par
这一机制的直接作用是：把寄存器态保护延伸到 VM 启动时刻。具体实现上，SEV-ES VM 的初始内存镜像与初始 CPU 寄存器状态都必须由 AMD Secure Processor（AMD-SP）加密后才能开始执行；启动过程中 AMD-SP 对二者做密码学度量，生成 launch receipt 供 attestation 使用，使 VM owner 在投放 secrets 前确认镜像与初始寄存器状态正确。\par
\noindent\textbf{（4）SNP / ownership 方向（证据边界声明）。}\par
这一方向的要点是进一步限制 hypervisor 对 guest page ownership 与完整性的攻击面：RMP（Reverse Map Table）一类 ownership 表与现代 SNP 叙事中的页验证/分配防止 host 把错误页映射给 guest。但必须强调：本地核验 PDF（2017 年 SEV-ES 白皮书）不含 RMP、VMPL 或 SNP ABI 的任何机制细节，本综述只能从 README 与产业背景确认其方向性；页粒度、VMPL 等级语义、PVALIDATE/RMPUPDATE 类指令行为均属本地材料未说明，须以 AMD64 APM Volume 2 与 SEV-SNP ABI 规范为准，此处不作展开。\par
\subsubsection{安全性与威胁边界}
从安全边界看，本地白皮书能够支撑的防护声明是：针对恶意 hypervisor 的寄存器读取（Figure 1 的 AES key 窃取场景）、寄存器写入（改写 RIP 跳过安全检查）与状态 replay 三类攻击，SEV-ES 通过加密、完整性校验值与原子 world switch 予以阻断；GHCB 把剩余必须共享的信息收敛到 guest 显式选择的字段。不能从本地材料推出的包括：SNP 对页级完整性/ownership 攻击的具体防护强度、性能开销、形式化安全性质，以及侧信道、中断完整性、设备 I/O（SEV-TIO 之前的 bounce buffer 问题见论文 39）等运行期攻击面。白皮书的证据等级是 industry evidence，应按产品方向材料使用。\par
\subsubsection{实验环境与证据}
本地材料无新实验，只有机制图示与文字说明：Figure 1（从 VM 窃取 AES key 的寄存器路径）、Figure 2（VMCB 两分结构）、Figure 3（NAE 示例流）。能够直接支撑的结论是 SEV-ES 的原子状态切换、寄存器态加密与完整性校验、AE/NAE 分类、\#VC/GHCB/VMGEXIT 协议与启动度量流程；不能支撑任何性能数字、SNP ABI 细节或与其他 CVM 方案的定量对比。\par
\subsubsection{实验结果与证据强度}
白皮书没有 benchmark。潜在开销来源可以从机制推断——\#VC/GHCB 的 guest 内处理、VMGEXIT 往返、页验证与（SNP 方向上的）RMP 检查——但真实开销取决于 workload、VMEXIT 频率、I/O emulation 比例与内核实现，本综述不编造数字。论文 39 的评估提供了相邻证据：SEV-SNP 式 SWIO（bounce buffer 加 hypervisor 介入）在其测试环境（AMD EPYC 7T83、100Gbps NIC）中牺牲 23\%$\sim$24\% 网络带宽，这是 I/O 路径的代价而非寄存器态保护本身的代价，二者不应混用。\par
\subsubsection{综合评价}
综合来看，AMD SEV-ES/SNP 这条材料线是理解现代 CVM 商业化的关键工业证据，但证据等级不同于 peer-reviewed 论文。它的优势在于贴近真实产品：把"hypervisor 不可信"从内存加密推进到寄存器态与启动度量，并用 AE/NAE 分类加 GHCB 协议给出了 guest 与 hypervisor 协作的最小接口，\#VC 集中 handler、驱动免改造、CPUID 缓存与请求合并等设计显示了对工程落地的考虑。它的局限在于白皮书粒度：本地材料止步于 SEV-ES，SNP 的 RMP/VMPL 细节、性能与形式化性质均需配套 AMD 架构手册、Linux/KVM 实现与后续安全分析另证。从部署角度看，SEV 系列已是云 CVM 生态核心路线；风险集中在 firmware（AMD-SP）、attestation 链与 guest 侧 \#VC 实现的正确性上。\par
\subsection{开发商安全实践建议}
本章四份材料覆盖 ISA 语义、I/O 检查电路、内存完整性元数据与 VM 态保护四个层次，对 CPU/MCU 设计商与固件厂商的可操作含义如下。\par
\noindent\textbf{实践一：启动期尽早用 PMP/Smepmp 把"默认拒绝"写死，而不是事后补救。} Privileged spec 的 PMP 默认语义是 S/U 无权限、M 全权，Smepmp 的 mseccfg.MMWP/MML 则允许把 M-mode 自身也纳入默认拒绝并锁定到复位\cite{riscv_privileged}。固件厂商应在 BootROM 阶段即置位这些字段并为安全 monitor、固件自身与 REE 划分 locked entry，利用 L 位"锁定到 hart 复位"的性质防止运行期篡改；PMP entry 有限（实现可选 0/16/64 个），区域规划必须在设计期完成，且修改 PMP 后必须执行规范要求的 SFENCE.VMA 同步序列。\par
\noindent\textbf{实践二：trap delegation 按最小权限配置，并把配置纳入审计。} medeleg/mideleg/hedeleg/hideleg 的每一位都决定一类事件落入可信还是不可信处理者\cite{riscv_privileged}。建议固件只下放 REE 必需处理的异常（如普通 page fault），保留安全相关异常与中断在 monitor；同时用 Smstateen 关闭未被 OS/hypervisor 感知的新扩展状态（规范以 AIA 的十个 supervisor CSR 为 covert channel 示例），用 Smcntrpmf 的 xINH 位阻断计数器向低特权级泄漏 privileged 执行信息。\par
\noindent\textbf{实践三：I/O 防护与 CPU 防护分层设计，检查放在设备数据通路上。} sIOPMP 的教训是 CPU 页表模型不适合 DMA：设备访问天然连续，区域式检查才匹配 workload\cite{feng2024siopmp}。SoC 设计商应采用区域式 IOPMP 类检查器（流水线加树形仲裁可把 1000+ entry 的检查压进时序预算，面积代价约 1.9\% LUT/FF），让 IOMMU 只做地址翻译；固件必须按 workload 正确维护设备冷热状态——论文实测状态误配在 1:10 请求比下浪费 85\% 热设备吞吐——并把 entry 修改期间的 per-SID 阻断纳入配置流程，杜绝新旧区域同时可访问的窗口。\par
\noindent\textbf{实践四：内存加密必须配套完整性与新鲜度，元数据布局按 tree-over-counter 设计。} BMT 的结论是：MAC 管 data、Merkle Tree 只管 counter，可以把完整性开销从两位数压到约 2\%（2007 年模拟数字），而纯内存加密不防 replay\cite{rogers2007bonsai}。CPU 设计商实现内存加密时应为 counter/MAC/树节点预留主存预算（128-bit MAC 下约 21.55\% 的历史数字提示成本量级），并保证 counter 完整性链条的根在片内；固件需处理 counter 溢出时的局部重加密与跨启动的唯一性（AISE 用片内非易失 GPC 解决），断电持久化与 crash consistency 则需在部署层另行设计。\par
\noindent\textbf{实践五：VM/机密执行场景的 CPU 状态切换要原子化、加密化，hypervisor 交互走协议化最小共享。} SEV-ES 的设计表明，仅加密内存会把攻击面挤到 VMEXIT 的寄存器保存路径\cite{amd_sev_snp}。CPU 设计商应把状态切换做成不可中断的原子操作、状态存取走加密加完整性校验值的路径；系统软件厂商实现 guest 侧 handler 时，必须把"共享哪些字段"做成按事件类型的显式最小集合，并对 hypervisor 写回的每个值做合法性校验后再接受——白皮书明确把校验责任放在 guest 侧。\par
\noindent\textbf{实践六：调试接口与不透明状态默认锁定，按证据链验收供应商声明。} Privileged spec 含独立 Debug Mode 语义（1.3 节），调试口一旦在生产件上开放即绕过前述全部权限语义\cite{riscv_privileged}；设备制造商应在出厂流程中熔断或认证化调试访问，并把 state-enable 覆盖范围列入安全评审。同时，对供应商材料要按证据等级使用：SEV-SNP 类白皮书可确认产品方向，但 RMP/VMPL 级细节与性能必须回到架构手册与可复现实验验收，本综述对本地未覆盖内容一律标注证据不足即是范例\cite{amd_sev_snp}。\par
\subsection{未来研究方向}
\noindent\textbf{（1）PMP/Smepmp 配置正确性的形式化验证与自动检查。} 问题是：规范定义了 64 entry、锁定规则、delegation 与 SFENCE 同步的语义，但"某张具体配置表是否实现了预期隔离"完全靠人工\cite{riscv_privileged}；现有工作未解决是因为语义散落于 221 页规范的多个章节，缺少可机读的统一模型。可行路线是从规范抽取 PMP/Smepmp/delegation 的状态机，用 Sail 或自研模型对 TEE monitor 的配置序列做属性验证（锁定不可绕过、S/U 默认拒绝、MPRV 组合语义），并自动生成 CSR misuse checker。需要的实验是在 Keystone/Penglai 类真实 monitor 上回放配置序列，注入位翻转与顺序错误，验证检测率与误报率。\par
\noindent\textbf{（2）I/O 隔离与 TEE-IO lifecycle 的统一建模。} 问题是：sIOPMP 只解决 DMA 访问控制一段，设备身份（SPDM/TDISP）、链路保护（PCIe IDE）、中断投递（AIA）与 attestation 各自为政，组合后的端到端保证无人证明\cite{feng2024siopmp}。现有工作未解决是因为各层标准分属不同组织、状态机各自定义。可行路线是把 sIOPMP 的 domain/entry 状态机与 CoVE-IO 的 TDI/TDM/DSM 生命周期统一成一张 ownership 迁移图，在 FireSim 原型上接入 SPDM/TDISP 栈做端到端评估。需要的实验覆盖设备热插拔、VF 动态创建、多冷设备并发切换（论文明确未覆盖）下的安全性与吞吐。\par
\noindent\textbf{（3）面向现代 CVM 的内存完整性元数据重评估。} 问题是：BMT 的 1.8\%/21.55\% 数字来自 2007 年单核模拟，现代场景新增多租户、CXL 内存池、持久内存 crash consistency 与侧信道变量\cite{rogers2007bonsai}；未解决是因为经典设计假设单一安全域与易失 DRAM。可行路线是把 tree-over-counter 映射到 CCA/CoVE/TDX/SEV-SNP 的 granule 模型，研究 counter/MAC/树的断电原子持久化与跨启动 rollback 防护，并评估树节点缓存与侧信道的互相影响。需要的实验是在现代模拟器（gem5/ChampSim）上用多核多租户负载重测 cache pollution 曲线，并对持久内存做故障注入。\par
\noindent\textbf{（4）guest 态保护协议的跨 ISA 形式化。} 问题是：SEV-ES 的 \#VC/GHCB 协议把校验责任放在 guest 侧，但"handler 是否对每个事件都校验了所有必要字段"没有形式化证据，本地白皮书更不含 SNP 的 RMP/VMPL 语义\cite{amd_sev_snp}；未解决是因为协议规范以白皮书和内核补丁形式存在，粒度不足以机器检验。可行路线是从 AMD APM 与 Linux/KVM 实现反向提取 GHCB 字段级协议模型，对返回值校验完备性做符号执行或模型检验，并与 Arm CCA 的 RSI、RISC-V CoVE 的 TSM 调用面做对照。需要的实验包括注入恶意 hypervisor 返回值（越界、重放、类型混淆）测量 guest 侧拒绝率。\par
\noindent\textbf{（5）可扩展检查电路的时序侧信道与 DoS 面分析。} 问题是：sIOPMP 的流水线与树形仲裁把安全检查变成多周期电路，检查延迟是否泄漏地址模式、CAM remapping 与 SID-missing 中断是否可被恶意设备用作 DoS 放大，论文显式将侧信道列为正交问题而未评估\cite{feng2024siopmp}。可行路线是在 RTL 级测量不同地址/权限组合下的检查延迟差，并建模冷设备切换风暴（伪造大量陌生 SID 触发 mount 中断）对 monitor 的占用。需要的实验是在 FireSim 上用恶意 DMA 发生器做延迟侧信道采样与中断洪泛，给出缓解（限速、批量 mount）后的开销。\par
\noindent\textbf{（6）计数器与调试通道的系统性泄漏枚举。} 问题是：Smcntrpmf 已确认 cycle/instret 跨特权级计数会泄漏 privileged 执行信息，Smstateen 以 AIA CSR 为例确认扩展状态可成 covert channel\cite{riscv_privileged}，但完整的"哪些微架构状态在 TEE 边界上可见"清单不存在；未解决是因为状态随扩展不断增加。可行路线是建立扩展状态与特权边界的登记册，对每项新扩展强制回答"谁可见、谁切换、谁清零"，并在模拟器上做跨特权级差分测量。需要的实验是对照 Smstateen 开关逐一测量各扩展状态的跨域可观测性，形成可回归的泄漏测试集。\par
\section{综合总结}

逐条阅读 PPT 选出的 41 个独立来源后可以看到，当前硬件安全研究的关键矛盾并不是“是否存在安全硬件”，而是分散在处理器、内存、设备、互连和验证方之间的机制能否在同一平台边界内正确组合。Arm CCA/RME/RMM、RISC-V CoVE/AP-TEE、SPDM、TDISP、IOMMU/IOPMP/SMMU、可信中断、加速器 TEE、attestation、MTE、PMP、能力架构和内存完整性分别缩小了不同层次的攻击面；它们不能彼此替代，也不能因为其中某一层成立就自动推出系统整体安全。TrustZone 漏洞谱系与 ret2ns 的研究从反面印证了这一点：隔离原语成立而组合失效，是真实漏洞的主要来源。

七个类别的逐篇分析呈现出三条贯穿性结论。第一，所有权状态机是所有方案的公共核心。CCA 的 granule、CoVE 的 MTT、SEV-SNP 的 RMP、SPEAR-V 的 page tag、sIOPMP 的访问域，本质上都在回答“这块物理资源此刻属于谁、谁能合法改变这个归属”；各章的失败模式分析也显示，安全问题集中出现在状态转换而非稳态。第二，设备路径是当前证据最薄弱的一段链。CPU 侧隔离已有架构规范、形式化验证和多个低开销原型，而设备身份、DMA 授权、中断投递、链路加密与 attestation 的端到端绑定仍由 draft 规范（CoVE-IO、AP-TEE）和早期原型（ACAI、Devlore、CAGE、TNIC、S-NIC、Hazel）支撑，量产证据稀缺。第三，证据成熟度本身构成研究瓶颈。本报告相当比例的性能数字来自 FVP、Gem5、SESC 或开发板回植；规范类材料中两份关键 Arm 文档本地不可核验，两份 RISC-V 规范尚未 ratified。把原型结论当作产品结论，是当前文献阅读中最需要警惕的推论方式。

从开发商视角看，各章建议可以收敛为一条共同主线：把安全机制的选择前移为平台对象模型的设计决策，而不是事后补丁。芯片与 SoC 厂商需要决定的是所有权检查在哪一层硬件路径上执行、TCB 固件有多小、调试与追踪接口如何在量产时锁定；固件与系统软件开发商需要决定的是状态转换如何审计、TCB 如何验证、证据如何绑定会话与生命周期；设备制造商（含网络设备厂商）需要决定的是设备身份如何注入、固件如何度量、接口暴露面如何版本化管理。各章“开发商安全实践建议”已按类别展开这些要求，其共同点是把论文中的攻击面分析转化为研发流程中的检查项，而非把某一篇论文的机制当作可直接采购的安全结论。

从论文证据走向工程落地，后续研究首先需要把 CPU、内存、设备、网络链路和持久存储的 ownership/lifecycle 状态机统一起来，使创建、共享、撤销和回收不再在不同层各自定义。其次，设备身份、DMA 权限、中断投递、链路保护和应用端点证明应被绑定到同一 workload，而不是由多个独立控制面松散拼接。再次，规范语义、形式化验证、开源实现和真实硬件实验需要形成可相互校验的证据链；对 draft、vendor、metadata-only 和 source-limited 材料，则应持续进行版本与可复现性审计。各章“未来研究方向”列出的具体问题——量产 CCA 硬件上的全栈性能复现、RMM 扩展组件的验证复用、设备 attestation 与 Realm 生命周期的端到端绑定、多租户加速器共享、I/O 隔离与 TEE-IO 生命周期的统一模型、内存完整性元数据在现代 CVM 下的重新评估等——都是从本报告的证据缺口中直接推导的，而非泛泛的方向罗列。

报告中保留的“论文未说明”和“证据不足”，正是后续深入工作的入口。若要将某一方向继续扩展为算法、图表、伪代码或复现实验级别的研究，应先回到 E1 原始论文和 E0 规范确认语义与威胁模型，再以 E2 Survey/SoK 校准分类，最后将 E3/E4/E5 作为标准化或工程状态证据。只有在这条证据链闭合后，具体方案的性能收益与安全边界才具有可比较、可复验的意义。
\bibliographystyle{IEEEtran}
\bibliography{reference}
\end{document}
